\documentclass[11pt, a4paper, titlepage, oneside]{scrbook}

\usepackage{../Tarc/Tarc}
\usepackage{../Tarc/Tarc-private}
\graphicspath{{immagini/}}

\title{Algebra Lineare e Geometria}
\author{Lorenzo Ferraro}
\date{\today}
\setDegree{Informatica}

\begin{document}

\frontmatter
\SapienzaCoverItalianoTriennale
\InfoSegnalazioni

\tableofcontents

\mainmatter

\chapter{Introduzione}

\section{Ricominciare da zero}

In queste prime pagine rivedremo la matematica da zero, senza dare nulla per scontato. 
Per fare ciò, è necessario descrivere le due operazioni fondamentali che tutti già conosciamo: la \textbf{somma} e il \textbf{prodotto}. \\
Preso un insieme $R$:

\bigskip

\begin{minipage}[c]{0.45\textwidth}
    \centering
    {\Large\textbf{Somma}:}
    \begin{align*}
        \hspace{0.4cm}R \times R &\longrightarrow R \\
        \hspace{0.4cm}(a,b) &\longmapsto a + b
    \end{align*}
\end{minipage}
\begin{minipage}[c]{0.45\textwidth}
    \centering
    {\Large\textbf{Prodotto}:}
    \begin{align*}
        R \times R &\longrightarrow R \\
        (a,b) &\longmapsto a \cdot b
    \end{align*}
\end{minipage}

\bigskip

Definite ora queste operazioni, vediamo alcune loro proprietà.

\begin{Proprieta}[Proprietà di Somma e Prodotto]
\label{prop-somma-prod}
\begin{enumerate}
    \item $\bm{0 + a = a} \enspace \forall a \in R$ \enspace (Esistenza dell'elemento neutro della somma)
    \item $\bm{a + (b + c) = (a + b) + c} \enspace \forall a,b,c \in R$ \enspace (Proprietà associativa della somma)
    \item $\forall a \in R \enspace \exists b \in R : \bm{a + b = 0}$ \enspace (Esistenza dell'elemento opposto)
    \item $\bm{a + b = b + a} \enspace \forall a,b \in R$ \enspace (Proprietà commutativa della somma)
    \item $\bm{1 \cdot a = a} \enspace \forall a \in R$ \enspace (Esistenza dell'elemento neutro del prodotto)
    \item $\bm{a \cdot (b \cdot c) = (a \cdot b) \cdot c} \enspace \forall a,b,c \in R$ \enspace (Proprietà associativa del prodotto)
    \item $\bm{a \cdot b = b \cdot a} \enspace \forall a,b \in R$ \enspace (Proprietà commutativa del prodotto)
    \item $\bm{a \cdot (b + c) = a \cdot b + a \cdot c} \enspace \forall a,b,c \in R$ \enspace (Proprietà distributiva)
\end{enumerate}
\end{Proprieta}

\begin{Osservazione}[]
\label{oss-prop-somma-prod}
\begin{enumerate}
    \item $\bm{\exists ! 0 \in R}:$ vale la proprietà 1), ossia l'elemento neutro della somma è unico
    \item $\bm{\forall a \in R \enspace \exists! b \in R}:$ vale la proprietà 3), ossia l'opposto di ogni elemento è unico, e diciamo che $\bm{b = -a}$
    \item $\bm{\exists! 1 \in R}:$ vale la proprietà 5), ossia l'elemento neutro del prodotto è unico
\end{enumerate}
\end{Osservazione}

\Dimostrazione 

\begin{enumerate}
    \item Poniamo $0_1, 0_2 \in R$ tali che $0_1 + a = a$ e $0_2 + a = a \enspace \forall a \in R$. Allora $0_1 + 0_2 = 0_1$,
    e per la proprietà 4) troviamo che $0_2 + 0_1 = 0_1$, ma sappiamo che $0_2 + 0_1  = 0_2$ per la definizione stessa di $0_2$. Dunque $0_2 = 0_1 = 0$. $\blacksquare$
    \item Poniamo $v,w_1,w_2 \in R$ tali che $v + w_1 = 0$ e $v + w _2 = 0$. Per la proprietà 1) sappiamo che $w_1 = w_1 + 0$, e sappiamo che $v + w_2 = 0$ per la definizione stessa di $v$ e $w_2$. 
    Troviamo quindi che $w_1 = w_1 + 0 = w_1 + (v + w_2)$ che è pari a $(w_1 + v) + w_2$ per la proprietà 2). Per la proprietà 4) troviamo quindi che $w_1 = (w_1 + v) + w_2 = 0 + w_2$ per la definizione stessa di $w_1$ e $v$ e dato che per la proprietà 1) $0 + w_2 = w_2$, troviamo che $w_1 = w_2 = w$. $\blacksquare$
    \item Ponaimo $1_1, 1_2 \in R$ tali che $1_1 \cdot a = a$ e $1_2 \cdot a = a$. Allora $1_2 = 1_1 \cdot 1_2$, 
    e per la proprietà 7) troviamo che $1_2 \cdot 1_1 = 1_2$, ma sappiamo che $1_2 \cdot 1_1 = 1_1$ per la definizione stessa di $1_1$. Dunque $1_2 = 1_1 = 1$. $\blacksquare$
\end{enumerate}

\section{Anelli e campi}

\begin{Definizione}[Anello]
\label{def-anello}
Un \textbf{anello commutativo con unità} (o più semplicemente \textbf{anello}) è un insieme $R$ con due operazioni ben definite, come possono essere le operazioni di somma e prodotto appena viste, che rispetta le Proprietà \ref{prop-somma-prod}.
\end{Definizione}

Anelli sono gli insiemi $\Z, \Q, \R$, con le operazioni classiche di somma e prodotto che tutti conoscimao, ma non l'insieme $\N$, dato che in questo insieme non è rispettata la proprietà 3), ossia l'esistenza dell'opposto.

\begin{Definizione}[Campo]
\label{def-campo}
Un anello $\K$ è un \textbf{campo} se $\forall a \neq 0 \ \exists b \in \K : \bm{a \cdot b = 1}$, ossia se esiste l'inverso moltiplicativo per ogni elemento diverso da $0$. 
\end{Definizione}

Campi sono gli anelli $\Q$ e $\R$, ma non l'anello $\Z$.

\begin{Definizione}[Insieme delle classi resto modulo $\bm{n}$]
\label{def-Z/(n)}
$\bm{\Z/(n)}$ è \textbf{l'insieme delle classi resto modulo} \bm{$n$}, ossia l'insieme delle classi di equivalenza
da $[0]$ a $[n-1]$ le quali rappresentano tutti gli interi che, divisi per $n$, hanno come resto quel numero.\\ 
Ad esempio, $[1]$ è la classe d'equivalenza dove ci sono tutti gli interi che divisi per $n$ hanno come resto $1$.
\end{Definizione}
È facile notare che quindi $[x]_n = [n-x]_n$, dove con $[x]_n$ rappresentiamo una classe resto modulo n.

Definiamo le operazioni di \textbf{somma} e \textbf{prodotto} su $\Z/(n)$.

\bigskip

\begin{minipage}[c]{0.45\textwidth}
    \centering
    {\Large\textbf{Somma}:}
    \begin{align*}
        \hspace{0.4cm}\Z/(n) \times \Z/(n) &\longrightarrow \Z/(n) \\
        \hspace{0.4cm}([a],[b]) &\longmapsto [a + b]
    \end{align*}
\end{minipage}
\begin{minipage}[c]{0.45\textwidth}
    \centering
    {\Large\textbf{Prodotto}:}
    \begin{align*}
        \Z/(n) \times \Z/(n) &\longrightarrow \Z/(n) \\
        ([a],[b]) &\longmapsto [a \cdot b]
    \end{align*}
\end{minipage}

\bigskip

Se ad esempio poniamo $n = 4$, avremo:

\bigskip

\begin{minipage}[c]{0.45\textwidth}
    \centering
    \textbf{Somma}
    \vspace{-5pt}
    \[
    \begin{array}{c|cccc}
        + & [0] & [1] & [2] & [3] \\ \hline
        [0] & [0] & [1] & [2] & [3] \\{}
        [1] & [1] & [2] & [3] & [0] \\{}
        [2] & [2] & [3] & [0] & [1] \\{}
        [3] & [3] & [1] & [2] & [3]  
    \end{array}
    \]
\end{minipage}
\begin{minipage}[c]{0.45\textwidth}
    \centering
    \textbf{Prodotto}
    \vspace{-5pt}
    \[
    \begin{array}{c|cccc}
        \cdot & [0] & [1] & [2] & [3] \\ \hline
        [0] & [0] & [0] & [0] & [0] \\{}
        [1] & [0] & [1] & [2] & [3] \\{}
        [2] & [0] & [2] & [0] & [2] \\{}
        [3] & [0] & [3] & [2] & [1]  
    \end{array}
    \]
\end{minipage}

\bigskip

Ora la domanda viene spontanea: questo insieme, con queste operazioni di somma e prodotto appena definite, è un anello? E se si, è un campo? 

È facile notare che rispetta tutte le 8 proprietà viste in Proprietà \ref{prop-somma-prod} qualunque sia il valore di $n$, dunque è facile concludere che \textbf{è un anello}.\\ 
Per essere però anche un campo, è necessario che, per ogni numero diverso da $0$, esista il suo \textit{inverso moltiplicativo}. Per verificarlo basta controllare se ogni riga (o colonna) che rappresenta una classe diversa da $[0]$, abbia il valore $[1]$. 
Riprendendo l'esempio di prima con $n = 4$, la classe $[1]$ notiamo subito che ha come invderso moltiplicativo sé stessa, e la classe $[3]$ anche ha come inverso moltiplicativo se stessa. La classe $[2]$, invece, non ha mai il valore $[1]$, e questo significa che non ha inverso moltiplicativo. Dato che $[2] \neq [0]$, possiamo concludere che $\Z/(4)$ \textbf{non è un campo}.
Questo però non significa che non sia mai un campo: se $n$ è un \textit{numero primo} (che indichiamo con $p$), allora $\Z/(p)$ \textbf{è un campo}. In questo caso, al posto di scrivere $\Z/(p)$, scriviamo $\bm{\F_p}$.

\section{Numeri complessi}
\begin{Definizione}[Numeri complessi]
\label{def-num-comp}
Si definisce \textbf{numero complesso} ogni numero esprimibile nella forma $\bm{z = a + bi}$, che prende il nome di \textbf{rappresentazione cartesiana}, dove $a,b \in \R$ e $\bm{i}$ è l'unità immaginaria tale che $i^2 = -1$. \\
Chiamiamo $\bm{a = Re(z)}$ la parte reale di $z$ e $\bm{b = Im(z)}$ la parte immaginaria di $z$. \\
$\bm{\C = \{z = a + bi \mid a,b \in \R\}}$ è l'insieme dei numeri complessi.
\end{Definizione}

Definiamo anche per i numeri complessi le operazioni di somma e prodotto.

\bigskip

\begin{minipage}[c]{0.45\textwidth}
    \centering
    {\hspace{-0.2cm}\Large\textbf{Somma}:}
    \begin{align*}
        \hspace{0.6cm} \C \times \C &\longrightarrow \C \\
        \hspace{0.6cm} (a + bi, c + di) &\longmapsto (a + c) + (b + d)i
    \end{align*}
\end{minipage}
\begin{minipage}[c]{0.45\textwidth}
    \centering
    {\hspace{-1.2cm}\Large\textbf{Prodotto}:}
    \begin{align*}
        \hspace{0.4cm} \C \times \C &\longrightarrow \C \\
        \hspace{0.4cm} (a + bi, c + di) &\longmapsto (ac - bd) + (ad + bc)i
    \end{align*}
\end{minipage}

\bigskip

Per capire velocemente se un insieme è un campo senza dover controllare tutte le 8 proprietà, basta verificare che:
\begin{enumerate}[label = \arabic*°]
    \item L'insieme è \textit{chiuso rispetto somma e prodotto}, ossia moltiplicando o sommando due elementi dell'insieme, troviamo sempre un elemento dell'insieme.
    \item Sono presenti gli \textit{elementi neutri} di somma e prodotto.
    \item Ogni elemento ha il proprio \textit{opposto}.
    \item Ogni elemento ha il proprio \textit{inverso moltiplicativo}.
\end{enumerate}
Dato che l'insieme $\C$ rispetta tutti questi requisiti, possiamo facilmente concludere che anch'esso è un \textbf{campo}.

\begin{Definizione}[Coniugato di un numero complesso]
\label{def-coniugato-num-comp}
Il coniugato di un numero complesso $z = a + bi$ è $\bm{\bar{z} = a - bi}$, ossia sé stesso ma con la parte immaginaria opposta.
\end{Definizione}

Si può notare come il prodotto tra un numero complesso e il proprio coniugato valga $a^2 + b^2$, che è sempre positivo $\forall z \in \C$.

\subsection{Geometria dei numeri complessi}
Geometricamente, per rappresentare i numeri complessi prendiamo un sistema di riferimento cartesiano sul piano e costruiamo dei vettori che partono dall'origine di tale piano. 
Sull'asse delle ascisse vi è la parte reale di $z$, mentre sulle ordinate la parte immaginaria.

\bigskip

\begin{minipage}[c]{0.35\textwidth}
\centering
\includegraphics[width=1\textwidth, valign=t]{Numeri complessi sul piano.png}
\end{minipage}\hfill
\begin{minipage}[c]{0.60\textwidth}
Il numero complesso $z = 3 + i$ viene qui rappresentato da un vettore lungo $\bm{\lvert z \rvert = \sqrt{a^2+b^2}} = 10$, che chiamiamo \textbf{modulo} di $z$, calcolato usando il classico teorema di Pitagora, e di angolo $\theta$, che chiamiamo \textbf{argomento di} $\bm{z}$ ($\bm{Arg(z) = \theta})$.
\end{minipage}

\begin{Definizione}[Rappresentazione polare]
\label{def-rapp-pol}
Possiamo scrivere tutti i numeri complessi come $\bm{z = \rho(\cos(\theta) + i \sin(\theta))}$, dove $\rho = \abs{z}$, che prende il nome \textbf{rappresentazione polare}.
\end{Definizione}

È facile capire il perché guardando la rappresentazione geometrica dei numeri complessi: si può infatti notare che $a = \abs{z} \cdot \cos(\theta)$, e che $b = \abs{z} \cdot \sin(\theta)$, dunque $z = a + bi = \abs{z} (\cos(\theta) + i\sin(\theta))$.

Vediamo ora invece le operazioni di somma e prodotto per i numeri complessi sul piano.

\bigskip

{\Large \textbf{Somma}} \nopagebreak

\bigskip

\begin{minipage}[c]{0.35\textwidth}
\centering
\includegraphics[width=1\textwidth, valign=t]{Somma di numeri complessi sul piano.png}
\end{minipage}\hfill
\begin{minipage}[c]{0.60\textwidth}
    La somma di due numeri complessi sul piano avviene come una classica somma tra vettori mediante il metodo del parallelogramma.
\end{minipage}

\bigskip

{\Large \textbf{Prodotto}} \nopagebreak

\bigskip

\begin{minipage}[c]{0.35\textwidth}
\centering
\includegraphics[width=1\textwidth, valign=t]{Prodotto di numeri complessi sul piano.png}
\end{minipage}\hfill
\begin{minipage}[c]{0.60\textwidth}
    Prendiamo $z_1 = \rho_1 (\cos(\theta_1) + i\sin(\theta_1))$ e \\ 
    $z_2 = \rho_2 (\cos(\theta_2) + i \sin(\theta_2))$. \\
    Vediamo come 
    \begin{align*}
    z_1 \cdot z_2 &= \rho_1 \rho_2 (\cos(\theta_1) + i\sin(\theta_1))(\cos(\theta_2) + i\sin(\theta_2))  \\
    &= \rho_1 \rho_2 (\cos(\theta_1) \cos(\theta_2) - \sin(\theta_1)\sin(\theta_2) + \\
    &+ i (\cos(\theta_1)\sin(\theta_2) + \cos(\theta_2)\sin(\theta_1)))
    \end{align*}
    Ricordando le formule di somma di seno e coseno, possiamo riscrivere tutto come \\ \bm{$z_1 \cdot z_2 = \rho_1 \rho_2 (\cos(\theta_1 + \theta_2) + i\sin(\theta_1)\sin(\theta_2))$}. 
\end{minipage}

\bigskip

Dunque il prodotto di due numeri complessi produce un altro numero complesso con modulo $|z_1 \cdot z_2| = \rho_1 \rho_2 = |z_1|  |z_2|$, e con angolo $\theta = \theta_1 + \theta_2$.

\subsection{Soluzioni complesse di un'equazione}

\begin{Proposizione}[Soluzioni complesse di un'equazione]
\label{propos-sol-comp}
Un'equazione di grado $\bm{n}$ avrà, nel campo dei numeri complessi $\C$, $\bm{n}$ \textbf{soluzioni}.
\end{Proposizione}

\Esempio

Trovare tutte le soluzioni complesse dell'equazione $\bm{z^n = 1}$. \nopagebreak

Cercare le soluzioni di questa equazione significa capire quali numeri complessi elevati alla $n$ hanno modulo pari a $1$ e angolo pari a $2k\pi$, con $k \in \Z$ (dato che seno e coseno sono funzioni con periodo pari a $2\pi$). In poche parole, possiamo dire che $\abs{z^n} = 1$ e che $Arg(z^n) = 2k\pi$.\\
Per quanto riguarda il modulo, sappiamo che $\abs{z^n} = \abs{z}^n = 1$, e che quindi $\abs{z} = 1$, mentre per quanto riguarda l'argomento, abbiamo visto prima come, preso $Arg(z) = \theta$, $Arg(z^n) = Arg(z \cdot z \ldots \cdot z) = \theta + \theta + \ldots + \theta = n \theta$, e dato che in questo caso $Arg(z^n) = 2k\pi$, troviamo che $nArg(z) = 2k\pi$ e che quindi $Arg(z) = \theta = \frac{2k\pi}{n}$. \\
Troviamo quindi che tutte le soluzioni complesse dell'equazione sono 
\[z = \abs{z} \cdot(\cos(\theta) + i \sin(\theta)) = \cos\left(\frac{2k\pi}{n}\right) + i \sin\left(\frac{2k\pi}{n}\right)\]

Nonostante sembri che abbiamo trovato un numero infinito di soluzioni, dato che $k \in \Z$, ma è facile notare che, non appena arriviamo a $k = n$, il valore di $z$ sarà pari al valore di quando $k = 0$. Troviamo quindi che le soluzioni dell'equazione sono in totale $n$, e sono  
\[\bm{z = \cos\left(\frac{2k\pi}{n}\right) + i \sin\left(\frac{2k\pi}{n}\right)} \textbf{ con } \bm{0 \leq k \leq n-1}\]

\begin{Proposizione}[Teorema fondamentale dell'algebra]
\label{propos-teo-fon-alg}
Definiamo $p(z) = c_0 z^n + c_1 z^{n-1} + \ldots + c_n$, con $c_0 \neq 0$ un polinomio di grado $n$ nel campo dei numeri complessi. È sempre possibile riscrivere tale polinomio come un \textbf{prodotto tra termini di primo grado}: $p(z) = c_0(z-\gamma_1)(z-\gamma_2)\ldots(z-\gamma_n)$, dove $\gamma_1,\gamma_2,\ldots,\gamma_n$ sono le radici complesse del polinomio.
\end{Proposizione}

\chapter{Spazi vettoriali}

\section{Cos'è uno spazio vettoriale?}

\begin{Definizione}[Segmento orientato]
\label{def-segm-orient}
Un \textbf{segmento orientato} è un segmento che unisce una coppia ordinata di punti $PQ$ in un piano.
\end{Definizione}

\begin{Definizione}[Segmenti equipollenti]
\label{def-segm-equip}
Due segmenti orientati $P_1Q_1$, $P_2Q_2$ si dicono \textbf{equipollenti} se le rette $\braket{P_1Q_1}$ e $\braket{P_2Q_2}$ sono parallele e anche le rette $\braket{P_1P_2}$ e $\braket{Q_1Q_2}$ lo sono.
\end{Definizione}

\begin{Definizione}[Vettore]
\label{def-vett}
Un \textbf{vettore} è un segmento orientato a meno di equipollenza.\\
\NB Per evitare confusione, dato che si parlerà spesso di scalari e vettori insieme, per definire lo $0$ vettoriale usiamo il simbolo $\bm{\undzero}$.
\end{Definizione}

Questo significa che due segmenti orientati equipollenti $P_1Q_1$ e $P_2Q_2$ rappresentano lo stesso vettore $\vv{PQ}$.

Definiamo le operazioni di \textbf{somma} e di \textbf{prodotto per scalare} per un insieme di vettori.

\bigskip

\begin{minipage}[c]{0.45\textwidth}
    \centering
    {\Large\textbf{Somma}:}
    \begin{align*}
        \hspace{0.4cm}\V \times \V &\longrightarrow \V \\
        \hspace{0.4cm}(v,w) &\longmapsto v + w
    \end{align*}
\end{minipage}
\begin{minipage}[c]{0.45\textwidth}
    \centering
    {\Large\textbf{Prodotto per scalare}:}
    \begin{align*}
        \K \times \V &\longrightarrow \V \\
        (\lambda,v) &\longmapsto \lambda v
    \end{align*}
\end{minipage}

\bigskip

Definite ora queste operazioni, vediamo alcune loro proprietà. \nopagebreak

\begin{Proprieta}[Proprietà di Somma e Prodotto per scalare]
\label{prop-somm-prod-scal}
\begin{enumerate}
    \item $\bm{\undzero + v = v} \enspace \forall v \in \V$ \enspace (Esistenza dell'elemento neutro della somma)
    \item $\bm{u + (v + w) = (u + v) + w} \enspace \forall u,v,w \in \V$ \enspace (Proprietà associativa della somma)
    \item $\forall v \in \V \enspace \exists w \in \V : \bm{v + w = \undzero}$ \enspace (Esistenza dell'elemento opposto)
    \item $\bm{v + w = w + v} \enspace \forall v,w \in \V$ \enspace (Proprietà commutativa della somma)
    \item $\bm{1 \cdot v = v} \enspace \forall v \in \V$ \enspace (Esistenza dell'elemento neutro del prodotto)
    \item $\bm{v \cdot (\lambda + \mu) = \lambda v + \mu v} \enspace \forall v \in \V \enspace \forall \lambda,\mu \in \K$ \enspace (Proprietà distributiva del prodotto rispetto agli scalari)
    \item $\bm{(v + w) \lambda = \lambda v + \lambda w} \enspace \forall v,w \in \V \enspace \forall \lambda \in \K$ \enspace (Proprietà distributiva del prodotto rispetto ai vettori)
    \item $\bm{(\lambda \cdot \mu) \cdot v = \lambda \cdot (\mu \cdot v)} \enspace \forall v \in V \enspace \forall \lambda,\mu \in \K$ \enspace (Proprietà associativa del prodotto)
\end{enumerate}
\end{Proprieta}

Valgono anche qui le Osservazioni \ref{oss-prop-somma-prod} viste precedentemente.

\begin{Definizione}[Spazio Vettoriale]
\label{def-spa-vet}
Uno \textbf{spazio vettoriale} su $\mathbb{K}$ è un insieme di vettori con le operazioni di somma e prodotto per scalare appena descritte che rispettano le Proprietà \ref{prop-somm-prod-scal}.
\end{Definizione}

Classici spazi vettoriali sono $\bm{\K[x]}$, ossia l'insieme dei polinomi in $x$ a coefficienti in $\K$, e anche $\bm{\K^n}$, ossia l'insieme delle tuple ordinate di $n$ elementi.

\Esempio

Verificare se $\bm{V(\E^2)} = \{\text{vettori di } \E^2\}$ sia uno \textbf{spazio vettoriale} o meno. \nopagebreak

Studiamo le operazioni di somma e prodotto per scalare di $V(\E^2)$.

\bigskip

\begin{minipage}[c]{0.35\textwidth}
\includegraphics[width=1\textwidth, valign=t]{Somma di vettori nel piano.png}
\end{minipage}\hfill
\begin{minipage}[c]{0.60\textwidth}
    {\centering
    {\Large\textbf{Somma}}
    \begin{align*}
        \hspace{0.4cm}V(\E^2) \times V(\E^2) &\longrightarrow V(\E^2) \\
        \hspace{0.4cm}(\vv{PQ},\vv{RS}) &\longmapsto \vv{PT}
    \end{align*}}
    La somma di due vettori avviene come una classica somma tra vettori mediante il metodo punta-coda.\\
    Dato che i due vettori, come in questo caso, potrebbero non toccarsi, per sommarli si prende un segmento orientato al secondo tale che l'inizio di questo segmento sia nello stesso punto in cui finisce il primo vettore.
\end{minipage}

\bigskip

\begin{minipage}[c]{0.35\textwidth}
\includegraphics[width=1\textwidth, valign=t]{Prodotto vettore per scalare nel piano.png}
\end{minipage}\hfill
\begin{minipage}[c]{0.60\textwidth}
    {\centering
    {\Large\textbf{Prodotto per scalare}}
    \begin{align*}
        \hspace{0.4cm}\R \times V(\E^2) &\longrightarrow V(\E^2) \\
        \hspace{0.4cm}(\lambda,\vv{PQ}) &\longmapsto \lambda\vv{PR}
    \end{align*}}
    Il prodotto tra un vettore e uno scalare avviene prendendo il vettore e modificandone la lunghezza. \\
    Se infatti il vettore $\vv{PQ}$ ha modulo $1$, il vettore $\vv{PR} = \lambda \cdot \vv{PQ}$ avrà come modulo $\lambda$. 
\end{minipage}

\bigskip

Dato che l'insieme $V(\E^2)$ rispetta tutte le 8 proprietà degli spazi vettoriali (facilmente verificabile), possiamo concludere che \textbf{è uno spazio vettoriale}. 

\begin{Proprieta}[Proprietà degli spazi vettoriali]
\label{prop-spaz-vett}
\begin{enumerate}[label=\Roman*)]
    \item $\bm{0 \cdot v = \undzero} \enspace \forall v \in \V$
    \item $\bm{\lambda \cdot \undzero = \undzero} \enspace \forall \lambda \in \K$
    \item $\bm{(-1) \cdot v + v = \undzero} \enspace \forall v \in \V$
\end{enumerate}
\end{Proprieta}

\Dimostrazione \nopagebreak

\begin{enumerate}[label=\Roman*)]
    \item La proprietà 1) ci dice che $0 \cdot v = (0 + 0) \cdot v$, e per la proprietà 6) $(0 + 0) \cdot v = 0 \cdot v + 0 \cdot v$. Troviamo quindi che $0 \cdot v = 0 \cdot v + 0 \cdot v$. Aggiungendo $-(0 \cdot v)$ ad entrambi i membri, troviamo che $0 \cdot v - (0 \cdot v) = 0 \cdot v + 0 \cdot v - 0 \cdot v$. 
    Per la proprietà 2) sappiamo che $0 \cdot v - (0 \cdot v) = \undzero$, dunque troviamo che $\undzero = 0 \cdot v + \undzero$, che per la proprietà 1) diventa $0 \cdot v = \undzero$. $\blacksquare$
    \item La proprietà 1) ci dice che $\lambda \cdot \undzero = \lambda \cdot (\undzero + \undzero)$, e per la proprietà 6) $\lambda \cdot (\undzero + \undzero) = \lambda \undzero + \lambda \undzero$. Troviamo quindi che $\lambda \cdot \undzero = \lambda \cdot \undzero + \lambda \cdot \undzero$. Aggiungendo $-(\lambda \cdot \undzero)$ ad entrambi i membri, troviamo che $\lambda \cdot \undzero - (\lambda \cdot \undzero) = \lambda \cdot \undzero + \lambda \cdot \undzero - (\lambda \cdot \undzero)$.
    Per la proprietà 2) sappiamo che $\lambda \cdot \undzero - (\lambda \cdot \undzero) = \undzero$, dunque troviamo che $\undzero = \lambda \cdot \undzero + \undzero$, che per la proprietà 1) diventa $\lambda \cdot \undzero = \undzero$. $\blacksquare$
    \item La proprietà 5) ci dice che $(-1) \cdot v + v = (-1) \cdot v + 1 \cdot v$, e per la proprietà 3) troviamo che $(-1) \cdot v + 1 \cdot v = v \cdot ((-1) + 1) = v \cdot 0$, e per la proprietà I) appena dimostrata sappiamo che $v \cdot 0 = \undzero$. Dunque $(-1) \cdot v + v = \undzero$. $\blacksquare$
\end{enumerate}

\section{Sottospazi vettoriali}

\begin{Definizione}[Sottospazi vettoriali]
\label{def-sott-vett}
Preso $\V$ spazio vettoriale, un \textbf{sottospazio vettoriale} $\W$ è un sottoinsieme di $\V$ con le seguenti proprietà:
\begin{itemize}
    \item Non è vuoto
    \item È chiuso rispetto alla somma
    \item È chiuso rispetto al prodotto per scalare
\end{itemize}
\end{Definizione}

\begin{Osservazione}[]
\label{oss-sott-vet}
Dato un sottospazio vettoriale $\W \subset \V$:
\begin{enumerate}[label=\alph*)]
    \item $\bm{\undzero \in \W}$
    \item $\forall v \in \W, \bm{-w \in \W}$
\end{enumerate}
\end{Osservazione}

\Dimostrazione \nopagebreak

\begin{enumerate}[label=\alph*)]
    \item Dato che $\W$ non è vuoto per definizione, sappiamo che $\exists v \in \W$, e dato che $\W$ è anche chiuso rispetto al prodotto per scalare, sicuramente $0 \cdot v = \undzero \in \W$. $\blacksquare$
    \item Dato che $\W$ non è vuoto per definizione, sappiamo che $\exists v \in \W$, e dato che $\W$ è anche chiuso rispetto al prodotto per scalare, sicuramente $(-1) \cdot v = -v \in \W$. $\blacksquare$
\end{enumerate}

\Esempio \nopagebreak

Verificare se $\K^n \supset \bm{\W := \{\X \in \K^n \mid \sum_{i = 1}^n x_i = 0\}}$ sia un \textbf{sottospazio vettoriale} o meno. \nopagebreak

Dato che $\K^n$ è uno spazio vettoriale, la domanda ha senso e dobbiamo quindi verificare le tre condizioni dei sottospazi vettoriali.

Dato che $\sum_{i = 1}^n 0 = 0$, possiamo concludere che $\undzero = (0,0,\ldots,0) \in \W$, e che quindi $\W$ non è vuoto. \\
Per verificare che $\W$ sia chiuso rispetto alla somma, prendiamo due tuple generiche $\X, \Y \in \W$. Sappiamo che $\sum_{i = 1}^n x_i = 0$ e che $\sum_{i = 1}^n y_i = 0$. Se prendiamo la somma delle due tuple, troviamo $\Z = \X + \Y = (x_1 + y_1, x_2 + y_2, \ldots, x_n + y_n)$. Troviamo quindi che $\sum_{i = 1}^n z_i = \sum_{i = 1}^n(x_i + y_i) = \sum_{i = 1}^n x_i + \sum_{i = 1}^n y_i = 0 + 0 = 0$, dunque $\Z \in \W$. Possiamo quindi concludere che $\W$ è chiuso rispetto alla somma. \\
Per verificare che $\W$ sia chiuso rispetto al prodotto per scalare, prendiamo una tupla generica $\X \in \W$ e prendiamo uno scalare generico $\lambda \in \K$. Se prendiamo il prodotto tra questa tupla e lo scalare, troviamo che $\Z = \lambda \cdot \X = (\lambda x_1, \lambda_x, \ldots, \lambda x_n)$. Troviamo quindi che $\sum_{i = 1}^n z_i = \sum_{i = 1}^n (\lambda x_i) = \lambda \sum_{i = 1}^n x_i = \lambda \cdot 0 = 0$, dunque $\Z \in \W$. Possiamo quindi concludere che $\W$ è chiuso rispetto al prodotto per scalare.

Avendo quindi dimostrato che $\W$ non è vuoto, è chiuso rispetto alla somma e rispetto al prodotto per scalare, possiamo concludere che $\bm{\W}$ \textbf{è un sottospazio vettoriale}.

\begin{Osservazione}[]
\label{oss-sott-vet-2}
\begin{enumerate}
    \item Se $\W$ è un sottospazio vettoriale, allora \textbf{è anche uno spazio vettoriale}.
    \item Se $\{\W_i\}_{i \in I}$ è un insieme di sottospazi vettoriali, allora anche $\bm{\bigcap\limits_{i \in I} \W_i}$ \textbf{è un sottospazio}.
\end{enumerate}
\end{Osservazione}

\Dimostrazione \nopagebreak

\begin{enumerate}
    \item Banalmente, se $\W$ è un sottospazio vettoriale allora questo rispetta tutte le Proprietà \ref{prop-spaz-vett}. $\blacksquare$
    \item Controlliamo se rispetta le tre proprietà dei sottospazi vettoriali. \\
    Sicuramente $\bigcap_{i \in I} \W_i$ non è vuoto, dato che per l'osservazione a) precedente $\undzero$ è presente in tutti i sottospazi vettoriali, e dunque anche nella loro intersezione. \\
    Per verificare che $\bigcap_{i \in I} \W_i$ sia chiuso rispetto alla somma, prendiamo due vettori \mbox{$v,w \in \bigcap_{i \in I} \W_i$}. Dato che questi due vettori sono presenti in tutti i sottospazi vettoriali dell'insieme, anche $v + w$ sarà presente in tutti i sottospazi, essendo questi per definizione chiusi rispetto alla somma. Dunque $v + w \in \bigcap_{i \in I} \W_i$, ossia $\bigcap_{i \in I} \W_i$ è chiuso rispetto alla somma. \\
    Per verificare che $\bigcap_{i \in I} \W_i$ sia chiuso rispetto al prodotto per scalare, prendiamo un vettore $v \in \bigcap_{i \in I} \W_i$ e prendiamo uno scalare generico $\lambda \in \K$. Dato che $v$ è presente in tutti i sottospazi vettoriali dell'insieme, anche $\lambda v$ sarà presente in tutti i sottospazi, essendo questi per definizione chiusi rispetto al prodotto per scalare. Dunque $\lambda v \in \bigcap_{i \in I} \W_i$, ossia $\bigcap_{i \in I} \W_i$ è chiuso rispetto al prodotto per scalare. \\
    Possiamo quindi concludere che $\bigcap_{i \in I} \W_i$ è un sottospazio vettoriale. $\blacksquare$
\end{enumerate} 

La prima osservazione è fondamentale: d'ora in poi, per dimostrare che un insieme è un sottospazio vettoriale, basterà dimostrare che è un sottospazio vettoriale, che ha solo 3 proprietà da dimostrare.

\begin{Definizione}[Somma di sottospazi vettoriali]
\label{def-somm-sott-vett}
Presi $\U,\W \subset \V$, $\bm{\U + \W = \{u + w \mid u \in \U \textbf{ e } w \in \W\}}$
\end{Definizione}

\section{Combinazioni lineari}

\begin{Definizione}[Vettori della base standard]
\label{def-vett-bas-can}
I vettori $e_1, e_2, \ldots, e_n \in \K$ sono chiamati \textbf{vettori della base standard} (vedremo più avanti cos'è la base standard, per ora ci interessano solo i suoi vettori) e sono così definiti: $\bm{e_1 = (1,0,\ldots,0), e_2 = (0,1,\ldots,0), \ldots, e_n = (0,0,\ldots,1)}$.
\end{Definizione}

\begin{Definizione}[Combinazione lineare]
\label{def-comb-line}
Una \textbf{combinazione lineare} dei vettori $v_1, \ldots, v_n$ è un vettore $v \in \V$ tale che \mbox{$\bm{v = \lambda_1 v_1 + \lambda_2 v_2 + \ldots + \lambda_n v_n}$}, con $\lambda_i \in \K$. \\
\NB $\undzero$ è considerato la combinazione lineare della lista vuota di vettori.
\end{Definizione}

\Esempio \nopagebreak

Dimostrare che ogni vettore di $\K^3$ è \textbf{combinazione lineare} di $\bm{e_1, e_2, e_3}$. \nopagebreak

Banalmente, prendendo una tupla qualsiasi $\X = (x_1, x_2, x_3) \in \K^3$, questa può essere riscritta come $\X = \bm{x_1 e_1 + x_2 e_2 + x_3 e_3}$.

\begin{Proposizione}[]
\label{propos-comb-lin}
Un sottoinsieme non vuoto $\W$ è un \textbf{sottospazio vettoriale} $\iff \forall v_1, v_2, \ldots, v_n \in \W$ \textbf{ogni combinazione lineare di tali vettori si trova in} $\bm{\W}$. 
\end{Proposizione}
In sostanza, verificando questa proprietà stiamo verificando contemporeanamente che $\W$ sia chiuso sia rispetto alla somma sia rispetto al prodotto per scalare.

\begin{Definizione}[SPAN]
\label{def-Span}
Preso $S$ sottoinsieme di uno spazio vettoriale $\V$ su $\mathbb{K}$, chiamiamo $\bm{\braket{S}}$ (o \textbf{SPAN($\bm{S}$)}) l'insieme delle combinazioni lineari formate dai vettori in $S$ prendendo coefficienti in $\mathbb{K}$. \\
Tutti i vettori $v_1, v_2, \ldots, v_n$ appartenenti ad $S$ sono chiamati \textbf{vettori generatori} di $\braket{S}$. \\
\NB Per alleggerire la notazione, al posto di scrivere $\braket{\{v_1, v_2, \ldots, v_n\}}$ scriveremo semplicemente $\braket{v_1, \ldots, v_n}$.
\end{Definizione}

\begin{Osservazione}[]
\label{oss-comb-lin}
Preso $S$ un sottoinsieme di uno spazio vettoriale $\V$:
\begin{enumerate}
    \item $\bm{S \subset \braket{S}}$
    \item $\bm{\braket{S}}$ è un \textbf{sottospazio vettoriale}
    \item Se $\W \subset \V$ è un sottospazio vettoriale che contiene $S$, allora $\bm{\braket{S} \subset \W}$
    \item $\bm{S = \braket{S}}$ solo se $S$ è un sottospazio vettoriale di $\V$ 
\end{enumerate}
\end{Osservazione}

\Dimostrazione \nopagebreak

\begin{enumerate}
    \item $\forall v \in \V$, $v = v \cdot 1$, dunque tutti i vettori di $S$ sono in $\braket{S}$. $\blacksquare$
    \item Controlliamo se rispetta le tre proprietà dei sottospazi vettoriali. \\
          $\undzero \in \braket{S}$ perché combinazione lineare della lista vuota, dunque $\braket{S}$ non è vuoto.\\
          Presi $v,w \in \braket{S}$, sappiamo che $v = \lambda_1 v_1 + \lambda_2 v_2 + \ldots + \lambda_n v_n$ e che $w = \mu_1 v_1 + \mu_2 v_2 + \ldots + \mu_n v_n$, dunque $v + w = v_1(\lambda_1 + \mu_1) + v_2 (\lambda_2 \mu_2) + \ldots + v_n (\lambda_n + \mu_n)$, dunque anche $v + w \in \braket{S}$. Possiamo quindi concludere che $\braket{S}$ è chiuso rispetto alla somma.\\
          Preso $v \in \braket{S}$ e $\mu \in \braket{S}$, sappiamo che $v = \lambda_1 v_1 + \lambda_2 v_2 + \ldots + \lambda_n v_n$, dunque \mbox{$\mu v = \mu \lambda_1 v_1 + \mu \lambda_2 v_2 + \ldots + \mu \lambda_n v_n$}, dunque anche $\mu v \in \braket{S}$. Possiamo quindi concludere che $\braket{S}$ è chiuso rispetto al prodotto per scalare.\\
          Possiamo quindi concludere che $\braket{S}$ è un sottospazio vettoriale. $\blacksquare$
    \item Preso $S = \{v_1, \ldots, v_n\}$, dato che $S$ contiene $S$, sappiamo che $v_1, \ldots, v_n \in \W$. Per la Proposizione \ref{propos-comb-lin} troviamo che ogni combinazione lineare di $v_1, v_2, \ldots, v_n$ è in $\W$, ossia $\braket{S} \subset \W$. $\blacksquare$
    \item Se $S$ è un sottospazio di $\V$, allora per la Proposizione \ref{propos-comb-lin} troviamo che ogni combinazione lineare dei vettori di $S$ si trova in $S$, ossia $\braket{S} \subset S$, e dato che l'osservazione 1) ci dice che $S \subset \braket{S}$, troviamo che $S = \braket{S}$. $\blacksquare$
\end{enumerate}

\begin{Definizione}[Spazio vettoriale finitamente generato]
\label{def-spa-vet-fin-gen}
Uno spazio vettoriale $\V$ è \textbf{finitamente generato} su $\mathbb{K}$ se esiste un insieme finito di vettori $S \subset \V$ tale che $\bm{\V = \braket{S}}$.
\end{Definizione}

\Esempi \nopagebreak

Dimostrare che $\bm{\K^n}$ \textbf{è finitamente generato} su $\K$. \nopagebreak

Dato che ogni vettore di $\K^n$ può essere riscritto come combinazione lineare di $e_1, e_2, \ldots, e_n$ (dimostrato in un esempio precedente con $n = 3$), troviamo che $\bm{\K^n = \braket{e_1, e_2, \ldots, e_n}}$.

\Separatore

Dimostrare che $\bm{K[x]}$ \textbf{non è finitamente generato} su $\K$. \nopagebreak

Dati $p_1, p_2, \ldots, p_n \in \K[x]$, se $p \in \braket{p_1, p_2, \ldots, p_n}$, il grado di $p$ sarà al più il massimo grado tra i polinomi $p_1, p_2, \ldots, p_n$. Dunque con il sottospazio $\braket{p_1, p_2 \ldots, p_n}$ è impossibile scrivere tutti i polinomi in $x$, dato che tutti quelli con grado superiore al massimo grado tra questi polinomi non saranno inclusi. Concludiamo quindi che $\K[x]$ \textbf{non è finitamente generato}. 

\begin{Osservazione}[]
\label{oss-comb-lin-2}
Preso $\W$ sottospazio vettoriale di $\V$ finitamente generato, anche $\W$ è \textbf{finitamente generato}.
\end{Osservazione}

\begin{Proposizione}[Corollario]
\label{propos-comb-lin-2}
L'insieme delle soluzioni di un sistema di $m$ equazioni lineari omogenee in $n$ incognite $\W \subset \K^n$ è \textbf{finitamente generato}.
\end{Proposizione}
Questo è vero perché $\W$ è un sottospazio vettoriale (facilmente dimostrabile), e dato che abbiamo dimostrato nell'esempio precedente che $\K^n$ è uno spazio vettoriale finitamente generato, allora per l'Osservazione \ref{oss-comb-lin-2} troviamo che anche $\W$ è finitamente generato.

\begin{Osservazione}[]
\label{oss-com-lin-3}
Per dimostrare che uno spazio vettoriale $\V$ su $\K$ sia chiuso rispetto alla somma e al prodotto per scalare, basta dimostrare che, presi $v, w \in \V$ e $\lambda, \mu \in \K$, $\bm{\lambda v + \mu w \in \V}$.
\end{Osservazione}

\section{Relazioni lineari}

\begin{Definizione}[Relazioni lineari]
\label{def-rel-lin}
Preso $\V$ uno spazio vettoriale, una \textbf{relazione lineare} tra i vettori $v_1, v_2, \ldots, v_n \in \V$ è un'uguaglianza del tipo
\[\bm{\lambda_1 v_1 + \lambda_2 v_2 + \ldots + \lambda_n v_n = \undzero} \enspace \text{con } \lambda_1, \lambda_2, \ldots, \lambda_n \in \K\]
Una relazione è detta \textbf{banale} se l'unica soluzione a tale equazione è $\lambda_1 = \lambda_2 = \ldots = \lambda_n = 0$.
\end{Definizione}

\begin{Definizione}[Dipendenza lineare]
\label{def-dip-lin}
I vettori $v_1, v_2, \ldots, v_n$ sono \textbf{linearmente dipendenti} se esiste una relazione lineare non banale tra loro. Altrimenti sono detti \textbf{linearmente indipendenti}.
\end{Definizione}

\Esempio \nopagebreak

Dimostrare che i vettori $\bm{e_1, e_2, \ldots, e_n} \in \K^n$ sono \textbf{linearmente indipendenti}. \nopagebreak

Banalmente \nopagebreak
\begin{align*}
 \undzero &= \lambda_1 e_1 + \lambda_2 e_2 + \ldots + \lambda_n e_n \\
 &= \lambda_1 (1,0,\ldots,0) + \lambda_2 (0,1,\ldots,0) + \ldots + \lambda_n( 0,0,\ldots,1) \\
 &= (\lambda_1, \lambda_2, \ldots, \lambda_n)
\end{align*}

è verificata se e solo se $\lambda_1 = \lambda_2 = \ldots = \lambda_n = 0$. Dunque i vettori $e_1, e_2, \ldots, e_n$ sono \textbf{linearmente indipendenti}.

\begin{Osservazione}[]
\label{oss-rel-lin}
I vettori $v_1, v_2, \ldots, v_n$ sono linearmente dipendenti $\iff \exists i \in \{1, 2, \ldots, n\}$ tale che \mbox{$\bm{v_i \in \braket{v_1, \ldots, v_{i-1}, v_{i+1}, \ldots, v_n}}$}, ossia se $v_i$ è combinazione lineare di tutti i vettori escluso sé stesso.
\end{Osservazione}

\Dimostrazione \nopagebreak

\begin{enumerate}[label = \arabic*°]
    \item $v_1, v_2 \ldots, v_n$ sono linearmente dipendenti $\implies$ $v_i$ è combinazione lineare di tutti i vettori escluso sé stesso.\\
    Dato che i vettori sono dipendenti, sappiamo che $\exists i \in \{1, 2, \ldots, n\}$ tale che $\lambda_i \neq 0$. \\ 
    Dunque
    \begin{align*}
    \lambda_1 v_1 + \ldots + \lambda_i v_i + \ldots + \lambda_n v_n &= \undzero \\
    \lambda_i v_i &= - \lambda_1 v_1 - \ldots - \lambda_{i-1} v_{i-1} - \lambda_{i+1} v_{i+1} - \ldots - \lambda_n v_n \\
    v_i &= -\frac{\lambda_1}{\lambda_i}v_1 - \ldots - \frac{\lambda_{i-1}}{\lambda_i} v_{i-1} - \frac{\lambda_{i+1}}{\lambda_i} v_{i+1} - \ldots - \frac{\lambda_n}{\lambda_i} v_n
    \end{align*}
    ossia $v_i$ è combinazione lineare di tutti i  vettori escluso sé stesso. $\blacksquare$
    \item $v_1, v_2, \ldots, v_n$ sono linearmente dipendenti $\impliedby$ $v_i$ è combinazione lineare di tutti i vettori escluso sé stesso.\\
    Dato che $v_i$ è combinazione lineare di tutti gli altri vettori, sappiamo che \mbox{$\exists \mu_1, \ldots, \mu_{i-1}, \mu_{i+1}, \ldots, \mu_n \in \K$} tale che \mbox{$v_i = \mu_1 v_1 + \ldots + \mu_{i-1} v_{i-1} + \mu_{i+1} v_{i+1} + \ldots + \mu_n v_n$}. \\
    Dunque
    \begin{align*}
    v_i &= \mu_1 v_1 + \ldots + \mu_{i-1} v_{i-1} + \mu_{i+1} v_{i+1} + \ldots + \mu_n v_n  \\
    \undzero &= \mu_1 v_1 + \ldots + \mu_{i-1} v_{i-1} + v_i + \mu_{i+1} v_{i+1} + \ldots + \mu_n v_n 
    \end{align*}
    Troviamo quindi che $\lambda_i = 1 \neq 0$, ossia i vettori sono linearmente dipendenti. $\blacksquare$
\end{enumerate}

\begin{Proposizione}[]
\label{propos-rel-lin}
Preso $\V$ uno spazio vettoriale su $\K$ finitamente generato da $m$ vettori.
\begin{enumerate}
    \item $w_1, w_2, \ldots, w_n \in \V$ sono \textbf{linearmente indipendenti} $\implies \bm{n \leq m}$
    \item $w_1, w_2, \ldots, w_n \in \V$ e $\bm{n > m} \implies w_1, w_2, \ldots, w_n$ sono \textbf{linearmente dipendenti}
\end{enumerate}
\end{Proposizione}

\Dimostrazione \nopagebreak

Queste due proposizioni, anche se scritte in modo differente, in realtà sono perfettamente identiche, dunque ci concentreremo a dimostrare solamente la prima proposizione.

Prendiamo $v_1, v_2, \ldots, v_m \in \V$ tali che $\braket{v_1, v_2, \ldots, v_m} = \V$. Possiamo quindi scrivere ogni $w_i$ come combinazione lineare di $v_1, v_2, \ldots, v_n$: $w_i = \sum_{j = 1}^m a_{ij} v_j$. Prendiamo ora $c_1, c_2, \ldots, c_n \in \K$. Troviamo che 
\[\sum_{i = 1}^n c_i w_i = \sum_{i = 1}^n c_i \sum_{j = 1}^m a_{ij} v_j = \sum_{j = 1}^m \left(\sum_{i = 1}^n a_{ij}  c_i\right)v_j\]
Questa combinazione lineare è nulla se, $\forall j \in \{1, 2, \ldots, m\}$, $\sum_{i = 1}^n a_{ij}c_i = 0$. Questo altro non è che un sistema lineare omogeneo di $m$ equazioni in $n$ incognite ($c_1, c_2, \ldots, c_n$). Dato che un sistema lineare omogeneo con più incognite che equazioni, ammette sempre una soluzione non banale (lo dimostreremo più avanti, per ora si può dare questa affermazione per vera), se supponiamo per assurdo $n > m$ troviamo proprio che $\exists i \in \{1,2,\ldots,n\}$ tale che $c_i \neq 0$, e dunque troviamo che $w_1, w_2, \ldots, w_n$ sono in realtà linearmente dipendenti, contraddicendo l'ipotesi. Concludiamo quindi che $n \leq m$. $\blacksquare$

\section{Basi di spazi vettoriali}

\begin{Definizione}[Base]
\label{def-base}
Una \textbf{base} di uno spazio vettoriale $\V$ è una lista di vettori $\B = \{v_1, v_2, \ldots, v_n\}$ tale che
\begin{itemize}
    \item $\bm{\braket{v_1, v_2, \ldots, v_n} = \V}$
    \item $\bm{v_1, v_2, \ldots, v_n}$ sono \textbf{linearmente indipendenti}
\end{itemize}
\end{Definizione}

\begin{Definizione}[Base standard]
\label{def-base-standard}
La base $\bm{S = \{e_1, e_2, \ldots, e_n\}}$ di $\K^n$ prende il nome di \textbf{base standard}.
\end{Definizione}
Dimostrare che $S$ sia effettivamente una base di $\K^n$ viene spontaneo, dato che abbiamo già dimostrato precedentemente che $\braket{e_1, e_2, \ldots e_n} = \K^n$ e che $e_1, e_2, \ldots, e_n$ sono linearmente indipendenti.

\Esempio \nopagebreak

Dimostrare che, presi due vettori $v_1, v_2 \in V(\E^2)$ tali che $v_1 \neq \lambda v_2 \enspace \forall \lambda \in \K$, allora $\B = \{v_1, v_2\}$ \textbf{è una base di} $\bm{V(\E^2)}$. \nopagebreak

Dobbiamo quindi verificare che $\braket{v_1, v_2} = V(\E^2)$ e che $v_1, v_2$ siano linearmente indipendenti.

Dire che $\braket{v_1, v_2} = V(\E^2)$ significa che $\forall v \in V(\E^2) \enspace \exists \lambda_1, \lambda_2 \in \K$ tali che $v = \lambda_1 v_1 + \lambda_2 v_2$. Ciò è geometricamente dimostrato dalla regola del parallelogramma.\\
$v_1$ e $v_2$ sono ovviamente linearmente indipendenti, perché se avessimo $\lambda_1 \neq 0$ allora dall'equazione $\lambda_1 v_1 + \lambda_2 v_2 = \undzero$ troveremmo $v_1 = -\frac{\lambda_2}{\lambda_1} v_2$, il che contraddice l'ipotesi che $v_1 \neq \lambda v_2$, dunque $\lambda_1 = 0$. Troviamo quindi che $\lambda_2 v_2 = \undzero$, che è vero se e solo se $\lambda_2 = 0$. Troviamo quindi che i vettori sono linearmente indipendenti.\\
Possiamo quindi concludere che $\B$ \textbf{è una base} di $V(\E^2)$.

\begin{Proposizione}[]
\label{propos-base}

Presi $v_1, v_2, \ldots, v_n$ vettori nello spazio vettoriale $\V$, definiamo l'applicazione $f$ nel seguente modo:
\begin{align*}
f : \K^n &\longrightarrow \V \\
(x_1, x_2, \ldots, x_n) &\longmapsto \sum_{i = 1}^n x_i v_i 
\end{align*}

\begin{enumerate}[label=\alph*)]
    \item $\bm{f}$ è \textbf{suriettiva} $\iff \bm{\braket{v_1, v_2, \ldots, v_n} = \V}$
    \item $\bm{f}$ è \textbf{iniettiva} $\iff \bm{v_1, v_2, \ldots, v_n}$ sono \textbf{linearmente indipendenti}
\end{enumerate}

\end{Proposizione}

\Dimostrazione \nopagebreak

\begin{enumerate}[label=\alph*)]
    \item \begin{enumerate}[label=\arabic*°]
        \item $f$ è suriettiva $\implies \braket{v_1, v_2, \ldots, v_n} = \V$ \\
        $f$ è suriettiva significa che $\forall v \in \V \enspace \exists \X = (x_1, x_2, \ldots, x_n) \in \K^n$ tale che $f(\X) = v$, e quindi $v = \sum_{i = 1}^n x_i v_i$. Questo significa che ogni vettore di $\V$ può essere riscritto come combinazione lineare di $v_1, v_2, \ldots, v_n$, ossia $\braket{v_1, v_2, \ldots, v_n} = \V$. $\blacksquare$
        \item $f$ è suriettiva $\impliedby \braket{v_1, v_2, \ldots, v_n} = \V$ \\
        $\braket{v_1, v_2, \ldots, v_n} = \V$ significa che $\forall v \in \V \enspace \exists \X = (x_1, x_2, \ldots, x_n) \in \K^n$ tale che \mbox{$v = x_1 v_1 + x_2 v_2 + \ldots x_n v_n = \sum_{i = 1}^n x_i v_i = f(\X)$}. Dato che questo vale $\forall v \in \V$, possiamo concludere che $f$ è suriettiva. $\blacksquare$ 
    \end{enumerate}
    \item \begin{enumerate}[label=\arabic*°]
        \item $f$ è iniettiva $\implies v_1, v_2, \ldots, v_n$ sono linearmente indipendenti \\
        Sappiamo che $f(0,0,\ldots,0) = \sum_{i = 1}^n 0 v_i = \undzero$. Dato che $f$ è iniettiva, troviamo che $f(\X) = \undzero$ se e solo se $\X = (0,0,\ldots,0)$, ossia $\sum_{i = 1}^n x_i v_i = \undzero$ se e solo se $x_i = 0 \enspace$ \mbox{$\forall i \in \{1,2, \ldots, n\}$}, ossia se i vettori $v_1, v_2, \ldots, v_n$ sono linearmente indipendenti. $\blacksquare$
        \item $f$ è iniettiva $\impliedby v_1, v_2, \ldots, v_n$ sono linearmente indipendenti \\
        Presi $\X = (x_1, x_2, \ldots, x_n), \Y = (y_1, y_2, \ldots, y_n) \in \K^n$, supponiamo che $f(\X) = f(\Y)$. Troviamo quindi che $\sum_{i = 1}^n x_i v_i = \sum_{i = 1} y_i v_i$, ossia che $\sum_{i = 1}^n (x_i - y_i) v_i = \undzero$. Dato che i vettori $v_1, v_2, \ldots, v_n$ sono linearmente indipendenti, troviamo che $x_i - y_i = 0 \enspace$ \mbox{$\forall i \in \{1, 2, \ldots, n\}$}, dunque $\X = \Y$, ossia $f$ è iniettiva. $\blacksquare$. 
    \end{enumerate}
\end{enumerate}

\begin{Proposizione}[Corollario]
\label{propos-corol-basi}
L'applicazione $f$ descritta nella Proposizione \ref{propos-base} è \textbf{biettiva} $ \iff \bm{v_1, v_2, \ldots, v_n}$ \textbf{è una base di} $\bm{\V}$. \\ 
Dunque se $\B = \{v_1, v_2, \ldots, v_n\}$ sono una base di $\V$, allora l'applicazione $f$ è biiettiva.
\end{Proposizione}

\begin{Definizione}[Applicazione inversa di $\bm{f}$]
\label{def-appl-inver-propos-base}
Dato che l'applicazione $f$ descritta nella Proposizione \ref{propos-base} è biiettiva se prendiamo $\B = \{v_1, v_2, \ldots, v_n\}$ è una base di $\V$, possiamo definire l'\textbf{applicazione inversa} di $f$:
\begin{align*}
\bm{X_\B} : \V &\longrightarrow \K^n \\
v &\longmapsto (x_1, x_2, \ldots, x_n)
\end{align*}
\end{Definizione}

\begin{Definizione}[Coordinate]
\label{def-coordinate}
Preso $v \in \V$, $\B$ una base di $\V$ e $X_\B$ l'applicazione vista in Definizione \ref{def-appl-inver-propos-base}, chiamiamo $\bm{X_\B(v)} = (x_1, x_2, \ldots, x_n)$ le \textbf{coordinate} di $v$ nella base $\B$.
\end{Definizione}

\Esempio \nopagebreak

Preso un vettore generico $v \in \V = \K^2$, calcolare le \textbf{coordinate} di $v$ nella base $\bm{\B = \{(1,1),(1,-1)\}}$.

Chiamiamo per comodità $v_1 = (1,1)$ e $v_2 = (1,-1)$. Prendendo $v = (a_1, a_2) \in \V$, sappiamo che $X_\B(v) = (x_1,x_2)$ tali che $v = x_1 v_1 + x_2 v_2 = x_1(1,1) + x_2(1,-1) = (x_1 + x_2, x_1 - x_2)$.\\
Dato che $v = (a_1, a_2)$, troviamo il sistema
\[
\begin{cases}
 x_1 + x_2 =& a_1 \\
 x_1 - x_2 =& a_2 
\end{cases}
\]

Dal quale troviamo che $\bm{X_\B(v) = (\frac{a_1+a_2}{2}, \frac{a_1-a_2}{2})}$.
 
\begin{Proposizione}[]
\label{propos-base-2}
\begin{enumerate}[label=\Roman*)]
    \item $\bm{\V}$ è \textbf{finitamente generato} $\implies$ $\V$ \textbf{ha una base}
    \item Preso $\V$ spazio vettoriale finitamente generato, $v_1,v_2, \ldots, v_n \in \V$ sono \textbf{linearmente indipendenti} $\implies \exists u_1, u_2, \ldots, u_l \in \V$ tali che $\bm{\B = \{v_1, v_2, \ldots, v_n, u_1, u_2, \ldots, u_l\}}$ è una \textbf{base di} $\bm{\V}$. 
\end{enumerate}
\end{Proposizione}


\Dimostrazione \nopagebreak

\begin{enumerate}[label=\Roman*)]
    \item Dato che $\V$ è finitamente generato, sappiamo che esistono $v_1, v_2, \ldots, v_n \in \V$ tali che $\braket{v_1, v_2, \ldots, v_n} = \V$. Ci sono ora due possibilità:
    \begin{enumerate}[label=\arabic*°]
        \item Se $v_1, v_2, \ldots, v_n$ sono linearmente indipendenti, allora banalmente \mbox{$\B = \{v_1, v_2, \ldots, v_n\}$} è una base di $\V$.
        \item Se invece sono linearmente dipendenti significa che $\exists i \in \{1, 2, \ldots, n\}$ tale che \mbox{$v_i \in \braket{v_1, \ldots, v_{i-1}, v_{i+1}, \ldots, v_n}$}. Questo in sostanza ci dice che $v_i$ è un vettore superfluo, dato che può essere riscritto come combinazione lineare degli altri vettori. 
        Da ciò concludiamo che $\braket{v_1, \ldots, v_{i-1}, v_{i+1}, \ldots, v_n}$. A questo punto, se questi vettori sono linearmente indipendenti allora rappresentano una base di $\V$, altrimenti si continua con questo procedimento finché non si arriva ad una lista di vettori linearmente indipendenti, che formeranno quindi una base di $\V$.
    \end{enumerate}
    Possiamo quindi concludere che $\V$ ha una base. $\blacksquare$
    \item Dato che $\V$ è finitamente generato, asappiamo che esistono $w_1, w_2, \ldots, w_m \in \V$ tali che $\braket{w_1, w_2, \ldots, w_m} = \V$. Dunque sicuramente anche $\braket{v_1, v_2, \ldots, v_n, w_1, w_2, \ldots, w_m} = \V$. Ci sono ora due possibilità:
    \begin{enumerate}[label=\arabic*°]
        \item Se $v_1, v_2, \ldots, v_n, w_1, w_2, \ldots, w_m$ sono linearmente indipendenti, allora banalmente \mbox{$\B = \{v_1, v_2, \ldots, v_n, w_1, w_2, \ldots, w_m\}$} è una base di $\V$.
        \item Se invece sono linearmente dipendenti significa che \mbox{$\exists \lambda_1, \lambda_2, \ldots, \lambda_n, \mu_1, \mu_2, \ldots, \mu_m \in \K$} tali che $\lambda_1 v_1 + \lambda_2 v_2 + \ldots + \lambda_n v_n + \mu_1 w_1 + \mu_2 w_2 + \ldots + \mu_m w_m = 0$ dove almeno uno tra i $\lambda_i$ e i $\mu_j$ è diverso da $0$, e dato che per ipotesi $\lambda_i \neq 0 \enspace \forall i \in \{1,2, \ldots, n\}$ essendo $v_1, v_2, \ldots, v_n$ linearmente indipendenti, troviamo che $\exists j \in \{1,2, \ldots, m\}$ tale che $\mu_j \neq 0$. 
        Possiamo quindi scrivere $w_j = \frac{\lambda_1}{\mu_j} v_1 + \ldots + \frac{\lambda_n}{\mu_j} v_n + \frac{\mu_1}{\mu_j} w_1 + \ldots + \frac{\mu_{j-1}}{\mu_j} w_{j-1} + \frac{\mu_{j+1}}{\mu_j} w_{j+1} + \ldots + \frac{\mu_m}{\mu_j} w_m$, ossia $\braket{v_1, \ldots, v_n, w_1, \ldots, w_{j-1}, w_{j+1}, \ldots, w_m} = \V$. 
        A questo punto, se questi vettori sono linearmente indipendenti allora rappresentano una base di $\V$, altrimenti si continua con questo procedimento finché non si arriva ad una lista di vettori linearmente indipendenti, che quindi formeranno una base di $\V$.
    \end{enumerate}
    Possiamo quindi concludere che $\exists u_1, u_2, \ldots, u_k \in \V$ tali che \\ $\B = \{v_1, v_2, \ldots, v_n, u_1, u_2, \ldots, u_l\}$ è una base di $\V$. $\blacksquare$
\end{enumerate}

\begin{Proposizione}[]
\label{propos-base-3}
Preso $\V$ uno spazio vettoriale finitamente generato, se $\B = \{v_1, v_2, \ldots, v_n\}$ e \mbox{$\Bs = \{w_1, w_2, \ldots, w_m\}$} sono basi di $\V$, allora $\bm{n = m}$. \\
Tutte le basi di uno spazio vettoriale $\V$ hanno quindi la stessa \textbf{cardinalità}.
\end{Proposizione}

\Dimostrazione \nopagebreak

Dato che $\B$ è una base di $\V$, sappiamo che $\braket{v_1, v_2, \ldots, v_n} = \V$. Siccome poi anche $\Bs$ è una base di $\V$, sappiamo che $w_1, w_2, \ldots, w_m$ sono linearmente indipendenti. Per la Proposizione \ref{propos-rel-lin}, troviamo che $m \leq n$. Ripetendo lo stesso procedimento invertendo $\B$ e $\Bs$, concludiamo che anche $n \leq m$, e che quindi $n = m$. $\blacksquare$

\section{Dimensione di uno spazio vettoriale}

\begin{Definizione}[Dimensione di uno spazio vettoriale]
\label{def-dim-spaz-vett}
Preso uno spazio vettoriale $\V$ finitamente generato su $\K$, la \textbf{dimensione} di $\V$ è la \textbf{cardinalità} di una sua base qualsiasi, e si indica scrivendo $\bm{\dim \V}$. \\
\NB La dimensione dello spazio vettoriale nullo, ossia lo spazio vettoriale che ha solo il vettore nullo al suo interno, è $0$: questo perché la lista vuota è a tutti gli effetti una base di $\V$, e questa ha cardinalità $0$. Dunque $\dim \{\undzero\} = 0$.
\end{Definizione}

La dimensione di uno spazio vettoriale \textit{può cambiare in base al campo su cui esso è definito}. Ad esempio, la dimensione dello spazio vettoriale $\C$ su $\C$ è $1$, mentre la dimensione di $\C$ su $\R$ è $2$. 

\begin{Osservazione}[]
\label{oss-dimen}
Preso $\V$ uno spazio vettoriale su $\K$ di dimensione $n$ e $\B = \{v_1, v_2, \ldots, v_n\}$ un'insieme di vettori di $\V$, se $\bm{\dim \V = \lvert \B \rvert}$ allora per dimostrare che $\B$ sia una \textbf{base} basta dimostrare che $\bm{\braket{v_1, v_2, \ldots, v_n} = \V}$ oppure che $\bm{v_1, v_2, \ldots, v_n}$ sono \textbf{linearmente indipendenti}. \\
In sostanza, se la cardinalità della lista di vettori è uguale alla dimensione dello spazio vettoriale, allora basta dimostrare \textbf{una sola delle due condizioni} necessarie affinché essa sia una base di $\V$.
\end{Osservazione}

\Dimostrazione \nopagebreak

\begin{enumerate}[label=\arabic*°]
    \item I vettori di $\B$ generano $\V$ \\
    Dobbiamo quindi dimostrare che i vettori di $\B$ sono linearmente indipendenti. Se per assurdo questi non lo fossero, allora esisterebbero \mbox{$i \in \{1,2, \ldots, n\}$} tali che \mbox{$v_i = \lambda_1 v_1 + \ldots + \lambda_{i-1} v_{i-1} + \lambda_{i+1} v_{i+1} + \ldots \lambda_n v_n$}, ossia troveremmo $\braket{v_1, v_2, \ldots, v_n} = \braket{v_1, \ldots, v_{i-1}, v_{i+1}, \ldots, v_n} = \V$. 
    Ma dato che $\dim \V = \abs\B = n$, è impossibile che lo spazio $\V$ sia generato da $n-1$ vettori. Questo perché, se fosse vero che $n-1$ vettori generano $\V$, allora potremmo trovare $n-1$ o meno vettori linearmente indipendenti che andrebbero quindi a formare una base di $\V$, ma dato che la dimensione di $\V$ è $n$, e dato che per la Proposizione \ref{propos-base-3} tutte le basi dello stesso spazio vettoriale hanno la stessa cardinalità, questo non è possibile. Concludiamo quindi che $v_1, v_2, \ldots, v_n$ sono linearmente indipendenti, ossia $\B$ è una base. $\blacksquare$
    \item I vettori di $\B$ sono linearmente indipendenti \\
    Dobbiamo quindi dimostrare che $\braket{v_1, v_2, \ldots, v_n} = \V$. Se per assurdo questo non fosse vero, allora $\exists v \in \V$ tale che $v \not\in \braket{v_1, v_2, \ldots, v_n}$. Potremmo quindi dire che $v_1, \ldots, v_n, v$ sono linearmente indipendenti. 
    Ma dato che $\dim \V = \abs\B = n$, è impossibile che nello spazio $\V$ vi siano $n+1$ vettori linearmente indipendenti. Infatti, $v, v_1, \ldots, v_n$ sono linearmente indipendenti se e solo se $x_1 v_1+ \ldots + x_n v_n + x_{n+1} v = \undzero$ ha come unica soluzione la soluzione banale $x_i = 0 \enspace \forall i \in \{1,2, \ldots, n+1\}$.
    Ma questa equazione così scritta altro non è che un sistema di $n$ equazioni in $n+1$ incognite, e come già accennato in nella dimostrazione della Proposizione \ref{propos-rel-lin}, un sistema di questo tipo ammette sempre almeno una soluzione non banale. Troviamo quindi che i vettori $v_1, \ldots, v_n, v$ sono linearmente dipendenti, arrivando quindi ad un assurdo. Dunque possiamo concludere che $\braket{v_1, v_2, \ldots, v_n} = \V$. $\blacksquare$
\end{enumerate}

\begin{Proposizione}[]
\label{propos-dim-spa-vet}
Preso uno spazio vettoriale $\V$ su $\K$ di dimensione $n$, se $\W$ è un sottospazio di $\V$ allora $\bm{\dim \W \leq n}$ e in particolare $\bm{\dim \W = n} \iff \bm{\W = \V}$. 
\end{Proposizione}

\Dimostrazione \nopagebreak

Sappiamo che $\W$ è finitamente generato per l'Osservazione \ref{oss-comb-lin-2}. Dunque ha una base \mbox{$\B = \{w_1, w_2, \ldots, w_m\}$}. Siccome $w_1, w_2, \ldots, w_m$ sono linearmente indipendenti, possono essere completati con altri vettori $u_1, u_2, \ldots, u_l \in \V$ per formare una base di $\V$, come visto nel punto II) della Proposizione \ref{propos-base-2}, formando la base $\Bs = \{w_1, w_2, \ldots, w_m, u_1, u_2, \ldots, u_l\}$. Dunque troviamo che $\dim \V = m + l = \dim \W + l$, e dato che $l \geq 0$ sempre, troviamo che $\dim \V \geq \dim \W$. Se poi prendiamo $l = 0$, significa che $\Bs$ è una base non solo di $\V$, ma anche di $\W$, dunque $\W = \V$. $\blacksquare$

\Esempi \nopagebreak

Calcolare la \textbf{dimensione} dello spazio vettoriale $\bm{\K^n}$. \nopagebreak

Sappiamo che la base standard $S = \{e_1, e_2, \ldots, e_n\}$ è una base di $\K^n$, dunque si può concludere che $\bm{\dim \K^n = n}$.

\Separatore

Calcolare la \textbf{dimensione} dello spazio vettoriale $\bm{\K[x]_{\leq d}}$. \nopagebreak

Con $d = 1$, lo spazio vettoriale diventa $\K[x]_{\leq 1} = \{p = a_1 x + a_0 \mid a_0,a_1 \in \K\}$. Una base di questo spazio vettoriale è, banalmente, la base $\B = \{1, x\}$. Infatti, $\braket{1,x} = \K[x]_{\leq 1}$ dato che, preso un qualsiasi $p \in \K[x]_{\leq 1}$, questo può essere riscritto come $1 \cdot a_0 + x \cdot a_1$, ed inoltre i due vettori sono anche linearmente indipendenti, dato che l'equazione $\lambda \cdot 1 + \mu \cdot x = \undzero$ è verificata solo se $\lambda = \mu = 0$. Dunque $\dim \K[x]_{\leq 1} = 2$. \\
Seguendo lo stesso procedimento, si può dimostrare per induzione che $\bm{\dim \K[x]_{\leq d} = d + 1}$.

\begin{Proposizione}[Formula di Grassmann]
\label{grassmann}
Preso $\V$ uno spazio vettoriale $\V$ su $\K$ finitamente generato e presi $\U, \W$ sottospazi di $\V$, $\bm{\dim (\U + \W) + \dim (\U \cap \W) = \dim \U + \dim \W}$.
\end{Proposizione}

\Dimostrazione \nopagebreak

Sia $k = \dim (\U \cap \W)$. Prendiamo ora la base $\B_{\U \cap \W} = \{v_1, v_2, \ldots, v_k\}$ di $\U \cap \W$. Estendiamo tale base a una base di $\U$, trovando $\B_{\U} = \{v_1, v_2, \ldots, v_k, u_1, u_2, \ldots, u_m\}$, e a una base di $\W$, trovando $\B_{\W} = \{v_1, v_2, \ldots, v_k, w_1, w_2, \ldots, w_p\}$. Dunque $\dim \U = k + m$ e $\dim \W = k + p$. A questo punto per dimostrare la formula di Grassmann non ci basta che dimostrare che $\B_{\U + \W} = \{v_1, v_2, \ldots, v_k, u_1, u_2, \ldots, u_m, w_1, w_2, \ldots, w_p\}$ è una base di $\U + \W$: questo perché a quel punto avremmo trovato che $\dim (\U + \W) + \dim (\U \cap \W) = k + m + p + k = \dim \U + \dim \W$, che è proprio la formula di Grassmann.

Innanzitutto cerchiamo di dimostrare che $\braket{v_1, v_2, \ldots, v_k, u_1, u_2, \ldots, u_m, w_1, w_2, \ldots, w_p} = \U + \W$.\\
Sappiamo che ogni elemento $z \in \U + \W$ può essere riscritto come $z = u + w$, come visto nella Definizione \ref{def-somm-sott-vett}. Poiché $u \in \braket{v_1, v_2, \ldots, v_k, u_1, u_2, \ldots, u_m}$ e \mbox{$w \in \braket{v_1, v_2, \ldots, v_k, w_1, w_2, \ldots, w_p}$}, troviamo che $z \in \braket{v_1, v_2, \ldots, v_k, u_1, u_2, \ldots, u_m, w_1, w_2, \ldots, w_p}$, e quindi troviamo che \mbox{$\braket{v_1, v_2, \ldots, v_k, u_1, u_2, \ldots, u_m, w_1, w_2, \ldots, w_p} = \U + \W$}.

Dobbiamo ora invece dimostrare che $v_1, v_2, \ldots, v_k, u_1, u_2, \ldots, u_m, w_1, w_2, \ldots, w_p$ sono linearmente indipendenti.\\
Vogliamo quindi dimostrare che l'equazione $\sum_{i = 1}^k a_i v_i + \sum_{j = 1}^m b_j u_j + \sum_{l = 1}^p c_l u_l = \undzero$ è verificata se e solo se $a_i = b_j = c_l = 0 \enspace \forall i \in \{1,2,\ldots,k\},  \forall j \in \{1,2,\ldots,m\},  \forall l \in \{1,2,\ldots,p\}$. Spostiamo ora tutti i vettori $w_l$ a destra, e troviamo $\sum_{i = 1}^k a_i v_i + \sum_{j = 1}^m b_j u_j = -\sum_{l = 1}^p c_l w_l$. Notiamo che il primo membro rappresenta un qualsiasi vettore di $\U$, mentre il secondo membro è sicuramente un vettore che si trova in $\W$. Se i due membri sono uguali, significa che il vettore a sinistra si trova in $\U \cap \W$. 
Troviamo quindi che $-\sum_{l = 1}^p c_l w_l = \sum_{i = 1}^k d_i v_i$, ossia $\sum_{l = 1}^p c_l w_l + \sum_{i = 1}^k d_i v_i = \undzero$. Dato che questi vettori formano la base $\B_{\W}$, sono linearmente indipendenti, dunque $c_l = 0 \enspace \forall l \in \{1,2,\ldots,p\}$. Tornando ora all'equazione $\sum_{i = 1}^k a_i v_i + \sum_{j = 1}^m b_i u_i = -\sum_{l = 1}^p c_l w_l$, sostituendo i valori dei vari $c_l$ troviamo che $\sum_{i = 1}^k a_i v_i + \sum_{j = 1}^m b_j u_j = \undzero$. Dato che questi vettori formano la base $\B_{\U}$, sono linearmente indipendenti, dunque $a_i = b_j = 0 \enspace \forall i \in \{1,2 \ldots, k\}, \forall j \in \{1,2, \ldots, m\}$. Abbiamo quindi dimostrato che i vettori sono linearmente indipendenti.

Avendo dimostrato che $\B_{\U \cap \W}$ è una base di $\U \cap \W$, possiamo concludere che \mbox{$\dim (\U + \W) + \dim (\U \cap \W) = \dim \U + \dim \W$}. $\blacksquare$

\begin{Proposizione}[Corollario]
\label{corol-grassman}
Preso $\V$ uno spazio vettoriale finitamente generato e $\U, \W$ sottospazi di $\V$, \mbox{$\bm{\dim (\U \cap \W) \geq \dim \U + \dim \W - \dim \V}$}.
\end{Proposizione}
Questo è vero perché, essendo $\U + \W$ un sottospazio vettoriale di $\V$, per la Proposizione \ref{propos-dim-spa-vet} \mbox{$\dim (\U + \W) \leq \dim \V$}, e per Grassmann invece sappiamo che \mbox{$\dim (\U \cap \W) = \dim \U + \dim \W - \dim (\U + \W)$}, dunque $\dim (\U \cap \W) \geq \dim \U + \dim \W - \dim \V$.


\begin{Proposizione}[]
\label{propos-dim-spaz-vett-2}
Preso lo spazio vettoriali $\V$ finitamente generato con $\dim \V = n$, e presi i sottospazi vettoriali di $\W$ $\W, \U$ tali che $\W \cap \U = \{\undzero\}$, $\bm{\dim \W + \dim \U = n} \iff \bm{\V = \W + \U}$. \\
In questo caso, diciamo che $\bm{\V = \W \oplus \U}$.
\end{Proposizione}

\Dimostrazione \nopagebreak

\begin{enumerate}[label=\arabic*°]
    \item $\dim \W + \dim \U = n \implies \V = \W + \U$ \\
    Dato che $\W \cap \U = \{\undzero\}$ anche $\dim (\W \cap \U) = 0$. Per la formula di Grassmann, troviamo che $\dim (\U + \W) + \dim (\U \cap \W) = \dim (\U + \W) = \dim \U + \dim \W = n$. Dato che $\W + \U$ è un sottospazio di $\V$, e $\dim (\W + \U) = \dim \V$, per la Proposizione \ref{propos-dim-spa-vet} troviamo che $\W + \U = \V$. $\blacksquare$ 
    \item $\dim \W + \dim \U = n \impliedby \V = \W + \U$ \\
    Dato che $\W \cap \U = \{\undzero\}$ anche $\dim (\W \cap \U) = 0$. Per la formula di Grassmann, troviamo che $\dim (\U + \W) + \dim (\U \cap \W) = \dim (\U + \W) = \dim \V = \dim \U + \dim \W = n$. $\blacksquare$
\end{enumerate}

\begin{Osservazione}[]
\label{oss-dim-spaz-vett}
Preso $\V$ uno spazio vettoriale finitamente generato su un campo $\K$ finito e di cardinalità $q$, $\bm{\dim \V  = d} \implies \bm{\lvert \V \rvert = q^d}$.
\end{Osservazione}

\Dimostrazione \nopagebreak

Dato che $\dim \V = d$, esiste sicuramente una base $\B = \{v_1, v_2, \ldots, v_d\}$ di $\V$. Dunque \mbox{$\V = \braket{v_1, v_2, \ldots, v_d} =\{v = \lambda_1 v_1 + \lambda_2 v_2 + \ldots + \lambda_k v_k \mid \lambda_1, \lambda_2, \ldots, \lambda_k \in \K\}$}. Per creare quindi un vettore $v \in \V$, abbiamo per ogni $\lambda_i$ $q$ possibilità, essendo la cardinalità di $\K$ $q$. Dunque, per creare un vettore di $\V$ dobbiamo scegliere $d$ volte tra $q$ opzioni, ossia $q^d$. Il totale di tutti i vettori in $\V$ è quindi $q^d$, ossia $\abs{\V} = q^d$. $\blacksquare$  

\section{Applicazioni lineari}

\begin{Definizione}[Applicazione lineare]
\label{def-app-lin}
Presi $\V, \W$ spazi vettoriali su $\K$, un'applicazione $f: \V \longrightarrow \W$ è detta \textbf{lineare} se:
\begin{itemize}
    \item $\bm{f(v_1 + v_2) = f(v_1) + f(v_2)} \enspace \forall v_1, v_2 \in \V$
    \item $\bm{f(\lambda v) = \lambda f(v)} \enspace \forall v \in \V, \forall \lambda \in \K$
\end{itemize}
\end{Definizione}


\Esempio \nopagebreak

Dimostrare che la seguente applicazione è \textbf{lineare}: \nopagebreak
\begin{align*}
\bm{f: \K^n} &\longrightarrow \bm{\K} \\
\bm{\X} &\longmapsto \bm{\sum_{i = 1}^n a_i x_i}
\end{align*}

Dobbiamo quindi verificare che, presi $\X, \Y \in \K^n$, $f(\X + \Y) = f(\X) + f(\Y)$ e che, preso $\X \in \K^n$ e $\lambda \in \K$, $f(\lambda \X) = \lambda f(\X)$.

Banalmente 
\[f(\X + \Y) = \sum_{i = 1}^n a_i(x_i + y_i) = \sum_{i = 1}^n a_i x_i + \sum_{i = 1}^n a_i y_i = f(\X) + f(\Y)\]
\[f(\lambda \X) = \sum_{i = 1}^n \lambda a_i x_i = \lambda \sum_{i = 1}^n a_i x_i = \lambda f(\X)\]

Possiamo quindi concludere che $f$ \textbf{è un'applicazione lineare}.

\begin{Osservazione}[]
\label{oss-app-lin}
Per dimostrare che un'applicazione $f: \V \longrightarrow \W$ sia lineare, basta dimostrare che, presi $\X, \Y \in \V$ e $\lambda, \mu \in \K$, $\bm{f(\lambda \X + \mu \Y) = \lambda f(\X) + \mu f(\Y)}$.
\end{Osservazione}

\begin{Proprieta}[Proprietà delle applicazioni lineari]
\label{prop-app-lin}
\begin{enumerate}
    \item $\bm{f(\lambda_1 v_1 + \lambda_2 v_2 + \ldots + \lambda_n v_n) = \lambda_1 f(v_1) + \lambda_2 f(v_2) + \ldots  + \lambda_n f(v_n)}$. \\ 
    In sostanza, \textit{l'immagine di una combinazione lineare è la combinazione lineare delle immagini}.
    \item $\bm{f(\undzero) = \undzero}$
    \item $\bm{f(-v) = -f(v)}$
\end{enumerate}
\end{Proprieta}

\Dimostrazione \nopagebreak

\begin{enumerate}
    \item La dimostrazione è molto semplice:
    \begin{align*}
    f(\lambda_1 v_1 + \lambda_2 v_2 + \ldots + \lambda_n v_n) &= f(\lambda_1 v_1) + f(\lambda_2 v_2 + \lambda_3 v_3 + \ldots + \lambda_n v_n)  \\
     &= \lambda_1 f(v_1) + f(\lambda_2 v_2) + f(\lambda_3 v_3 + \lambda_4 v_4 + \ldots + \lambda_n v_n) \\
     &= \lambda_1 f(v_1) + \lambda_2 f(v_2) + f(\lambda_3 v_3) + f(\lambda_4 v_4 + \ldots + \lambda_n v_n)
    \end{align*}
    Ripetendo questo stesso procedimento per tutti i $\lambda_i v_i$, troviamo che \mbox{$f(\lambda_1 v_1 + \lambda_2 v_2 + \ldots + \lambda_n v_n) = \lambda_1 f(v_1) + \lambda_2 f(v_2) + \ldots + \lambda_n f(v_n)$}. $\blacksquare$
    \item $f(\undzero) = f(0 \cdot \undzero) = 0 \cdot f(\undzero) = \undzero$. $\blacksquare$
    \item $f(-v) = f(-1 \cdot v) = -1 \cdot f(v) = -f(v)$. $\blacksquare$
\end{enumerate}

\begin{Osservazione}[]
\label{oss-app-lin-2}
Tutte le applicazioni lineari $f: \K^n \longrightarrow \K$ sono scritte nella forma $\bm{f(\X) = \sum_{i = 1}^n a_i x_i}$.
\end{Osservazione}

\Dimostrazione \nopagebreak

Sappiamo che ogni vettore $\X \in \K^n$ può essere riscritto come $\X = \sum_{i = 1}^n x_i e_i$, essendo \mbox{$S = \{e_1, e_2, \ldots, e_n\}$} la base standard di $\K^n$. Per la proprietà 1) vista in Proprietà \ref{prop-app-lin}, sappiamo che $f(\X) = f(x_1 e_1 + x_2 e_2 + \ldots + x_n e_n) = x_1 f(e_1) + x_2 f(e_2) + \ldots + x_n f(e_n)$. Ponendo quindi $a_i = f(e_i)$, troviamo che $f(\X) = a_1 x_1 + a_2 x_2 + \ldots + a_n x_n = \sum_{i = 1}^n a_i x_i$. $\blacksquare$

\begin{Definizione}[Insieme delle applicazioni lineaari]
\label{def-insieme-app-lin}
Presi gli spazi vettoriali $\V$ e $\W$ su $\K$, chiamiamo $\bm{\mathcal{L}(\V, \W)}$ l'insimee delle applicazioni lineari da $\V$ a $\W$.
\end{Definizione}

Definiamo le operazioni di \textbf{somma} e \textbf{prodotto per scalare} per $\mathcal{L}(\V, \W)$.

\bigskip

\begin{minipage}[c]{0.45\textwidth}
    \centering
    {\Large\textbf{Somma}:}
    \begin{align*}
        \hspace{0.4cm}\mathcal{L}(\V, \W) \times \mathcal{L}(\V, \W) &\longrightarrow \mathcal{L}(\V, \W) \\
        \hspace{0.4cm}(f,g) &\longmapsto f+g \\
        \hspace{0.4cm}\text{Dove } (f + g)(v) &= f(v) + g(v)
    \end{align*}
\end{minipage}
\begin{minipage}[c]{0.45\textwidth}
    \centering
    {\Large\textbf{Prodotto per scalare}:}
    \begin{align*}
        \K \times \mathcal{L}(\V, \W) &\longrightarrow \mathcal{L}(\V, \W) \\
        (\lambda,f) &\longmapsto \lambda f \\
        \text{Dove } (\lambda f)(v) &= \lambda \cdot f(v)
    \end{align*}
\end{minipage}

\bigskip

\begin{Proposizione}[]
\label{propos-insieme-app-lin-spaz-vett}
Con le operazioni di somma e prodotto appena descritte, $\mathcal{L}(\V, \W)$ è uno \textbf{spazio vettoriale} su $\K$.
\end{Proposizione}

\Dimostrazione \nopagebreak

Per dimotrare che $\mathcal{L}(\V, \W)$ è uno spazio vettoriale, la prima osservazione vista in Osservazione \ref{oss-sott-vet-2} ci dice che basta dimostrare che questo abbia l'elemento nullo e che sia chiuso rispetto alla somma e al prodotto. \\
L'elemento nullo è presente dato che, presi $v_1, v_2 \in \V$ e $\lambda_1, \lambda_2 \in \K$, $\undzero(\lambda_1 v_1 + \lambda_2 v_2) = \undzero = \lambda_1 \undzero(v_1) + \lambda_2 \undzero(v_2)$, dove $\undzero$ è l'applicazione nulla che manda tutto in $\undzero$.\\ 
Se poi prendiamo $f, g \in  \mathcal{L}(\V, \W)$, $v_1, v_2 \in \V$ e $\lambda_1, \lambda_2 \in \K$, allora 
\begin{align*}
 (f + g)(\lambda_1 v_1 + \lambda_2 v_2) &= f(\lambda_1 v_1 + \lambda_2 v_2) + g(\lambda_1 v_1 + \lambda_2 v_2)   \\
 &= \lambda_1 f(v_1) + \lambda_2 f(v_2) + \lambda_1 g(v_1) + \lambda_2 g(v_2) \\
 &= \lambda_1 (f + g)(v_1) + \lambda_2 (f + g)(v_2)
\end{align*}
Dunque $\mathcal{L}(\V, \W)$ è chiuso rispetto alla somma.\\
Se infine prendiamo $f \in \mathcal{L}(\V, \W)$, $v_1, v_2 \in \V$ e $\lambda, \lambda_1, \lambda_2 \in \K$, allora
\begin{align*}
 (\lambda f)(\lambda_1 v_1 + \lambda_2 v_2) &= \lambda (\lambda_1 f(v_1) + \lambda_2 f(v_2))  \\
 &= \lambda \lambda_1 f(v_1) + \lambda \lambda_2 f(v_2) \\
 &= \lambda_1 (\lambda f(v_1)) + \lambda_2 (\lambda f(v_2))
\end{align*}
Dunque $\mathcal{L}(\V, \W)$ è chiuso rispetto al prodotto. \\
Possiamo quindi concludere che $\mathcal{L}(\V, \W)$ è uno spazio vettoriale. $\blacksquare$

\section{Proprietà delle applicazioni lineari}

\begin{Definizione}[Nucleo]
\label{def-nucleo}
Il \textbf{nucleo} di un'applicazione lineare $f: \V \longrightarrow \W$ è così definito:
$\bm{\ker f := \{v \in \V \mid f(v) = \undzero\}}$, ossia l'insieme degli elementi di $\V$ la cui immagine è $\undzero$.
\end{Definizione}

\begin{Osservazione}[]
\label{oss-nucleo-sottospazio}
Il \textbf{nucleo} di un'applicazione lineare è un \textbf{sottospazio del dominio}.
\end{Osservazione}

\Dimostrazione \nopagebreak

Presa un'applicazione lineare $f: \V \longrightarrow \W$, per dimostrare che $\ker f$ è un sottospazio di $\V$ dobbiamo dimostrare che non è vuoto e che, presi $v_1, v_2 \in \ker f$ e $\lambda_1, \lambda_2 \in \K$, allora $\lambda_1 v_1 + \lambda_2 v_2 \in \ker f$. \\
È facile notare che non sia vuoto, dato che $f(\undzero) = \undzero$ e dunque $\undzero \in \ker f$. \\
Dato che $f$ è lineare, $f(\lambda_1 v_1 + \lambda_2 v_2) = \lambda_1 f(v_1) + \lambda_2 f(v_2) = \lambda_1 \cdot \undzero + \lambda_2 \cdot \undzero = \undzero$, perché $v_1, v_2 \in \ker f$. Dunque $\lambda_1 v_1 + \lambda_2 v_2 \in \ker f$.\\
Possiamo quindi concludere che $\ker f$ è un sottospazio di $\V$. $\blacksquare$ 

\begin{Proposizione}[]
\label{propos-applicazione-lineare-iniettiva-nucleo}
Presa un'applicazione lineare $f$, $\bm{f}$ \textbf{è iniettiva} $\iff \bm{\ker f = \{\undzero\}}$, ossia il nucleo è il più piccolo possibile. 
\end{Proposizione}

\vspace{2cm}

\Dimostrazione \nopagebreak

\begin{enumerate}[label=\arabic*°]
    \item $f$ è iniettiva $\implies \ker f = \{\undzero\}$ \\
    Se $f$ è iniettiva, allora c'è un solo elemento del dominio che abbia come immagine $\undzero$, e abbiamo già dimostrato che in ogni applicazione lineare $f(\undzero) = \undzero$. Dunque $\ker f = \{\undzero\}$. $\blacksquare$
    \item $f$ è iniettiva $\impliedby \ker f = \{\undzero\}$ \\
    Prendiamo due elementi del dominio $v_1$ e $v_2$ tali che $f(v_1) = f(v_2)$. Essendo $f$ lineare, $f(v_1 - v_2) = f(v_1) - f(v_2) = \undzero$. Dunque $v_1 - v_2 \in \ker f$, ma dato che $\ker f = \{\undzero\}$, troviamo che $v_1 = v_2$, ossia $f$ è iniettiva. $\blacksquare$
\end{enumerate}

\begin{Osservazione}[]
\label{oss-immagine-sottospazio}
L'\textbf{insieme immagine} di un'applicazione lineare è un \textbf{sottospazio del codominio}.
\end{Osservazione}

\Dimostrazione \nopagebreak

Presa un'applicazione lineare $f: \V \longrightarrow \W$, per dimostrare che $Im(f)$ è un sottospazio di $\W$ dobbiamo dimostrare che non è vuoto e che, presi $w_1, w_2 \in Im(f)$ e $\lambda_1, \lambda_2 \in \K$, \mbox{$\lambda_1 w_1 + \lambda_2 w_2 \in Im(f)$}. \\
È facile notare che non sia vuoto, dato che $f(\undzero) = \undzero$ e dunque $\undzero \in Im(f)$. \\
Dato che $w_1, w_2 \in Im(f)$, sappiamo che esistono $v_1, v_2 \in \V$ tali che $f(v_1) = w_1$ e $f(v_2) = w_2$. Dunque $\lambda w_1 + \lambda w_2 = \lambda f(v_1) + \lambda f(v_2) = f(\lambda_1 v_1 + \lambda_2 v_2)$, essendo $f$ lineare. Quindi \mbox{$\lambda_1 w_1 + \lambda_2 w_2 \in Im(f)$}, essendo l'immagine di $\lambda_1 v_1 + \lambda_2 v_2$.\\
Possiamo quindi concludere che $Im(f)$ è un sottospazio di $\W$. $\blacksquare$

\begin{Proposizione}[]
\label{propos-dim-dominio-dim-ker-dim-im}
Presa un'applicazione lineare $f: \V \longrightarrow \W$, con $\V$ e $\W$ finitamente generati, \mbox{$\bm{\dim \V = \dim \ker f + \dim Im(f)}$}.
\end{Proposizione}

\Dimostrazione \nopagebreak

Prendiamo una base del $\ker f$ $\B = \{v_1, \ldots, v_a\}$, e estendiamola ad una base di $\V$, essendo il nucleo di $f$ un sottospazio vettoriale del suo dominio $\V$, formando la base $\Bs = \{v_1, \ldots, v_a, w_1, \ldots, w_b\}$. Dunque $\dim \ker f = a$ e $\dim \V = a + b$. Dobbiamo quindi semplicemente dimostrare che $\dim Im(f) = b$. \\
Per farlo, dimostriamo che $\D = \{f(w_1), \ldots, f(w_b)\}$ è una base di $Im(f)$.
\begin{itemize}
    \item $\braket{f(w_1), \ldots, f(w_b)} = Im(f)$ \\
    Preso $u = f(v) \in Im(f)$, sappiamo che $v = \lambda_1 v_1 + \ldots + \lambda_a v_a + \mu_1 w_1 + \ldots + \mu_b w_b$, dunque $u = f(v) = f(\lambda_1 v_1 + \ldots + \lambda_a v_a + \mu_1 w_1 + \ldots + \mu_b w_b)$. Essendo $f$ lineare, troviamo che \mbox{$u = f(\lambda_1 v_1 + \ldots + \lambda_a v_a) + f(\mu_1 w_1 + \ldots + \mu_b w_b) = \undzero + f(\mu_1 w_1 + \ldots + \mu_b w_b)$}, dato che i vettori $v_1, \ldots, v_a$ formano una base del nucleo di $f$. 
    Dunque $u = \mu_1 f(w_1) + \ldots + \mu_b f(w_b)$, ossia $\braket{f(w_1), \ldots, f(w_b)} = Im(f)$.
    \item $f(w_1), \ldots, f(w_b)$ sono linearmente indipendenti \\
    Partendo dall'equazione $\mu_1 f(w_1) + \ldots \mu_b f(w_b)$, essendo $f$ lineare troviamo \mbox{$f(\mu_1 w_1 + \ldots + \mu_b w_b) = \undzero$}, ossia \mbox{$\mu_1 w_1 + \ldots + \mu_b w_b$}. Troviamo quindi l'equazione \mbox{$\mu_1 w_1 + \ldots + \mu_b w_b = \lambda_1 v_1 + \ldots + \lambda_a v_a$}. Se spostiamo tutti i membri a sinistra, troviamo \mbox{$\lambda_1 v_1 + \ldots + \lambda_a v_a - \mu_1 w_1 - \ldots - \mu_b w_b = \undzero$}. Dato che $\Bs = \{v_1, \ldots, v_a, w_1, \ldots, w_b\}$ è una base di $\V$, sappiamo che questi vettori sono linearmente indipendenti. Dunque l'equazione che abbiamo appena trovato ha come unica soluzione la soluzione banale $\lambda_1 =  \ldots = \lambda_a = \mu_1 = \ldots = \mu_b = 0$. Dunque anche \mbox{$\mu_1 f(w_1) + \ldots + \mu_b f(w_b) = \undzero$} ha come unica soluzione $\mu_1 = \ldots = \mu_b = 0$, e quindi $f(w_1), \ldots, f(w_b)$ sono linearmente indipendenti.
\end{itemize}

Essendo quindi $\D$ una base di $Im(f)$, troviamo che $\dim Im(f) = b$ e quindi concludimao che $\dim \V = \dim \ker f + \dim Im(f)$. $\blacksquare$

\begin{Proposizione}[Corollario]
\label{corollario-dim-dominio-nucleo-im}
Presa un'applicazione lineare $f: \V \longrightarrow \W$, con $\V$ e $\W$ finitamente generati, \mbox{$\bm{\dim \ker f \geq \dim \V - \dim \W}$}.
\end{Proposizione}

Questo è vero dato perché $Im(f)$ è un sottospazio di $\W$, dunque per la Proposizione \ref{propos-dim-spa-vet} $\dim Im(f) \leq \dim \W$, e dato che per la Proposizione \ref{propos-dim-dominio-dim-ker-dim-im} $\dim \ker f = \dim \V - \dim Im(f)$, troviamo che $\dim \ker f \geq \dim \V - \dim \W$.

\begin{Proposizione}[]
\label{propos-app-biunivoca-nucleo-nullo}
Presa un'applicazione lienare $f: \V \longrightarrow \W$, con $\V$ e $\W$ finitamente generati tali che $\dim \V = \dim \W$, \mbox{$\bm{f}$ \textbf{è biunivoca} $\iff \bm{\ker f = \{\undzero\}}$}.
\end{Proposizione}

\Dimostrazione \nopagebreak

\begin{enumerate}[label=\arabic*°]
    \item $f$ è biunivoca $\implies \ker f = \{\undzero\}$ \\
    Se $f$ è biunivoca allora è anche iniettiva, e per la Proposizione \ref{propos-applicazione-lineare-iniettiva-nucleo} possiamo dire che $\ker f = \{\undzero\}$. $\blacksquare$
    \item $f$ è biunivoca $\impliedby \ker f = \{\undzero\}$ \\
    Se $\ker f = \{\undzero\}$ allora $f$ è iniettiva sempre per la stessa proposizione. Dato che \mbox{$\dim \V = \dim \ker f + \dim Im(f) = \dim Im(f)$} essendo $\ker f = \{\undzero\}$, e dato che $\dim \V = \dim \W$ per ipotesi, troviamo che $\dim \W = \dim Im(f)$. Essendo $Im(f)$ un sottospazio di $\W$, significa che $\W = Im(f)$, ossia $f$ è suriettiva. Dunque $f$ è biunivoca. $\blacksquare$
\end{enumerate}

\begin{Proposizione}[Combinazione di applicazioni lineari]
\label{propos-comb-app-lin}
Prese due applicazioni lienari $g: \U \longrightarrow \V$ e $f: \V \longrightarrow \W$, $\bm{f \circ g: \U \longrightarrow \W}$ è \textbf{lineare}.
\end{Proposizione}

\Dimostrazione \nopagebreak

Per dimostrare che $f \circ g$ è lineare prendiamo $u_1, u_2 \in \U$ e $\lambda_1, \lambda_2 \in \K$ e dimostriamo che \mbox{$f \circ g (\lambda_1 u_1 + \lambda_2 u_2) = \lambda_1 f \circ g (u_1) + \lambda_2 f \circ g(u_2)$}. \\
Sappiamo che $f \circ g (\lambda_1 u_1 + \lambda_2 u_2) = f(g(\lambda_1 u_1 + \lambda_2 u_2))$. Essendo $g$ lineare, troviamo che \mbox{$f(g(\lambda_1 u_1 + \lambda_2 u_2)) = f(\lambda_1 g(u_1) + \lambda_2 g(u_2))$}. Essendo però anche $f$ lineare, troviamo che \mbox{$f(\lambda_1 g(u_1) + \lambda_2 g(u_2)) = \lambda_1 f(g(u_1)) + \lambda_2 f(g(u_2)) = \lambda_1 f \circ g (u_1) + \lambda_2 f \circ g (u_2)$}. Possiamo quindi concludere che $f \circ g$ è lineare. $\blacksquare$

\begin{Osservazione}[]
\label{oss-id-app-lin}
L'applicazione
\vspace{-15pt}
\begin{align*}
id_\V : \V &\longrightarrow \V \\
v &\longmapsto v 
\end{align*}
è un'\textbf{applicazione lineare}.
\end{Osservazione}

\Dimostrazione \nopagebreak

Banalmente, presi $v_1, v_2 \in \V$ e $\lambda_1, \lambda_2 \in \K$, \mbox{$id_\V(\lambda_1 v_1 + \lambda_2 v_2) = \lambda_1 v_1 + \lambda_2 v_2 = \lambda_1 id_\V(v_1) + \lambda_2 id_\V(v_2)$}. $\blacksquare$

\section{Isomorfismi}

\begin{Definizione}[Isomorfismo]
\label{def-isomorfismo}
Presa un'applicazione lineare $f: \V \longrightarrow \W$, questa è un \textbf{isomorfismo} se esiste un'altra applicazione lineare $g: \W \longrightarrow \V$ tale che $g \circ f = Id_\V$ e $f \circ g = id_\W$, ossia se esiste \textbf{l'applicazione lineare inversa}.
\end{Definizione}

\begin{Osservazione}[]
\label{oss-def-isomorfismo}
Un'isomorfismo è un'\textbf{applicazione lineare biunivoca} e viceversa.
\end{Osservazione}

Infatti un isomorfismo per definizione è un'applicazione lineare, e dato che ammette un'inversa è per forza anche biunivoca. Questa osservazione ci sarà molto utile in futuro quando dovremo dimostrare l'isomorfismo di alcune applicazioni. 

\begin{Definizione}[Isomorfi]
\label{def-isomorfi}
Due spazi vettoriali $\V$ e $\W$ sono detti \textbf{isomorfi} ($\bm{\V \cong \W}$) se esiste un isomorfismo $f: \V \longrightarrow \W$.
\end{Definizione}

\begin{Osservazione}[]
\label{oss-basi-isomorfismi}
Preso uno spazio vettoriale $\V$ e una sua base $\B = \{v_1, \ldots, v_n\}$, se $f: \V \longrightarrow \W$ è un \textbf{isomorfismo}, allora $\bm{\B' = \{f(v_1), \ldots, f(v_n)\}}$ \textbf{è una base} di $\bm{\W}$.
\end{Osservazione}

\Dimostrazione \nopagebreak

\begin{itemize}
    \item $\braket{f(v_1), \ldots, f(v_n)} = \W$ \\
    Vogliamo quindi dimostrare che preso $w \in \W$, $w = \lambda_1 f(v_1) + \ldots + \lambda_n f(v_n)$. Essendo $f$ lineare, $w = f(\lambda_1 v_1 + \ldots + \lambda_n v_n)$, e dato che $\B = \{v_1, \ldots, v_n\}$ è una bse di $\V$, $\lambda_1 v_1 + \ldots + \lambda_n v_n = v \in \V$. Dunque troviamo $w = f(v)$. Essendo $f$ un isomorfismo, sappiamo che esiste l'applicazione lineare inversa di $f$, e che quindi questa è biiettiva e dunque suriettiva. Dunque abbiamo trovato che $\forall w \in \W \enspace \exists v \in \V$ tale che $f(v) = w$, e dato che $f(v) = \lambda_1 f(v_1) + \ldots + \lambda_n f(v_n)$, troviamo che $\braket{f(v_1), \ldots, f(v_n)} = \W$.
    \item $f(v_1), \ldots, f(v_n)$ sono linearmente indipendenti \\
    Vogliamo quindi dimostrare che l'equazione $\lambda_1 f(v_1) + \ldots + \lambda_n f(v_n) = \undzero$ ha come unica soluzione la soluzione banale. Essendo $f$ lineare, possiamo riscriverla come $f(\lambda_1 v_1 + \ldots + \lambda_n v_n) = \undzero$. Essendo $f$ un isomorfismo, questa è un'applicazione biiettiva e quindi iniettiva, dunque sappiamo che se $f(v) = \undzero$ allora $v = \undzero$. Dato che $v_1, \ldots, v_n$ sono linearmente indipendenti, l'unica soluzione all'equazione $\lambda_1 v_1 + \ldots + \lambda_n v_n = \undzero$ è la soluzione banale $\lambda_1 = \ldots = \lambda_n = 0$. Dunque $f(v_1), \ldots, f(v_n)$ sono linearmente indipendenti.
\end{itemize}
Possiamo quindi concludere che $\B'$ è una base di $\W$. $\blacksquare$

\begin{Proprieta}[Proprietà degli isomorfismi]
\label{prop-isomorfismi}
Presi $\U, \V, \W$ spazi vettoriali su $\K$
\begin{enumerate}
    \item $\bm{\V \cong \V}$
    \item $\bm{\V \cong \W} \implies \bm{\W \cong \V}$
    \item $\bm{\U \cong \V}$ e $\bm{\V \cong \W} \implies \bm{\U \cong \W}$
\end{enumerate}
\end{Proprieta}

\Dimostrazione \nopagebreak

\begin{enumerate}
    \item Banalmente, essendo $id_\V : \V \longrightarrow \V$ un isomorfismo, troviamo che $\V$ è isomorfo con sè stesso. $\blacksquare$
    \item Se $f: \V \longrightarrow \W$ è un isomorfismo, allora esiste l'appliazione lineare inversa $g: \W \longrightarrow \V$. Ma anche questa è un isomorfismo, che ha come applicazione lineare inversa $f$. Dunque anche $\W$ e $\V$ sono isomorfi. $\blacksquare$
    \item Se $f: \U \longrightarrow \V$ è un isomorfismo, e anche $g: \V \longrightarrow \W$ lo è, allora sappiamo che esistono $f^{-1}: \V \longrightarrow \U$ inversa di $f$ e $g^{-1}: \W \longrightarrow \V$ inversa di $g$. Per la proprietà 2) appena vista, sappiamo che anche $f^{-1}$ e $g^{-1}$ sono isomorfismi. Cerchiamo quindi di dimostrare che $h := g \circ f: \U \longrightarrow \W$ è un isomorfismo, la cui inversa è $h^{-1} := f^{-1} \circ g^{-1} : \W \longrightarrow \U$. 
    \begin{itemize}
        \item $h \circ h^{-1} = id_\W$ \\
        Preso $w \in \W$, ponendo $v := g^{-1}(w) \in \V$ e $u := f^{-1}(v) \in \U$ troviamo che \nopagebreak 
        \begin{align*}
         h \circ h^{-1} (w) &= h(h^{-1}(w)) \\
         &= g \circ f (f^{-1} \circ g^{-1}(w)) \\
         &= g(f(f^{-1}(g^{-1}(w)))) \\
         &=  g(f(f^{-1}(v))) \\
         &= g(f(u)) \\
         &= g(v) = w
        \end{align*}
        Dunque $h \circ h^{-1} = id_\W$.
        \item $h^{-1}\circ h = id_\U$ \\
        Preso $u \in \U$, ponendo $v := f(u) \in \V$ e $w := g(v) \in \W$ troviamo che \nopagebreak
        \begin{align*}
         h^{-1} \circ h (u) &= h^{-1}(h(u)) \\
         &= f^{-1} \circ g^{-1} (g \circ f(u)) \\
         &= f^{-1}(g^{-1}(g(f(u)))) \\
         &= f^{-1}(g^{-1}(g(v))) \\
         &= f^{-1}(g^{-1}(w)) \\
         &= f^{-1}(v) = u
        \end{align*}
        Dunque $h^{-1} \circ h = id_\U$.
    \end{itemize}
    Possiamo quindi concludere che $h$ è un isomorfismo, e che quindi $\U$ e $\W$ sono isomorfi. $\blacksquare$
\end{enumerate}

\begin{Osservazione}[]
\label{oss-dim-isomorfismi}
Presi due spazi vettoriali $\V$ e $\W$, $\bm{\V \cong \W} \iff \bm{\dim \V = \dim \W}$.
\end{Osservazione}

\Dimostrazione \nopagebreak

\begin{enumerate}[label=\arabic*°]
    \item $\V \cong \W \implies \dim \V = \dim \W$ \\
    Banalmente, presa una base $\B = \{v_1, \ldots, v_n\}$ di $\V$, sappiamo per l'Osservazione \ref{oss-basi-isomorfismi} che $\B' = \{f(v_1), \ldots, f(v_n)\}$ è una base di $\W$, dunque $\dim \V = \dim \W$. $\blacksquare$
    \item $\V \cong \W \impliedby \dim \V = \dim \W$ \\
    Se $\dim \V = \dim \W$, allora possiamo prendere due basi $\B = \{v_1, \ldots, v_n\}$ e $\Bs = \{w_1, \ldots, w_n\}$ rispettivamente di $\V$ e di $\W$. Prendiamo le applicazioni $X_\B: \V \longrightarrow \K^n$ e $X_\Bs: \W \longrightarrow \K^n$. Essendo queste due applicazioni lineari e biunivoche, per la Proposizione \ref{propos-corol-basi}, per l'Osservazione \ref{oss-def-isomorfismo} queste sono anche degli isomorfismi. Troviamo quindi che $\V \cong \K^n$ e $\W \cong K^n$. Per la proprietà 2) appena vista nelle Proprietà \ref{prop-isomorfismi}, troviamo anche che $\K^n \cong \W$, e per la proprietà 3) possiamo concludere che $\V \cong \W$. $\blacksquare$
\end{enumerate}

\chapter{Matrici}

\section{Definizioni}

\begin{Definizione}[Matrice]
\label{def-matrice}
Una \textbf{matrice} $m \times n$, ossia formata da $m$ righe e da $n$ colonne, è un'applicazione del tipo
\begin{align*}
\{(i,j) \mid 1 \leq i \leq m, 1 \leq j \leq n\} & \longrightarrow \K \\
(i,j) &\longmapsto a_{ij}
\end{align*}
Viene rappresentata così:
\[
\bm{A = (a_{ij}) = }
\begin{bmatrixb}
a_{11} & \dots & a_{1n}   \\
\vdots & \ddots & \vdots \\
a_{n1} & \dots & a_{nn}
\end{bmatrixb}
\]
L'insieme delle matrici $m \times n$ con entrate in $\K$ lo indichiamo così: $\bm{M_{mn}(\K)}$.
\end{Definizione}

\begin{Definizione}[Righe e colonne di una matrice]
\label{def-righ-col-matr}
Presa una matrice $A \in M_{mn}(\K)$, indichiamo con $\bm{A^i}$ la riga i-esima di $A$ e con $\bm{A_j}$ la colonna j-esima di $A$.
\end{Definizione}

È quindi possibile riscrivere una qualunque matrice $m \times n$ in questi due modi:
\[A = 
\begin{bmatrix}
A^1 \\
\vdots \\
A^m
\end{bmatrix}
=
\begin{bmatrix}
A_1 & \dots & A_n 
\end{bmatrix}
\]

Definiamo le operazioni di \textbf{somma} e \textbf{prodotto per scalare} per $M_{mn}(\K)$.

\bigskip

\begin{minipage}[c]{0.45\textwidth}
    \centering
    {\Large\textbf{Somma}:}
    \begin{align*}
        \hspace{0.4cm}M_{mn}(\K) \times M_{mn}(\K) &\longrightarrow M_{mn}(\K) \\
        \hspace{0.4cm}((a_{ij}),b_{ij}) &\longmapsto (c_{ij}) \\
        \hspace{0.4cm}\text{Dove } c_{ij} = a_{ij} + b_{ij}
    \end{align*}
\end{minipage}
\begin{minipage}[c]{0.45\textwidth}
    \centering
    {\Large\textbf{Prodotto per scalare}:}
    \begin{align*}
        \K \times M_{mn}(\K) &\longrightarrow M_{mn}(\K) \\
        (\lambda,(a_{ij})) &\longmapsto (b_{ij}) \\
        \text{Dove } b_{ij} = \lambda a_{ij}
    \end{align*}
\end{minipage}

\bigskip

\begin{Proposizione}[]
\label{propos-insieme-matrici-spazio-vettoriale}
Con le operazioni di somma e prodotto appena descritte, $\bm{M_{mn}(\K)}$ è uno \textbf{spazio vettoriale} su $\K$, la cui dimensione è pari a $\bm{m \cdot n}$.
\end{Proposizione}

\begin{Definizione}[Matrice unità]
\label{def-matrice-unita}
La \textbf{matrice unità} $\bm{1_n}$ è una matrice $n \times n$ che ha sulla diagonale principale come valori solo $1$ e tutti $0$ per il resto della matrice.
\[\bm{1_n} = 
\begin{bmatrixb}
1 & 0 & \dots & 0 \\
0 & 1 & \dots & 0 \\
\vdots & \vdots & \ddots & \vdots \\
0 & 0 & \dots & 1 \\
\end{bmatrixb}
\]
\end{Definizione}

\begin{Definizione}[Matrice scalare]
\label{def-matrice-scalare}
Una \textbf{matrice scalare} $\bm{\lambda 1_n}$ è una matrice $n \times n$ che ha sulla diagonale principale come valori solo $\lambda$ e tutti $0$ per il resto della matrice.
\[\bm{\lambda1_n} = 
\begin{bmatrixb}
\lambda & 0 & \dots & 0 \\
0 & \lambda & \dots & 0 \\
\vdots & \vdots & \ddots & \vdots \\
0 & 0 & \dots & \lambda \\
\end{bmatrixb}
\]
\end{Definizione}

\begin{Definizione}[Matrice diagonale]
\label{def-matrice-diagonale}
Una \textbf{matrice diagonale} $\bm{\lambda_i \delta_{ij}}$ è una matrice $n \times n$ che ha tutti i valori fuori dalla diagonale principale pari a $0$.
\[\bm{\lambda_i \delta_{ij}}
\begin{bmatrixb}
a_{11} & 0 & \dots & 0 \\
0 & a_{22} & \dots & 0 \\
\vdots & \vdots & \ddots & \vdots \\
0 & 0 & \dots & a_{nn} \\
\end{bmatrixb}
\]
\end{Definizione}

\subsection{Prodotto tra matrici}

\begin{Definizione}[Prodotto tra matrici]
\label{def-prodotto-tra-matrici}
Prese $A \in M_{mn}(\K)$ e $\B \in M_{np}(\K)$, $A \cdot B = C \in M_{mp}$, dove $\bm{c_{ih} = \sum_{j = 1}^n a_{ij} b_{jh}}$, ossia l'entrata $c_{ih}$ è la somma del prodotto scalare tra la riga i-esima di $A$ e la colonna $h$-esima di $B$.
\end{Definizione}

Prese $A \in M_{mn}(\K)$ e $B \in M_{np}(\K)$, possiamo vedere il prodotto tra matrici in questo modo:
\[A \cdot B = \begin{bmatrix}
    A^1 \\ \vdots \\ A^m
\end{bmatrix} \cdot \begin{bmatrix}
    B_1 & \dots & B_p
\end{bmatrix} = \begin{bmatrix}
    A^1B_1 & \dots & A^1 B_p \\
    \vdots & \ddots & \vdots \\
    A^m B_1 & \dots & A^m B_p
\end{bmatrix}\]

Da questa definizione, è facile capire che:
\begin{itemize}
    \item Si possono moltiplicare due matrici solo se \textit{il numero di colonne della prima è uguale al numero di righe della seconda}.
    \item Il prodotto tra matrici \textit{non è commutativo}.
\end{itemize}

\Esempio \nopagebreak

Calcolare il \textbf{prodotto} tra le matrici \nopagebreak
\[A =
\begin{bmatrixb}
2 & 1  \\
0 & -1
\end{bmatrixb}
\enspace 
B =
\begin{bmatrixb}
2 & 1 & 1 \\
1 & 5 & -1
\end{bmatrixb}
\]

Dato che $A$ ha $2$ colonne e $B$ ha $2$ righe, ha senso calcolare $A \cdot B$. 

Sappiamo che la matrice $A \cdot B$ è una matrice $2 \times 3$, dunque dobbiamo calcolarci i valori da $c_{11}$ fino a $c_{23}$.

\begin{align*}
&c_{11} = \sum_{j = 1}^2 a_{1j} b_{j1} = 2 \cdot 2 + 1 \cdot 1 = 5  \qquad c_{12} = \sum_{j = 1}^2 a_{1j} b_{j2} = 2 \cdot 1 + 1 \cdot 5 = 7 \\
&c_{13} = \sum_{j = 1}^2 a_{1j} b_{j3} = 2 \cdot 1 + 1 \cdot (-1) = 1 \qquad c_{21} = \sum_{j = 1}^2 a_{2j} b_{j1} = 0 \cdot 2 + (-1) \cdot 1 = -1 \\
&c_{22} = \sum_{j = 1}^2 a_{2j} b_{j2} = 0 \cdot 1 + (-1) \cdot 5 = -5 \qquad c_{23} = \sum_{j = 1}^2 a_{2j} b_{j3} = 0 \cdot 1 + (-1) \cdot (-1) = 1
\end{align*}

Concludiamo quindi che
\[C = \bm{A \cdot B =}
\begin{bmatrixb}
5 & 7 & 1 \\
-1 & -5 & 1 
\end{bmatrixb}
\]

\begin{Osservazione}[]
\label{oss-matrice-per-matrice-unità}
Presa una matrice $A \in M_{mn}(\K)$, $\bm{A \cdot 1_n = 1_m \cdot A = A}$.
\end{Osservazione}

\Dimostrazione \nopagebreak

Sappiamo che 
\[
(1_n)_{ik} =
\begin{cases}
 1 & \text{se } i = k \\
 0 & \text{se } i \neq k
\end{cases}
\]

Guardiamo prima $A \cdot 1_n$: dato che $(A \cdot 1_n)_{ik} = \sum_{j=1}^n a_{ij}(1_n)_{jk} = a_{ik}$ $\forall i \in \{1,2,\ldots,m\} \\ \forall k \in \{1,2,\ldots,n\}$, troviamo che $A \cdot 1_n = A$. \\
Guardiamo ora $1_m \cdot A$: dato che $(1_m \cdot A)_{ik} = \sum_{j=1}^m (1_m)_{ij}a_{jk} = a_{ik}$ $\forall i \in \{1,2,\ldots,m\} \\ \forall k \in \{1,2,\ldots,n\}$, troviamo che $1_m \cdot A = A$. \\
Possiamo quindi concludere che $A \cdot 1_n = 1_m \cdot A = A$. $\blacksquare$.

\begin{Proprieta}[Proprietà del prodotto tra matrici]
\label{prop-prodotto-matrici}
Prese $A \in M_{mn}, B \in M_{np}, C \in M_{pq}$, e $\lambda \in \K$
\begin{enumerate}
    \item $\bm{A \cdot \lambda 1_n = \lambda 1_m \cdot A = \lambda A }$.
    \item $\bm{A \cdot (B \cdot C) = (A \cdot B) \cdot C}$ (il prodotto tra matrici è \textit{associativo})
    \item $\bm{A \cdot (B + C) = A \cdot B + A \cdot C}$ (il prodotto tra matrici è \textit{distributivo})
\end{enumerate}
\end{Proprieta}

\Dimostrazione \nopagebreak

\begin{enumerate}
    \item Come già visto nella dimostrazione precedente dell'Osservazione \ref{oss-matrice-per-matrice-unità}
    \[
    (1_n)_{ik} =
    \begin{cases}
     1 & \text{se } i = k \\
     0 & \text{se } i \neq k
    \end{cases}
    \]
    Guardiamo prima $A \cdot \lambda1_n$: dato che $(A \cdot \lambda1_n)_{ik} = \sum_{j = 1}^n a_{ij}(\lambda 1_n)_{jk} = \lambda a_{ik}$ $\forall i \in \{1,2,\ldots,m\} \\ \forall k \in \{1,2,\ldots,n\}$, troviamo che $A \cdot \lambda 1_n = \lambda A$. \\
    Guardiamo ora $\lambda 1_m \cdot A$: dato che $(\lambda 1_m \cdot A)_{ik} = \sum_{j = 1}^m (\lambda 1_m)_{ij} a_{jk} = \lambda a_{ik}$ $\forall i \in \{1,2,\ldots,m\} \\ \forall k \in \{1,2,\ldots,n\}$, troviamo che $\lambda 1_m \cdot A = A \lambda = \lambda A$. \\
    Possiamo quindi concludere che $A \cdot \lambda 1_n = \lambda 1_m \cdot A  = \lambda A$. $\blacksquare$.
    \item Banalmente
    \begin{align*}
    &(A \cdot (B \cdot C))_{il} = \sum_{j = 1}^n a_{ij} (B \cdot C)_{jl} = \sum_{j = 1}^n a_{ij} \sum_{k = 1}^p b_{jk}c_{kl} = \sum_{j = 1}^n \sum_{k = 1}^p a_{ij} b_{jk} c_{kl} \\
    &((A \cdot B) \cdot C)_{il} = \sum_{k = 1}^p (A \cdot B)_{ik} c_{kl} = \sum_{k = 1}^p \left(\sum_{j = 1}^n a_{ij} b_{jk}\right) c_{kl} = \sum_{k = 1}^p \sum_{j = 1}^n a_{ij} b_{jk} c_{kl}
    \end{align*}
    Essendo tutte le entrate di entrambi i prodotti identiche, possiamo concludere che \mbox{$A \cdot (B \cdot C) = (A \cdot B) \cdot C$}. $\blacksquare$
    \item Banalmente
    \begin{align*}
    (A \cdot (B + C))_{ik} &= \sum_{j = 1}^n a_{ij}(B + C)_{jk} \\ 
    &= \sum_{j = 1}^n a_{ij}(b_{jk} + c_{jk}) \\
    &= \sum_{j = 1}^n a_{ij} b_{jk} + \sum_{j = 1}^n a_{ij} c_{jk} \\
    &= (A \cdot B)_{ik} + (A \cdot C)_{ik}
    \end{align*}
    Essendo tutte le entrate di entrambi i prodotti identiche, possiamo concludere che \mbox{$A \cdot (B + C) = A \cdot B + A \cdot C$}. $\blacksquare$
\end{enumerate}

\section{Matrici}

\begin{Definizione}[Applicazione lineare associata ad una matrice]
\label{def-app-lin-associata-matrice}
Presa una matrice $A \in M_{mn}(\K)$, l'\textbf{applicazione lineare associata a} $\bm{A}$ $\bm{L_A}$ è l'applicazione
\begin{align*}
\bm{L_A: \K^n} &\longrightarrow \bm{\K^m} \\
\bm{\X} &\longmapsto \bm{A \cdot \X} 
\end{align*}
Dove $\X$ è un vettore colonna, ossia $X = \begin{bmatrixb}
    x_1 \\ \vdots \\ x_n
\end{bmatrixb}$.
\end{Definizione}

Dimostriamo velocemente che questa applicazione è lineare. \\
Presi $\X, \Y \in \K^n$ e $\lambda, \mu \in \K$, verifichiamo che $L_A(\lambda \X + \mu \Y) = \lambda L_A(\X) + \mu L_A(\Y)$. \\ Banalmente
\[L_A(\lambda \X + \mu \Y) = A \cdot (\lambda \X + \mu Y) = A \lambda \X + A \mu Y = \lambda A \X + \mu A \Y = \lambda L_A(\X) + \mu L_A(\Y) \quad \blacksquare\]

\Esempio \nopagebreak

Presa una matrice $A \in M_{mn}(\K)$, calcolare la \textbf{dimensione} dell'insieme $\bm{\{\X \in \K^n \mid A \cdot \X = \undzero\}}$, ossia l'insieme delle soluzioni di un sistema di $m$ equazioni lienari omogenee in $n$ incognite che ha come coefficienti le entrate di $A$. \nopagebreak

Possiamo notare che l'insieme appena descritto altro non è che il nucleo dell'appliazione lineare associata alla matrice $A$, ossia il nucleo di $L_A$. Dunque $\{\X \in \K^n \mid A \cdot \X = \undzero\} = \ker L_A$. Per la Proposizione \ref{corollario-dim-dominio-nucleo-im} troviamo che $\dim \ker L_A \geq \dim \K^n - \dim \K^m$, ossia $\bm{\dim \ker L_A \geq n - m}$. \\Questo significa che se $n > m$, allora $\dim \ker L_A \geq 1$ e dunque il sistema ha almeno una soluzione non banale. Questa è una affermazione che avevamo utilizzato senza dimostrarla, che è stata quindi dimostrata tramite questo esempio.

\begin{Proposizione}[Lemmi]
\label{lemmi-matrici-e-app-lin}
Presi due spazi vettoriali $\V$ e $\W$ su $\K$, con $\V$ finitamente generato:
\begin{enumerate}[label=\Roman*)]
    \item Presi i vettori $v_1, \ldots, v_n \in \V$ linearmente indipendenti, $\bm{\forall w_1, \ldots, w_n \in \W} \enspace \exists \bm{f}$ \textbf{lineare} tale che $\bm{f(v_i) = w_i} \enspace \forall i \in \{1,\ldots,n\}$
    \item $\bm{\V = \braket{v_1, \ldots, v_n}}$ e $\bm{f,g}$ \textbf{applicazioni lineari} da $\V$ a $\W$ e $\bm{f(v_i) = g(v_i)} \enspace$ \mbox{$\forall i \in \{1,\ldots,n\}$} $\implies \bm{f = g}$
\end{enumerate}
\end{Proposizione}

\Dimostrazione \nopagebreak

\begin{enumerate}[label=\Roman*)]
    \item Essendo $v_1, \ldots, v_n$ linearmente indipendenti, questi possono formare o meno una base di $\V$. Dividiamo per questo due possibili casi:
    \begin{enumerate}[label=\arabic*°]
        \item $\B = \{v_1, \ldots, v_n\}$ è una base di $\V$ \\
        Sapendo che $\forall v \in \V$ possiamo scrivere $v = \sum_{i = 1}^n x_i v_i$, dove quindi $(x_1, \ldots, x_n) = X_\B(v)$, definiamo $f$ in questo modo:
        \begin{align*}
            f: \V &\longrightarrow \W \\
            v &\longmapsto \sum_{i = 1}^n x_i w_i
         \end{align*}
         Dimostriamo ora che $f$ è lineare, verificando che presi $u,v \in \V$ e $\lambda, \mu \in \K$, \mbox{$f(\lambda u + \mu v) = \lambda f(u) + \mu f(v)$}. \\
         Banalmente, se poniamo $u = \sum_{i = 1}^n x_i v_i$ e $v = \sum_{i = 1}^n y_i v_i$, troviamo che
         \[f(\lambda u + \mu v) = \sum_{i = 1}^n (\lambda (x_i w_i) + \mu (y_i w_i)) = \lambda \sum_{i = 1}^n x_i w_i + \mu \sum_{i = 1}^n y_i w_i = \lambda f(u) + \mu f(v) \blacksquare\]
    \item $\B = \{v_1, \ldots, v_n\}$ non è una base di $\V$.
    In questo caso, estendiamo $\B$ ad una base $\Bs = \{v_1, \ldots, v_n, u_1, \ldots, u_m\}$ di $\V$. A questo punto definiamo $f$ nello stesso modo in cui l'abbiamo definita prima, e, perla dimostrazione precedente, troviamo che $f$ è lineare. $\blacksquare$
    \end{enumerate}
    \item Sia $v = \sum_{i = 1}^n x_i v_i \in \V$. Per la proprietà 1) vista nelle Proprietà \ref{prop-app-lin}, sappiamo che \mbox{$f(x_1 v_1 + \ldots + x_n v_n) = x_1 f(v_1) + \ldots + x_n f(v_n) = \sum_{i = 1}^n x_i f(v_i)$}. Per lo stesso ragionamento, troviamo che $g(v) = \sum_{i = 1}^n x_i g(v_i)$. Dato che per ipotesi $f(v_i) = g(v_i) \enspace \forall i \in \{1, \ldots, n\}$, troviamo che $f = g$. $\blacksquare$ 
\end{enumerate}

\begin{Proposizione}[Corollario]
\label{corollario-applicazioni-lineari-matrici}
Presi gli spazi vettoriali $\V$ e $\W$ su $\K$, con $\V$ finitamente generato, e presa la base $\B = \{v_1, \ldots, v_n\}$ di $\V$ e i vettori $w_1, \ldots, w_m \in \W$, \textbf{esiste un'unica applicazione lineare} $\bm{f}$ che va da $\V$ a $\W$ tale che $\bm{f(v_i) = w_i} \enspace \forall i \in \{1, \ldots, n\}$.
\end{Proposizione}

I due lemma visti in precedenza erano necessari proprio per arrivare a questa conclusione.

\begin{Osservazione}[]
\label{oss-L_A(e_i) = A_i}
Presa una matrice $A \in M_{mn}(\K)$, $\bm{L_A(e_i) = A \cdot e_i =  A_i} \enspace \forall i \in \{1, \ldots, n\}$.
\end{Osservazione}

\Dimostrazione \nopagebreak

Sappiamo che $L_A(\X) = A \cdot \X$, e che $A = \begin{bmatrix}
    A_1 & \dots & A_n
\end{bmatrix}$, dove in particolare $A_i$ è la colonna i-esima di $A$. Dato che $\X$ abbiamo detto è un vettore colonna, troviamo che 
\[L_A(e_i) = A \cdot e_i = \begin{bmatrix}
    A_1 & \dots & A_i & \dots & A_n
\end{bmatrix} \cdot \begin{bmatrix}
    0 \\ \vdots \\ 1 \\ \vdots \\ 0
\end{bmatrix} = 0 \cdot A_1 + \ldots + 1 \cdot A_i + \ldots + 0 \cdot A_n = A_i \quad \blacksquare\] 

\begin{Proposizione}[]
\label{propos-unica-appl-lin-ass-a-matrice}
Presa l'applicazione lineare $f: \K^n \longrightarrow \K^m$, \textbf{esiste un'unica matrice} $\bm{A \in M_{mn}(\K)}$ tale che $\bm{f = L_A}$. \\
In sostanza, per ogni applicazione lineare esiste un'unica matrice a cui questa è associata.
\end{Proposizione}

\Dimostrazione \nopagebreak

Dimostriamo prima che $A$ sicuro esiste, e poi dimostriamo che è unica. \\
Per dimostrare che $A$ esiste, poniamo $A_j := f(e_j)$. Per l'Osservazione \ref{oss-L_A(e_i) = A_i}, sappiamo che $A_j = L_A(e_j)$. Dato che $L_A$ e $f$ sono lineari, e dato che $\K^n = \braket{e_1, \ldots, e_n}$, per la proposizione II) vista in Proposizione \ref{lemmi-matrici-e-app-lin}, possiamo concludere che $L_A = f$. \\
Per dimostrare che $A$ è unica, supponiamo che $L_A = L_B = f$. Allora $L_A(e_j) = L_B(e_j)$, ossia $A_j = B_j$, e dato che questo vale $\forall j \in \{1, \ldots, n\}$, troviamo che $A = B$. \\
Possiamo quindi concludere che esiste un'unica matrice $A$ tale che $f = L_A$. $\blacksquare$

\begin{Definizione}[Matrice associata]
\label{def-matrice-associata}
Presa un'applicazione lineare $f: \V \longrightarrow \W$, una base $\B = \{v_1, \ldots, v_n\}$ di $\V$ e una base $\Bs = \{w_1, \ldots, w_m\}$ di $\W$, la \textbf{matrice associata} di $f$ $\bm{M_\Bs^\B(f)}$ è la matrice che ha come colonne le coordinate in base $\Bs$ dei vettori di $\B$, ossia 
\[
\bm{M_\Bs^\B(f)}
\begin{bmatrixb}
X_\Bs(f(v_1)) & \dots & X_\Bs(f(v_n))
\end{bmatrixb}
\]
È facile capire quindi come il numero di righe di $M_\Bs^B(f)$ è pari al numero di vettori di $\Bs$ e invece il numero di colonne il numero di vettori di $\B$, dunque $\bm{M_\Bs^\B(f) \in M_{\dim \W, \dim \V}}$.
\end{Definizione}

\Esempi \nopagebreak

Presa l'applicazione
\begin{align*}
\bm{f: \R[x]_{\leq 2}} &\longrightarrow \bm{\R[x]_{\leq 2}} \\
 p &\longmapsto p + p'
\end{align*}
trovare una matrice associata $\bm{M_\Bs^\B(f)}$. \nopagebreak

Dato che dominio e codominio sono lo stesso spazio vettoriale, possiamo scegliere una stessa base $\B = \Bs = \{1, x, x^2\}$. Calcoliamo ora le coordinate delle immagini dei vettori di $\B$ nella base stessa $\B$.
\begin{align*}
&f(1) = 1 \longrightarrow X_\Bs(f(1)) = X_\Bs(1) = X_\Bs(1 \cdot 1 + 0 \cdot x + 0 \cdot x^2) = (1,0,0) \\
&f(x) = x+1 \longrightarrow X_\Bs(f(x)) = X_\Bs(x+1) = X_\Bs(1 \cdot 1 + 1 \cdot x + 0 \cdot x^2) = (1,1,0) \\
&f(x^2) = x^2+2x \longrightarrow X_\Bs(f(x^2)) = X_\Bs(x^2+2x) = X_\Bs(0 \cdot 1 + 2 \cdot x + 1 \cdot x^2) = (0,2,1) \\
\end{align*} 

Dunque 
\[ \bm{M_\Bs^B(f)} = 
\begin{bmatrixb}
1 & 1 & 0  \\
0 & 1 & 2 \\
0 & 0 & 1 
\end{bmatrixb}
\]

\Separatore

Presa l'applicazione $\bm{p_\theta : V(\E^2)} \longrightarrow \bm{V(\E^2)}$, che rappresenta la rotazione $\theta$ di un vettore nel piano euclideo, trovare una matrice associata $\bm{M_\Bs^\B(p_\theta)}$. \nopagebreak


\begin{minipage}[c]{0.60\textwidth}
Prima di cercare la matrice associata, cerchiamo di capire graficamente cosa rappresenta la funzione $p_\theta$. Preso un vettore $\vv{AB}$ sul pinao, rappresentato qui a destra in blu, il vettore $p_\theta(\vv{AB})$ è lo stesso vettore traslato di un angolo $\theta$, rappresentato qui a destra in rosso. 
\end{minipage}\hfill
\begin{minipage}[c]{0.35\textwidth}
\centering
\includegraphics[width=1\textwidth, valign=t]{Esercizio base ortonormale.png}
\end{minipage}

Dato che anche qui dominio e codominio coincidono, possiamo scegliere una stessa base $\B = \Bs = \{\underline{i}, \underline{j}\}$, che non è una base qualunque ma è una \textit{base ortonormale}. Una base ortonormale, in breve, è una base dove tutti i vettori sono perpendicolari tra loro e hanno come modulo $1$. \\
Per calcolare $p_\theta(\underline{i})$ e $p_\theta(\underline{j})$, ci è utile un'altra rappresentazione grafica. 

\begin{minipage}[c]{0.60\textwidth}
Ruotando il vettore $\underline{i}$ di un angolo $\theta$, troviamo che \mbox{$p_\theta(\underline{i}) = \underline{i} \cdot \cos(\theta) + \underline{j} \cdot \sin(\theta)$}. \\
Invece, ruotando il vettore $\underline{j}$ di un angolo $\theta$, troviamo che \mbox{$p_\theta(\underline{j}) = \underline{j} \cdot \cos(\theta) - \underline{i} \cdot \sin(\theta)$}.
\end{minipage}\hfill
\begin{minipage}[c]{0.35\textwidth}
\centering
\includegraphics[width=1\textwidth, valign=t]{Esercizio base ortonormale 2.png}
\end{minipage}

Troviamo quindi che
\begin{align*}
&X_\Bs(p_\theta(\underline{i})) = X_\Bs(\underline{i} \cdot \cos(\theta) + \underline{j} \cdot \sin(\theta)) = (\cos(\theta), \sin(\theta)) \\
&X_\Bs(p_\theta(\underline{j})) = X_\Bs(\underline{j} \cdot \cos(\theta) - \underline{i} \cdot \sin(\theta)) = (-\sin(\theta), \cos(\theta)) 
\end{align*}

Dunque
\[ \bm{M_\Bs^\B(f)} = 
\begin{bmatrixb}
\cos(\theta) & - \sin(\theta) \\
\sin(\theta) & \cos(\theta)
\end{bmatrixb}
\]

\begin{Proposizione}[]
\label{propos-matrice-associata-prodotto}
Presi due spazi vettoriali $\V$ e $\W$, un vettore $v \in \V$, un'applicazione lineare $f: \V \longrightarrow \W$, una base $\B$ di $\V$ e una base $\Bs$ di $\W$, $\bm{X_\Bs(f(v)) = M_\Bs^\B(f) \cdot X_\B(v)}$. 
\end{Proposizione}

\Dimostrazione \nopagebreak

Poniamo $\B = \{v_1, \ldots, v_n\}, \Bs = \{w_1, \ldots, w_m\}$ e, per comodità, $M_\Bs^\B(f) = A$. \\
Sappiamo che $X_\B(v) = (x_1, \ldots, x_n)$ tali che $v = \sum_{i = 1}^n x_i v_i$. Dunque $f(v) = f(\sum_{i = 1}^n x_i v_i)$, ed essendo $f$ lineare troviamo che $f(v) = \sum_{i = 1}^n x_i f(v_i)$. 
Poiché \mbox{$A_j = X_\Bs(f(v_j)) \enspace \forall j \in \{1, \ldots, n\}$}, possiamo scrivere $f(v_j) = \sum_{i = 1}^m a_{ij} w_i$. Dunque troviamo che
\begin{align*}
f(v) &= x_1 \sum_{i = 1}^m a_{i1} w_i + \ldots + x_n \sum_{i = 1}^m a_{in} w_i  \\
 &= w_1 \sum_{i = 1}^n a_{1i} x_i + \ldots + w_m \sum_{i = 1}^n a_{mi} x_i \\
 &= w_1 A^1 \cdot \X + \ldots + w_m A^m \cdot \X
\end{align*}
Dunque $X_\Bs(f(v)) = \begin{bmatrix}
    A^1 \cdot \X & \dots & A^m \cdot \X
\end{bmatrix} = A \cdot \X = M_\Bs^\B(f) \cdot X_\B(v)$. $\blacksquare$

\begin{Proposizione}[]
\label{propos-spaz-vet-app-lin-matrici}
Presi due spazi vettoriali $\V$ e $\W$ su $\K$ finitamente generati, con $\dim \V = n$ e $\dim \W = m$, e prese una base $\B$ di $\V$ e una base $\Bs$ di $\W$, allora l'applicazione
\begin{align*}
\bm{M_\Bs^\B: \mathcal{L(\V, \W)}} &\longrightarrow \bm{M_{mn}(\K)} \\
\bm{f} &\longmapsto \bm{M_\Bs^\B(f)}
\end{align*}
è un \textbf{isomorfismo}.
\end{Proposizione}

\Dimostrazione \nopagebreak

Vogliamo quindi dimostrare che $M_\Bs^\B$ è un'applicazione lineare biiettiva, come già visto nell'Osservazione \ref{oss-def-isomorfismo}.

\begin{enumerate}[label=\arabic*°]
    \item $M_\Bs^\B$ è lineare \\
    Affinché $M_\Bs^\B$ sia lineare è necessario che, prese $f_1, f_2 \in \mathcal{L}(\V, \W)$ e presi $\lambda_1, \lambda_2 \in \K$, $M_\Bs^\B(\lambda_1 f_1 + \lambda_2 f_2) = \lambda_1 M_\Bs^\B(f_1) + \lambda_2 M_\Bs^\B(f_2)$. Sappiamo che la colonna j-esima della matrice associata ad $f$ è pari a $X_\Bs(f(v_j))$. Dunque ci basta dimostrare che, $\forall j \in \{1, \ldots, n\}$, \mbox{$X_\Bs(\lambda_1 f_1(v_j) + \lambda_2 f_2(v_j)) = \lambda_1 X_\Bs(f_1(v_j)) + \lambda_2 X_\Bs(f_2(v_j))$}. Questo è vero essendo $X_\Bs$ un'applicazione lineare, dunque anche $M_\Bs^\B$ è lineare.
    \item $M_\Bs^\B$ è biiettiva \\
    Per dimostrare che $M_\Bs^\B$ sia biiettiva, per la Proposizione \ref{propos-app-biunivoca-nucleo-nullo} ci basta dimostrare che $\dim \mathcal{L}(\V, \W) = \dim M_{mn}(\K)$ e che $\ker M_\Bs^\B = \{\undzero\}$. \\
    Sappiamo già per la Proposizione \ref{propos-insieme-matrici-spazio-vettoriale} che $\dim M_{mn}(\K) = m \cdot n$. Per calcolare la dimensione di $\mathcal{L}(\V, \W)$, invece, prendiamo la base $\B = \{v_1, \ldots, v_n\}$ di $\V$ e la base $\Bs = \{w_1, \ldots, w_m\}$ di $\W$. 
    Per ogni coppia $(i,j)$, con $i \in \{1, \ldots, m\}$ e $j \in \{1, \ldots, n\}$, per la Proposizione \ref{corollario-applicazioni-lineari-matrici} esiste un'unica applicazione lineare $f_{ij}: \V \longrightarrow \W$ tale che 
    \[
    f_{ij}(v_k)=
    \begin{cases}
     w_i & \text{se } k = j \\
     \undzero & \text{se } k \neq j 
    \end{cases}
    \]
    Per dimostrare che $\dim \mathcal{L}(\V, \W) = n \cdot m$ è sufficiente quindi dimostare che \mbox{$\D = \{f_{ij} \mid i \in \{1, \ldots, m\} \text{ e } j \in \{1, \ldots, n\}\}$} è una sua base.
    \begin{itemize}
        \item $\braket{f_{11}, \ldots, f_{mn}} = \mathcal{L}(\V, \W)$ \\
        Preso $f \in \mathcal{L}(\V, \W)$, $f(v_j) \in \W$ sicuramente, dunque possiamo riscriverlo come $f(v_j) = \sum_{i = 1}^m a_{ij} w_i$ essendo $\Bs = \{w_1, \ldots, w_m\}$ una base di $\W$.
        Poniamo ora \mbox{$g(v_k):= \sum_{i = 1}^m \sum_{j = 1}^n a_{ij} f_{ij}(v_k)$}, e cerchiamo di dimostrare che $f = g$, così da dimostrare che una qualsiasi applicazione lineare $f: \V \longrightarrow \W$ è scrivibile come combinazione lineare di $f_{11}, \ldots, f_{mn}$. \\
        Abbiamo visto prima che $f_{ij}(v_k) = \undzero$ se $j \neq k$, dunque troviamo che \mbox{$g(v_k) = \sum_{i = 1}^m \sum_{j = 1}^n  a_{ij} f_{ij}(v_k) = \sum_{i = 1}^m a_{ik} f_{ik}(v_k) = \sum_{i = 1}^m a_{ik} w_i = f(v_k)$}, questo perché $f_{ik}(v_k) = w_k$. Quindi $\forall k \in \{1, \ldots, n\} \enspace g(v_k) = f(v_k)$. Per il lemma II) visto in Proposizione \ref{lemmi-matrici-e-app-lin}, troviamo che $g = f$.
        \item $f_{11}, \ldots, f_{mn}$ sono linearmente indipendenti \\
        Prendiamo l'equazione $\sum_{i = 1}^m \sum_{j = 1}^n a_{ij} f_{ij} = \undzero$. Se prendiamo ora un $v_k$ generico, troviamo che $\sum_{i = 1}^m \sum_{j = 1}^n a_{ij} f_{ij}(v_k) = \undzero(v_k) = \undzero$. Dato che $f_{ij}(v_k) = 0$ se $j \neq k$ e invece $f_{ij}(v_k) = w_k$ se $j = k$, troviamo che $\sum_{i = 1}^m a_{ik} w_i = \undzero$. Dato che $w_1, \ldots, w_m$ formano la base $\Bs$, sappiamo che questi sono linearmente indipendenti, dunque $a_{ij} = 0 \enspace \forall i \in \{1, \ldots, m\}, \forall j \in \{1, \ldots, m\}$. Dunque anche $f_{11}, \ldots, f_{mn}$ sono linearmente indipendenti. 
    \end{itemize}
    Abbiamo quindi dimostrato che $\dim \mathcal{L}(\V, \W) = n \cdot m$. Non ci manca che dimostrare che $\ker M_\Bs^\B = \{\undzero\}$. \\
    Ipotizziamo che $M_\Bs^\B(f) = \undzero$, ossia la matrice con soli $0$. Significa che $\forall j \in \{1, \ldots, n\} \enspace X_\Bs(f(v_j)) = \undzero$. Dunque, preso $v \in \V$, dato che $\B = \{v_1, \ldots, v_n\}$ è una base di $\V$ possiamo riscrivere $v = a_1 v_1 + \ldots + a_n v_n$, ed essendo $f$ lineare troviamo che $f(v) = a_1 f(v_1) + \ldots + a_n f(v_n) = a_1 \undzero + \ldots + a_n \undzero = \undzero$. Dunque $f$ è per forza l'applicazione nulla, e per questo troviamo che $\ker M_\Bs^\B(f) = \{\undzero\}$.\\
    Possiamo quindi concludere che $M_\Bs^\B$ è biiettiva.   
\end{enumerate}
Essendo $M_\Bs^\B$ sia lineare che biiettiva, troviamo che questa è un isomorfismo. $\blacksquare$

\begin{Proposizione}[]
\label{propos-matrice-associata-composizione}
Presi tre spazi vettoriali $\U, \V$ e $\W$ su $\K$ finitamente generati, prese le applicazioni lineari $g: \U \longrightarrow \V$ e $f: \V \longrightarrow \W$ e prese la base $\A$ di $\U$, la base $\B$ di $\V$ e la base $\Bs$ di $\W$, \mbox{$\bm{M_\Bs^\A(f \circ g) = M_\Bs^\B(f) \cdot M_\B^\A(g)}$}.
\end{Proposizione}

\Dimostrazione \nopagebreak

Preso $u \in \U$, $X_\Bs(f \circ g(u)) = X_\Bs(f(g(u))) = M_\Bs^\B(f) \cdot X_\B(g(u))$ per la Proposizione \ref{propos-matrice-associata-prodotto}. Sempre per questa proposizione troviamo che $X_\B(g(u)) = M_\B^\A(g) \cdot X_\A(u)$. Dunque troviamo che $X_\Bs(f \circ g(u)) = M_\Bs^\B(f) \cdot M_\B^\A(g) \cdot X_\A(u)$. Per la stessa proposizione già vista sappiamo però che $X_\Bs(f \circ g(u)) = M_\Bs^\A(f \circ g) \cdot X_\A(u)$, e da questo troviamo che $X_\Bs(f \circ g(u)) = M_\Bs^\A(f \circ g) \cdot X_\A(u) = M_\Bs^\B(f) \cdot M_\B^A(g) \cdot X_\A(u)$. Concludiamo quindi che $M_\Bs^\A(f \circ g) = M_\Bs^\B(f) \cdot M_\B^\A(g)$. $\blacksquare$

\Esempio \nopagebreak

Prese le applicazoni $\bm{p_\alpha: V(\E^2)} \longrightarrow \bm{V(\E^2)}$ e $\bm{p_\beta: V(\E^2)} \longrightarrow \bm{V(\E^2)}$, dove $p_\alpha \circ p_\beta = p_{\alpha + \beta}$, trovare una matrice associata $\bm{M_\Bs^\A(p_\alpha \circ p_\beta)}$.  \nopagebreak

Dato che in questo caso dominio e codominio di entrambe le applicazioni coincidono, possiamo scegliere una stessa base $\A = \B = \Bs = \{\underline{i}, \underline{j}\}$. Dalla Proposizione \ref{propos-matrice-associata-composizione} sappiamo che $M_\Bs^\A(p_\alpha \circ p_\beta) = M_\Bs^\B(p_\alpha) \cdot M_\B^\A(p_\beta)$. Dall'esempio precedente sappiamo già che
\[ M_\Bs^\B(p_\alpha) = 
\begin{bmatrix}
 \cos(\alpha) & -\sin(\alpha) \\
 \sin(\alpha) & \cos(\alpha)
\end{bmatrix}
\]
e che
\[ M_\B^\A(p_\beta) = 
\begin{bmatrix}
 \cos(\beta) & -\sin(\beta) \\
 \sin(\beta) & \cos(\beta)
\end{bmatrix}
\]
Dunque troviamo che 
\[\bm{M_\Bs^\A(p_\alpha \circ p_\beta)} = 
\begin{bmatrixb}
\cos(\alpha) \cdot \cos(\beta) - \sin(\alpha) \sin(\beta) & -\cos(\alpha) \sin(\beta) - \sin(\alpha) \cos(\beta)  \\
\sin(\alpha) \cos(\beta) + \cos(\alpha) \sin(\beta) & - \sin(\alpha)\sin(\beta) + \cos(\alpha) \cos(\beta) 
\end{bmatrixb}
\]

Dato che però $p_\alpha \circ p_\beta = p_{\alpha + \beta}$, troviamo anche che 
\[ M_\Bs^\A(p_\alpha \circ p_\beta) = M_\Bs^\A(p_{\alpha + \beta}) = 
\begin{bmatrix}
 \cos(\alpha + \beta) & -\sin(\alpha + \beta)  \\
 \sin(\alpha + \beta) & \cos(\alpha + \beta)
\end{bmatrix}
\]

Questi risultati ci permettono quindi di definire formalmente le formule di addizione del seno e del coseno, ossia
\begin{align*}
\sin(\alpha + \beta) &= \sin(\alpha) \cos(\beta) + \cos(\alpha) \sin(\beta) \\
\cos(\alpha + \beta) &= \cos(\alpha)  \cos(\beta) - \sin(\alpha)\sin(\beta) 
\end{align*}

\section{Trovare una base del nucleo e dell'immagine di un'applicazione lineare}

\subsection{Base dell'immagine di un'applicazione lineare}

\begin{Definizione}[Matrice a scala per colonne]
\label{def-matrice-scala-per-colonne}
Una matrice \textbf{a scala per colonne} è una matrice $A \in M_{mn}(\K)$ che ha il primo elemento non nullo di ogni colonna (chiamato \textit{pivot}) almeno una riga più in basso rispetto alla colonna precedente. Se si arriva ad una colonna nulla, tutte le successive colonne anche dovranno essere nulle. \\
Una matrice a scala per colonne è una matrice del tipo
\[
\begin{bmatrix}
p_1 & 0 & 0 & 0 & 0  \\
* & p_2 & 0 & 0 & 0 \\
* & * & 0 & 0 & 0 \\
* & * & p_3 & 0 & 0 \\
* & * & * & 0 & 0 
\end{bmatrix}
\]
\end{Definizione}

\begin{Definizione}[Rango]
\label{def-rango}
Il \textbf{rango} $\bm{rg(f)}$ di un'applicazione lineare $f:\V \longrightarrow \W$, dove $\V$ e $\W$ sono spazi vettoriali su $\K$ finitamente generati, è la \textbf{dimensione della sua immagine}.
\end{Definizione}

\NB Per comodità scriveremo quasi sempre $Im(A)$ al posto di $Im(L_A)$, oppure anche $\ker (A) = \ker (L_A)$. È quasi scontanto che si parli ovviamente di $L_A$ e non della matrice stessa $A$ perché questi concetti hanno senso se riferiti ad un'applicazione lineare (basta rivedersi la definizione di rango o nucleo per notare subito che sono riferiti ad un'applicazione lineare).  

\begin{Proposizione}[]
\label{propos-immagine-combinazione-lineare-colonne}
Presa una matrice $A \in M_{mn}(\K)$, $\bm{Im(A) = \braket{A_1, \ldots, A_n}}$.
\end{Proposizione}

\Dimostrazione \nopagebreak

Per dimostrare questa proposizione, dimostriamo prima che $Im(A) \subset \braket{A_1, \ldots, A_n}$, e poi il viceversa. Questa doppia dimostrazione implica per forza di cose che $Im(A) = \braket{A_1, \ldots, A_n}$ (infatti, se è vero che $A$ sia un sottoinsieme di $B$ e anche $B$ è un sottoinsieme di $A$, allora l'unica soluzione plausibile è che appunto $A$ sia uguale a $B$).
\begin{itemize}
    \item $Im(A) \subset \braket{A_1, \ldots, A_n}$ \\
    Prendiamo $\Y \in Im(A)$. Sappiamo che esiste quindi $\X \in \K^n$ tale che $L_A(\X) = A \cdot \X = \Y$. Essendo $S = \{e_1, \ldots, e_n\}$ la base canonica di $\K^n$, troviamo che $\X = e_1 x_1 + \ldots + e_n x_n$, dunque $\Y = A \cdot (e_1 x_1 + \ldots + e_n x_n) = x_1 \cdot (A e_1) + \ldots + x_n \cdot (A e_n) = x_1 A_1 + \ldots x_n A_n$. Questo significa che tutti gli elementi in $Im(A)$ sono scrivibili come combinazione lineare delle colonne di $A$, ossia $Im(A) \subset \braket{A_1, \ldots, A_n}$.
    \item $Im(A) \subset \braket{A_1, \ldots, A_n}$ \\
    Prenidamo $\Y \in \braket{A_1, \ldots, A_n}$. Sappiamo che esistono qunidi $\lambda_1, \ldots, \lambda_n \in \K$ tali che $\Y = \lambda_1 A_1 + \ldots + \lambda_n A_n = \lambda_1 \cdot (A e_1) + \ldots + \lambda_n \cdot A(e_n) = A (\lambda_1 e_1 + \ldots + \lambda_n e_n)$. Dato che $S = \{e_1, \ldots, e_n\}$ è la base canonica di $\K^n$, troviamo che $\lambda_1 e_1 + \ldots + \lambda_n e_n \in \K^n$. Definiamo quindi $\X = \lambda_1 e_1 + \ldots + \lambda_n e_n$. Troviamo quindi che $\Y = A \cdot \X = L_A(\X)$, ossia $\Y \in Im(A)$. 
\end{itemize}
Possiamo quindi concludere che $Im(A) = \braket{A_1, \ldots, A_n}$. $\blacksquare$

\begin{Proposizione}[]
\label{propos-base-immagine-lista-colonna-non-nulle}
Presa una matrice $A \in M_{mn}(\K)$ a \textit{scala per colonne}, una base di $Im(A)$ è $\bm{\B = \{A_i \mid A_i \neq 0\}}$, ossia la lista delle colonne non nulle di $A$. \\
Troviamo quindi che $\bm{rg(A) = \{\textbf{numero di colonne non nulle}\}}$.
\end{Proposizione}

\Dimostrazione \nopagebreak

Prima di passare alla dimostrazione della proposizione vera e propria, dimostriamo che questa proposizione non è vera se la matrice non è a scala per colonne. \\
Prendiamo ad esempio la matrice 
\[ A =
\begin{bmatrix}
1 & 2 \\
1 & 2
\end{bmatrix}
\]
È evidente che questa non sia a scala per colonne. In questo caso $\B = \{A_1, A_2\} = \{(1,2), (1,2)\}$ è la lista delle colonne non nulle di $A$, ma questa non può mai essere una base, dato che le due colonne sono uguali e quindi linearmente dipendenti. 

Passiamo ora alla dimostrazione vera e propria. \\
Prendiamo una matrice a scala per colonne
\[A = 
\begin{bmatrix}
p_1 & 0 & \dots & 0 & 0 & 0 \\ * & p_2 & \dots & 0 & 0 & 0\\ \vdots & \vdots & \ddots & \vdots &  \vdots &  \vdots \\ * & * & \dots & p_k & \dots & 0
\end{bmatrix}
\]
Da questa, prendiamo la matrice che ha solo le colonne non nulle
\[A' = 
\begin{bmatrix}
p_1 & 0 & \dots & 0  \\ * & p_2 & \dots & 0 \\ \vdots & \vdots & \ddots & \vdots \\ * & * & \dots & p_k  
\end{bmatrix}
\]
Sappiamo per la Proposizione \ref{propos-immagine-combinazione-lineare-colonne} che $Im(A) = \braket{A_1, \ldots, A_n}$, e dato che i vettori nulli non influiscono nella combinazione lineare, possiamo dire che $Im(A) = \braket{A_1, \ldots, A_k}$. \\
Per dimostrare che questi vettori sono linearmente indipendenti, dobbiamo dimostrare che l'unica soluzione all'equazione $\lambda_1 A_1 + \ldots + \lambda_k A_k = \undzero$ sia $\lambda_1 = \ldots = \lambda_k = 0$. Riscriviamoci per esteso questa equazione:
\[
\sum_{j = 1}^k \lambda_j A_j = 
\begin{bmatrix}
    \lambda_1 p_1 \\
    \lambda_1 * \\
    \vdots \\
    \lambda_1 * \\
\end{bmatrix}
+ \begin{bmatrix}
    0 \\
    \lambda_2 p_2 \\
    \vdots \\
    \lambda_2 *
\end{bmatrix}
+ \dots + 
\begin{bmatrix}
    0 \\
    0 \\
    \vdots \\
    \lambda_k p_k
\end{bmatrix}
= 
\begin{bmatrix}
    \lambda_1 p_1 \\
    \lambda_1 * + \lambda_2 p_2 \\
    \vdots \\
    \sum_{j = 1}^k \lambda_j p_j
\end{bmatrix}
=
\begin{bmatrix}
0  \\
0 \\
\vdots \\
0 
\end{bmatrix}
\] 

Alla prima riga troviamo l'equazione $\lambda_1 p_1 = 0$. Essendo $p_1 \neq 0$ per definizione, troviamo che $\lambda_1 = 0$. Alla seconda riga, dato che $\lambda_1 = 0$, troviamo che anche $\lambda_2 = 0$. Continuando così fino ad arrivare a $\lambda_k$, troviamo che tutti i coefficienti $\lambda$ sono pari a $0$. Dunque $A_1, \ldots, A_k$ sono linearmente indipendenti, e quindi possiamo concludere che $\B = \{A_1, \ldots, A_k\}$ è una base di $Im(A)$. $\blacksquare$ 

\begin{Proprieta}[Operazioni elementari sulle colonne di una matrice]
\label{prop-operazioni-elementari-colonne}
Presa una matrice $A = [A_1, \ldots, A_n] \in M_{mn}(\K)$ 
\begin{enumerate}
\item \textit{Scambio di due colonne}: presi $1 \leq i < j \leq n$, \\
$\bm{A = [A_1, \ldots, A_n] \leadsto B = [A_1, \ldots, A_{i-1}, A_j, A_{i+1}, \ldots, A_{j-1}, A_i, A_{j+1}, \ldots, A_n]}$
\item \textit{Aggiungere a una colonna il multiplo di un'altra}: presi $j,k \in \{1, \ldots, n\}$ tali che $j \neq k$,\\
$\bm{A = [A_1, \ldots, A_n] \leadsto B = [A_1, \ldots, A_j + \lambda A_k, \ldots, A_n]}$, con $\lambda \in \K$
\item \textit{Moltiplicare una colonna per uno scalare non nullo}: preso $j \in \{1, \ldots, n\}$, \\
$\bm{A = [A_1, \ldots, A_n] \leadsto B = [A_1, \ldots, \lambda A_j, \ldots, A_n]}$, con $\lambda \in \K$ tale che $\lambda \neq 0$
\end{enumerate}
\end{Proprieta}

\begin{Proposizione}[]
\label{propos-arrivare-matrice-a-scala-per-colonne}
Presa una matrice $A \in M_{mn}(\K)$, è sempre possibile, attraverso una serie di operazioni elementari sulle colonne viste in Proprietà \ref{prop-operazioni-elementari-colonne}, passare da $A$ a una matrice a scala per colonne $S$ \textbf{senza modificare l'insieme generato dalla matrice}, ossia $\bm{Im (A) = Im (S)}$.
\end{Proposizione}

\Dimostrazione \nopagebreak

Concentriamoci prima a passare da una matrice $A$ a una matrice a scala per colonne $S$, per poi dimostrare che $Im(A) = Im(S)$.

Presa $A \in M_{mn}(\K)$, ci sono due possibilità:
\begin{enumerate}[label=\arabic*°]
    \item $A^1 = \undzero$, e dunque passiamo direttamente alla seconda riga lasciando perdere la prima
    \item $A^1 \neq \undzero$
\end{enumerate} 
Se $A^1 \neq \undzero$, allora sappiamo che almeno una colonna è non nulla. Prendiamo questa colonna e la scambiamo con la prima colonna, eseguendo un'operazione elementare 1). Ora, per la riga 1 dobbiamo eliminare tutti gli elementi su questa linea di tutte le colonne che non siano la prima. \\
Per ogni colonna $i$ aggiungiamo un certo $\lambda A_1$ tale che $\lambda_1 a_{11} = - a_{1i}$. In questo modo, tutta la riga $1$ sarà uguale a $0$ trannè che per $a_{11}$. \\
Ora se guardaimo la sotto matrice $(m - 1) \times (n - 1)$ togliendo la prima riga e la prima colonna, ripetendo lo stesso passaggio troviamo che anche la seconda riga è tutta nulla dalla terza entrata in poi. Continuando così con tutte le rige, troviamo una nuova matrice $S$ a scala per colonne.

Dobbiamo ora verificare che $Im(A) = Im(S)$. Se vediamo le colonne di $A$ come dei vettori $v_1, \ldots, v_n$, possiamo riscriverci tutte e 3 le operazioni elementari utilizzando questi vettori. Quindi
\begin{enumerate}
    \item Presi $1 \leq i \leq j \leq n$, $\{v_1, \ldots, v_n\} \leadsto \{v_1, \ldots, v_{i-1}, v_j, v_{i+1}, \ldots, v_{j-1}, v_i, \ldots, v_{j+1}, \ldots, v_n\}$
    \item Presi $j,k \in \{1, \ldots, n\}$ tali che $j \neq k$, $\{v_1, \ldots, v_n\} \leadsto \{v_1, \ldots, v_j + \lambda v_k, \ldots, v_n\}$, con $\lambda \in \K$
    \item Preso $j \in \{v_1, \ldots, v_n\} \leadsto \{v_1, \ldots, \lambda v_j, \ldots, v_n\}$
\end{enumerate}

Dobbiamo dimostrare che, in tutti e tre i casi, lo SPAN generato dai vettori a sinistra è identico allo SPAN generato dai vettori a destra (questo perché per la Proposizione \ref{propos-immagine-combinazione-lineare-colonne}).
\begin{enumerate}
    \item $\braket{v_1, \ldots, v_n} = \braket{v_1, \ldots, v_{i-1}, v_j, v_{i+1}, \ldots, v_{j-1}, v_i, v_{j+1}, \ldots, v_n}$ \\ 
    Questa è la più banale, dato che modificare l'ordine dei vettori in una lista dei vettori, lo SPAN di essi non cambia mai.
    \item  $\braket{v_1, \ldots, v_n} = \braket{v_1, \ldots, v_j + \lambda v_k, \ldots, v_n}$ \\
    Dimostriamo prima che $\braket{v_1, \ldots, v_n} \subset \braket{v_1, \ldots, v_j + \lambda v_k, \ldots, v_n}$ e poi dimostriamo che $\braket{v_1, \ldots, v_n} \supset \braket{v_1, \ldots, \lambda v_j, \ldots, v_n}$ per dimostrare che questi sono uguali.
    \begin{itemize}
        \item $\braket{v_1, \ldots, v_n} \subset \braket{v_1, \ldots, v_j + \lambda v_k, \ldots, v_n}$ \\
        Tutti i generatori dello SPAN a sinistra si trovano tra i generatori dello SPAN di destra fatta eccezione per $v_j$, ma è facile dimostrare che anche lui appartiene allo SPAN. Se infatti nella combinazione lineare $\lambda_1 v_1 + \ldots + \lambda_j (v_j + \lambda v_k) + \lambda_n v_n$ poniamo \mbox{$\lambda_i = 0 \enspace \forall i \in \{1, \ldots, n\}$} tali che $i \neq j$ e $i \neq k$, e poniamo poi $\lambda_j = 1$ e $\lambda_k = -\lambda$, troviamo proprio l'ultimo generatore che ci mancava $v_j$. 
        \item $\braket{v_1, \ldots, v_n} \supset \braket{v_1, \ldots, v_j + \lambda v_k, \ldots, v_n}$ \\
        Anche qui notiamo che tutti i generatori dello SPAN a destra si trovano tra i generatori dello SPAN di sinistra fatta eccezione per $v_j + \lambda v_k$, ma è facile dimostrare che anche lui appartiene allo SPAN. Se infatti nella combinazione lineare $\lambda_1 v_1 + \ldots + \lambda_n v_n$ poniamo \mbox{$\lambda_i = 0 \enspace \forall i \in \{1, \ldots, n\}$} tali che $i \neq j$ e $i \neq k$, e poniamo poi $\lambda_j = 1$ e $\lambda_k = \lambda$, troviamo proprio l'ultimo generatore che ci mancava $v_j + \lambda v_k$. 
    \end{itemize}
    Possiamo quindi concludere che $\braket{v_1, \ldots, v_n} = \braket{v_1, \ldots, v_j + \lambda v_k, \ldots, v_n}$.
    \item $\braket{v_1, \ldots, v_n} = \braket{v_1, \ldots, \lambda v_j, \ldots, v_n}$ \\
    Anche qui cerchiamo di dimostrare sia che $\braket{v_1, \ldots,v_n} \subset \braket{v_1, \ldots, \lambda v_j, \ldots, v_n}$ e poi dimostriamo che $\braket{v_1, \ldots, v_n} \supset \braket{v_1, \ldots, \lambda v_j, \ldots, v_n}$. 
    \begin{itemize}
        \item $\braket{v_1, \ldots, v_n} \subset \braket{v_1, \ldots, \lambda v_j, \ldots, v_n}$ \\
        Tutti i generatori dello SPAN a sinistra si trovano tra i generatori dello SPAN di destra fatta eccezione per $v_j$, ma è facile dimostrare che anche lui appartiene allo SPAN. Se infatti nella combinazione lineare $\lambda_1 v_1 + \ldots + \lambda_j (\lambda v_j) + \ldots + \lambda_n v_n$ poiniamo \mbox{$\lambda_i = 0 \enspace \forall i \in \{1, \ldots, n\}$} tali che $i \neq j$, e poniamo $\lambda_j = \frac{1}{\lambda}$, troviamo proprio l'ultimo generatore che ci mancava $v_j$. 
        \item $\braket{v_1, \ldots, v_n} \supset \braket{v_1, \ldots, \lambda v_j, \ldots, v_n}$ \\
        Anche qui notiamo che tutti i generatori dello SPAN a destra si trovano tra i generatori dello SPAN di sinistra fatta eccezione per $\lambda v_j$, ma è facile dimostrare che anche lui appartiene allo SPAN. Se infatti nella combinazione lineare $\lambda_1 v_1 + \ldots + \lambda_n v_n$ poiniamo \mbox{$\lambda_i = 0 \enspace \forall i \in \{1, \ldots, n\}$} tali che $i \neq j$, e poniamo $\lambda_j = \lambda$, troviamo proprio l'ultimo generatore che ci mancava $\lambda_j v_j$. 
    \end{itemize}
    Possiamo quindi concludere che $\braket{v_1, \ldots, v_n} = \braket{v_1, \ldots, \lambda v_j, \ldots, v_n}$.
\end{enumerate}

Dato che tutte e tre le operazioni elementari sulle colonne non modificano lo SPAN generato dai vettori, anche lo SPAN generato dalle colonne della matrice non si modifica, dunque $Im(A) = Im(S)$. $\blacksquare$


\Esempio \nopagebreak

Calcolare il \textbf{rango} della matrice \nopagebreak

\[ A = 
\begin{bmatrixb}
3 & 2 & 1 & 5  \\
1 & 0 & 2 & 11 \\
2 & -3 & -1 & -1 \\
6 & -1 & 2 & 15 
\end{bmatrixb}
\]

Per trovare il rango di $A$, che è uguale alla dimensione di $Im(A)$, ci basta trasformare $A$ in una matrice a scala per colonne $S$. Infatti, dato che $Im(A) = Im(S)$ per la Proposizione \ref{propos-arrivare-matrice-a-scala-per-colonne}, e dato che $rg(S)$ è pari al numero di colonne non nulle per la Proposizione \ref{propos-base-immagine-lista-colonna-non-nulle}, trovare il rango di $S$ è facilissimo, e dato che $rg(S) = \dim Im(S) = \dim Im(A) = rg(A)$, troviamo con questo procedimento anche il rango di $A$.

Seguendo il procedimento base qui dovremmo cercare, attraverso operazioni elementari sulle colonne di $A$ di ricavarci la riga $[3,0,0,0]$, ma essendoci un $1$ nella prima riga, ci conviene invece scambiare la colonna con l'$1$ con la prima colonna, così da ricavarsi la riga $[1,0,0,0]$, che è più semplice da ricavare perché banalmente non lavoriamo con le frazioni.

Con un'operazione elementare 1) facciamo quindi
\[ A = 
\begin{bmatrix}
    3 & 2 & 1 & 5 \\
    1 & 0 & 2 & 11 \\
    2 & -3 & -1 & -1 \\
    6 & .1 & 2 & 15
\end{bmatrix}
\leadsto
B = \begin{bmatrix}
    1 & 2 & 3 & 5 \\
    2 & 0 & 1 & 11 \\
    -1 & -3 & 2 & -1 \\
    2 & -1 & 6 & 15
\end{bmatrix}
\]
Cerchiamo ora di ricavare la riga $[1,0,0,0]$. Troviamo un'altra matrice $C$ così definita:
\begin{align*}
C_1 &= B_1  \\
C_2 &= B_2 + (-2 B_1) \\
C_3 &= B_3 + (-3 B_1) \\
C_4 &= B_4 + (-5 B_1)
\end{align*}

Troviamo quindi che 
\[
C = \begin{bmatrix}
    1 & 0 & 0 & 0 \\
    2 & -4 & -5 & 1 \\
    -1 & -1 & 5 & 4 \\
    2 & -5 & 0 & 5
\end{bmatrix}
\]

Ripetendo questo procedimento per tutte le righe, arriviamo alla matrice 
\[ S =
\begin{bmatrix}
    1 & 0 & 0 & 0 \\
    2 & 1 & 0 & 0 \\
    -1 & 4 & 25 & 0 \\
    2 & 5 & 25 & 0 
\end{bmatrix}
\]
Questa è una matrice a scala per colonne, e dato che ha 3 colonne non nulle, troviamo che $rg(S) = \bm{rg(A) = 3}$.

\subsection{Base del nucleo di un'applicazione lineare}

\begin{Definizione}[Trasposta]
\label{def-trasp}
Presa una matrice $A \in M_{mn}(\K)$, la \textbf{trasposta} di $A$ è la matrice $\bm{A^t} \in M_{nm}(\K)$ tale che \mbox{$\bm{a_{ij}^t = a_{ji}} \enspace \forall i \in \{1, \ldots, n\}, j \in {1, \ldots, m}$}. \\
In poche parole, $A^t$ è la matrice che ha come righe le colonne di $A$. 
\end{Definizione}

\begin{Definizione}[Simmetrica]
\label{def-simmetrica}
Presa una matrice $A \in M_{nn}(\K)$, questa è \textbf{simmetrica} se e solo se $\bm{A = A^t}$. \\
In poche parole, la riga i-esima di $A$ è uguale alla colonna i-esima di $A$.  
\end{Definizione}

Per esempio la matrice 
\[ A =
\begin{bmatrix}
    2 & 3 & 4 \\
    5 & 6 & 7
\end{bmatrix}
\]
non è simmetrica (essendo una matrice rettangolare questo è evidente), e la sua trasposta è
\[ A^t =
\begin{bmatrix}
    2 & 5 \\ 
    3 & 6 \\
    4 & 7
\end{bmatrix}\]

\begin{Proposizione}[Trasposta del prodotto tra matrici]
\label{propos-trasposta-prodotto}
Prese $n$ matrici $A_1, \ldots, A_n$ il cui prodotto $A_1 \cdot \ldots \cdot A_n$ sia sensato, la trasposta di tale prodotto è pari al prodotto delle trasposte prese in ordine inverso, ossia \mbox{$\bm{(A_1 \cdot \ldots \cdot A_n)^t = A_n^t \cdot \ldots A_1^t}$}. 
\end{Proposizione}

\Dimostrazione \nopagebreak

Dimostriamo questa proposizione per induzione.
\begin{itemize}
    \item Caso base: $n = 2$ \\
    In questo caso abbiamo quindi due matrici $A, B$. Sappiamo che $A = (a_{ij}) \in M_{mp}$ e che $B = (b_{ij}) \in M_{pn}$. Sappiamo anche che $(AB)_{ij} = \sum_{k = 1}^p a_{ik} b_{kj}$ e che \mbox{$(A_B)^t_{ij} = (AB)_{ji} = \sum_{k = 1}^p a_{jk} b_{ki} = \sum_{k = 1}^p b_{ki} a_{jk}$}. Dato che $a_{jk} = a_{kj}^t$ e $b_{ki} = b^t_{ik}$, troviamo che $(AB)_{ij}^t = \sum_{k = 1}^p b_{ik}^t a_{kj}^t$, e quindi $(AB)^t = B^t A^t$. $\blacksquare$
    \item Ipotesi induttiva \\
    Assumiamo la proposizione vera per un prodotto tra $n-1$ matrici.
    \item Passo induttivo \\
    Prendiamo questa volta $n$ matrici $A_1, \ldots, A_n \in M_{nn}(\K)$. Sappiamo ovviamente che \mbox{$(A_1 \cdot \ldots \cdot A_{n-1} \cdot A_n) = ((A_1 \cdot \ldots \cdot A_{n-1}) \cdot A_n)^t$}, e dato che $A_1 \cdot \ldots \cdot A_{n-1}$ è sempre una matrice, possiamo vedere questo prodotto come un prodotto tra due matrici, e avendo già dimostrato il caso base possiamo dire che la trasposta del prodotto tra queste due matrici è il prodotto delle trasposte in ordine ivnerso, ossia \mbox{$((A_1 \cdot \ldots \cdot A_{n-1}) \cdot A_n)^t = A_n^t \cdot (A_1 \cdot \ldots \cdot A_{n-1})^t$}. Per ipotesi induttiva, sappiamo che \mbox{$(A_1 \cdot \ldots \cdot A_{n-1})^t =  A_{n-1}^t \cdot \ldots \cdot A_1$}, dunque troviamo che \mbox{$(A_1 \cdot \ldots A_n)^t = A_n^t \cdot \ldots \cdot A_1^t$.} $\blacksquare$
\end{itemize}

\begin{Definizione}[Matrice a scala per righe]
\label{def-matr-scala-per-righe}
Una matrice \textbf{a scala per righe} è una matrice la cui trasposta è una matrice a scala per colonne.
\end{Definizione}

\begin{Proposizione}[]
\label{propos-dimensione-nucleo-(n-r)}
Presa una matrice $S \in M_{mn}(\K)$ a scala per righe, con pivot sulle colonne $j_1, \ldots, j_r$, dove $1 \leq j_1 \leq \ldots \leq j_r \leq n$, l'applicazione
\begin{align*}
\bm{\phi S: \ker S} &\longrightarrow \bm{\K^{n-r}} \\
\bm{(x_1, \ldots, x_n)} &\longmapsto \bm{(x_1, \ldots, x_{j_1 - 1}, x_{j_1 + 1}, \ldots, x_{j_r - 1}, x_{j_r + 1}, \ldots, x_n)}
\end{align*}
che prende un vettore nel nucleo di $S$ e ne ritorna un'altro con gli stessi vettori tolti quelli sugli indici $j_1, \ldots, j_r$, è un \textbf{isomorfismo}. 
\end{Proposizione}

\Dimostrazione \nopagebreak

Per dimostrare che $\phi S$ è un isomorfismo, per l'Osservazione \ref{oss-def-isomorfismo} ci basta dimostrare che questa è lineare e biiettiva. È evidente che questa sia lineare, dunque ci concentreremo a dimostrare solo che questa è biiettiva, ossia che è iniettiva e suriettiva.
\begin{itemize}
    \item $\phi S$ è iniettiva \\
    Preso $\X \in \ker S$, supponiamo che $\phi S (\X) = \undzero$. Dato che $\X \in \ker S$, sappiamo che $\X = (x_1, \ldots, x_n)$ è soluzione dell'equazione $S \cdot \X = \undzero$, ossia è soluzione del sistema
    \[
    \begin{cases}
    p_1 x_{j_1} + s_{1 j_1} x_{j_1+1} + \ldots + s_{1n} x_n &= 0 \\
    p_2 x_{j_2} + s_{2 j_2} x_{j_2 + 1} + \ldots + s_{2n} x_n &= 0 \\
    \vdots \\
    p_r x_{j_r} + s_{r j_r + 1} x_{j_r + 1} \ldots + s_{rn} x_n &= 0 \\
    0 = 0 \\
    \vdots \\
    0 = 0 
    \end{cases}
    \]
    Dato che $\phi S (\X) = \undzero$, tutte le variabili $x_i$ tali che $i \neq j_k \enspace \forall k \in \{1, \ldots, r\}$ sono uguali a $0$. Partendo dal basso troviamo quindi che $p_r x_{jr} = 0$, e dato che per definizione $p_r \neq 0$ essendo il pivot, troviamo che $x_{jr} = 0$. Risalendo tutto il sistema, troviamo che $x_{j_k} = 0 \enspace \forall k \in \{1, \ldots, r\}$, ossia troviamo $\X = \undzero$. Dunque $\ker \phi S = \{\undzero\}$. Per la Proposizione \ref{propos-applicazione-lineare-iniettiva-nucleo} troviamo quindi che $\phi S$ è iniettiva.
    \item $\phi S$ è suriettiva \\
    Prendiamo un vettore $\Y = (y_1, \ldots, y_{n-r}) \in \K^{n-r}$ e cerchiamo di dimostrare che esiste un $\X \in \ker S$ tale che $\phi S (\X) = \Y$. Vogliamo quindi dimostrare che $\Y$ è $\X$ ma senza tutte le entrate in corrispondenza degli indici dei pivot. \\
    Riprendiamo il sistema già visto prima:
    \[
    \begin{cases}
    p_1 x_{j_1} + s_{1 j_1} x_{j_1+1} + \ldots + s_{1n} x_n &= 0 \\
    p_2 x_{j_2} + s_{2 j_2} x_{j_2 + 1} + \ldots + s_{2n} x_n &= 0 \\
    \vdots \\
    p_r x_{j_r} + s_{r j_r + 1} x_{j_r + 1} \ldots + s_{rn} x_n &= 0 \\
    \end{cases}
    \]
    Risolvendo dal basso verso l'altro, troviamo che $p_r x_{j_r} = -(\text{somma di costanti note})$, ossia troviamo $x_{j_r} = k$ dove $k$ è un termine noto. Risalendo tutto il sistema troviamo tutti i valori di $x_{j_k}$, ma tutti i restanti valori $x_i$ tali che $i \neq j_k \enspace \forall k \in \{1, \ldots, r\}$ non sono in alcun modo vincolati dal sistema. Possiamo quindi porre questi valori pari a $y_i$, arrivando quindi a dire che $\phi S(\X) = \Y$. Dunque $\phi S$ è suriettiva. 
\end{itemize}
Essendo quindi $\phi S$ lineare, iniettiva e suriettiva, e quindi biettiva, possiamo concludere che è un isomorfismo. $\blacksquare$

\begin{Proposizione}[Corollario]
\label{corol-propos-dimensione-nucleo}
Per la Proposizione \ref{propos-dimensione-nucleo-(n-r)}, presa una matrice $S \in M_{mn}(\K)$ a scala per righe, troviamo che $\bm{\dim (\ker S) = n -r}$, dove $r$ è il numero di pivot di $S$.
\end{Proposizione}

Questo corollario è evidente grazie all'Osservazione \ref{oss-dim-isomorfismi}.

\begin{Osservazione}[]
\label{oss-soluzione-non-banale-sistema}
Presa una matrice $A \in M_{mn}(\K)$, $\bm{n > m} \implies $ esiste $\X \neq \undzero$ tale che $\bm{A \cdot \X = \undzero}$. \\
In poche parole, per un sistema lineare omogeneo con più variabili che equazioni \textbf{esiste una soluzione non banale}. È un'affermazione che avevamo dimostrato con un esempio precedente, ma che ora cerchiamo di dimostrare in modo più rigoroso.
\end{Osservazione}

\Dimostrazione \nopagebreak

Se $n > m$, allora $n > r$, dato che $m \geq r$. Dunque $n-r \geq n- m \geq 1$, ossia esiste una variabile libera, che si trova in una colonna senza pivot. Questo singifica che quella variabile non deve essere per forza $0$, e che quindi $\X \neq 0$. $\blacksquare$

\begin{Proposizione}[]
\label{propos-arrivare-matrice-a-scala-per-righe}
Presa una matrice $A \in M_{mn}(\K)$, è sempre possibile, attraverso una serie di operazioni elementari sulle righe (che sono identiche a quelle sulle colonne viste in Proprietà \ref{prop-operazioni-elementari-colonne} ma applicate sulle righe), passare da $A$ a una matrice a scala per righe $S$ \textbf{senza modificare l'insieme generato dalle righe della matrice}.
\end{Proposizione}
La dimostrazione è pressoché identica alla dimostrazione per la Proposizione \ref{propos-arrivare-matrice-a-scala-per-colonne}, perciò è omessa.


\begin{Proposizione}[Corollario]
\label{corol-ker(A)=ker(S)}
Prese due matrici $A, S$, dove $S$ è una matrice a scala per righe, $\bm{A \leadsto S} \implies \bm{\ker A = \ker S}$.
\end{Proposizione}

\Esempio \nopagebreak

Trovare una \textbf{base} di $\bm{\ker A}$ dove \nopagebreak
\[
A = 
\begin{bmatrixb}
    3 & 1 & 2 & 6 \\
    2 & 0 & -3 & -1 \\
    1 & 2 & -1 & 2 \\
    5 & 11 & -1 & 15
\end{bmatrixb}
\]

Con una serie di operazioni elementari sulle righe di $A$ troviamo la matrice a scala per righe
\[ 
S = \begin{bmatrix}
    1 & 2 & -1 & 2 \\
    0 & 1 & 4 & 5 \\
    0 & 0 & 1 & 1 \\
    0 & 0 & 0 & 0
\end{bmatrix}
\] 
Per il Corollario \ref{corol-ker(A)=ker(S)} sappiamo che $\ker A = \ker S$. Per trovare una base di $A$ ci basta quindi trovare una base di $S$. Inoltre, per la Proposizione \ref{propos-dimensione-nucleo-(n-r)} troviamo che $\dim \ker A = 1$, dato che $S$ è una matrice con $3$ pivot e $4$ colonne. \\
Dato che stiamo cercando una base di $\ker S$, cerchiamo di risolvere l'equazione $S \cdot \X = \undzero$, che si traduce nel sistema
\[
\begin{cases}
x_1 + 2x_2 - x_3 + 2x_4 &= 0 \\
x_2 + 4x_3 +  5x_4 &= 0 \\
x_3 + x_4 &= 0 
\end{cases}
\implies
\begin{cases}
    x_3 = - x_4 \\
    x_2 = -x_4 \\
    x_1 = -x_4
\end{cases}
\]

Troviamo quindi che l'unica variabile libera è $x_4$, e quindi \mbox{$\ker S = \{(-x_4, -x_4, -x_4, x_4) \mid x_4 \in \K\} = \{t (1,1,1,-1)\mid t \in \K\}$}, e diventa quindi evidente che una base di $\ker S$, e quindi di $\ker A$, è $\bm{\B = \{(1,1,1,-1)\}}$ (confermiamo quindi che $\dim \ker A = 1$). 

\begin{Proposizione}[]
\label{propos-rango-uguale-trasposta}
Presa una matrice $A \in M_{mn}(\K)$, $\bm{rg(A) = rg(A^t)}$. \\
In poche parole, il numero di colonne linearmente indipendenti di $A$ è uguale al numero di righe linearmente indipendenti di $A$.
\end{Proposizione}

\Dimostrazione \nopagebreak

Presa l'applicazione $L_A: \K^n \longrightarrow \K^m$, sappiamo per la Proposizione \ref{propos-dim-dominio-dim-ker-dim-im} che $\dim \K^n = n = \dim \ker A + \dim Im(A) = \dim \ker A + rg(A)$, ossia $rg(A) = n - \dim \ker A$. Ora da $A$ riconduciamoci ad una matrice a scala per righe $S$. Sappiamo quindi per il corollario \ref{corol-ker(A)=ker(S)} che $\ker A = \ker S$, e quindi, per il Corollario \ref{corol-propos-dimensione-nucleo} $\dim \ker A = \dim \ker S = n -r$. Dunque troviamo che $rg(A) = r$, ossia che il numero di pivot di $S$, ossia il numero di righe non nulle di $S$, è uguale al rango di $A$. Dato che per definizione il numero di righe non nulle è il rango di $A^t$, troviamo che $rg(A) = rg(A^t)$. $\blacksquare$ 

\section{Eliminazione di Gauss}

\begin{Definizione}[Eliminazione di Gauss]
\label{def-elim-di-gauss}
L'\textbf{eliminazione di Gauss} è un algoritmo che ci aiuta a trovare, se esistono, tutte le soluzioni di sistemi del tipo 
\[\bm{A \cdot \X = B}\]
dove $A \in M_{mn}(\K), \X \in M_{n1}(\K)$ e $B \in M_{m1}(\K)$. \\
L'idea è quella di trovare, attraverso operazioni elementari sulle righe di $A$, una matrice a scala per righe $S \in M_{mn}(\K)$, dato che l'\textbf{insieme delle soluzioni del sistema} $\bm{A \cdot \X = B}$ è pari all' \textbf{insieme delle soluzioni del sistema} $\bm{S \cdot \X = C}$, dove $C$ è la matrice che si ottiene ripetendo le stesse identiche operazioni elementari effettuate su $A$ su $B$.
\end{Definizione}
L'utilità di questo algoritmo sta nel fatto che il sistema $S \cdot \X = C$ è un sistema molto semplice da risolvere, essendo $S$ a scala per righe.

\Esempio \nopagebreak

Risolvere il sistema \nopagebreak
\[
\begin{cases}
\bm{3x_1 - x_2 + x-3= 1} \\
\bm{x_1 - x_2 + 5x_3= 2} \\
\bm{-x_1 - 2x_2 + x_3= 0}
\end{cases}
\]

Risolvere questo sistema equivale a risolvere l'equazione $A \cdot \X = B$, dove 
\[A = \begin{bmatrix}
    3 & -1 & 1 \\
    1 & -1 & 5 \\
    -1 & -2 & 1
\end{bmatrix}
\quad \text{e} \quad
B = \begin{bmatrix}
    1 \\ 2 \\ 0
\end{bmatrix}
\]

Usando l'eliminazione di Gauss, prendiamo $[A | B]$ e trasformiamo $A$ in $S$ e $B$ in $C$. Dunque
\[
[A|B] = \left[\begin{array}{ccc|c}
    3 & -1 & 1 & 1 \\
    1 & -1 & 5 & 2 \\
    -1 & -2 & 1 & 0
\end{array}\right] \leadsto [S|C] = \left[ \begin{array}{ccc|c}
    1 & -1 & 5 & 2 \\
    0 & 2 & -14 & -5 \\
    0 & 0 & -15 & -\frac{11}{2}
\end{array}\right]
\]

Troviamo quindi un sistema molto più semplice da risolvere:
\[
\begin{cases}
x_1 - x_2 + 5x_3 = 2 \\
x_2 - 14 x_3 = -5  \\
-15 x_3 = -\frac{11}{2}
\end{cases}
\implies
\begin{cases}
    \bm{x_3 = \frac{11}{30}} \\
    \bm{x_2 = \frac{1}{15}} \\
    \bm{x_1 = \frac{7}{30}}
\end{cases}
\]

\section{Equazioni cartesiane di un sottospazio vettoriale}
\begin{Definizione}[]
\label{def-utile-solo-per-la-dimostrazione}
Preso uno spazio vettoriale $\V$ e una base $\B = \{v_1, \ldots, v_n\}$ di $\V$, \mbox{$\bm{v_i^*}: \V \longrightarrow \K$} è l'unica applicazione lineare tale che
\[
\bm{v_i^*(v_j) = \delta_{ij}} = \begin{cases}
    \bm{1} & \text{ se} \bm{j = i} \\
    \bm{0} & \text { se} \bm{j \neq i}
\end{cases}
\]
\NB Queste applicazioni hanno senso se e solo se $v_i$ si trova in una base di $\V$.
\end{Definizione}
Questa definizione ci sarà utile solamente per una dimostrazione che faremo tra poco.

\begin{Definizione}[Equazioni cartesiane di un sottospazio]
\label{def-eq-cart-sott}
Preso uno spazio vettoriale $\V$ e un suo sottospazio vettoriale $\W$, esistono $\bm{f_1, \ldots, f_d}$ applicazioni lienari da $\V$ a $\K$ tali che \mbox{$\bm{\W = \{v \in \V \mid f_1(v) = \ldots = f_d(v) = \undzero\}}$}, ossia \mbox{$\bm{\W = \ker f_\lambda}$} dove
\begin{align*}
\bm{f_\lambda: \V} &\longrightarrow \bm{\K^d} \\
\bm{v}  &\longmapsto \bm{(f_1(v), \ldots, f_d(v))}
\end{align*}
\end{Definizione}

\begin{Osservazione}[]
\label{oss-equazioni-cartesiane}
Preso uno spazio vettoriale $\W$ e un suo sottospazio vettoriale $\W$, è \textbf{sempre} possibile trovare un sistema di equazioni lineari il cui insieme delle soluzioni è uguale a $\W$.
\end{Osservazione}

\Dimostrazione \nopagebreak

Sia $\B = \{w_1, \ldots, w_m\}$ una base di $\W$. Estendiamo questa base a \mbox{$\Bs = \{w_1, \ldots, w_m, v_1, \ldots, v_d\}$}, base di $\V$. Troviamo quindi che $d = \dim \V - m$. Poniamo ora $f_i = v_i^* \enspace \forall i \in \{1, \ldots, d\}$ (ha senso perché $v_1, \ldots, v_d$ fanno parte di una base di $\V$). Definiamo ora il sottospazio vettoriale \mbox{$\U = \{v \in \V \mid f_1(v) = \ldots = f_d(v) = \undzero\}$}. Sappiamo che $w_i \in \ker f_j$ dato che $w_i \neq v_j$ e $f_j = v_j^*$, dunque notiamo subito come $\W \subset \U$, dato che tutti i suoi generatori si trovano in $\U$. Non ci manca quindi che dimostrare che $\U \subset \W$. \\
Prendiamo un vettore $u \in \U$. Essendo $\Bs$ una base di $\V$, ed essendo $\U$ un sottospazio vettoriale di $\V$, possiamo dire che $u = \sum_{i = 1}^m x_i w_i + \sum_{k = 1}^d y_k v_k$. Dunque \mbox{$f_j(u) = v_j^*(u) = \sum_{i = 1}^m x_i v_j^*(w_i) + \sum_{k = 1}^d y_k v_j^*(v_k)$}. La prima sommatoria, per quello detto prima, sappiamo che è pari a $0$, e della seconda sommatoria l'unico termine che rimane è $y_j$, dato che $v_j^*(v_j) = 1$. Dunque troviamo che $f_j(u) = y_j$, e dato che $u \in \U$, troviam oche $y_j = 0$. Dato che questo vale $\forall j \in \{1, \ldots, d\}$, troviamo che $u = \sum_{i = 1}^m x_i w_i$, ossia $u \in \W$. Dunque $\U \subset \W$. \\
Possiamo quindi concludere che $\W = \U = \{v \in \V \mid f_1(v)  = \ldots = f_d(v) = \undzero\}$. 

\Esempio \nopagebreak

Preso il sottospazio vettoriale $\R^4 \supset \bm{\W = \braket{(-1,0,3,2),(1,-1,1,-1),(1,-2,5,0)}}$, trovare delle \textbf{equazioni cartesinae di} $\bm{\W}$. \nopagebreak

Prima di inizare a risolvere l'esercizio, è bene notare una cosa: nel testo è richiesto di trovare \textit{delle} equazioni cartesiane di $\W$, non \textit{le} equazioni cartesiane di $\W$. Questo perché le equazioni cartesiane di un sottospazio vettoriale sono effettivamente infinite. Vedremo il perché alla fine dell'esercizio non appena le avremo trovate.

Cerchiamo quindi una funzione
\begin{align*}
f_\lambda: \R^4 &\longrightarrow \R \\
\X &\longmapsto \lambda_1 x_1 + \lambda_2 x_2 + \lambda_3 x_3 + \lambda_4 x_4
\end{align*}
tale che $\ker f_\lambda = \W$, ossia $f_\lambda (\W) = \{\undzero\}$. Questo è vero se e solo se preso un vettore $w \in \W$ allora $f_\lambda(w) = \undzero$. Sappiamo che $w = \mu_1 (-1,0,3,2) + \mu_2 (1,-1,1,-1) + \mu_3 (1,-2,5,0)$. Essendo $f_\lambda$ lineare, per dimostrare che $\ker f_\lambda = \W$ ci basta dimostrare che \mbox{$f_\lambda(-1,0,3,2) = f_\lambda(1,-1,1,-1) = f_\lambda(1,-2,5,0) = \undzero$}. Questo si traduce nel sistema
\[
\begin{cases}
-\lambda_1 + 3\lambda_2 + 2\lambda_4 = 0 \\
\lambda_1 - \lambda_2 + \lambda_3 - \lambda_4 = 0 \\
\lambda_1 - 2\lambda_22 + 5\lambda_3 = 0
\end{cases}
\]
Risolvere questo sistema equivale a risolvere l'equazione $A \cdot \Delta = \undzero$, dove 
\[ A =
\begin{bmatrix}
    -1 & 0 & 3 & 2  \\
    1 & -1 & 1 & -1 \\
    1 & -2 & 5 & 0
\end{bmatrix}
\quad \text{e} \quad 
\Delta = \begin{bmatrix}
    \lambda_1 \\
    \lambda_2 \\
    \lambda_3 \\
    \lambda_4
\end{bmatrix}
\]

Usando l'eliminazione di Gauss trasformiamo $A$ nella matrice a scala per righe
\[S = 
\begin{bmatrix}
    -1 & 0 & 3 & 2 \\
    0 & -1 & 4 & 1 \\
    0 & 0 & 0 & 0
\end{bmatrix}
\]

Dunque troviamo il sistema
\[
\begin{cases}
-\lambda_1 + 3\lambda_2 + 2\lambda_4 = 0 \\
-\lambda_2 + 4\lambda_3 + \lambda_4 = 0 \\
0 = 0 
\end{cases}
\implies
\begin{cases}
    \lambda_1 = 3 \lambda_3 + 2 \lambda_4 \\
    \lambda_2 = 4\lambda_3 +  \lambda_4
\end{cases}
\]

Troviamo quindi che 
\begin{align*}
(\lambda_1, \lambda_2, \lambda_3, \lambda_4) &= (3\lambda_3 + 2 \lambda_4, 4\lambda_3 + \lambda_4, \lambda_3, \lambda_4) \\
 &= (3t_1 + 2t_2, 4t_1 + t_2, t_1, t_2) \\
 &= t_1 (3,4,1,0) + t_2 (2,1,0,1)
\end{align*}
Dunque non abbiamo trovato solo una $f_\lambda$, ma ben due:
\begin{align*}
f_{\lambda_1} &= 3x_1 + 4x_2 + x_3  \\
f_{\lambda_2} &= 2x_2 + x_2 + x_4 
\end{align*}

In questo caso $\W \subset \ker f_{\lambda_1} \cap \ker f_{\lambda_2}$, ossia 
\[ \bm{\W} = 
\left\{\bm{\X} \in \bm{\R^4} \;\middle|\; \begin{cases}
\bm{3 x_1 + 4x_2 + x_3 = 0} \\
\bm{2x_1 + x_2 + x_4 = 0}
\end{cases}
\! \! \! \!
\right\}
\]

\section{Inversa di una matrice}

\begin{Definizione}[Inversa di una matrice]
\label{def-inversa}
Presa una matrice $A \in M_{nn}(\K)$, la sua \textbf{inversa} è la matrice $\bm{A^{-1}}$ tale che $\bm{A \cdot A^{-1} = A^{-1} \cdot A = 1_n}$. \\
Se una matrice $A$ è invertibile allora diciamo che $\bm{A \in AL_n(\K)}$, altrimenti si dice che la matrice è \textbf{singolare}.
\end{Definizione}

\begin{Proposizione}[Invertibilità di una matrice]
\label{propos-invertibilità-matrice}
Presa una matrice $A \in M_{nn}(\K)$: 
\begin{itemize}
    \item \textbf{è invertibile} $\iff \bm{rg(A) = n}$
    \item \textbf{è invertibile} $\iff \bm{\dim \ker A = 0}$
\end{itemize}
\end{Proposizione}

\Dimostrazione \nopagebreak

È facile notare come le due affermazioni siano in realtà equivalenti, dato che per la Proposizione \ref{propos-dimensione-nucleo-(n-r)} sappiamo che $\dim \ker A = n - r$, dove $r$ è il numero di pivot della matrice $S$ a scala per righe ottenuta da $A$, che corrisponde al $rg(A^t)$ che è pari al $rg(A)$. Dunque se $rg(A) = n$ troviamo proprio che $\dim \ker A = 0$. Per questo, ci limiteremo a dimostrare che \mbox{$A$ è invertibile $\iff \dim \ker A = 0$}. 
\begin{enumerate}[label=\arabic*°]
    \item $A$ è invertibile $\implies \dim \ker A = 0$ \\
    Prendiamo $v \in \ker A$. Sappiamo quindi che $A \cdot v = \undzero$. Poiché $A$ è invertibile, sappiamo che esiste $A^{-1}$, dunque ha senso moltiplicare entrambi i lati dell'equazione per $A^{-1}$. Troviamo quindi $A^{-1} (A \cdot v) = A^{-1} \cdot \undzero = \undzero$. Dato che $A^{-1} (A \cdot v) = (A^{-1} \cdot A) v$, e dato che $A^{-1} \cdot A = 1_n$, troviamo che $1_n v = \undzero$, ossia $v = \undzero$. Dunque $\ker A = \{\undzero\}$, ossia $\dim \ker A = 0$. $\blacksquare$
    \item $A$ è invertibile $\impliedby \dim \ker A = 0$ \\
    Se $\dim \ker A = 0$, allora sappiamo che $L_A$ è iniettiva. Dato che $\dim \K^n = \dim \ker A + \dim Im(A) = \dim Im(A)$, troviamo che $\dim Im(A) = n$, ed essendo $Im(A)$ un sottospazio vettoriale di $\K^n$, per la Proposizione \ref{propos-dim-spa-vet} troviamo che $Im(A) = \K^n$. Dunque $L_A$ è anche suriettiva, e quindi è biiettiva. Essendo biiettiva e lienare, $L_A$ è un isomorfismo, che per definizione ammette sempre un'applicazione lineare inversa $L_A^{-1}$. La matrice associata a $L_A^{-1}$ è proprio la matrice inversa di $A$, ossia $A^{-1}$. Dunque $A$ è invertibile. $\blacksquare$
\end{enumerate}

\begin{Proposizione}[Trovare l'inversa di una matrice]
\label{propos-trovare-inversa}
Presa una matrice $A \in M_{nn}(\K)$ invertibile, per trovare la sua inversa $\bm{A^{-1}}$ prendiamo $\bm{[A|1_n]}$, e trasofrmiamo con delle operazioni sulle righe $A$ in una matrice a scala per righe $S$, e ripetiamo le stesse operazioni su $1_n$. Troviamo quindi $\bm{[S|B]}$, e, facendo altre operazioni su $S$ per trasofrmarla in $1_n$, e ripetendo le stesse operazioni anche su $B$, troviamo $\bm{[1_n | A^{-1}]}$, trovando quindi $A^{-1}$.
\end{Proposizione}

\Dimostrazione \nopagebreak

Stiamo cercando una matrice $\X \in M_{nn}(\K)$ tale che $A \cdot \X = 1_n$. Sappiamo che $A = [A_1, \ldots, A_n]$, $\X = [\X_1, \ldots, \X_n]$ e  $1_n = [e_1, \ldots, e_n]$. Dunque dobbiamo risolvere il sistema
\[
\begin{cases}
A \cdot x_1 = e_1 \\
\vdots \\
A \cdot x_n  = e_n
\end{cases}
\]

Per poter risolvere questo sistema, possiamo usare l'eliminazione di Guass: prendiamo $[A|e_1]$ e trasformiamo $A$ in una matrice a scala per righe $S$, trovando $[S|B_1]$. Dato che però questo proceidmento lo dobbiamo ripetere per tutti i vettori $e_i$, possiamo farlo una sola volta trasformando $[A|e_1, \ldots, e_n] = [A|1_n]$ in $[S|B]$. \\
Essendo ora $A$ invertibile, noi sappiamo che $S$ può essere ritrasformata in $1_n$. Questo perché $S$ ha $n$ pivot su $n$ righe, che si trovano quindi lungo la diagonale, essendo $S$ a scala per righe, dunque basta eliminare tutti gli elementi sopra la diagonale e dividere ogni riga per il proprio pivot per trovare $1_n$. Dunque da $[S|B]$ possiamo arrivare a $[1_n | C]$. Dal sistema originale troviamo quindi
\[
\begin{cases}
1_n \cdot x_1 = c_1 \\
\vdots \\
1_n \cdot x_n = c_n
\end{cases}
\]
Dunque troviamo $\X = C$, e dato che $A \cdot \X = 1_n$, troviamo che $C = A^{-1}$. Siamo quindi passati da $[A|1_n]$ a $[1_n|A^{-1}]$. $\blacksquare$

\Esempio \nopagebreak

Trovare l'\textbf{inversa} della matrice \nopagebreak

\[ \bm{A} = 
\begin{bmatrixb}
2 & 1 & 3 \\
-1 & 0 & 1 \\
3 & 2 & 8 
\end{bmatrixb}
\]

Dimostriamo innanzitutto che $A^{-1}$ sia invertibile. Riduciamo quindi $A$ a una matrice a scala per righe per studiare il rango di $A$.
\[
A = \begin{bmatrix}
    2 & 1 & 3 \\
    -1 & 0 & 1 \\
    3 & 2 & 8
\end{bmatrix}
\leadsto 
S = \begin{bmatrix}
    -1 & 0 & 1 \\
    0 & 1 & 5 \\
    0 & 0 & 1
\end{bmatrix}
\] 
Dato che i pivot sono $3$, troviamo che $rg(A) = 3$, e dato che $A$ è una matrice $3 \times 3$, troviamo che $A$ è invertibile.

Per trovare ora $A^{-1}$, dobbiamo passare da $[A|1_n]$ a $[1_n|A^{-1}]$.
\[[A|1_n] = \left[\begin{array}{ccc|ccc}
    2 & 1 & 3 & 1 & 0 & 0 \\
    -1 & 0 & 1 & 0 & 1 & 0 \\
    3 & 2 & 8 & 0 & 0 & 1
\end{array}\right]\leadsto [1_n | A^{-1}]
\left[ \begin{array}{ccc|ccc}
    1 & 0 & 0 & -2 & -2 & 1 \\
    0 & 1 & 0 & 11 & 7 & .5 \\
    0 & 0 & 1 & -2 & -1 & 1
\end{array}\right]
\]

Troviamo quindi che
\[
\bm{A^{-1}} = \begin{bmatrixb}
    -2 & -2 & 1 \\
    11 & 7 & -5 \\
    -2 & -1 & 1
\end{bmatrixb}
\]

\NB Se si vuole verificare che questo risultato sia corretto, basta verificare che $A \cdot A^{-1} = 1_n$.

\begin{Osservazione}[]
\label{oss-inversa}
Presa una matrice $A \in M_{nn}(\K)$, 
$\bm{A}$ \textbf{invertibile} $\iff \bm{\{A_1, \ldots, A_n\}}$ \textbf{sono linearmente indipendenti}.
\end{Osservazione}

\vspace{2cm}

\Dimostrazione \nopagebreak

\begin{enumerate}[label=\arabic*°]
    \item $A$ è invertibile $\implies$ $\{A_1, \ldots, A_n\}$ sono linearmente indipendenti \\
    Per verificare che i vettori colonna di $A$ siano linearmente indipendenti dobbiamo dimostrare che l'equazione $\lambda_1 A_1 + \ldots + \lambda_n A_n = A \cdot \Lambda = \undzero$ ha come unica soluzione la soluzione banale. Dato che $A$ è invertibile, sappiamo che esiste $A^{-1}$. Se molitplichiamo entrambi i membri dell'equazione per $A^{-1}$ troviamo l'equazione $A^{-1}(A \cdot \Lambda) = A^{-1} \cdot \undzero = \undzero$. Dato che $A^{-1}(A \cdot \Lambda) = (A^{-1} \cdot A)\Lambda$, e dato che $A^{-1} \cdot A = 1_n$, troviamo che $1_n \Lambda = \undzero$, ossia $\Lambda = \undzero$, ossia trovimao che l'unica soluzione all'equazione iniziale è la soluzione banale, dunque i vettori colonna di $A$ sono linearmente indipendenti. $\blacksquare$
    \item $A$ è invertibile $\impliedby$ $\{A_1, \ldots, A_n\}$ sono linearmente indipendenti \\
    Affinché $A$ sia invertibile è necessario che esista una matrice $A^{-1}$ tale che $A \cdot A^{-1} = 1_n$. Sappiamo che $A^{-1} = [A_1^{-1}, \ldots, A_n^{-1}]$ e che $1_n = [e_1, \ldots, e_n]$, dunque l'equazione $A \cdot A^{-1} = 1_n$ possiamo riscriverla come il sistema
    \[
    \begin{cases}
    A \cdot A_1^{-1} = e_1 \\
    \vdots \\
    A \cdot A_n^{-1} = e_n 
    \end{cases}
    \]
    Dato che $A_1, \ldots, A_n$ sono linearmente indipendenti, per l'Osservazione \ref{oss-dimen} questi sono una base di $\K^n$, dunque $\braket{A_1, \ldots, A_n} = \K^n$. Per la Proposizione \ref{propos-immagine-combinazione-lineare-colonne} sappiamo però che $\braket{A_1, \ldots, A_n} = Im(A)$, dunque troviamo che $\K^n = Im(A)$, ossia troviamo che $A$ è suriettiva. Esiste quindi sicuramente una matrice $A^{-1}$ che risolva il sistema appena visto, ossia $A$ è invertibile. $\blacksquare$
\end{enumerate}

\begin{Proposizione}[Inversa del prodotto tra matrici]
\label{propos-inversa-prodotto}
Prese $n$ matrici $A_1, \ldots, An$ il cui prodotto $A_1 \cdot \ldots \cdot A_n$ sia sensato e invertibile, l'inversa di tale prodotto è pari al prodotto delle inverse prese in ordine inverso, ossia $\bm{(A_1 \cdot \ldots \cdot A_n)^{-1} = A_n^{-1} \cdot \ldots \cdot A_1^{-1}}$.
\end{Proposizione}
La dimostrazione qui è omessa ma segue lo stesso procedimento visto per la dimostrazione della Proposizione \ref{propos-trasposta-prodotto}.

\section{Applicazioni diagonalizzabili}

\begin{Definizione}[Matrice del cambiamento di base]
\label{def-matrice-cambiamento-base}
Preso uno spazio vettoriale $\V$ e un'applicazione lineare $f: \V \longrightarrow \V$, la \textbf{matrice del cambiamento di base} è la matrice $\bm{M_\Bs^\B(Id_\V)}$.
\end{Definizione}
Questa matrice ci permette di trovare le coordinate di $\Bs$ di un qualsiasi vettore $v \in \V$ a partire dalle sue coordinate nella base $\B$. Infatti, per la Proposizione \ref{propos-matrice-associata-prodotto} troviamo che $X_\Bs(Id_\V(v)) = M_\Bs^\B(Id_V) \cdot X_\B(v)$, ma dato che $Id_\V(v) = v$, troviamo che $X_\Bs(v) = M_\Bs^\B(Id_\V) \cdot X_\B(v)$. \\
È possibile anche notare che $M_\Bs^\B(Id_\V) = [X_\Bs(v_1), \ldots, X_\Bs(v_n)]$. 

\begin{Osservazione}[]
\label{oss-invertibilita-matrice-cambiamento-base}
La matrice $M_\Bs^\B(Id_\V)$ è invertibile, e in particolare $\bm{M_\Bs^\B(Id_\V)^{-1} = M_\B^\Bs(Id_\V)}$. 
\end{Osservazione}

\Dimostrazione \nopagebreak

Sappiamo per la Proposizione \ref{propos-matrice-associata-composizione} che $M_\Bs^\B(Id_\V) \cdot M_\B^\Bs(Id_\V) = M_\Bs^\Bs(Id_\V \circ Id_\V) = M_\Bs^\Bs(Id_\V)$. È facile notare che la j-esima colonna di questa matrice non è altro che $X_\Bs(v_j)$, con $v_j \in \Bs$. Dunque \mbox{$M_\Bs^\Bs(Id_\V)_j = X_\Bs(v_j) = (0,\ldots,1,\ldots,0)$}, dove l'$1$ si trova alla j-esima colonna. Rappresentando $M_\Bs^\Bs(Id_\V)$ notiamo quindi che 
\[
M_\Bs^\Bs(Id_\V) = \begin{bmatrix}
    1 & 0 & \dots & 0 \\
    0 & 1 & \dots & 0 \\
    \vdots & \vdots & \ddots & \vdots \\
    0 & 0 & \dots & 1
\end{bmatrix}
= 1_n
\]
Dato che $M_\Bs^\Bs(Id_\V) = M_\Bs^\B(Id_\V) \cdot M_\B^\Bs(Id_\V)$, troviamo che $M_\Bs^\B(Id_\V)^{-1} = M_\B^\Bs(Id_\V)$. $\blacksquare$

\begin{Proposizione}[]
\label{propos-matrice-cambiamento-base-base-standard}
Presa la base $\{B_1, \ldots, B_n\}$ di $\K^n$ e la matrice $B = [B_1, \ldots, B_n]$, $\bm{B = M_S^\B(Id_\V)}$.
\end{Proposizione}

\Dimostrazione \nopagebreak

Banalmente, dato che la colonna j-esima di $M_S^\B(Id_\V)$ è $X_S(Id_\V(B_j)) = X_S(B_j)$, e dato che $X_S(B_j) = B_j$, troviamo prorpio che $M_S^\B(Id_\V) = B$. $\blacksquare$ 

\begin{Osservazione}[]
\label{oss-matrice-cambiamento-base-passare-per-due-basi}
$\bm{M_\Bs^\B(Id_\V) = M_\Bs^S(Id_\V) \cdot M_S^\B(Id_\V) = M_S^\Bs(Id_\V)^{-1} \cdot M_S^\B(Id_\V)}$.
\end{Osservazione}
Questa osservazione è ovvia ma molto importante: infatti, abbiamo visto come le matrici $M_S^\B(Id_\V)$ siano tutte incredibilmente facili da calcolare, dato che le loro colonne sono i vettori della base $\B$, dunque abbiamo appena trovato un modo molto semplice per calcolare una qualsiasi matrice del cambiamento di base $M_\Bs^\B(Id_\V)$. 

\Esempio \nopagebreak

Prese le basi $\B = \{(2,1,0), (1,0,-1), (0,-1,2)\}$ e $\Bs = \{(2,-1,3), (1,0,2), (3,1,8)\}$ di $\Q^3$, calcolare $\bm{M_\Bs^\B(Id_\V)}$. \nopagebreak

Dato che $M_\Bs^\B(Id_\V) = M_S^\Bs(Id_\V)^{-1} \cdot M_S^\B(Id_\V)$, ci basta cercare queste due matrici. Sappiamo che 
\[
M_S^\Bs(Id_\V) = \begin{bmatrix}
    2 & 1 & 3 \\ 
    -1 & 0 & 1 \\
    3 & 2 & 8
\end{bmatrix}
\quad \text{e} \quad
M_S^\B(Id_\V) = \begin{bmatrix}
    2 & 1 & 0 \\
    1 & 0 & -1  \\
    0 & -1 & 2
\end{bmatrix}
\]

Con un po di calcoli è facile trovare che 
\[M_\Bs^S(Id_\V) =  M_S^\Bs(Id_\V)^{-1}
\begin{bmatrix}
    -2 & 11 & -2 \\
    -2 & 7 & -1 \\
    1 & -5 & 1
\end{bmatrix}
\]

Dunque troviamo che
\[
\bm{M_\Bs^\B(Id_\V)} = 
\begin{bmatrixb}
    -6 & -3 & 4 \\
    29 & 16 & -17 \\
    -5 & -3 & 3
\end{bmatrixb}
\]

\begin{Osservazione}[]
\label{oss-matrice-diagonalizzata-M_C^C(f) = M M M}
Preso uno spazio vettoriale $\V$ e un endomorfismo $f$ di $\V$, e prese due basi $\B$ e $\Bs$ di $\V$, $\bm{M_\Bs^\Bs(f) = M_\Bs^\B(Id_\V) \cdot M_\B^\B(f) \cdot M_\B^\Bs(Id_\V)}$.
\end{Osservazione}

\Dimostrazione \nopagebreak

Sappiamo già che $M_\Bs^\B(Id_\V) \cdot M_\B^\B(f) \cdot M_\B^\Bs(f) = M_\Bs^\Bs(Id_\V \circ f \circ Id_\V)$. Dato che \mbox{$(Id_\V \circ f \circ Id_\V)(v)= Id_\V(f(Id_\V(v))) = Id_\V(f(v)) = f(v)$}, troviamo che $Id_\V \circ f \circ Id_\V$. Dunque $M_\Bs^\Bs(Id_\V \circ f \circ Id_\V) = M_\Bs^\Bs(f) = M_\Bs^\B(Id_\V) \cdot M_\B^\B(f) \cdot M_\B^\Bs(Id_\V)$. $\blacksquare$

\begin{Definizione}[Matrici coniugate]
\label{def-matrici-coniugate}
Prese due matrici $A, B \in M_{nn}(\K)$, queste si dicono \textbf{coniugate} (o \textbf{simili}) se esiste una matrice $G$ invertibile tale che $\bm{A = G^{-1} \cdot B \cdot G}$ ($\bm{A \sim B}$).
\end{Definizione}

\begin{Definizione}[Endomorfismi coniugati]
\label{def-endomorfismi-coniugati}
Preso uno spazio vettoriale $\V$ e due endomorfismi $f,g$ su $\V$, questi si dicono \textbf{coniugati} se esistono due basi $\B$ e $\Bs$ di $\V$ tali che $\bm{M_\B^\B(f) = M_\Bs^\Bs(g)}$.
\end{Definizione}

\begin{Osservazione}[]
\label{oss-matrici-coniugate}
Le matrici $M_\B^\B(f)$ e $M_\Bs^\Bs(f)$ sono coniugate.
\end{Osservazione}
Questo è evidente dato che $M_\Bs^\Bs(f) = M_\Bs^\B(Id_\V) \cdot M_\B^\B(f) \cdot M_\B^\Bs(Id_\V) = M_\B^\Bs(Id_\V)^{-1} \cdot M_\B^\B(f) \cdot M_\B^\Bs(Id_\V)$

\begin{Proposizione}[]
\label{propos-dimensione-coniugate}
Preso uno spazio vettoriale $\V$ e due endomorfismi $f,g$ su tale spazio vettoriale $\V$, \mbox{$\bm{f \sim g } \implies \bm{\dim \ker f = \dim \ker g}$}.
\end{Proposizione}

\Dimostrazione \nopagebreak

Sappiamo che $f$ e $g$ sono coniugate se esistono due basi $\B$ e $\Bs$ di $\V$ tali che $M_\B^\B(f) = M_\Bs^\Bs(g) = A$. Dato che $\dim \V = \dim \ker f + \dim Im(f) = \dim \ker f + rg(M_\B^\B(f)) = \dim \ker f + rg(A)$ (la dimensione dell'immagine di un endomorfismo è pari al rango di una sua matrice associata), e dato che allo stesso modo $\dim \V = \dim \ker g + \dim Im(g) = \dim \ker g + rg(M_\Bs^\Bs(g)) = \dim \ker g + rg(A)$, troviamo che $\dim \ker f = \dim \ker g$. $\blacksquare$

\begin{Osservazione}[]
\label{oss-coniugio-rel-equiv}
La relazione di coniugio è una \textbf{relazione d'equivalenza}.
\end{Osservazione}

\Dimostrazione \nopagebreak

Per essere una relazione d'equivalenza la relazione di coniugio deve essere riflessiva, simmetrica e transitiva. 
\begin{itemize}
    \item Riflessività \\
    $A \sim A$ è sempre vera dato che $A = 1_n^{-1} \cdot A \cdot 1_n = 1_n \cdot A = A$. 
    \item Simmetria \\
    Preso $A \sim B$, dimostriamo che $B \sim A$. $A \sim B$ significa che $A = G^{-1} \cdot B \cdot G$. Moltiplicando ambo i memrbi per $G$ e poi per $G^{-1}$, troviamo che $G \cdot A \cdot G^{-1} = B$. Siccome $G = (G^{-1})^{-1}$, troviamo che $B \sim A$.
    \item Transitività \\
    Preso $A \sim B$ e $B \sim C$ dimostriamo che $A \sim C$. $A \sim B$ significa che $A = G^{-1} \cdot B \cdot G$ e $B \sim C$ significa che $B = H^{-1} \cdot C \cdot H$. Troviamo quindi che $A = G^{-1} \cdot B \cdot G = G^{-1} \cdot H^{-1} \cdot C \cdot G \cdot H$. Presa la matrice $H \cdot G$, la Proposizione \ref{propos-inversa-prodotto} ci dice che l'inversa di questa è esattamente $G^{-1} \cdot H^{-1}$. Dunque troviamo $A = (H \cdot G)^{-1} \cdot C \cdot (H \cdot G)$, ossia $A \sim C$.
\end{itemize}
Possiamo quindi concludere che la relazione di coniugio è una relazione d'equivalenza. $\blacksquare$

\begin{Osservazione}[]
\label{oss-B = M_S^S(L_B)}
Presa una matrice $B \in M_{mn}(\K)$, $\bm{B = M_S^S(L_B)}$.
\end{Osservazione}

\Dimostrazione \nopagebreak

Banalmente, la colonna j-esima di $M_S^S(L_B)$ altro non è che $X_S(L_B(e_j)) = X_S(B_j) = B_j$, dunque $M_S^S(L_B) = B$. $\blacksquare$

\Esempio \nopagebreak

Presa la base $\B = \{(1,1), (1,-1)\}$ di $\K^2$ e la matrice \nopagebreak

\[ \bm{B} = 
\begin{bmatrixb}
    1 & 2 \\
    2 & 1
\end{bmatrixb}
\]
trovare $\bm{M_\B^\B(L_B)}$. \nopagebreak

Ricordiamo che per l'Osservazione \ref{oss-B = M_S^S(L_B)} $B = M_S^S(L_B)$, e che per l'Osservazione \ref{oss-matrice-diagonalizzata-M_C^C(f) = M M M} \mbox{$M_B^B(L_B) = M_\B^S(Id) \cdot M_S^S(L_B) \cdot M_S^\B(Id)$} (per comodità spesso da ora in avati scriveremo $Id$ per indicare l'applicazione identità). Sappiamo inoltre per l'Osservazione \ref{oss-invertibilita-matrice-cambiamento-base} che $M_\B^S(Id) = M_S^\B(Id)^{-1}$, e che per la Proposizione \ref{propos-matrice-cambiamento-base-base-standard} 
\[M_S^\B(Id) = \begin{bmatrix}
    1 & 1 & \\
    1 & -1
\end{bmatrix}\]
Con un po' di calcoli troviamo che 
\[M_\B^S(Id) = M_S^\B(Id)^{-1} = \begin{bmatrix}
    \frac{1}{2} & \frac{1}{2} \\
    \frac{1}{2} & -\frac{1}{2}
\end{bmatrix}\]
Dunque
\[
\bm{M_\B^\B(L_B)} = M_\B^S(Id) \cdot M_S^S(L_B) \cdot M_S^S(L_B) \cdot M_S^\B(Id) = \begin{bmatrix}
\frac{1}{2} & \frac{1}{2} \\
\frac{1}{2} & -\frac{1}{2}    
\end{bmatrix} 
\cdot \begin{bmatrix}
    1  & 2 \\
    2 & 1
\end{bmatrix}
\cdot \begin{bmatrix}
    1 & 1 \\ 
    1 & -1
\end{bmatrix}
= \begin{bmatrixb}
    3 & 0 \\
    0 & -1
\end{bmatrixb}
\]

\subsection{Cercare la matrice più semplice possibile}
Preso un endomorfismo $f: \V \longrightarrow \V$, cerchiamo ora un metodo semplice per trovare una base $\B$ di $\V$ tale che $M_\B^\B(f)$ sia la matrice più semplice possibile, ossia la matrice \textit{diagonale}. 

\begin{Definizione}[Endomorfismo diagonalizzabile]
\label{def-endomorfismo-diagonalizzabile}
Preso uno spazio vettoriale $\V$, un endomorfismo $f$ su $\V$ e una base $\B$ di $\V$, questa \textbf{diagonalizza} $f$ se $\bm{M_\B^\B(f)}$ è \textbf{diagonale}.\\
$f$ è detta \textbf{diagonalizzabile} se esiste una base di $\V$ che diagonalizzi $f$. \\
Una matrice $A \in M_{mn}(\K)$ è detta \textbf{diagonalizzabile} se $L_A$ è diagonalizzabile.
\end{Definizione}

\begin{Osservazione}[]
\label{oss-matrice-diagonale-f(vi)= lambda i vi}
$\bm{M_\B^\B(f)}$ \textbf{diagonale}  $\iff \bm{f(v_i) = \lambda_i v_i} \enspace \forall i \in \{1, \ldots, n\}$, dove $v_1, \ldots, v_n$ sono i vettori della base $\B$.  
\end{Osservazione}

\Dimostrazione \nopagebreak

\begin{enumerate}[label=\arabic*°]
    \item $M_\B^\B(f)$ diagonale $\implies f(v_i) = \lambda_i v_i \enspace \forall i \in \{1, \ldots, n\}$\\
    Sappiamo che la colonna i-esima di $M_\B^\B(f)$ è $X_\B(f(v_i))$. Per ipotesi sappiamo che $M_\B^\B(f)$ è diagonale, dunque la sua colonna i-esima ha un solo elemento non nullo che si trova proprio sulla riga i-esima. Dunque $X_\B(f(v_i)) = \lambda_i e_i$. Dato che $v_i$ è l'i-esimo vettore della base $\B$, troviamo che $X_\B(v_i) = e_i$, ed essendo $X_\B$ un'applicazione lineare, troviamo che $X_\B(\lambda_i v_i) = \lambda_i X_\B(v_i) = \lambda_i e_i$. Dunque $X_\B(\lambda_i v_i) = X_\B(f(v_i))$. Essendo $X_\B$ l'applicazione inversa dell'applicazione $f$ descritta nella Proposizione \ref{propos-base}, troviamo che $X_\B$ è iniettiva, dunque $f(v_i) = \lambda_i v_i$. Questo vale $\forall i \in \{1, \ldots, n\}$. $\blacksquare$
    \item $M_\B^\B(f)$ diagonale $\impliedby f(v_i) = \lambda_i v_i \enspace \forall i \in \{1, \ldots, n\}$\\
    Abbiamo già visto come la colonna i-esima di $M_\B^\B(f)$ sia $X_\B(f(v_i))$. Per ipotesi sappiamo che $f(v_i) = \lambda_i v_i$, dunque $X_\B(f(v_i)) = X_\B(\lambda_i v_i)$. Abbiamo già visto come $X_\B(\lambda_i v_i) = \lambda_i e_i$, e che questo valga $\forall i \in \{1, \ldots, n\}$. Rappresentando graficamente $M_\B^\B(f)$, notiamo subito che 
    \[
    M_\B^\B(f) = \begin{bmatrix}
        \lambda_1 & 0 & \dots & 0 \\
        0 & \lambda_2 & \dots & 0 \\
        \vdots & \vdots & \ddots & \vdots \\
        0 & 0 & \dots & \lambda_n
    \end{bmatrix}
    \]
    Dunque $M_\B^\B(f)$ è diagonale. $\blacksquare$
\end{enumerate}

\begin{Proposizione}[Corollario]
\label{corol-capire-matrice-diagonalizzabile}
Presa una matrice $A \in M_{mn}(\K)$, per capire se questa sia \textbf{diagonalizzabile} bisogna verificare che esista una base $\B$ di $\K^n$ tale che 
\[
\begin{cases}
\bm{L_A(v_1) = \lambda_1 v_1} \\
\vdots \\
\bm{L_A(v_n) = \lambda_1 v_n} \\
\end{cases}
\]
\end{Proposizione}

\Esempio \nopagebreak

Verificare se la matrice \nopagebreak

\[\bm{A} = 
\begin{bmatrixb}
0 & 1 \\
0 & 0
\end{bmatrixb} \in M_{22}(\Q)
\]
sia \textbf{diagonalizzabile}. \nopagebreak

Cerchiamo quindi una base $\B = \{v_1, v_2\}$ di $\Q^2$ che risolva il sistema
\[
\begin{cases}
L_A(v_1) = A \cdot v_1 = \lambda_1 \cdot v_1 \\
L_A(v_2) = A \cdot v_2 = \lambda_2 \cdot v_2
\end{cases}
\]

Questo sistema possiamo vederlo come un'equazione del tipo $A \cdot \X = \Lambda \cdot \X$, dove 
\[\X = 
\begin{bmatrix}
    v_1 & v_2
\end{bmatrix} = \begin{bmatrix}
    x_1 & x_2 \\
    y_1 & y_2
\end{bmatrix}
\quad \text{e} \quad
\Lambda = \begin{bmatrix}
    \lambda_1 & 0 \\
    0 & \lambda_2
\end{bmatrix}
\]
Troviamo quindi che
\[A \cdot \X = \begin{bmatrix}
    0 & 1 \\
    0 & 0
\end{bmatrix} \cdot \begin{bmatrix}
    x_1 & x_2 \\
    y_1 & y_2
\end{bmatrix} = \begin{bmatrix}
y_1 & y_2 \\
0 & 0
\end{bmatrix} =\X \cdot \Lambda = \begin{bmatrix}
    \lambda_1 x_1 & \lambda_2 x_2 \\
    \lambda_1 y_1 & \lambda_2 y_2
\end{bmatrix}\]

Da qui troviamo il sistema
\[
\begin{cases}
y_1 = \lambda_1 x_1 \\
y_2 = \lambda_2 x_2 \\
0 = \lambda_1 y_1 \\
0 = \lambda_2 y_2
\end{cases}
\]

Questo sistema si risolve solo ponendo $y_1 = y_2 = 0$. Troviamo quindi che $v_1 = (x_1, 0)$ e $v_2 = (x_2, 0)$. Questi due vettori non possono essere una base di $\Q^2$, dunque possiamo concludere che $A$ \textbf{non è diagonalizzabile}.

\section{Autospazi, autovalori e autovettori}

\begin{Definizione}[Autospazi, autovalori e autovettori]
\label{def-autospazi-autovalori-autovettori}
Preso uno spazio vettoriale $\V$ e un endomorfismo $f$ su $\V$, chiamiamo \mbox{$\bm{V_\lambda(f) = \ker (\lambda Id_\V - f)} = \{v \in \V \mid f(v) = \lambda v\}$} l'\textbf{autospazio} associato all'\textbf{autovalore} $\lambda$ se $\bm{V_\lambda(f) \neq \{\undzero\}}$. \\
Tutti i vettori dell'autospazio $\V_\lambda(f)$ diversi dal vettore nullo sono detti \textbf{autovettori di autovettore} $\bm{\lambda}$. 
\end{Definizione}

\begin{Osservazione}[]
\label{oss-autospazio-lambda-fuori-matrice-diagonale}
Se $\lambda_1, \ldots, \lambda_n$ sono i valori della matrice diagonale $M_\B^\B(f)$, allora \mbox{$\bm{\forall \lambda \notin \{\lambda_1, \ldots, \lambda_n\}} \implies \bm{V_\lambda = \{\undzero\}}$}.
\end{Osservazione}

\Dimostrazione \nopagebreak

Se prendiamo $M_\B^\B(\lambda Id_\V - f)$, sappiamo che $M_\B^\B(\lambda Id_\V - f) = M_\B^\B(\lambda Id_\V) - M_\B^\B(f)$ (questo perché, l'applicazione che prende un'applicazione lineare e due basi del suo dominio e restituisce la matrice associata è lineare), e che quindi
\[M_\B^\B(\lambda Id_\V - f) = 
\begin{bmatrix}
    \lambda & 0 & \dots & 0 \\
    0 & \lambda & \dots & 0 \\
    \vdots & \vdots & \ddots & \vdots \\
    0 & 0 & \dots & \lambda
\end{bmatrix}
- \begin{bmatrix}
    \lambda_1 & 0 & \dots & 0 \\
    0 & \lambda_2 & \dots & 0 \\
    \vdots & \vdots & \ddots & \vdots \\
    0 & 0 & \dots & \lambda_n
\end{bmatrix}
=
\begin{bmatrix}
    \lambda - \lambda_1 & 0 & \dots & 0\\
    0 & \lambda - \lambda_2 & \dots & 0 \\ 
    \vdots & \vdots & \ddots & \vdots \\
    0 & 0 & \dots & \lambda - \lambda_n
\end{bmatrix}
\] 

Prendiamo $v \in \V$. Sappiamo che 
\[
M_\B^\B(\lambda Id_\V - f) \cdot v = \begin{bmatrix}
    \lambda - \lambda_1 & 0 & \dots & 0\\
    0 & \lambda - \lambda_2 & \dots & 0 \\ 
    \vdots & \vdots & \ddots & \vdots \\
    0 & 0 & \dots & \lambda - \lambda_n
\end{bmatrix} \cdot \begin{bmatrix}
    x_1 \\
    x_2 \\
    \vdots \\
    x_n
\end{bmatrix} = \begin{bmatrix}
    x_1 (\lambda - \lambda_1) \\
    x_2 (\lambda - \lambda_2) \\
    \vdots \\
    x_n (\lambda - \lambda_n)
\end{bmatrix}
\]

Dato che $\lambda \notin \{\lambda_1, \ldots, \lambda_n\}$, sappiamo che affinché $M_\B^\B(\lambda Id_\V - f) \cdot v = \undzero$ è necessario che $x_1 = \ldots = x_n = 0$, ossia $v = \undzero$. Concludiamo quindi che $\V_\lambda = \ker (\lambda Id_\V - f) = \{\undzero\}$. $\blacksquare$

\begin{Proposizione}[Corollario]
\label{corol-autovalori-matrice}
Tutti gli \textbf{autovalori} di $f$ sono i valori della matrice diagonale $M_\B^\B(f)$.
\end{Proposizione}

\begin{Osservazione}[]
\label{oss-diagonalizzabile-esiste-base}
$\bm{f}$ \textbf{diagonalizzabile} $\iff$ \textbf{esiste una base} $\bm{\B}$ \textbf{di} $\bm{\V}$ \textbf{i cui vettori sono autovettori di} $\bm{f}$. \\
\NB $M_\B^\B(f)$ è una matrice diagonale le cui entrate sono gli autovalori corrispondenti agli vettori di $\B$.
\end{Osservazione}

\Dimostrazione \nopagebreak

\begin{enumerate}[label=\arabic*°]
    \item $f$ diagonalizzabile $\implies$ esiste una base $\B$ di $\V$ i cui vettori sono autovettori di $f$ \\
    Se $f$ è diagonalizzabile, allora sappiamo che esiste una base $\B = \{v_1, \ldots, v_n\}$ di $\V$ tale che 
    \[
    M_\B^\B(f) = \begin{bmatrix}
        \lambda_1 & 0 & \dots & 0 \\
        0 & \lambda_2 & \dots & 0 \\
        \vdots & \vdots & \ddots & \vdots \\
        0 & 0 & \dots & \lambda_n
    \end{bmatrix}
    \]

    Sappiamo che la j-esima colonna di $M_\B^\B(f)$ è pari a $X_\B(f(v_j))$, dunque troviamo che $f(v_j) = \lambda_j v_j$. Dunque $v_j$ è un autovettore di $f$. Dato che questo vale $\forall j \in \{1, \ldots, n\}$, troviamo che la base $\B$ è una base di $\V$ con tutti autovettori di $f$. $\blacksquare$
    \item $f$ diagonalizzabile $\impliedby$ esiste una base $\B$ di $\V$ i cui vettori sono autovettori di $f$ \\
    Sia $\B = \{v_1, \ldots, v_n\}$ la base di autovettori di $\V$. Sappiamo quindi che $\forall j \in \{1, \ldots, n\} \enspace f(v_j) = \lambda_j v_j$. Dunque $X_\B(f(v_j)) = (0,\ldots,0,\lambda_j,0,\ldots,0)$. Dunque
    \[
    M_\B^\B(f) = [X_\B(f(v_1)), \ldots, X_\B(f(v_n))] = \begin{bmatrix}
        \lambda_1 & 0 & \dots & 0 \\
        0 & \lambda_2 & \dots & 0 \\
        \vdots & \vdots & \ddots & \vdots \\
        0 & 0 & \dots & \lambda_n
    \end{bmatrix}
    \] 
    che è proprio una matrice diagonale. $\blacksquare$
\end{enumerate}

\begin{Proposizione}[]
\label{propos-matrice-diag-matrice-coniugata-diagonale}
Presa una matrice $A \in M_{mn}(\K)$, $\bm{A}$ \textbf{diagonalizzabile} $\iff$ $\bm{A}$ \textbf{coniugata alla matrice diagonale} $\bm{\Lambda}$.
\end{Proposizione}

\Dimostrazione \nopagebreak

\begin{enumerate}[label=\arabic*°]
    \item $A$ diagonalizzabile $\implies$ $A$ coniugata con la matrice diagonale $\Lambda$ \\
    Dato che $A$ è diagonalizzabile, sappiamo che esiste una base $\B = \{v_1, \ldots, v_n\}$ tale che $M_\B^\B(L_A)$ sia diagonale. Presa ora la base standard $S$, sappiamo che $A = M_S^S(L_A)$, e per l'Osservazione \ref{oss-matrice-diagonalizzata-M_C^C(f) = M M M} troviamo che $A = M_S^\B(Id) \cdot M_\B^\B(L_A) \cdot M_B^S(Id)$. Sappaimo che $M_\B^\B(L_A) = \Lambda$, e che $M_\B^S(Id) = M_S^\B(Id)^{-1}$, dunque, se poniamo $M_\B^S(Id) = \X$, troviamo che $A = \X^{-1} \cdot \Lambda \cdot \X$, ossia $A$ è coniugata a $\Lambda$. $\blacksquare$
    \item $A$ diagonalizzabile $\impliedby$ $A$ coniugata con la matrice diagonale $\Lambda$ \\
    Dato che $A$ è coniugata a $\Lambda$, sappiamo che esiste una matrice $\X$ tale che $A = \X \cdot \Lambda \cdot X^{-1}$. Poiché $\X$ è invertibile, sappiamo che le sue colonne sono vettori linearmente indipendenti, ed essendo queste $n$ colonne in uno spazio vettoriale di dimensione $n$, per l'Osservazione \ref{oss-dimen} questi vettori colonna formano una base. Dunque $\B = \{\X_1, \ldots, \X_n\}$ è una base di $\K^n$. Cerchiamo ora $M_\B^\B(L_A)$. \\
    Sappiamo che $M_\B^\B(L_A) = M_\B^S(L_A) \cdot M_S^S(L_A) \cdot M_S^\B(L_A) = M_S^\B(L_A)^{-1} \cdot A \cdot M_S^\B(L_A)$. Dato che $\X$ è la matrice che ha come colonne i vettori di $\B$, per la Proposizione \ref{propos-matrice-cambiamento-base-base-standard} che $\X = M_S^\B(L_A)$. Troviamo quindi che \mbox{$M_\B^\B(L_A) = \X^{-1} \cdot A \cdot \X = \X^{-1} \cdot (\X \cdot \Lambda \cdot \X^{-1}) \cdot \X = \Lambda$}, ossia $A$ è diagonalizzabile. $\blacksquare$
\end{enumerate}

\begin{Osservazione}[]
\label{oss-A diagonalizzabile con Lambda rende potenza facile}
Presa una matrice $A \in M_{nn}(\K)$ diagonalizzabile, con $A = G^{-1} \cdot \Lambda \cdot G$, $\bm{A^n = G^{-1} \cdot \Lambda^n \cdot G}$.
\end{Osservazione}

\Dimostrazione \nopagebreak

Dimostriamo questa proposizione per induzione.
\begin{itemize}
    \item Caso base: $n = 2$ \\
    In questo vogliamo quindi dimostrare che  $A^2 = G^{-1} \cdot \Delta^2 \cdot G$. Banalmente \mbox{$A^2 = (G^{-1} \cdot \Delta \cdot G)^2 = (G^{-1} \cdot \Delta \cdot G) \cdot (G^{-1} \cdot \Delta \cdot G) = G^{-1} \cdot \Delta \cdot (G \cdot G^{-1}) \cdot \Delta \cdot G = G^{-1} \cdot \Delta^2 \cdot G$}. $\blacksquare$
    \item Ipotesi induttiva \\
    Assumiamo l'osservazione vera per una potenza di $n-1$.
    \item Passo induttivo \\
    Per ipotesi induttiva \mbox{$A^n  = A^{n-1} \cdot (G^{-1} \cdot \Delta \cdot G) = (G^{-1} \cdot \Delta^{n-1} \cdot G) \cdot (G \cdot \Delta \cdot G^{-1})$}. Dato che \mbox{$(G^{-1} \cdot \Delta^{n-1} \cdot G) \cdot (G \cdot \Delta \cdot G^{-1}) = G^{-1} \cdot \Delta^{n-1} \cdot (G \cdot G^{-1}) \cdot \Delta \cdot G = G^{-1} \cdot \Delta \cdot G$}. $\blacksquare$
\end{itemize}

\Esempio \nopagebreak

Data la matrice \nopagebreak
\[ \bm{A} = 
\begin{bmatrixb}
1 & -1  \\
2 & 4 
\end{bmatrixb} \in M_{22}(\Q)
\]
trovare una \textbf{base} di $\Q$ che diagonalizzi $A$ e calcolare $\bm{A^{100}}$. \nopagebreak

Prendiamo la base $\B = \{(1,-1), (1,-2)\} = \{v_1, v_2\}$ di $\Q$. Sappiamo che per l'Osservazione \ref{oss-diagonalizzabile-esiste-base} (che vale anche per le matrici, dato che una matrice $A$ è diagonalizzabile solo se $L_A$ lo è) che se $v_1$ e $v_2$ sono autovettori di $A$, allora $A$ è diagonalizzabile e in particolare $M_\B^\B(\Q)$ è diagonale. Cerchiamo quindi di dimostrare che $v_1$ e $v_2$ sono autovettori di $A$. 
\[L_A(v_1) = L_A((1,-1)) = A \cdot \begin{bmatrix}
        1 \\ -1
    \end{bmatrix} = \begin{bmatrix}
    1 & -1 \\
    2 & 4    
    \end{bmatrix} \cdot \begin{bmatrix}
        1 \\
        -1
    \end{bmatrix}
    =
    \begin{bmatrix}
        2 \\ -2
    \end{bmatrix}= 2 v_1 \]

\[L_A(v_2) = L_A((1,-2)) = A \cdot \begin{bmatrix}
        1 \\ -2
    \end{bmatrix} = \begin{bmatrix}
    1 & -1 \\
    2 & 4    
    \end{bmatrix} \cdot \begin{bmatrix}
        1 \\
        -2
    \end{bmatrix}
    =
    \begin{bmatrix}
        3 \\ -6
    \end{bmatrix}= 3 v_2 \]

Dunque entrambi i vettori sono autovettori di $A$ ($v_1$ con autovalore $2$ e $v_2$ con autovalore $3$), e quindi la base $\bm{\B = \{(1,-1), (1,-2)\}}$ diagonalizza $A$. In particolare, troviamo che 
\[M_\B^\B(L_A) = 
\begin{bmatrix}
    2 & 0 \\
    0 & 3
\end{bmatrix}
\]

Per calcolare $A^{100}$ ci conviene ora ricordare la formula $A = M_S^S(L_A) = M_S^\B(Id) \cdot M_\B^\B(L_A) \cdot M_\B^\S(Id)$ vista nell'Osservazione \ref{oss-matrice-diagonalizzata-M_C^C(f) = M M M}. Conosciamo già $M_\B^\B(L_A)$, e sappiamo che 
\[
M_S^\B(Id) = [v_1, v_2] = \begin{bmatrix}
    1 & 1 \\
    -1 & -2
\end{bmatrix}
\]
Per calcolare ora $M_\B^S(Id)$ non ci basta che ricordare che $M_\B^S(Id) = M_S^\B(Id)$. Seguendo il procedimento descritto nella Proposizione \ref{propos-trovare-inversa}, troviamo che 
\[
M_\B^S(Id) = \begin{bmatrix}
    2 & 1\\
    -1 & -1
\end{bmatrix}
\]

Dunque
\[
A = \begin{bmatrix}
    1 & -1 \\
    2 & 4
\end{bmatrix} = \begin{bmatrix}
    1 & 1 \\
    -1 & -2
\end{bmatrix} \cdot \begin{bmatrix}
    2 & 0 \\
    0 & 3
\end{bmatrix} \cdot \begin{bmatrix}
    2 & 1 \\
    -1 & -1
\end{bmatrix} = G^{-1} \cdot \Lambda \cdot G
\]
Abbiamo visto nell'Osservazione \ref{oss-A diagonalizzabile con Lambda rende potenza facile} che $A^{100} = G^{-1} \cdot \Lambda^{100} \cdot G$, dunque troviamo che 
\[
\bm{A^{100}} = \begin{bmatrix}
    1 & 1 \\
    -1 & -2
\end{bmatrix} \cdot \begin{bmatrix}
    2^{100} & = \\
    0 & 3^{100}
\end{bmatrix}
\cdot
\begin{bmatrix}
    2 & 1 \\
    -1 & -1
\end{bmatrix} = \begin{bmatrixb}
    2^{101} - 3^{100} & 2^{100} - 3^{100} \\
    -2^{101} + 2 \cdot 3^{100} & -2^{100} + 2 \cdot 3^{100}
\end{bmatrixb}
\]

\begin{Definizione}[]
\label{def-e^A}
Presa una matrice $A \in M_{nn}(\R)$,
\[
\bm{e^A := 1_n + A + \frac{A^2}{2!} + \ldots = \sum_{n = 0}^\infty} \frac{A^n}{n!}
\]
Questa serie è la serie di Taylor dell'esponenziale senza resto applicata alle matrici.
\end{Definizione}

\begin{Proposizione}[]
\label{propos-e^A+B}
Prese due matrici $A,B \in M_{nn}(\K)$, $\bm{e^{A+B} = e^A \cdot e^B} \iff \bm{AB = BA}$.
\end{Proposizione}

\Dimostrazione \nopagebreak

Abbiamo visto nella Definizione \ref{def-e^A} che $e^{A+B} = \sum_{n = 0}^\infty \frac{(A+B)^n}{n!}$, ma per comodità in questa dimostrazione ci fermeremo a dire che $e^{A+B} =  1_n + (A + B) + \frac{(A+B)^2}{2!}$. 
Dato che $(A+B)^2 = A^2 + B^2 + AB + BA$, troviamo che $e^{A+B} = 1_n + A + B + \frac{A^2}{2} + \frac{B^2}{2} + \frac{AB}{2} + \frac{BA}{2}$. Dato che $e^A = 1_n + A + \frac{A^2}{2}$ e $e^B = 1_n + B + \frac{B^2}{2}$, troviamo che $e^A \cdot e^B = 1_n + B + \frac{B^2}{2} + A + AB + \frac{AB^2}{2} + \frac{A^2}{2} + \frac{A^2 B}{2} + \frac{A^2B^2}{4}$. Dato che però stiamo considerando solo i termini fino al secondo grado, possiamo troncare quelli di terzo e quarto trovando che $e^A \cdot e^B = 1_n + A + B + AB + \frac{A^2}{2} + \frac{B^2}{2}$.
Dunque, affinché $e^{A+B} = e^A \cdot e^B$, è necessario che $\frac{AB}{2} + \frac{BA}{2} = AB$, e questo è vero se e solo se $AB = BA$. $\blacksquare$

\begin{Proposizione}[]
\label{propos-e^A-diagonalizzazione}
Presa una matrice $A \in M_{nn}(\K)$ diagonalizzabile tale che $A = G^{-1} \cdot \Delta \cdot G$, $\bm{e^A = G^{-1} \cdot e^\Delta \cdot G}$.
\end{Proposizione}

\Dimostrazione \nopagebreak

Anche in questa dimostrazione per comodità ci fermeremo a dire che $e^A = 1_n + A + \frac{A^2}{2}$. Dato che $A = G^{-1} \cdot \Delta \cdot G$, troviamo che $e^A = G^{-1} \cdot 1_n \cdot G + G^{-1} \cdot \Delta \cdot G + \frac{G^{-1} \cdot \Delta^2 \cdot G}{2}$, dato che $1_n = G^{-1} \cdot \Delta \cdot G$, e dunque $e^A = G^{-1} \cdot (1_n + \Delta + \frac{\Delta^2}{2}) \cdot G$. Dato che $e^\Delta = 1_n + \Delta + \frac{\Delta^2}{2}$, concludiamo che $e^A = G^{-1} \cdot e^\Delta \cdot G$. $\blacksquare$  

\Esempio \nopagebreak

Presa la matrice \nopagebreak
\[ \bm{A} = 
\begin{bmatrixb}
    1 & -1 \\
    2 & 4
\end{bmatrixb}
\]
trovare $\bm{e^A}$. \nopagebreak

Questa matrice è la stessa dell'esempio precedente, dunque già sappiamo che
\[
A = \begin{bmatrix}
    1 & 1 \\
    -1 & -2
\end{bmatrix} \cdot \begin{bmatrix}
    2 & 0 \\
    0 & 3
\end{bmatrix} \cdot \begin{bmatrix}
    2 & 1 \\
    -1 & -1
\end{bmatrix} = G^{-1} \cdot \Delta \cdot G
\]
È facile quindi concludere che 
\[
\bm{e^A} = G^{-1} \cdot e^\Delta \cdot G = \begin{bmatrixb}
    2e^2 - e^3 & e^2 - e^3 \\
    -2e^2 + 2e^3 & -e^2 + 2e^3
\end{bmatrixb}
\]

\chapter{Determinante}

\section{Definizioni}

\begin{Definizione}[Sottomartice complementare]
\label{def-sottomatrice-complementare}
Presa una matrice $A \in M_{nn}(\K)$ e presi gli indici $i,j \in \{1, \ldots, n\}$, la \textbf{sottomatrice complementare} associata a $(i,j)$ è la sottomatrice $\bm{A^i_j} \in M_{(n-1)(n-1)}$ ottenuta da $A$ eliminando la riga $i$ e la colonna $j$.
\end{Definizione}

\begin{Definizione}[Determinante]
\label{def-determinante}
Preso un anello commutativo con unità $R$, definiamo $\forall n \in \N_+$ l'applicazione $\bm{\Det_n: M_nn(R) \longrightarrow R}$ così definita:

\[
\Det_n(A) =
\begin{cases}
\bm{a} & \text{se } \bm{n = 1} \\
\bm{\sum_{j = 1} (-1)^{n+j} \cdot a_{nj} \cdot Det_{n-1}(A^n_j)} & \text{se } \bm{n \neq 1} 
\end{cases}
\]

\NB Il determinante di $A$ può essere indicato anche con i simboli $\bm{\Det(A)}, \bm{\det A}$ o $\bm{|A|}$.
\end{Definizione}

\Esempio \nopagebreak

Calcolare il \textbf{determinante di una matrice} $\bm{2 \times 2}$. \nopagebreak

Una matrice $2 \times 2$ è una matrice del tipo 
\[ A = 
\begin{bmatrix}
    a_{11} & a_{12} \\
    a_{21} & a_{22}
\end{bmatrix}
\]

Dunque $\bm{\det A} = -a_{21} \cdot a_{12} + a_{22} \cdot a_{11} = \bm{a_11 \cdot a_{22} - a_{12} \cdot a_{21}}$.

\section{Proprietà del determinante}

\begin{Proprieta}[Proprietà algebriche del determinante]
\label{prop-determinante}
Presa una matrice $A \in M_{nn}(R)$
\begin{enumerate}
    \item Se $A$ ha \textbf{una colonna nulla} allora $\bm{\det A = 0}$
    \item Se $a_{ij} = 0 \enspace \forall i < j$, ossia tutti i valori sopra la diagonale principale sono $0$, allora 
    \vspace{-0.3cm}
    \[\bm{\det A = \prod_{i = 1}^n a_{ii}} = a_{11} \cdot \ldots \cdot a_{nn}\]
    \item $\bm{\det 1_n = 1}$
    \item Presi $\lambda, \mu \in R$ e presi $B,C \in M_{n1}(R)$, 
    \vspace{-0.3cm}
    \[\bm{\det [A_1, \ldots, \lambda B + \mu C, \ldots, A_n] = \lambda \det[A_1, \ldots, B, \ldots, A_n] + \mu \det[A_1, \ldots, C, \ldots, A_n]}\]
    \vspace{-0.9cm}
    \item Il determinante è \textbf{alternante nelle colonne}, ossia se $A$ ha \textbf{due colonne uguali} allora $\bm{\det A = 0}$.
\end{enumerate}
\end{Proprieta}

\Dimostrazione \nopagebreak

\begin{enumerate}
    \item Dimostriamo questa proprietà per induzione.
    \begin{itemize}
        \item Caso base: $n = 1$ \\
        Allora $A = (0)$ e dunque, per la stessa definizione del determinante, $\det A = 0$. $\blacksquare$
        \item Ipotesi induttiva \\
        Assumiamo la proprietà vera per tutte le matrici $(n-1) \times (n-1)$.
        \item Passo induttivo \\
        Sappiamo che $A = [A_1, \ldots, \undzero, \ldots, A_n]$, dove la colonna $j_0$ di $A$ è proprio la colonna nulla. Se proviamo a calcolare il determinante di $A$, troviamo che \mbox{$\det A = \sum_{j = 1}^n (-1)^{n+j} \cdot a_{nj} \cdot \det(A_j^n)$}. Sappiamo che $\forall j \neq j_0$, $(-1)^{n+j} \cdot a_{nj} \cdot \det(A_j^n) = 0$, dato che $A_j^n$ non elimina la colonna nulla ($j \neq j_0$), e quindi, per ipotesi induttiva, $\det A_j^n = 0$. Dunque troviamo che $\det A = (-1)^{n+j_0} \cdot a_{n j_0} \cdot \det(A^n_{j_0})$, e dato che $a_{n j_0} = 0$, troviamo che $\det A = 0$. $\blacksquare$
    \end{itemize}
    \item Dimostriamo questa proprietà per induzione.
    \begin{itemize}
        \item Caso base: $n = 1$ \\
        Allora $A = (a)$ e quindi per la stessa definizione del determinante $\det A = a$. $\blacksquare$
        \item Ipotesi induttiva \\
        Assumiamo la proprietà vera per tutte le matrici $(n-1) \times (n-1)$
        \item Passo induttivo \\
        Sappiamo che $\det A = \sum_{j = 1}^n (-1)^{n+j} \cdot a_{nj} \cdot \det(A^n_j)$. Guardando $A_j^n$, è facile notare che $\forall j \neq n$ questa matrice abbia sempre una colonna nulla, ossia l'ultima. Dunque, per la proprietà 1) possiamo dire che $\det A_j^n = 0 \enspace \forall j \neq n$. Dunque rimane \mbox{$\det A = (-1)^{n+n} \cdot a_{nn} \cdot \det A_n^n$}. Dato che $A_n^n$ è una matrice $(n-1) \times (n-1)$, per ipotesi induttiva possiamo dire che $\det A_n^n = a_{11} \cdot \ldots \cdot a_{(n-1)(n-1)}$, dunque concludiamo che \mbox{$\det A = a_{11} \cdot \ldots \cdot a_{nn}$}. $\blacksquare$
    \end{itemize}
    \item Dato che $1_n$ ha ovviamente tutti i valori sopra la diagonale principale $0$, per 2) troviamo che \mbox{$\det 1_n = 1 \cdot \ldots \cdot 1 = 1$}. $\blacksquare$
    \item Questa proprietà verrà dimostrata quando verranno introdotte le applicazioni multilineari.
    \item Questa proprietà verrà dimostrata quando verranno introdotte le applicazioni multilineari.
\end{enumerate}


\begin{Proprieta}[Operazoni elementari sulle colonne e determinante]
\label{prop-oper-elem-determinante}
\begin{enumerate}
    \item \textit{Scambio tra due colonne}: il determinante \textbf{cambia segno}
    \item \textit{Aggiungere a una colonna il multiplo di un'altra}: il determinante \textbf{non cambia}
    \item \textit{Moltipicare una colonna per uno scalare non nullo}: il determinante \textbf{viene moltiplicato per lo stesso scalare}
\end{enumerate}
\end{Proprieta}

\Dimostrazione \nopagebreak

\begin{enumerate}
    \item Per comodità vedremo lo scambio tra le prime due colonne, ma il discorso è identico per qualsiasi coppia di colonne di $A$. Per la proprietà 4) sappiamo che
    \begin{align*}
    \det [A_1 + A_2, A_2 + A_1, \ldots, A_n] &= \det [A_1 + A_2, A_2, \ldots, A_n] \\
     &+ \det [A_1 + A_2, A_1, \ldots, A_n]  \\
     &= \det [A_1, A_2, \ldots, A_n] + \det [A_2, A_2, \ldots, A_n] \\ 
     &+ \det [A_1, A_1, \ldots, A_n] + \det [A_2, A_1, \ldots, A_n]
    \end{align*}
    Per la proprietà 5) però sappiamo anche che $\det [A_1 + A_2, A_2 + A_1, \ldots, A_n] = 0$  dato che ha due colonne uguali, e per lo stesso motivo anche \mbox{$\det[A_2,A_2, \ldots, A_n] = \det [A_1, A_1, \ldots, A_n] = 0$}. Dunque troviamo che \mbox{$\det[A_1, A_2, \ldots, A_n] = - \det[A_2, A_1, \ldots, A_n]$}. $\blacksquare$
    \item Per comodità anche qui vediamo questa operazione eseguita sulla prima colonna aggiungendo la seconda colonna, ma il discorso è identico per qualsiasi coppia di colonne di $A$.  Per la proprietà 4) sappiamo che \mbox{$\det [A_1 + \lambda A_2, A_2, \ldots, A_n] = \det [A_1, A_2, \ldots, A_n] + \lambda \det [A_2, A_2, \ldots, A_n]$}, e per la proprietà 5) sappiamo che $\det [A_2, A_2, \ldots, A_n] = 0$. Dunque concludiamo che \mbox{$\det [A_1, A_2, \ldots, A_n] = \det [A_1 + \lambda A_2, \ldots, A_n]$}. $\blacksquare$
    \item Per la proprietà 4) troviamo che \mbox{$\det [A_1, \ldots, \lambda A_j, \ldots, A_n] = \lambda \det [A_1, \ldots, A_j, \ldots, A_n]$. $\blacksquare$}
\end{enumerate}

\begin{Osservazione}[]
\label{oss-det diverso da 0 A invertibile}
Presa una matrice $A \in M_{nn}(\K)$, $\bm{\det A \neq 0} \iff \bm{A} \textbf{ è invertibile}$. \\
\NB È importante che $A$ abbia entrate in un campo e non in un semplice anello perché affinché abbia senso dire che $A$ è invertibile, è necessario trovarsi in un insieme che ammetta l'inverso moltiplicativo di ogni suo elemento. 
\end{Osservazione}

\Dimostrazione \nopagebreak

\begin{enumerate}[label=\arabic*°]
    \item $\det A \neq 0 \implies A$ è invertibile \\
    Partiamo col trasformare $A$ in una matrice a scala per colonne. Sappiamo che per farlo dobbiamo fare delle operazioni elementari sulle sue colonne, ma come abbiamo visto queste operazioni possono al massimo cambiare il segno del determinante o moltiplicarlo per uno scalare. Dunque troviamo che $\det A = \mu \det S$, con $\mu \neq 0$. Dato che $\det A \neq 0$, anche $\det S \neq 0$. Dato che $S$ è una matrice a scala per colonne, sicuramente tutti gli elementi sopra la diagonale principale sono diversi da $0$, dunque $\det S = s_{11} \cdot \ldots \cdot s_{nn}$. Dato che $\det S \neq 0$, troviamo che $s_{ii} \neq 0 \enspace \forall i \in \{1, \ldots, n\}$. Dunque $S$ è una matrice $n \times n$ con $n$ pivot, ossia $rg(S) = n$. Dato che $rg(A) = rg(S) = n$, per la Proposizione \ref{propos-invertibilità-matrice} troviamo che $A$ è invertibile. $\blacksquare$
    \item $\det A \neq 0 \impliedby A$ è invertibile \\
    Se $A$ è invertibile, possiamo fare diverse operazioni elementari sulle colonne di $A$ per arrivare alla matrice unità $1_n$. Abbiamo già detto che le operazioni elementari sulle colonne moltiplicano il determinante di $A$ per uno scalare $\mu \neq 0$. Dunque $\det 1_n = \mu \det A$. Dato che, per la proprietà algebrica 3), $\det 1_n = 1$, troviamo che $\det A = \frac{1}{\mu} \neq 0$. $\blacksquare$
\end{enumerate}

\begin{Proposizione}[Corollario]
\label{corol-det-invertibilita}
Presa una matrice $A \in M_{nn}(\K)$, $\bm{\det A = 0} \iff \bm{\ker A \neq \{\undzero\}} \iff \bm{rg(A) \neq n} \iff$ \mbox{\textbf{i vettori colonna di} $\bm{A}$ \textbf{sono linearmente dipendenti}}.
\end{Proposizione}

\section{Applicazioni multilineari}
\begin{Definizione}[Applicazioni multilineari]
\label{def-app-multilineari}
Presi gli spazi vettoriali $\V_1, \ldots, \V_m$ su $\K$, l'applicazione
\begin{align*}
\bm{\Phi: \V_1 \times \ldots \times \V_m} &\longrightarrow \bm{\K}  \\
\bm{(v_1, \ldots, v_m)} &\longmapsto \bm{\Phi(v_1, \ldots, v_m)}
\end{align*}
è detta \textbf{multilineare} se questa risulta lineare rispetto a ciascuno degli $m$ valori in input. In altre parole, se \textit{congeliamo} $(\bar{v}_1, \ldots, \bar{v}_{j-1}, \bar{v}_{j+1}, \ldots, \bar{v}_n)$ e prendiamo l'applicazione
\begin{align*}
f:V_j &\longrightarrow \K \\
v_j &\longmapsto \Phi(\bar{v}_1, \ldots, \bar{v}_{j-1}, v_j, \bar{v}_{j+1}, \ldots, \bar{v}_n) 
\end{align*}
questa è \textbf{lineare} $\bm{\forall j \in \{1, \ldots, m\}}$.
\end{Definizione}

Vediamo degli esempi con quest'applicazione dove viene preso come spazio vettoriale $\K$ tre volte:
\begin{align*}
\Phi: \K^3 &\longrightarrow \K \\
(x_1, x_2, x_3) & \longmapsto \Phi(x_1, x_2, x_3)
\end{align*}

\begin{itemize}
    \item Se $\Phi(x_1, x_2, x_3) = x_1^3$, quest'applicazione \textbf{non è multilineare}. Infatti, se prendiamo l'applicazione
    \begin{align*}
    f:V_2 &\longrightarrow \K \\
    (1, x_2, 1) &\longmapsto \Phi(1, x_2, 1) = 1
    \end{align*}
    quest'applicazione non è lineare, dato che $f(\undzero) \neq \undzero$.
    \item Se prendiamo invece l'applicazione $\Phi(x_1, x_2, x_3) = 5 \cdot x_1 \cdot x_2 \cdot x_3$, quest'applicazione è \textbf{multilineare}. Infatti, se blocchiamo $\bar{x}_2$ e $\bar{x}_3$, troviamo che $\Phi(x_1, \bar{x}_2, \bar{x}_3) = x_1 \cdot 5\bar{x}_2 \bar{x}_3 = Cx_1$, dove $C = 5 \bar{x}_2 \bar{x}_3$ è una costante. È facile notare come quest'applicazione è lineare, e questo proceidmento può essere ripetuto anche per le variabili $x_2$ e $x_3$. Dunque $\Phi$ è effettivamente multilineare.  
\end{itemize}

\begin{Proposizione}[]
\label{propos-applicaz-multilineari-tutte-uguali}
Presi gli spazi vettoriali $\K^d$ e $\K^e$, tutte le applicazioni su questi due spazi vettoriali sono del tipo:
\begin{align*}
\bm{\Phi: \K^d \times \K^e} &\longrightarrow \bm{\K} \\
\bm{(X, Y)} &\longmapsto \bm{\sum_{i = 1}^d \sum_{j = 1}^e a_{ij} x_i y_i} \quad \text{con } a_{ij} \in \K
\end{align*}
\end{Proposizione}

\Dimostrazione \nopagebreak

Dimostriamo prima che questa applicazione sia multilinare e poi dimostriamo che tutte le applicazioni multilineari su $\K^d$ e $\K^e$ cadano in questo caso.\\ 
Dobbiamo quindi dimostrare che qunado blocchiamo la $Y$ l'applicazione $f(X) = \sum_{i = 1}^d x_i (\sum_{j = 1}^e a_{ij} \bar{y}_j)$ è lineare e che allo stesso modo $f(Y) = \sum_{j = 1}^e y_j (\sum_{i = 1}^d a_{ij} \bar{x}_i)$ è lineare. Dato che $\sum_{j = 1}^e a_{ij} \bar{y}_j$ e $\sum_{i = 1}^d a_{ij} \bar{x}_i$ sono delle somme di costanti, dobbiamo semplciemente dimostrare che $f(X) = \sum_{i = 1}^d C_i x_i$ e che $f(Y) = \sum_{j = 1}^e D_j y_j$ sono lineari, e questo è evidente. Dunque $\Phi$ è un'applicazione multilineare. 

Per dimostrare che tutte le applicazioni multilineari su $\K^d$ e $\K^e$ cadono in questo caso, dobbiamo dimostrare che questo è l'unico tipo di applicazioni multilineari. \\
Se prendiamo una base $\B = \{v_1, \ldots, v_d\}$ di $\K^d$ e una base $\Bs = \{w_1, \ldots, w_e\}$ di $\K^e$, troviamo che $X = \sum_{i = 1}^d x_i v_i$ e $Y = \sum_{j = 1}^e y_j w_j$. Dunque $\Phi(X, Y) = \Phi(\sum_{i = 1}^d x_i v_i, \sum_{j = 1}^e y_j w_j)$. Dobbiamo quindi dimostrare che $\Phi(X, Y) = \sum_{i = 1}^d \sum_{j = 1}^e x_i y_j \Phi(v_i, w_j)$. \\
Espandendo $X$ troviamo che $\Phi(X,Y) = \Phi(x_1 v_1 + \ldots + x_d v_d, Y) =  \sum_{i = 1}^d x_i \Phi(v_i, Y)$. Ripetendo lo stesso passaggio espandendo $Y$, troviamo proprio che \mbox{$\Phi(X, Y) = \sum_{i = 1}^d \sum_{j = 1}^e x_i y_j \Phi(v_i, w_j)$}, come volevamo dimsotrare.

Possiamo dunque concludere che qualsiasi applicazione multilineare può essere ricondotta alla forma $\sum_{i = 1}^d \sum_{j = 1}^e a_{ij} x_i y_j$. $\blacksquare$

\begin{Proposizione}[Corollario]
\label{corol-app-multilin}
Tutte le applicazioni multilineari su $m$ spazi vettoriali $\K^{d_1}, \ldots, \K^{d_m}$ sono del tipo
\begin{align*}
\bm{\Phi: \K^{d_1} \times \ldots \times \K^{d_m}} &\longrightarrow \bm{\K} \\
\bm{(x_1, \ldots, x_m)} &\longmapsto \bm{\sum_{i_1, \ldots, i_m} a_{i_1, \ldots, i_m} x_1(i_1) \ldots x_m(i_m)}
\end{align*}

Nonostante sembri una formula assurda, in realtà questo è semplicemente ciò che abbiamo visto con due spazi vettoriali ma con $m$ spazi vettoriali. È importante anche notare come tutti i vettori si trovino solo al primo grado.\\
\NB $(i_1), \ldots, (i_m)$ rappresentano gli indici dei vettori $x$ da prendere.
\end{Proposizione}

\begin{Proposizione}[]
\label{propos-multilineareità-di-un-applicazione}
Presa un'applicazione $\Phi$ su $m$ spazi vettoriali, \\
\[
\bm{\Phi} \textbf{ è multilineare} \iff
\begin{aligned}
\bm{\Phi(x_1, \ldots, \alpha x + \beta y, \ldots, x_m)} &= \bm{\alpha \Phi(x_1, \ldots, x, \ldots, x_m)} \\
&+ \bm{\beta \Phi(x_1, \ldots, y, \ldots, x_m)}
\end{aligned}
\]
\end{Proposizione}

\vspace{2cm}

\Dimostrazione \nopagebreak

\begin{enumerate}[label=\arabic*°]
    \item \leavevmode
    $\Phi \textbf{ è multilineare} \implies
    \begin{aligned}[t]
    \Phi(x_1, \ldots, \alpha x + \beta y, \ldots, x_m) &= \alpha \Phi(x_1, \ldots, x, \ldots, x_m) \\
    &+ \beta \Phi(x_1, \ldots, y, \ldots, x_m)
    \end{aligned}$
    \vspace{0.2cm}

    Dato che $\Phi$ è multilineare, sappiamo che se prendiamo un indice $k \in \{1, \ldots, m\}$ e congeliamo tutti i vettori di entrata diversa dall'indice $k$, troviamo un'applicazione $f$ lineare. Dato che $f(\alpha x + \beta y) = \alpha f(x) + \beta f(y)$, possiamo concludere che
    \mbox{$\Phi(x_1, \ldots, \alpha x + \beta y, \ldots, x_m) = \alpha \Phi(x_1, \ldots, x, \ldots, x_m) + \beta \Phi(x_1, \ldots, y, \ldots, x_m)$}. $\blacksquare$

    \item \leavevmode
    $\Phi \textbf{ è multilineare} \impliedby
    \begin{aligned}[t]
    \Phi(x_1, \ldots, \alpha x + \beta y, \ldots, x_m) &= \alpha \Phi(x_1, \ldots, x, \ldots, x_m) \\
    &+ \beta \Phi(x_1, \ldots, y, \ldots, x_m)
    \end{aligned}$
    \vspace{0.2cm}

    Riprendendo la stessa applicazione $f$, sappiamo quindi che $f(\alpha x + \beta y) = \alpha f(x) + \beta f(y)$. Questa altro non è che la definizione di applicazione lineare, dunque $f$ è lineare, e dunque $\Phi$ è multilineare. $\blacksquare$
\end{enumerate}

\begin{Definizione}[Insieme delle applicazioni multilineari]
\label{def-insieme-app-multilineari}
Chiamiamo $\bm{\Multilin (\V_1 \times \ldots \times \V_m, \K)}$ \textbf{l'insieme delle applicazioni multilineari} che vanno da $\V_1, \ldots, \V_m$ a $\K$.
\end{Definizione}

Definiamo ora le operazioni di somma e prodtto per scalare per questo insieme.

\bigskip

    {\Large\textbf{Somma}:}
    \begin{align*}
        \hspace{0.4cm}\Phi + \Psi: \V_1 \times \ldots \times \V_m &\longrightarrow \K \\
        \hspace{0.4cm}(v_1,\ldots,v_m) &\longmapsto \Phi(v_1, \ldots, v_m) + \Psi(v_1, \ldots, v_m)
    \end{align*}


    {\Large\textbf{Prodotto per scalare}:}
    \begin{align*}
        \lambda \Phi: \V_1 \times \ldots \times \V_m &\longrightarrow \K \\
        (v_1,\ldots,v_m) &\longmapsto \lambda \Phi(v_1, \ldots, v_m)
    \end{align*}

\bigskip

\begin{Proposizione}[]
\label{propos-multilin-spazio-vett}
Con le operazioni di somma e prodotto per scalare appena descritte, $\bm{\Multilin}$ è uno \textbf{spazio vettoriale} su $\K$.
\end{Proposizione}
 
\Dimostrazione \nopagebreak

Sappiamo già che $\Multilin$ è chiuso rispetto a somma e prodotto per la Proposizione \ref{propos-multilineareità-di-un-applicazione}, dunque dobbiamo solo verificare che $\undzero$ sia in $\Multilin$. \\
Per dimostrare che $\undzero \in \Multilin$ dobbiamo dimostare che 
\[
\undzero(v_1, \ldots, \alpha x + \beta y, \ldots, v_m) = \alpha \undzero (v_1, \ldots, x, \ldots, v_m) + \beta \undzero(v_1, \ldots, y, \ldots, v_m)
\]
Ma questo è evidente, dato che l'applicazione nulla $\undzero$ manda tutto in $\undzero$, dunque questa uguaglianza si semplifica in $\undzero = \undzero$. Dunque $\undzero \in \Multilin$. \\
Possiamo quindi concludere che $\Multilin$ è uno spazio vettorialie. $\blacksquare$

\begin{Definizione}[Applicazioni multilineari alternanti]
\label{def-app-multilin-altern}
Un'applicazione multilineare è detta \textbf{alternante} se l'applicazione \textbf{si annulla se due dei suoi argomenti sono uguali}.
\end{Definizione}

\begin{Definizione}[Insieme delle applicazioni multilineari alternanti]
\label{def-insieme-app-multilin-altern}
Chiamiamo $\bm{\Alt(\V^m, \K)} \subset \Multilin(\V^m, \K)$ \textbf{l'insieme delle applicazioni multilineari alternanti} che vanno da $\V^m$ a $\K$.
\end{Definizione}

\begin{Osservazione}[]
\label{oss-inversa-app-multil-alt}
Presa un'applicazione multilineare alternante $\Phi$ che va da $\V_1 \times \ldots \times \V_m$ a $\K$, \mbox{$\bm{\Phi(v_1, \ldots, v_m) = -\Phi(v_2, v_1, \ldots, v_m)}$}.
\end{Osservazione}

\Dimostrazione \nopagebreak

Dato che $\Phi$ è alternante, sappiamo che $\Phi(v_1 + v_2, v_2 + v_1, \ldots, v_m) = \undzero$. Essendo anche multilineare, troviamo che 
\begin{align*}
\Phi(v_1 + v_2, v_2 + v_1, \ldots, v_m) &= \Phi(v_1, v_2 + v_1, \ldots, v_m) + \Phi(v_2, v_2 + v_1, \ldots, v_m) \\
 &= \Phi(v_1, v_2, \ldots, v_m) + \Phi(v_1, v_1, \ldots, v_m) \\
 &+ \Phi(v_2, v_2, \ldots, v_m) + \Phi(v_2, v_1, \ldots, v_m) 
\end{align*}

\vspace{-0.5cm}
Dato che $\Phi$ è alternante, però, $\Phi(v_1, v_1, \ldots, v_m) = \Phi(v_2, v_2, \ldots, v_m) = \undzero$, dunque troviamo che 
\begin{align*}
&\Phi(v_1, v_2, \ldots, v_m) + \Phi(v_2, v_1, \ldots, v_m) = 0 \\
&\Phi(v_1, v_2, \ldots, v_m) = - \Phi(v_2, v_1, \ldots, v_m) \quad \blacksquare
\end{align*}

\begin{Proposizione}[]
\label{propos-alt-spaz-vett}
Con le operazioni di somma e prodotto per scalare descritte per lo spazio vettoriale $\Multilin$, anche $\bm{\Alt(\V^m, \K)}$ è uno \textbf{spazio vettoriale} su $\K$.
\end{Proposizione}

La dimostrazione è pressoché identica alla dimostrazione per la Proposizione \ref{propos-multilin-spazio-vett}, perciò è omessa.

\section{Determinante come applicazione multilineare}

\begin{Proposizione}[]
\label{propos-determinante-app-multilineare}
Il \textbf{determinante} è un'\textbf{applicazione multilineare alternante} ($\det A \in \Alt((\K^n)^n, \K)$).
\end{Proposizione}

\Dimostrazione \nopagebreak

Per dimostrare questa proposizione dimostriamo prima che $\det A \in \Multilin((\K^n)^n, \K)$ e poi che $\det A \in \Alt((\K^n)^n, \K)$.

\begin{enumerate}[label=\arabic*°]
    \item $\det A \in \Multilin((\K^n)^n, \K)$ \\
    Per la Proposizione \ref{propos-multilineareità-di-un-applicazione}, dobbiamo dimostrare che
    \begin{align*}
    \det [A_1, \ldots, \lambda B_j + \mu C_j, \ldots, A_n] &= \lambda \det [A_1, \ldots, B_j, \ldots, A_n] + \mu \det [A_1, \ldots, C_j, \ldots, A_n] 
    \end{align*}
    Dimostriamolo per induzione.
    \begin{itemize}
        \item Caso base: $n = 1$ \\
        $\det [\lambda b + \mu c] = \lambda b + \mu c = \lambda \det [b] + \mu \det[c]$. $\blacksquare$
        \item Ipotesi induttiva \\
        Assumiamo la proposizione vera per tutte le matrici $(n-1) \times (n-1)$.
        \item Passo induttivo \\
        Prendiamo la matrice $A = [A_1, \ldots, \lambda B_j + \mu C_j, \ldots, A_n]$ e calcoliamone il determinante:
        \begin{align*}
        \det A &= \sum_{i = 1}^n (-1)^{n+1} \cdot a_{ni} \cdot \det A^n_i  \\
         &= (-1)^{n + j} \cdot a_{nj} \cdot \det A^n_j + \sum_{i \neq j} (-1)^{n+i} \cdot a_{ni} \cdot \det A^n_i \\
         &= (-1)^{n+j} \cdot (\lambda b_j + \mu c_j) \det A^n_j + \sum_{i \neq j} (-1)^{n+i} \cdot a_{ni} \cdot \det A^n_i \\
        \end{align*}
        Dato che $A^n_i$ è una matrice $(n-1) \times (n-1)$, possiamo applicare l'ipotesi induttiva e trovare quindi che $\det A^n_i = \lambda \det B^n_i + \mu \det C^n_i$, dove $B$ e $C$ non sono altro che la matrice $A$ ma con la colonna j-esima pari rispettivamente a $B$ e $C$. Troviamo quindi che
        \begin{align*}
        \det A &= (-1)^{n+j} \cdot (\lambda b_j + \mu c_j) \det A^n_j + \sum_{i \neq j} (-1)^{n+i} \cdot a_{ni} \cdot (\lambda \det B^n_i + \mu \det C^n_i) \\
         &= (-1)^{n+j} \cdot (\lambda b_j) \det A^n_j + \sum_{i \neq j} (-1)^{n+i} \cdot a_{ni} \cdot \lambda \det B^n_i \\
         &+ (-1)^{n+j} \cdot (\mu c_j) \det A^n_j + \sum_{i \neq j} (-1)^{n+i} \cdot a_{ni} \cdot \mu \det C^n_i \\
         &= \lambda \det B + \mu \det C \quad \blacksquare 
        \end{align*}
    \end{itemize}

    \item $\det A \in \Alt((\K^n)^n, \K)$ \\
    Per comodità, dimostreremo che se due colonne consecutive sono uguali allora il determinante di tale matrice è nullo, ma il discorso è analogo con qualsiasi coppia di indici. Facciamolo per induzione.
    \begin{itemize}
        \item Caso base: $n = 2$ \\
        In questo caso abbiamo quindi
        \[ A = 
        \begin{bmatrix}
        a & b \\
        a & b
        \end{bmatrix}
        \]
        Dunque $\det A = ab - ab = 0$. $\blacksquare$
        \item Ipotesi induttiva \\
        Assumiamo la proposizione vera per tutte le matrici $(n-1) \times (n-1)$.
        \item Passo induttivo \\
        Prendiamo la matrice $A$ le cui colonne $A_j$ e $A_{j+1}$ sono uguali e calcoliamone il determinante. 
        \begin{align*}
        \det A &= \sum_{i = 1}^n (-1)^{n+i} \cdot a_{ni} \cdot \det A^n_i \\
         &= (-1)^{n+j} \cdot a_{nj} \cdot \det A^n_{j+1} + (-1)^{n+j+1} \cdot a_{n(j+1)} \cdot \det A^n_(j+1)  \\
         &+ \sum_{i \neq j, j+1} (-1)^{n+i} \cdot a_{ni} \cdot \det A^n_i
        \end{align*}
        Dato che $A_j = A_{j+1}$, sappiamo che $a_{nj} = a_{n(j+1)}$ e che $\det A^n_j = \det A^n_{j+1}$, e dato che $(-1)^{n+j} = - (-1)^{n+j+1}$, troviamo che
        \[
        (-1)^{n+j} \cdot a_{nj} \cdot \det A^n_j + (-1)^{n+j+1} \cdot a_{n(j+1) \cdot \det A^n_{j+1} = 0}
        \] 
        Infine, dato che $A^n_i$ è una matrice $(n-1) \times (n-1)$, troviamo che $\det A^n_i = 0$ sempre, dunque $\det A = 0$. $\blacksquare$
    \end{itemize}
\end{enumerate}

\begin{Proposizione}[]
\label{propos-ev_B isomorfismo}
Preso uno spazio vettoriale $\V$ finitamente generato su $\K$, con $\dim \V = n$ e con base \mbox{$\B = \{\bar{v}_1, \ldots, \bar{v}_n\}$}, l'applicazione
\begin{align*}
\bm{ev_\B: \Alt(\V^n, \K)} &\longrightarrow \bm{\K} \\
\bm{\Phi} &\longmapsto \bm{\Phi(\bar{v}_1, \ldots, \bar{v}_n)} 
\end{align*}
è un \textbf{isomorfismo di spazi vettoriali}.
\end{Proposizione}

\Dimostrazione \nopagebreak

Per l'Osservazione \ref{oss-def-isomorfismo}, ci basta dimostrare che $ev_\B$ è un'applciazione lineare e biunivoca.
\begin{enumerate}[label=\arabic*°]
    \item $ev_\B$ è lineare \\
    Presi $\lambda, \mu \in \K$ e $\Phi, \Psi \in \Alt(\V^n, \K)$, 
    \vspace{-0.4cm}
    \begin{align*}
    ev_\B(\lambda \Phi, \mu \Psi) &= (\lambda \Phi + \mu \Psi)(\bar{v}_1, \ldots, \bar{v}_n) = \lambda \Phi(\bar{v}_1, \ldots, \bar{v}_n) + \mu \Psi(\bar{v}_1, \ldots, \bar{v}_n) \\
    &= \lambda ev_\B (\Phi) + \mu ev_\B(\Psi)
    \end{align*}
    \vspace{-1cm}

    Dunque $ev_\B$ è lineare. $\blacksquare$

    \item $ev_\B$ è biunivoca \\
    Per dimostrare che $ev_\B$ è biunivoca, dimostriamo che è iniettiva e suriettiva.
    \begin{enumerate}[label=\arabic*°]
        \item $ev_\B$ è iniettiva \\
        Essendo $ev_\B$ lineare, per la Proposizione \ref{propos-applicazione-lineare-iniettiva-nucleo} basta dimostrare che $\ker ev_\B = \{\undzero\}$, dove $\undzero$ è l'applicazione nulla che manda tutto in $\undzero$. Affinché un'applicazione mulilineare si trovi nel nucleo di $ev_\B$, è necessario che $\Phi(\bar{v}_1, \ldots, \bar{v}_n) = \undzero$. \\
        Dato che $\B$ è una base di $\V$, presi i vettori $v_1, \ldots, v_n \in \V$ possiamo riscriverli come combinazione lineare dei vettori di questa base. Dunque \mbox{$\Phi(v_1, \ldots, v_n) = \Phi(\sum_{j = 1}^n a_{1j} \bar{v}_j, \ldots, \sum_{j = 1}^n a_{nj} \bar{v}_j)$}. Ora, applicando tante volte la multilineareità di $\Phi$, possiamo arrivare a dire che \mbox{$\Phi(v_1, \ldots, v_n) = \sum_{1 \leq j_i \leq n} a_{1 j_1} \cdot \ldots \dots a_{nj_n} \Phi(\bar{v}_{j_1}, \ldots, \bar{v}_{j_n})$}. Dato che $\Phi$ è anche alternante, se esistono due vettori identici in questa base, allora sappiamo già che $\Phi(\bar{v}_1, \ldots, \bar{v}_n) = \undzero$. Dunque, nella sommatoria precedente, rimangono solo i termini dove $j_1, \ldots, j_n$ sono tutti distinti (ma in ordine comunque sparso), e dunque diventa \mbox{$\Phi(v_1, \ldots, v_n) = \sum_{\sigma \in \mathrm{Perm}_{(n)}} a_{1 \sigma_{(1)}} \cdot \ldots \cdot a_{n \sigma_{(n)}} \cdot \Phi(\bar{v}_{\sigma_{(1)}}, \ldots, \bar{v}_{\sigma_{(n)}})$}, dove $\sigma$ rapprsenta quindi tutte le possibili permutazioni di $j_1, \ldots, j_n$ e $\sigma_{(i)}$ rappresenta l'i-esimo elemento della permutazione. 
        Ora, se analizziamo $\Phi(\bar{v}_{\sigma_{(1)}}, \ldots, \bar{v}_{\sigma_{(n)}})$, possiamo notare che per passare da questo a $\Phi(\bar{v}_1, \ldots, \bar{v}_n)$ non vengono fatte altro che una composizione di trasposizioni (ossia una composizione di scambi di vettori). Quando scambiamo due vettori, per l'Osservazione \ref{oss-inversa-app-multil-alt} sappiamo che $\Phi$ cambia segno. Dunque, in base al numero di vettori scambiati, troviamo che \mbox{$\Phi(\bar{v}_{\sigma_{(1)}}, \ldots, \bar{v}_{\sigma_{(n)}}) = \pm \Phi(\bar{v}_1, \ldots, \bar{v}_n) = \undzero$}. Dunque $\Phi(v_1, \ldots, v_n) = \undzero$, e quindi $\ker ev_\B = \{\undzero\}$.
        \item $ev_\B$ è suriettiva \\
        Dato che $ev_\B$ è iniettiva, troviamo che $\dim \Alt(\V^n, \K) = \dim \ker (ev_\B) + \dim Im (ev_\B)$, e dato che $\ker (ev_\B) = \{\undzero\}$, troviamo che $\dim \ker (ev_\B) = 0$ e dunque \mbox{$\dim \Alt(\V^n, \K) = \dim Im(ev_\B)$}. Dunque, per dimostrare che $ev_\B$ è suriettiva, basta dimostrare che $\dim Alt(\V^n, \K) = \dim \K = 1$. Per fare ciò, ci basta dimostrare che esiste un'applicazione multilineare alternante $\Phi$ tale che $ev_\B(\Phi) \neq 0$. Infatti, dato che $ev_\B$ è iniettiva, sappiamo che $\dim \Alt(\V^n, \K) \leq \dim \K = 1$, dunque ci basta dimostare che $\dim \Alt(\V^n, \K) \geq 1$, ossia trovare una $\Phi$ tale che $ev_\B(\Phi) \neq 0$. Un'applicazione del genere è $\Phi(v_1, \ldots, v_n) = \det [X_\B(v_1), \ldots, X_\B(v_n)]$, che ha senso dato che il determinante è un'applicazione multilineare alternante. Se calcoliamo ora $ev_\B(\Phi)$, troviamo che \mbox{$ev_\B(\Phi) = \det [X_\B(\bar{v}_1), \ldots, X_\B(\bar{v}_n)] = \det 1_n = 1$}, dato che $X_\B(\bar{v}_i) = e_i$. 
    \end{enumerate}
    Dunque $ev_\B$ è biunivoca. $\blacksquare$
\end{enumerate}

\section{Binet, Laplace e Cramer}

\subsection{Binet}

\begin{Proposizione}[Teorema di Binet]
\label{teorema-binet}
Prese due matrici $A, M \in M_{nn}(\K)$, $\bm{\det (A \cdot M) = \det A \cdot \det M}$.
\end{Proposizione}

\Dimostrazione \nopagebreak

Partiamo definendo l'applicazione
\begin{align*}
f: M_{nn}(\K) &\longrightarrow \K \\
M &\longmapsto \det (A \cdot M)
\end{align*}
e cerchiamo di dimostrare che tale applicazione è multilineare alternante nelle colonne di $M$.

\begin{enumerate}[label=\arabic*°]
    \item $f$ è multilineare nelle colonne \\
    Per dimostrare che $f$ è multilineare nelle colonne di $M$, dobbiamo verificare che $f$ sia lineare in tutte le colonne di $M$. Dato che il proceidmento è identico per ogni colonna di $M$, lo dimostreremo solo per la prima colonna.  \\
    Prese due vettori colonne $C,D \in M_{n1}(\K)$ e $\lambda, \mu \in \K$, vogliamo dimostrare che \mbox{$f([\lambda C + \mu D, \ldots, M_n]) = \lambda f([C, \ldots, M_n]) + \mu f([D, \ldots, M_n])$}, ossia
    \[ \det
    \begin{bmatrix}
        A^1 (\lambda C + \mu D) & \dots & A^1 M_n \\
        \vdots & \ddots & \vdots \\
        A^n(\lambda C + \mu D) & \dots & A^n M_n
    \end{bmatrix} = \lambda \det \begin{bmatrix}
        A^1 C & \dots & A^1 M_n \\
        \vdots & \ddots & \vdots \\
        A^n C & \dots & A^n M_n
    \end{bmatrix}
    + \mu \det \begin{bmatrix}
        A^1 D & \dots & A^1 M_n \\
        \vdots & \ddots & \vdots \\
        A^n C & \dots & A^n M_n
    \end{bmatrix}
    \]

Dato che $A^i (\lambda C + \mu D) = \lambda A^i C + \mu A^i D$ troviamo che

    \[ \det
    \begin{bmatrix}
        A^1 (\lambda C + \mu D) & \dots & A^1 M_n \\
        \vdots & \ddots & \vdots \\
        A^n(\lambda C + \mu D) & \dots & A^n M_n
    \end{bmatrix} = \det \begin{bmatrix}
        \lambda A^1 C & \dots & \lambda A^1 M_n \\
        \vdots & \ddots & \vdots \\
        \lambda A^n C & \dots & \lambda A^n M_n
    \end{bmatrix}
    \]

Essendo poi $\det$ multilineare, concludiamo che
\[
 \det \begin{bmatrix}
        \lambda A^1 C & \dots & \lambda A^1 M_n \\
        \vdots & \ddots & \vdots \\
        \lambda A^n C & \dots & \lambda A^n M_n
    \end{bmatrix} = 
    \lambda \det \begin{bmatrix}
        A^1 C & \dots & A^1 M_n \\
        \vdots & \ddots & \vdots \\
        A^n C & \dots & A^n M_n
    \end{bmatrix}
    + \mu \det \begin{bmatrix}
        A^1 D & \dots & A^1 M_n \\
        \vdots & \ddots & \vdots \\
        A^n C & \dots & A^n M_n
    \end{bmatrix}
\]

Dunque anche $f$ è multilineare.

\item $f$ è alternante nelle colonne di $M$ \\
Banalmente

\[
f([M_1, M_1, M_3, \ldots, M_n]) = \det \begin{bmatrix}
    A^1 M_1 & A^1 M_1 & \dots & A^1 M_n \\
    \vdots & \vdots & \ddots & \vdots \\
    A^n M_1 & A^n M_1 & \dots & A^n M_n
\end{bmatrix} = 0
\]
dato che $\det$ è un'applicazione multilineare e alternante. Dunque anche $f$ è alternante.
\end{enumerate}

Abbiamo quindi dimostrato che $f$ è un'applicazione multilineare e alternante. \\
Prendiamo ora invece $M = [e_1, \ldots, e_n]$. Troviamo che $f(M) = \det A$. Definiamo ora l'applicazione
\begin{align*}
g: M_{nn}(\K) &\longrightarrow \K \\
M &\longmapsto \det A \cdot \det M
\end{align*}

Dato che $A$ è una matrice fissata, $\det A$ è una costante, dunque $g \in \Alt(\V^n, \K)$, essendo un multiplo di $f$. Essendo $S = \{e_1, \ldots, e_n\}$ una base di $\K^n$, ed essendo l'applicazione $ev_\B$ un'isomorfismo, come visto nella Proposizione \ref{propos-ev_B isomorfismo}, dato che $g(e_1, \ldots, e_n) = f(e_1, \ldots, e_n)$, per il Corollario \ref{corollario-applicazioni-lineari-matrici} troviamo che $f = g$, ossia \mbox{$\det (A \cdot M) = \det A \cdot \det M$}. $\blacksquare$

\begin{Proposizione}[Proposizione]
\label{propos-matrici-coniugate-stesso-determinante}
Due matrici coniugate \textbf{hanno determinante uguale}.
\end{Proposizione}

\Dimostrazione \nopagebreak

Prese due matrici $A,B \in M_{nn}(\K)$ coniugate, ossia tali che $A = G^{-1} \cdot B \cdot G$, troviamo che $\det A = \det (G^{-1} \cdot B \cdot G)$. Proprio per Binet, troviamo che \mbox{$\det A = \det G^{-1} \cdot \det B \cdot \det G = \det B \cdot \det (G^{-1} \cdot G) = \det B$}. $\blacksquare$

\begin{Proposizione}[Corollario]
\label{corollario-det M B^B(f) = det M C^C(f)}
Preso un endomorfismo $f$ di uno spazio vettoriale $\V$ finitamente generato su $\K$, e prese due basi $\B, \Bs$ di $\V$, $\bm{\det M_\B^\B(f) = \det M_\Bs^\Bs(f)}$.
\end{Proposizione}

\begin{Definizione}[Determinante di un'applicazione lineare]
\label{def-determinante app lineare}
Preso un endomorfismo $f$ di uno spazio vettoriale $\V$ finitamente generato su $\K$, $\bm{\det f = \det M_\B^\B(f)}$, dove $\B$ è una qualsiasi base di $\V$. \\
Questa definizione ha senso dato che per il Corollario \ref{corollario-det M B^B(f) = det M C^C(f)} troviamo che tutte le matrici associate ad una base di $\V$ di $f$ hanno lo stesso determinante. 
\end{Definizione}

\subsection{Laplace}

\begin{Definizione}[Sviluppo di Laplace secondo riga i]
\label{def-sviluppo di Laplace riga}
Presa una matrice $A \in M_{nn}(\K)$ e una riga $i$ di $A$, $\bm{\det A = \sum_{j = 1}^n (-1)^{i+j} \cdot a_{ij} \cdot \det A_j^i}$. \\
\NB Questa è una generalizzazione dello sviluppo che abbiamo usato fino ad ora, dove prendavamo sempre la riga n-esima. 
\end{Definizione}

\begin{Proposizione}[]
\label{propos-det A = det A^t}
Presa una matrice $A \in M_{nn}(\K)$, $\bm{\det A = \det A^t}$.
\end{Proposizione}

\Dimostrazione \nopagebreak

Dimostriamolo per induzione. 
\begin{itemize}
    \item Caso base: $n = 1$ \\
    $A = (a)$, dunque $A = A^t$ e quindi anche $\det A = \det A^t$. $\blacksquare$ 
    \item Ipotesi induttiva \\
    Assumiamo la proposizione vera per tutte le matrici $(n-1) \times (n-1)$. 
    \item Passo induttivo \\
    Prendiamo la matrice $A$. Calcolandoci il determinante di $A$ troviamo che \mbox{$\det  A = \sum_{j = 1}^n (-1)^{n+j} \cdot a_{nj} \cdot \det A_j^n = (-1)^{2n} \cdot a_{nn} \cdot \det A^n_n + \sum_{j = 1}^{n-1} (-1)^{n+j} \cdot a_{nj} \cdot A_j^n$}.
    Invece, per la matrice trasposta, $\det A^t = \sum_{j = 1}^n (-1)^{n+j} \cdot a_{nj}^t \cdot \det (A^t)_j^n$. Sappiamo che $a^t_{nj} = a_{jn}$, e che per ipotesi induttiva $\det (A^t)^n_j = \det A^j_n$. Dunque troviamo che \mbox{$\det A^t = \sum_{j = 1}^n (-1)^{j+n} \cdot a_{jn} \cdot \det A^j_n = (-1)^{2n} \cdot a_{nn} \cdot \det A_n^n + \sum_{j = 1}^{n-1} (-1)^{j+n} \cdot a_{jn} \cdot \det A^j_n$}. Dato che i primi termini di entrambi i determinanti sono uguali, ci basta dimostrare che $\sum_{j = 1}^{n-1} (-1)^{j+n} \cdot a_{nj} \cdot \det A_j^n = \sum_{j = 1}^{n-1} (-1)^{j+n} \cdot a_{jn} \cdot \det A^j_n$. Se sviluppiamo $\det A^n_j$ e $\det A_n^j$ secondo riga i, troviamo che \mbox{$\sum_{j = 1}^{n-1} \sum_{j = 1}^{n-1} (-1)^{j+i-1} \cdot a_{nj} \cdot a_{in} \cdot \det (A^n_j)^i_n = \sum_{j = 1}^{n-1} \sum_{i = 1}^{n-1} (-1)^{i+j-1} \cdot a_{jn} \cdot a_{ni} \cdot \det (A^j_n)^n_i$}. È facile notare ora come $a_{nj} \cdot a_{in} = a_{jn} \cdot a_{ni}$ dato che $i$ e $j$ sono entrambe variabili che prendono tutti i valori da $1$ a $n-1$, e per lo stesso motivo anche $\det (A_j^n)^i_n  = \det (A_n^j)^n_i$. Dunque troviamo che $\det A = \det A^t$. $\blacksquare$
\end{itemize}

\begin{Definizione}[Sviluppo di Laplace secondo colonna j]
\label{def-sviluppo di Laplace colonna j}
Presa una matrice $A \in M_{nn}(\K)$ e una colonna j di $A$, $\bm{\det A = \sum_{i = 1}^n (-1)^{i+j} \cdot a_{ij} \cdot \det A^i_j}$. \\
\NB Questo sviluppo ha senso ora che abbiamo dimostrato che $\det A = \det A^t$. Infatti, lo possiamo vedere come lo sviluppo secondo riga i di $A^t$. 
\end{Definizione}

\subsection{Matrice di Vandermonde}

\begin{Definizione}[Matrice di Vandermonde]
\label{def-matrice-vandermonde}
Presa una matrice $A \in M_{nn}(\K)$, questa è una \textbf{matrice di Vandermonde} se è nella forma 
\[ A = 
\begin{bmatrixb}
1 & 1 & \dots & 1 \\
x_1 & x_2 & \dots & x_n \\
x_1^2 & x_2^2 & \dots & x_n^2 \\
\vdots & \vdots & \ddots & \vdots \\
x_1^{n-1} & x_2^{n-2} & \dots & x_n^{n-1}
\end{bmatrixb}
\]
\end{Definizione}

\begin{Proposizione}[]
\label{propos-determinante vandermonde}
Presa una matrice $A \in M_{nn}(\K)$ di Vandermonde, $\bm{\det A = \prod_{i \leq j} (x_j - x_i)}$, con \mbox{$i,j \in \{1, \ldots, n\}$}.
\end{Proposizione} 

\Dimostrazione \nopagebreak

Dimostriamolo per induzione. 
\begin{itemize}
    \item Caso base: $n = 1$ \\
    Banalmente $\det A = \det [1] = 1$. $\blacksquare$
    \item Ipotesi induttiva \\
    Assumiamo la proposizione vera per tutte le matrici di Vandermonde $(n-1) \times (n-1)$. 
    \item Passo induttivo \\
    Prendiamo la matrice $A$ e aggiungiamo a ogni riga i-esima la riga precedente moltiplicata per $x_1$, così da trasformare la prima colonna in una colonna con un solo $1$ e tutti $0$. 
    \[ A = 
    \begin{bmatrix}
        1 & 1 & \dots & 1 \\
        x_1 & x_2 & \dots & x_n \\
        x_1^2 & x_2^2 & \dots & x_n^2 \\
        \vdots & \vdots & \ddots & \vdots \\
        x_1^{n-1} & x_2^{n-1} & \dots & x_n^{n-1}
    \end{bmatrix}
    \leadsto  B = 
    \begin{bmatrix}
        1  & 1 & \dots & 1 \\
        0 & x_2 - x_1 & \dots & x_n - x_1 \\
        0 & x_2^2 - x_1 x_2 & \dots & x_n^2 - x_1 x_n^2 \\
        \vdots & \vdots & \ddots & \vdots \\
        0 & x_2^{n-1} - x_1 x_2^{n-2} & \dots & x_n^{n-1} - x_1 x_n^{n-2}
    \end{bmatrix}    
    \]

    Calcoliamo ora il determinante di $A$ con lo sviluppo di Laplace secondo la prima colonna. \mbox{$\det A = \det B = \sum_{i = 1}^n(-1)^{1+i} \cdot b_{i1} \cdot \det B^i_1$}. Dato che $b_{i1} = 0$ se $i \neq 1$, troviamo che \mbox{$\det A = \det B_{11}$}. Dato che
    \[
    B_1^1 = \begin{bmatrix}
        x_2 - x_1 & \dots & x_n - x_1 \\
        x_2 ( x_2 - x_1) & \dots & x_n (x_n - x_1) \\ 
        \vdots & \ddots & \vdots \\
        x_2^{n-2} ( x_2 - x_1) & \dots & x_n^{n-2} (x_n - x_1)
    \end{bmatrix}
    \]

Da ogni colonna possiamo raccogliere il termine $(x_i - x_1)$. Troviamo quindi che \mbox{$\det A = (x_2 - x_1) \dots (x_n - x_1) \det C$}, dove
\[
C = \begin{bmatrix}
    1 & \dots & 1 \\
    x_2 & \dots & x_n \\
    \vdots & \ddots & \vdots \\
    x_2^{n-2} & \dots & x_n^{n-2}
\end{bmatrix}
\]
\end{itemize}


Questa matrice è un'altra matrice di Vandermonde, e per ipotesi induttiva concludiamo che \mbox{$\det A = (x_2 - x_1) \cdot \ldots \cdot (x_n - x_1) \cdot \det C = \prod_{i \leq j} (x_j - x_i) \enspace \forall i,j \in \{1, \ldots, n\}$}. $\blacksquare$


\subsection{Cramer}

\begin{Definizione}[Cofattore]
\label{def-cofattore}
Presa una matrice $A \in M_{nn}(\K)$, il \textbf{cofattore} $\bm{(i,j)}$ di $A$ è $\bm{A_{ij} := (-1)^{i+j} \cdot \det A_j^i}$.
\end{Definizione}

\begin{Definizione}[Matrice dei cofattori]
\label{def-matrice dei cofattori}
Presa una matrice $A \in M_{nn}(\K)$, la \textbf{matrice dei cofattori} di $A$ è la matrice \mbox{$\bm{A^c} := (a_{ij}^c) \in M_{nn}(\K)$} dove $\bm{a_{ij}^c = A_{ji}}$, ossia è la trasposta della matrice che ha come entrata $(i,j)$ $A_{ij}$. 
\end{Definizione}

\begin{Proposizione}[Formula di Cramer]
\label{Formula-di-Cramer}
Presa una matrice $A \in M_{nn}(\K)$, $\bm{A \cdot A^c = A^c \cdot A = \det A \cdot 1_n}$. 
\end{Proposizione}

\Dimostrazione \nopagebreak

Dimostreremo solo che $A \cdot A^c = \det A \cdot 1_n$, dato che il procedimento per dimostrare che anche $A^c \cdot A = \det A \cdot 1_n$ è pressoché identico.

Prendiamo $B = A \cdot A^c$. Troviamo che $b_{ij} = A^i \cdot A^c_j$, dove $A^i =[a_{i1}, \ldots, a_{in}]$ e 
\[
A^c_j = \begin{bmatrix}
    a_{1j}^c \\
    \vdots \\
    a_{nj}^c
\end{bmatrix} =  \begin{bmatrix}
    A_{j1} \\
    \vdots \\
    A_{jn}
\end{bmatrix}
\]

Dunque $b_{ij} = \sum_{s = 1}^n a_{is} \cdot A_{js}$. Dato che $A_{js} = (-1)^{j+s} \cdot \det A_s^j$, troviamo proprio che \mbox{$b_{ij} = \sum_{s = 1}^n a_{is} \cdot (-1)^{j+s} \cdot \det A_s^j$}. Distinguiamo ora due casi:
\begin{itemize}
    \item $i = j$ \\
    In questo caso $b_{ii} = A^i \cdot A^c_i = \sum_{s = 1}^n a_{is} \cdot (-1)^{j+s} \cdot \det A^i_s = \det A$, calcolato secondo lo sviluppo di Laplace secondo riga i.
    \item $i \neq j$ \\
    Prendiamo ora la matrice $C \in M_{nn}(\K)$ ottenuta da $A$ sostituendo la riga j di $A$ con la sua riga i. Dato che in $C$ ci sono due righe uguali, per la proprietà 4) vista nelle Proprietà \ref{prop-determinante} sappiamo che $\det C = 0$, dato che per la Proposizione \ref{propos-det A = det A^t} sappiamo che $\det C = \det C^t$, e $C^t$ ha proprio due colonne uguali. Sviluppiamo ora il determinante di $C$ secondo la riga j: \mbox{$\det C = \sum_{s = 1}^n (-1)^{j+s} \cdot c_{js} \cdot \det C_s^j$}. Dato che $C^j = A^i$ e $c_{js} = a_{is} \enspace \forall s \in \{1, \ldots, n\}$ e anche $C^j_s = A^j_s$, troviamo che \mbox{$\det C = \sum_{s = 1}^n (-1)^{i+s} \cdot a_{is} \cdot \det A^j_s = b_{ij}$}. Troviamo quindi che $b_{ij} = 0$ se $i \neq j$.
\end{itemize}

Dunque possiamo concludere che
\[
B = A \cdot A^c = \begin{bmatrix}
    \det A & 0 & \dots & 0 \\
    0 & \det A & \dots & 0 \\
    \vdots & \vdots & \ddots & \vdots \\
    0 & 0 & \dots & \det A
\end{bmatrix} = \det A \cdot 1_n \quad \blacksquare
\]

\begin{Proposizione}[]
\label{proposizione-invertibilità-cramer}
Presa una matrice $A \in M_{nn}(\K)$, $\bm{A}$ \textbf{è invertibile} $\iff \bm{\det A \neq 0}$. \\
In tal caso si ha:
\begin{itemize}
    \item $\bm{\det A  = (\det A^{-1})^{-1}}$ 
    \item $\bm{A^{-1} = (\det A)^{-1} \cdot A^c}$
\end{itemize}

\end{Proposizione}

\Dimostrazione \nopagebreak

\begin{enumerate}[label=\arabic*°]
    \item $A$ è invertibile $\implies \det A \neq 0$ \\
    Applicando il teorema di Binet troviamo che $\det (A \cdot A^{-1}) = \det A \cdot \det A^{-1}$. Dato che $A \cdot A^{-1} = 1_n$, troviamo che $\det A \cdot \det A^{-1} = \det 1_n = 1$, ossia $\det A = (\det A^{-1})^{-1} \neq 0$, dato che $\det A \cdot \det A^{-1} = 1 \neq 0$. $\blacksquare$
    \item $A$ è invertibile $\impliedby \det A \neq 0$ \\
    Dalla formula di Cramer sappiamo che $A \cdot A^c = \det A \cdot 1_n$. Sappiamo che $\det A \neq 0$, dunque possiamo dividere tutto per $\det A$, trovando che $A \cdot \frac{A^c}{\det A} = 1_n$, dunque troviamo che $A^{-1} = \frac{A^c}{\det A} = (\det A)^{-1} \cdot A^c$. $\blacksquare$ 
\end{enumerate}

\Esempi \nopagebreak

Trovare il \textbf{determinante di una matrice} $\bm{2 \times 2}$ \textbf{invertibile} usando la formula di Cramer. \nopagebreak

Prendiamo una matrice $A \in M_{22}(\K)$ invertibile. Abbiamo quindi
\[
A = \begin{bmatrix}
    a &  b \\
    c & d
\end{bmatrix}
\enspace \text{e} \enspace A^c = \begin{bmatrix}
    A_{11} & A_{21} \\
    A_{12} & A_{22}
\end{bmatrix} = \begin{bmatrix}
    d & -b \\
    -c & a
\end{bmatrix}
\]
Troviamo quindi che
\[
A \cdot A^c = \begin{bmatrix}
    ad - bc & 0 \\
    0 & ad-bc
\end{bmatrix} = \det A \cdot 1_n
\]

Dunque $\bm{\det A = ad -bc}$. 

\Separatore

Presa la matrice invertibile

\[
\bm{A =} \begin{bmatrixb}
    1 & 1 & 1 \\
    1 & -1 & 2 \\
    1 & 1 & 4
\end{bmatrixb}
\]

trovare $\bm{A^{-1}}$.

Per la Proposizione \ref{proposizione-invertibilità-cramer} sappiamo che $A^{-1} = (\det A)^{-1} \cdot A^c$. Per calcolare il determinante di $A$, riduciamo la matrice a scala per righe:
\[
A = \begin{bmatrix}
    1 & 1 & 1 \\
    1 & -1 & 2\\
    1 & 1 & 4
\end{bmatrix}
\leadsto S = \begin{bmatrix}
    1 & 1 & 1 \\
    0 & -2 & 1 \\
    0 & 0 & 3
\end{bmatrix}
\]

Troviamo quindi che $\det A = \det S = 1 \cdot (-2) \cdot 3 = -6$, e quindi $(\det A)^{-1} =  - \frac{1}{6}$. Possiamo quindi concludere che
\[
\bm{A^{-1}} = (\det A)^{-1} \cdot A^c = -\frac{1}{6} \cdot \begin{bmatrixb}
    A_{11} & A_{21} & A_{31} \\
    A_{12} & A_{22} & A_{32} \\
    A_{13} & A_{23} & A_{33}
\end{bmatrixb} = -\frac{1}{6} \begin{Bmatrix}
    -6 & -3 & 3 \\
    -2 & 3 & -1 \\
    2 & 0 & -2
\end{Bmatrix} = \begin{bmatrixb}
    1 & \frac{1}{2} & -\frac{1}{2} \\
    \frac{1}{3} & -\frac{1}{2} & \frac{1}{6} \\
    -\frac{1}{3} & 0 & \frac{1}{3}
\end{bmatrixb}
\]

\section{Determinante e permutazioni}

\begin{Definizione}[Matrice di una permutazione]
\label{def-matrice-permutazione}
Presa una permutazione $\sigma \in S_n = \{\text{insieme delle permutazioni di} \{1, \ldots, n\} = [n]\}$, dove quindi $\sigma: [n] \longrightarrow [n]$ è un'applicazione biunivoca, chiamiamo $\bm{M_\sigma = (m_{ij}) \in M_{nn}(\Z)}$ la \textbf{matrice della permutazione} $\bm{\sigma}$, dove 
\[
\bm{m_{ij}}
\begin{cases}
\bm{1} & \enspace \textbf{se } \bm{i = \sigma(j)}  \\
\bm{0} & \enspace \textbf{altrimenti} 
\end{cases}
\]
\end{Definizione}

Dunque $M_\sigma$ è una matrice formata da tutti $0$ dove, in ogni colonna j, è presente uno e un solo $1$ alla riga $\sigma(j)$.

\begin{Osservazione}[]
\label{oss-matrice-permutazioni-1}
Presa una matrice $A \in M_{nn}(\Z)$ con $a_{ij} \in \{0, 1\}$, esiste una permutazione $\sigma \in S_n$ tale che $A = M_\sigma$ $\iff$ su ogni riga e colonna di $A$ è presente uno e un solo $1$.
\end{Osservazione}

Se ad esempio prendiamo una trasposizione $\sigma$ (ossia una permutazione dove solo due elementi vengono scambiati tra loro), dove vengono scambiate le colonne $j$ e $j+1$, troviamo che
\[
M_\sigma = [(1_n)_1, \ldots, (1_n)_{j+1}, (1_n)_j, \ldots, (1_n)_n] = \begin{bmatrix}
    1 & 0 & \dots & 0 & 0 & \dots & 0 \\
    0 & 1 & \dots & 0 & 0 & \dots & 0 \\
    \vdots & \vdots & \ddots & \vdots & \vdots & \ddots & \vdots \\
    0 & 0 & \dots & 0 & 1 & \dots & 0 \\
    0 & 0 & \dots & 1 & 0 & \dots & 0 \\
    \vdots & \vdots & \ddots & \vdots & \vdots & \ddots & \vdots \\
    0 & 0 & \dots & 0 & 0 & \dots & 1
\end{bmatrix}
\] 

\begin{Osservazione}[]
\label{oss-matrice-permutazione-2}
$\bm{\det M_\sigma} \in \{1, -1\}$.
\end{Osservazione}

\Dimostrazione \nopagebreak

Banalmnente, non essendo $M_\sigma$ altro che una permutazione delle colonne di $1_n$, troviamo proprio che $\det M_\sigma = \pm \det 1_n = \pm 1$. $\blacksquare$  \\
In particolare, $\det M_\sigma = 1$ se il numero di scambi tra le colonne di $M_\sigma$ per arrivare a $1_n$ è pari, ed invece $-1$ se il numero è dispari.

\begin{Definizione}[Segno di una permutazione]
\label{def-segno-permutazione}
Il \textbf{segno} di una permutazione $\bm{\sigma}$ è il \textbf{determinante di} $\bm{M_\sigma}$. Per indicare il segno di una permutazione $\sigma$ scriviamo $\bm{\epsilon(\sigma)}$ (dunque $\epsilon(\sigma) = \det M_\sigma$).
\end{Definizione}

\begin{Osservazione}[]
\label{oss-matrice-permutazione-3}
Prese due permutazioni $\sigma, \tau \in S_n$, $\bm{\epsilon(\sigma \circ \tau) = \epsilon(\sigma) \cdot \epsilon(\tau)}$.
\end{Osservazione}

\Dimostrazione \nopagebreak

Sappiamo che $\epsilon(\sigma \circ \tau) = \det M_{\sigma \circ \tau}$, e che $\epsilon(\sigma) \cdot \epsilon(\tau) = \det M_\sigma \cdot \det M_\tau$, dove \mbox{$M_{\sigma \circ \tau}, M_\sigma, M_\tau \in M_{nn}(\Z)$}. Vogliamo quindi dimostrare che $\det M_{\sigma \circ \tau} = \det M_\sigma \cdot \det M_\tau$. Per farlo, dimostriamo semplicemente che $M_{\sigma \circ \tau} = M_\sigma \cdot M_\tau$, dato che, per Binet, abbiamo che \mbox{$\det M_{\sigma \circ \tau} = \det (M_\sigma \cdot M_\tau) = \det M_\sigma \cdot M_\tau$}. \\
Per dimostrare che due matrici sono uguali, basta dimostrare che le loro applicazioni lineari associate lo sono, ossia che $L_{M_{\sigma \circ \tau}} = L_{M_\sigma \cdot M_\tau} = L_{M_\sigma} \circ L_{M_\tau}$. Per il lemma 2) visto nella Proposizione \ref{lemmi-matrici-e-app-lin}, sappiamo che ci basta dimostrare che $L_{M_{\sigma \circ \tau}}(v_i) = L_{M_\sigma} \circ L_{M_\tau} (v_i)$ \mbox{$\forall i \in \{1, \ldots, n\}$}, dove $v_i$ sono i vettori di una base canonica di $Z^n$. Prendendo la base standard $\{e_1, \ldots, e_n\}$, sappiamo che $L_{M_\sigma}(e_i) = M_\sigma \cdot e_i = (M_\sigma)_i = e_{\sigma_i}$. Dunque troviamo che $L_{M_{\sigma \circ \tau}}(e_i) = e_{\sigma \circ \tau(i)} = e_{\sigma(\tau(i))}$ e $L_{M_\sigma} \circ L_{M_\tau}(e_i) = L_{M_\sigma}(e_{\tau(i)}) = e_{\sigma(\tau(i))}$, dunque concludiamo che $L_{M_{\sigma \circ \tau}} = L_{M_\sigma} \circ L_{M_\tau}$ e quindi $M_{\sigma \circ \tau} = M_\sigma \cdot M_\tau$. $\blacksquare$

\begin{Osservazione}[]
\label{oss-matrice-permutazione-4}
Dato che una permutazione $\sigma$ è una composizione di trasposizioni ($\sigma = \tau_1 \circ \ldots \circ \tau_d$), troviamo che $\epsilon(\sigma) = \epsilon (\tau_1) \cdot \ldots \cdot \epsilon(\tau_d)$. Dato che ogni $\epsilon(\tau_i)$ cambia il segno di $\epsilon(\sigma)$, troviamo che
\[
\bm{\epsilon(\sigma)} = 
\begin{cases}
\bm{1} & \enspace \textbf{se } \bm{d} \textbf{ è pari} \\
\bm{-1} & \enspace \textbf{se } \bm{d} \textbf{ è dipari} \\
\end{cases}
\]
\end{Osservazione}

Ci si potrebbe chiedere perché non abbiamo direttamente definito il segno di una permutazione in questo modo. Non era possibile farlo perché semplicemente non sapevamo se il numero di trasposizioni con cui si può scrivere una permutazione è sempre pari o dispari: avremmo ad esempio potuto trovare una permutazione $\sigma = \tau_1 \circ \tau_2 = \tau_1' \circ \tau_2' \circ \tau_3'$. In questo caso, il segno di $\sigma$ sarebbe $\pm 1$, il che non ha senso. Questa osservazione è quindi la dimostrazione che \textit{ogni combinazioni di trasposizioni che forma una permutazione $\sigma$ con segno $1$ avrà un numero pari di trasposizioni, altrimenti se il segno è $-1$ allorà avrà un numero dispazi di traspoziioni}.


\begin{Proposizione}[]
\label{propos-matrice-permutazioni}
Presa una matrice $A \in M_{nn}(\K)$, $\bm{\det A = \sum_{\sigma \in S_n} \epsilon(\sigma) \cdot a_{1 \sigma(1)} \cdot \ldots \cdot a_{n\sigma(n)}}$, dove $\abs{S_n} = n!$.
\end{Proposizione}

Se prendiamo ad esempio $n = 2$, abbiamo solo due permutazioni $\sigma$: $\sigma_1 = \{1,2\}$ e $\sigma_2 = \{2,1\}$. Dunque \mbox{$\det A = a_{1\sigma_1(1)} \cdot a_{2 \sigma_1 (2)} - a_{1\sigma_2(1)} a_{2\sigma_2(2)} = a_{11}a_{22} - a{12}a_{21}$}, come già sapevamo.

\section{Polinomio caratteristico}

\begin{Definizione}[Polinomio caratteristico]
\label{def-polinomio-caratteristico}
Preso un endomorfismo $f$ su uno spazio vettoriale $\V$, il \textbf{polinomio caratteristico} di $f$ è $\bm{P_f = \det (\lambda Id_\V - f)}$.

Allo stesso modo, presa una matrice $A \in M_{nn}(\K)$, il \textbf{polinomio caratteristico} di $A$ è $\bm{P_A = \det (\lambda 1_n - A)}$.
\end{Definizione}

\begin{Osservazione}[]
\label{oss-polinomio-caratteristico-è-monico}
Il polinomio caratteristico è un polinomio in $\lambda$ a coefficienti in $\K$ \textbf{monico di grado} $\bm{n}$, ossia è un polinomio in $\lambda$ di grado $n$ con il coefficiente di $\lambda^n$ pari a $1$.
\end{Osservazione}

\Dimostrazione \nopagebreak

Ponendo $B = \lambda 1_n - A$, troviamo che $\det (\lambda 1_n - A) = \det B = \sum_{\sigma \in S_n}\epsilon(\sigma) \cdot b_{1\sigma(1)} \cdot \ldots \cdot b_{n\sigma(n)}$, come visto nella Proposizione \ref{propos-matrice-permutazioni}. Sappiamo che 
\[
b_{ij}=
\begin{cases}
\lambda - a_{ii} & \enspace \text{se } i = j \\
-a{ij} & \enspace \text{se } i \neq j 
\end{cases}
\]

Ora, prendiamo la permutazione $\sigma$ che non sposta alcun elemento. In questa permutazione, \mbox{$\epsilon(\sigma) \cdot b_{1\sigma(1)} \cdot \ldots \cdot b_{n\sigma(n)} = b_{11} \cdot \ldots \cdot b_{nn} = \lambda^n - (\sum_{i = 1}^n a_{ii})\lambda^{n-1} \cdot \ldots$}, dunque troviamo che per questa permutazione il coefficiente di $\lambda^n$ è $1$ e il coefficiente di $\lambda^{n-1}$ è $-\sum_{i = 1}^n a_{ii}$ (questo ci sarà utile per una dimostrazione futura). Se guardiamo anche tutte le altre permutazioni, è facile notare che queste devono muovere almeno due elementi (dato che ci deve essere almeno uno scambio). Dunque vi sono almeno due indici $k_1$ e $k_2$ tali che $k_1 \neq \sigma(k_1)$ e $k_2 \neq \sigma(k_2)$. Dunque per tutte queste altre permutazioni, troviamo al più $n-2$ elementi di $B$ che vengono moltiplicati, e quindi il grado massimo raggiungibilie di $\lambda$ è $n-2$. Possiamo quindi concludere che il coefficiente di $\lambda^n$ è $1$. $\blacksquare$

\begin{Proposizione}[]
\label{propos-polinomio-autovalori-radici}
Preso un polinomio caratteristico $P_f$ (può essere anche $P_A$, è indifferente), \mbox{$\bm{\lambda_0 \in \K}$ \textbf{è un autovalore di} $\bm{f} \iff \bm{\lambda_0}$ \textbf{è una radice di} $\bm{P_f}$.}
\end{Proposizione}

\vspace{2cm}

\Dimostrazione \nopagebreak

\begin{enumerate}[label=\arabic*°]
    \item $\lambda_0 \in K$ è un autovalore di $f \implies \lambda_0$ è una radice di $P_f$\\ 
    Sappiamo che $\lambda_0$ è un autovalore di $f$ $\iff \exists v \neq 0$ tale che $f(v) = \lambda_0 v$. Dunque $\ker(\lambda_0 Id_\V - f) \neq \{\undzero\}$, dunque $\lambda_0 Id_\V -  f$ non è iniettiva, e dunque non è invertibile. Per il Corollario \ref{corol-det-invertibilita}, sappiamo quindi che $\det(\lambda_0 Id_\V - f) = 0$. Dato che $P_f = \det(\lambda Id_\V - f)$, e dato che $\det (\lambda_0 Id_\V - f) = 0$, troviamo che $P_f(\lambda_0) = 0$, ossia $\lambda_0$ è una radice di $P_f$. $\blacksquare$
    \item $\lambda_0 \in K$ è un autovalore di $f \impliedby \lambda_0$ è una radice di $P_f$\\ 
    Sappiamo che $P_f(\lambda_0) = 0$, ossia $\det (\lambda_0 Id_\V - f) = 0$. Dunque, per lo stesso corollario visto prima, troviamo che $\ker (\lambda_0 Id_\V - f) \neq \{\undzero\}$ ossia $\lambda_0$ è un autovalore di $f$. $\blacksquare$
\end{enumerate}


\Esempio \nopagebreak

Trovare una base di $\R^2$ che diagonalizzi la matrice \nopagebreak

\[\bm{A} = 
\begin{bmatrixb}
    1 & 5 \\
    -\frac{1}{5} & -2
\end{bmatrixb}
\]

Prendiamo il polinomio caratteristico di $A$:
\[
P_A = \begin{vmatrix}
    \lambda - 1 & -5 \\
    \frac{1}{5} & \lambda + 2
\end{vmatrix} = (\lambda -1)(\lambda + 2) + \frac{5}{4}
\]
Ponendo $P_A = 0$, troviamo l'equazione $\lambda^2 + \lambda - \frac{3}{4} = 0$. Le radici di questa equazione sono $\lambda_1 = \frac{1}{2}$ e $\lambda_2 = -\frac{3}{2}$, che per la Proposizione \ref{propos-polinomio-autovalori-radici} sono quindi gli autovalori di $A$. \\
Ora che abbiamo gli autovalori, ci servono i loro corrispettivi autospazi per poter diagonalizzare $A$: infatti, per il Corollario \ref{corol-autovalori-matrice}, sappiamo che gli autovalori di $A$ sono proprio i valori della matrice diagonale $M_\B^\B(L_A)$, e dunque $\B = \{v_1, v_2\}$ è una base che diagonalizza $A$ e $v_1$ e $v_2$ sono due autovettori di autovalore rispettivamente $\lambda_1$ e $\lambda_2$. \\
Vogliamo quindi calcolare gli autospazi $\V_\lambda = \ker (\lambda 1_n - A)$.
\begin{itemize}
    \item $\V_{\lambda_1} = \ker(\lambda_1 1_n - A)$ \\
    \[
    (\lambda_1 1_n - A)  \cdot \begin{pmatrix}
        x \\ y
    \end{pmatrix} = \begin{bmatrix}
        -\frac{1}{2} & -5 \\
        \frac{1}{4} & \frac{5}{2}
    \end{bmatrix} \cdot \begin{pmatrix}
        x \\ y
    \end{pmatrix} = \begin{pmatrix}
        -\frac{1}{2}x  -5y \\
        \frac{1}{4}x + \frac{5}{2} y
    \end{pmatrix} = \begin{pmatrix}
        0 \\ 0
    \end{pmatrix} = \undzero
     \implies
     \begin{cases}
        x = -10y \\
        0 = 0
     \end{cases}
    \]
    Possiamo quindi prendere come autovettore il vettore $v_1 = (-10,1)$, che avrà quindi come autovalore $\lambda_1 = \frac{1}{2}$.
    \item $\V_{\lambda_2} = \ker(\lambda_2 1_n - A)$ \\
    \[
    (\lambda_2 1_n - A) \cdot \begin{pmatrix}
        x \\ y
    \end{pmatrix} = \begin{bmatrix}
        -\frac{5}{2} & - 5 \\
        \frac{1}{4} & \frac{1}{2}
    \end{bmatrix} \cdot \begin{pmatrix}
        x \\ y
    \end{pmatrix} = \begin{pmatrix}
        -\frac{5}{2} x - 5y \\
        \frac{1}{4} x + \frac{1}{2}y
    \end{pmatrix} = \begin{pmatrix}
         0 \\ 0
    \end{pmatrix} = \undzero
    \implies
    \begin{cases}
        x = -2y \\
        0 = 0 
    \end{cases}
    \]
    Possiamo quindi prendere come autovettore il vettore $v_2 = (2,-1)$, che avrà quindi come autovalore $\lambda_2 = -\frac{3}{2}$.
\end{itemize}

Troviamo quindi che $M_\B^\B(L_A)$ è una matrice diagonale, e quindi $\bm{\B = \{(-10,1), (2,-1)\}}$ è una base che diagonalizza $A$.

\begin{Osservazione}[]
\label{oss-PA = PB}
Se due matrici $A,B \in M_{nn}(\K)$ sono coniugate, allora $\bm{P_A = P_B}$.
\end{Osservazione}

\Dimostrazione \nopagebreak

Essendo $A$ e $B$ coniugate, sappiamo che $\exists G \in GL_n(\K)$ (ossia una matrice $G$ invertibile) tale che $A = G^{-1} \cdot B \cdot G$. Dunque \mbox{$P_A = \det (\lambda 1_n - A) = \det(\lambda 1_n - G^{-1} \cdot B \cdot G)$}. Dato che $\lambda 1_n = G^{-1} \cdot \lambda 1_n \cdot G$, troviamo che $P_A = \det(G^{-1} \cdot \lambda 1_n \cdot G - G^{-1} \cdot B \cdot G) = \det (G^{-1} \cdot (\lambda 1_n - B)\cdot G)$. Per Binet, troviamo che $P_A = \det G^{-1} \cdot \det(\lambda 1_n - B) \cdot \det G = \det (\lambda 1_n - B) = P_B$. $\blacksquare$

\begin{Definizione}[Invariante per coniugio]
\label{def-invariante-per-coniugio}
Presa un'applicazione lineare $f$, questa è detta \textbf{invariante per coniugio} se, prese due matrici $A,B \in M_{nn}(\K)$ coniugate, $\bm{f(A) = f(B)}$. \\
Ad esempio, il \textit{polinomio caratteristico} è un'applicazione invariante per coniugio, come dimostrato nell'Osservazione \ref{oss-PA = PB}
\end{Definizione}

\begin{Definizione}[Traccia]
\label{def-traccia}
Presa una matrice $A \in M_{nn}(\K)$, la sua \textbf{traccia} è $\bm{\mathrm{Tr}A := \sum_{i = 1}^n a_{ii}}$, ossia la somma di tutti i suoi elementi sulla diagonale principale.
\end{Definizione}

\begin{Proposizione}[]
\label{propos-matrici-coniugate-traccia-uguale}
Prese due matrici $A,B \in M_{nn}(\K)$, $A$ e $B$ \textbf{sono coniugate} $\implies \bm{\mathrm{Tr} A = \mathrm{Tr} B}$.
\end{Proposizione}

\Dimostrazione \nopagebreak

Nella dimostrazione dell'Osservazione \ref{oss-polinomio-caratteristico-è-monico} abbiamo visto come il coefficiente del termine $\lambda^{n-1}$ del polinomio caratteristico sia $-\sum_{i = 1}^n m_{ii}$ per una generica matrice $M$. Dato che $A$ e $B$ sono coniugate, per l'Osservazione \ref{oss-PA = PB} troviamo che $P_A = P_B$, e dunque $-\sum_{i = 1}^n a_{ii} = -\sum_{i = 1}^n b_{ii}$, ossia quindi $\mathrm{Tr}A = \mathrm{Tr}B$. $\blacksquare$ 

\begin{Proposizione}[]
\label{propos-autovettori con autovalori distinti indipendenti}
Preso un endomorfismo $f$ dello spazio vettoriale $\V$, e presi $v_1, \ldots, v_n$ autovettori di $f$ con autovalori distinti, $\bm{v_1, \ldots, v_n}$ \textbf{sono linearmente indipendenti}.
\end{Proposizione}

\Dimostrazione \nopagebreak

Per assurdo, poniamo $v_1, \ldots, v_n$ linearmente dipendenti. Sappiamo quindi che esiste un minimo $m$ tale che $1 \leq i_1 \leq \ldots \leq i_m \leq n$, con $v_{i_1}, \ldots, v_{i_m}$ linearmente dipendenti. Quindi, sappiamo che se prendiamo $m-1$ vettori tra questi, sicuramente saranno linearmente indipendenti, essendo $m$ il minimo. \\
Dunque esisono $\mu_1, \ldots, \mu_m \in K$ tali che $\mu_1 v_{i_1}+ \ldots + \mu_m v_{i_m} = \undzero$ con almeno un $\mu_j \neq 0$. Essendo però $m$ il minimo, in realtà possiamo dire che tutti i $\mu_j$ sono diversi da $0$, dato che se questo non fosse vero, allora avremmo un numero di vettori linearmente dipendenti minore di $m$, il che è impossibile per definizione di $m$. \\
Chiamiamo ora $\lambda_k$ l'autovalore di $v_{i_k}$. Dunque sappiamo che $f(v_{i_k}) = \lambda_k v_{i_k}$. Dunque troviamo che \mbox{$f(\undzero) = f(\mu_1 v_{i_1} + \ldots + \mu_m v_{i_m}) = \mu_1 \lambda_1 v_{i_1} + \ldots + \mu_m \lambda_m v_{i_m} = \undzero$}. Dato che $\mu_1 v_{i_1} +\ldots + \mu_m v_{i_m} = \undzero$, troviamo che sicuramente anche $\lambda_m (\mu_1 v_{i_1} + \ldots + \mu_m v_{i_m}) = \undzero$. Dunque 
\begin{align*}
 \mu_1 \lambda_1 v_{i_1} + \ldots + \mu_m \lambda_m v_{i_m} - \lambda_m (\mu_1 v_{i_1} + \ldots + \mu_m v_{i_m}) &= \undzero \\
 \mu_1 (\lambda_1 - \lambda_m) v_{i_1} + \ldots + \mu_{m-1} (\lambda_{m-1} - \lambda_m) v_{i_{m-1}} = \undzero
\end{align*}

Dato che $\mu_1, \ldots, \mu_{m-1}$ sono tutti diversi da $0$, e dato che $\lambda_1, \ldots, \lambda_{m-1}$ sono tutti distinti, troviamo che questa è una relazione di dipendenza lineare. Abbiamo quindi trovato $m-1$ vettori linearmente dipendenti, ma essendo $m$ il minimo di vettori linearmente dipendenti questo è impossibile. \\
Essendo quindi arrivati ad un assurdo, concludiamo che $v_1, \ldots, v_n$ sono vettori linearmente indipendenti. $\blacksquare$

\begin{Proposizione}[]
\label{propos-f-con-n-autovalori-diagonalizzabile}
Sia $f$ un endomorfismo di uno spazio vettoriale $\V$, con $\dim \V = n$. $\bm{f}$ \textbf{ha} $\bm{n}$ \textbf{autovalori distinti in} $\bm{\K} \implies f$ \textbf{è diagonalizzabile}. 
\end{Proposizione}

\Dimostrazione \nopagebreak

Siano $\lambda_1, \ldots, \lambda_n$ gli autovalori di $f$ e $v_1, \ldots, v_n$ gli autovettori corrispondenti. Sappiamo per la Proposizione \ref{propos-autovettori con autovalori distinti indipendenti} che $\B = \{v_1, \ldots, v_n\}$ è una base di $\V$, dato che $\abs{\B} = \dim \V$ e per l'Osservazione \ref{oss-dimen} questo implica che $\B$ sia una base. Dato che $f(v_i) = \lambda_i v_i$, troviamo che la colonna i-esima di $M_\B^\B(f)$ sarà $X_\B(f(v_i)) = X_\B(\lambda_i v_i) = (0,\ldots, \lambda_i, \ldots, 0)$, con $\lambda_i$ alla riga i-esima. Dunque 
\[
M_\B^\B(f) = \begin{bmatrix}
    \lambda_1 & \dots & 0 \\
    \vdots & \ddots & \vdots \\
    0 & \dots & \lambda_n
\end{bmatrix}
\]
Concludiamo quindi che $f$ è diagonalizzabile. $\blacksquare$

\begin{Definizione}[Molteplicità Algebrica]
\label{def-molteplicità}
Preso uno scalare $\lambda_0 \in K$ di un'endomorfismo $f$, la \textbf{molteplicità algebrica} di $\lambda_0$, indicata con il simbolo $\bm{\mathrm{mult}_{\lambda_0}(P_f)}$ è il \textbf{massimo intero} $\bm{m \geq 1}$ tale che $\bm{(\lambda - \lambda_0)^m}$ \textbf{divida il polinomio caratteristico} $\bm{P_f}$. 
\end{Definizione}
È facile capire come se $\mathrm{mult}_{\lambda_0}(P_f) \geq 1$ allora $\lambda_0$ è un autovalore di $f$, dato che questo è a tutti gli effetti una radice di $P_f$, e per la Proposizione \ref{propos-polinomio-autovalori-radici} sappiamo che questo lo rende un autovalore di $f$.

\begin{Proposizione}[]
\label{propos-dim V < mult}
Preso un'endomorfismo $f$ di uno spazio vettoriale $\V$, $\bm{\lambda_0}$ \textbf{autovalore di} $\bm{f} \implies$ \mbox{$\bm{\dim \V_{\lambda_0}(f) \leq \mathrm{mult}_{\lambda_0}(P_f)}$}.
\end{Proposizione}

\Dimostrazione \nopagebreak

Sia $d = \dim \V_{\lambda_0}(f)$ e $\Bs = \{v_1, \ldots, v_d\}$ una base di $V_{\lambda_0}(f)$. Estendiamo $\Bs$ ad una base $\B = \{v_1, \ldots, v_d, v_{d+1}, \ldots, v_n\}$ di $\V$, ponendo dunque $\dim \V = n$. Sappiamo che le prime $d$ colonne di $M_\B^\B(f)$ saranno del tipo $(0,\ldots, \lambda_0, \ldots, 0)$ con $\lambda_0$ che per la i-esima colonna si troverà alla i-esima riga. Dunque
\[
M_\B^\B(f) = \left[
\begin{array}{ccc|c}
    \lambda_0 & \dots & 0 &\\
    \vdots & \ddots & \vdots & M \\
    0 & \dots & \lambda_0 & \\ \hline
    0 & \dots & 0 & Q
\end{array}
\right]
\]
Dove $M$ e $Q$ sono matrici che non ci interessano particolarmente. 

Troviamo quindi che 
\[
P_f = \det M_\B^\B(\lambda Id - f) = \left|
\begin{array}{ccc|c}
    \lambda - \lambda_0 & \dots & 0 & \\
    \vdots & \ddots & \vdots & -M \\
    0 & \dots & \lambda - \lambda_0 & \\ \hline
    0 & \dots & 0 & \lambda 1_{n-d}  - Q 
\end{array}
\right|
\]
Se calcoliamo il determinante usando lo sviluppo di Laplace secondo la 1° colonna troviamo che
\[
P_f = (\lambda - \lambda_0) \left|
\begin{array}{ccc|c}
    \lambda - \lambda_0 & \dots & 0 & \\
    \vdots & \ddots & \vdots & -M' \\
    0 & \dots & \lambda - \lambda_0 & \\ \hline
    0 & \dots & 0 & \lambda 1_{n-d}  - Q 
\end{array}
\right|
\]
Dove $M'$ è la matrice $M$ senza la prima riga. Ripetendo questo procedimento d-volte, troviamo che $P_f = (\lambda - \lambda_0)^d \cdot \det(\lambda 1_{n-d} - Q)$, e dunque troviamo che $P_f$ è divisibile almeno d-volte per $\lambda - \lambda_0$, e dato che $\dim \V_{\lambda_0}(P_f) = d$, troviamo proprio che \mbox{$\mathrm{molt}_{\lambda_0}(P_f) \geq \dim \V_{\lambda_0}(f)$}. $\blacksquare$

\begin{Proposizione}[Corollario]
\label{corollario-f diagonalizzabile Pf radici}
Preso un endomorfismo $f$ su uno spazio vettoriale $\V$, con $\dim \V = n$, \mbox{$\bm{f}$ \textbf{è diagonalizzabile} $\iff$} $\bm{P_f}$ \textbf{ha} $\bm{n}$ \textbf{radici in} $\bm{\K}$ \textbf{e per ogni autovalore} $\bm{\lambda_0}$ \textbf{di} $\bm{f}$ \textbf{si ha} $\bm{\dim \V_{\lambda_0}(f) = \mathrm{mult}_{\lambda_0}(P_f)}$.
\end{Proposizione}

\chapter{Spazi affini}
\section{Introduzione agli spazi affini}

Abbiamo già visto lo spazio vettoriale $V(\E^2)$ che racchiude tutti i vettori geometrici nello spazio euclideo $\E^2$ nel primo capitolo. Ora, riprendendo questo concetto, cerchiamo di trovare una relazione tra $\E^3$ e $V(\E^3)$ (lavoriamo con lo spazio euclideo piuttosto che con il piano euclideo). 

Possiamo innanzitutto definire un'applicazione del tipo
\begin{align*}
\E^3 \times V(\E^3) &\longrightarrow \E^3 \\
(P, v) &\longmapsto P + v
\end{align*}

Ma $P + v$ è una somma tra un punto nello spazio e un vettore di tale spazio, non ha molto senso. Definiamo quindi la somma $P+v$ come il punto $Q \in \V^3$ tale che $v = \vv{PQ}$. Tale definizione ha perfettamente senso perché esiste un unico punto $Q$ tale che $v = \vv{PQ}$. 

Grazie a questa definizione di somma, possiamo elencare 3 proprietà importanti.

\begin{Proprieta}[Proprietà della somma tra punti e vettori]
\label{prop-somma-punto-vettore}
\begin{enumerate}
    \item $\forall P \in \E^3 \enspace \bm{P + \undzero = P}$
    \item $\forall P \in \E^3$ e $\forall v,w \in V(\E^3) \enspace \bm{P + (v+w) = (P+v) + w}$
    \item $\forall P,Q \in \E^3 \enspace \exists ! v \in V(\E^3)$ tale che $\bm{P + v = Q}$ ($v = \vv{PQ}$ e quindi $\vvbm{PQ} \bm{= P - Q}$) 
\end{enumerate}
\end{Proprieta}

\begin{Definizione}[Spazio affine]
\label{def-spazio-affine}
Uno \textbf{spazio affine con gruppo delle traslazioni} $\bm{\V}$ (che è uno spazio vettoriale su $\K$) è un \textbf{insieme non vuoto} $\bm{\S}$ con un'applicazione
\begin{align*}
\S \times \V &\longrightarrow \S \\
(P,v) &\longmapsto P + v
\end{align*}
con le Proprietà \ref{prop-somma-punto-vettore}. \\
\NB Spesso indicheremo il gruppo di traslazioni dello spazio affine $\S$ con il simbolo $\bm{V(\S)}$.
\end{Definizione}

È facile capire come lo \textit{spazio euclideo} sia uno affine con gruppo delle traslazioni $V(\E^3)$.

\Esempio \nopagebreak

Prese due matrici $A \in M_{mn}(\K)$ e $B \in M_{m1}(\K)$, dimostrare che l'insieme delle soluzioni del sistema di equazioni lineari non omogeneo $AX = B$ $\bm{\S = \{X \in M_{n1}(\K) \mid AX = B\} \neq \varnothing}$ \textbf{sia uno spazio affine}. \nopagebreak

Sia $\V = \{X \in M_{n1}(\K) \mid AX = \undzero\} = \ker A$ l'insieme delle soluzioni del sistema di equazioni lineari omogeneo associato al sistema $AX = B$. Se $X \in \S$ e $Y \in \V$, allora $(X+Y) \in \S$, dato che $A(X + Y) = AX + AY = B + \undzero = B$. Dunque ha senso definire l'applicazione
\begin{align*}
\S \times \V &\longrightarrow \S \\
(X,Y) &\longmapsto X + Y
\end{align*}
Verifichiamo che quest'applicazione rispetti le 3 proprietà:
\begin{enumerate}
    \item $\forall X \in \S \enspace X + \undzero = X$ è vero, dato che $X$ è un vettore normale di $\K^n$ e qualsiasi vettore in questo spazio vettoriale sommato al vettore nullo $\undzero$ ritorna sè stesso.
    \item $\forall X \in \S$ e $\forall Y_1, Y_2 \in \V \enspace X + (Y_1 + Y_2) = (X + Y_1) + Y_2$ è vero, dato che $X,Y_1$ e $Y_2$ sono dei semplici vettori di $\K^n$.
    \item $\forall X,Z \in \S \enspace \exists ! Y \in \V: X + Y = Z$ è vero: infatti, troviamo che $Y = Z - X$ ha senso, dato che $Z - X \in \V$ perché $AX = AZ = B$ significa che $A \cdot (Z - X) = AZ - AX = \undzero$ e dunque $Z - X \in \V$.
\end{enumerate}
Possiamo quindi concludere che $\bm{\S}$ \textbf{è uno spazio affine}.

\begin{Proprieta}[Proprietà spazi affini]
\label{prop-spazi-affini}
\begin{enumerate}[label=\alph*)]
    \item $\forall P \in \S \enspace \vvbm{PP} = \bm{\undzero}$ 
    \item $\forall P,Q,R \in \S \enspace \vvbm{PQ} + \vvbm{QR} = \vvbm{PR}$
    \item $\forall P,Q \in \S \enspace \vvbm{PQ} = - \vvbm{QP}$
\end{enumerate}
\end{Proprieta}

\Dimostrazione \nopagebreak

Nello spazio euclideo queste proprietà ci sembrano ovvie, ma dimostriamo che valgono per tutti gli spazi affini.

\begin{enumerate}[label=\alph*)]
    \item Per la proprietà 1) delle Proprietà \ref{prop-somma-punto-vettore} sappiamo che $P + \undzero = P$, dunque $\undzero = \vv{PP}$. $\blacksquare$
    \item Per dimostrare questa proprietà dobbiamo dimostrare che $P + (\vv{PQ} + \vv{QR}) = R$. Per la proprietà 2) sappiamo che $P + (\vv{PQ} + \vv{QR}) = (P + \vv{PQ}) + \vv{QR}$, e dato che per la proprietà 3) troviamo che $(P + \vv{PQ}) + \vv{QR} = Q + \vv{QR} = R$.
    \item Per la proprietà b) sappiamo che $\vv{PQ} + \vv{QP} = \vv{PP}$, e per la proprietà a) sappiamo che $\vv{PP} = \undzero$. Dunque $\vv{PQ} = - \vv{QP}$.
\end{enumerate}

\begin{Definizione}[Dimensione di uno spazio affine]
\label{def-dimensione-spazio-affine}
La \textbf{dimensione di uno spazio affine} $\S$ è la dimensione del suo gruppo di traslazioni $V(\S)$. Dunque $\bm{\dim \S = \dim V(\S)}$.
\end{Definizione}

\begin{itemize}
    \item Una \textbf{retta} è uno spazio affine di dimensione $\bm{1}$, e per questo motivo uno spazio affine di dimensione $1$ prende il nome di \textit{retta affine}.
    \item Un \textbf{piano} è uno spazio affine di dimensione $\bm{2}$, e per questo motivo uno spazio affine di dimensione $2$ prende il nome di \textit{piano affine}.
    \item Un \textbf{solido} è uno spazio affine di dimensione $\bm{3}$, e per questo motivo uno spazio affine di dimensione $3$ prende il nome di \textit{solido affine}.
\end{itemize}

\section{Combinazione affine di punti}

Abbiamo già visto come sommare un punto dello spazio affine $\S$ con un vettore del suo gruppo di traslazioni. Cerchiamo ora di definire invece una "somma" tra punti. 

Prendiamo il caso della ricerca di un punto medio $M$ tra due punti. Il problema qui è che, per trovare il punto medio tra due punti, è necessario avere un'origine da cui far partire i vettori verso i due punti e trovare così il punto medio.

\begin{minipage}[c]{0.35\textwidth}
\centering
\includegraphics[width=0.8\textwidth, valign=t]{Punto medio.png}
\end{minipage}\hfill
\begin{minipage}[c]{0.60\textwidth}
È però possibile notare che, qualsiasi punto $O \in \S$ noi prendiamo come origine, il punto medio tra $P_0$ e $P_1$ sarà sempre lo stesso (cosa abbastanza ovvia, dato che la metà di un segmento è sempre la stessa). Il grafico a sinistra ci aiuta molto a capire questa cosa: sia che prendiamo $Q$ come orgiine, sia che prendiamo $R$, il punto medio $M$ è sempre nella stessa identica posizione.
\end{minipage}

L'idea è quindi quella di dimostrare che, per calcolare il punto medio, non sia necessario prendere un'origine fissa, dato che un qualsiasi punto nello spazio può funzionare come origine. 

\begin{Proposizione}[]
\label{propos-somma-punti-spazio-affine}
Preso uno spazio affine $\S$ su $\K$ (ossia con il proprio gruppo delle traslazioni $V(\S)$ su $\K$), i coefficienti $\lambda_0, \ldots, \lambda_d \in \K$ tali che $\sum_{i = 0}^d \lambda_i = 1$ e i punti $Q, R, P_0, \ldots, P_d \in \S$, allora
\[
\bm{Q + \sum_{i = 0}^d \lambda_i} \vvbm{QP_i} \bm{= R + \sum_{i = 0}^d \lambda_i} \vvbm{RP_i}
\]
\end{Proposizione}

\Dimostrazione \nopagebreak

Per dimostrare questa proposizione ci basta dimostrare che $Q + \sum_{i = 0}^d \lambda_i \vv{QP_i} - \sum_{i = 0}^d \lambda_i \vv{RP_i} = R$, dato che $R + \sum_{i = 0}^d \lambda_i \vv{RP_i} - \sum_{i = 0}^d \lambda_i \vv{RP_i} = R + \undzero = R$. \\
Con un po' di manipolazione algebrica possiamo concludere proprio che
\begin{align*}
Q + \sum_{i = 0}^d \lambda_i \vv{QP_i} - \sum_{i = 0}^d \lambda_i \vv{RP_i} &= Q + \sum_{i = 0}^d \lambda_i (\vv{QP_i} - \vv{RP_i})  \\
 &= Q + \sum_{i = 0}^d \lambda_i (\vv{QP_i} + \vv{P_iR}) \\
 &= Q + \sum_{i = 0}^d \lambda_i \vv{QR} \\
 &= Q + \vv{QR} \sum_{i = 0}^d \lambda_i \\
 &= R \quad \blacksquare
\end{align*}

\begin{Definizione}[Combinazione affine di punti]
\label{def-combinazione-affine-punti}
Preso uno spazio affine $\S$ su $\K$, i coefficienti $\lambda_0, \ldots, \lambda_d \in \K$ tali che $\sum_{i = 0}^d \lambda_i = 1$ e i punti $P_0, \ldots, P_d \in \S$. allora
\[
\bm{\sum_{i = 0}^d \lambda_i P_i = Q + \sum_{i = 0}^d \lambda_i} \vvbm{QP_i}
\]
Dove $Q$ è un qualsiasi punto appartenente allo spazio affine $\S$.
\end{Definizione}

Abbiamo così definito la combinazione affine di più punti in uno spazio affine. Se torniamo al caso precedente, nel cercare il punto medio tra $P_0$ e $P_1$, troviamo che $M = \frac{1}{2}P_0 + \frac{1}{2}P_1 = Q + \sum_{i = 0}^1 \frac{1}{2} \vv{QP_i} = Q + (\frac{1}{2} \vv{QP_0} + \frac{1}{2} \vv{QP_1})$.

\section{Dipendenza lineare tra punti}

Quando parliamo di vettori linearmente dipendenti, parliamo di vettori $v_1, \ldots, v_n \in \V$ i quali, presi i coefficienti $\lambda_1, \ldots, \lambda_n \in \K$ non tutti nulli, risolvono l'equazione $\lambda_1 v_1 + \ldots + \lambda_n v_n = \undzero$. C'è però un altro modo per definire la dipendenza lineare tra diversi vettori.

Prendiamo ora l'applicazione lineare 
\begin{align*}
L: \K^n &\longrightarrow \V \\
(\lambda_1, \ldots, \lambda_n)
\end{align*}

Questa prende un vettore di $\K^n$ con $n$ coefficienti e restituisce la combinazione lineare tra i vettori $v_1, \ldots, v_n$ con i coefficienti del vettore. Se quindi prendiamo un vettore $\Lambda = (\mu_1, \ldots, \mu_n) \in \ker L$, allora sappiamo che $\mu_1 v_1 + \ldots + \mu_n v_n = \undzero$. Dato che $L$ è un'applicazione lineare, per la Proposizione \ref{propos-applicazione-lineare-iniettiva-nucleo} sappiammo che questa è iniettiva se e solo se $\ker L = \{\undzero\}$, ossia se l'unico vettore di coefficienti tale che $\lambda_1 v_1 + \ldots + \lambda_n v_n = \undzero$ è il vettore $(\lambda_1, \ldots, \lambda_n) = (0,\ldots,0)$. Possiamo quindi dare una nuova definizione di vettori linearmente dipendenti.

\begin{Definizione}[Dipendenza lineare tra vettori]
\label{def-dipendenza-lineare-vettori}
I vettori $v_1, \ldots, v_n$ sono \textbf{linearmente dipendenti} $\iff$ l'applicazione $\bm{L}$ \textbf{non è iniettiva}.
\end{Definizione}

\begin{Definizione}[Dipendenza lineare tra punti]
\label{def-dipendenza-lineare-tra-punti}
Preso uno spazio affine $\S$ su $\K$, i punti $P_0, \ldots, P_d \in \S$ e l'applicazione lineare
\begin{align*}
\varphi_{P_0, \ldots, P_d}: \{(\lambda_0, \ldots, \lambda_d) \in \K^d \mid \sum_{i = 0}^d \lambda_i = 1\} &\longrightarrow \S  \\
(\lambda_1, \ldots, \lambda_n) &\longmapsto \sum_{i = 0}^d \lambda_i P_i 
\end{align*}
$P_0, \ldots, P_d$ sono \textbf{linearmente dipendenti} $\iff$ $\bm{\varphi}$ \textbf{non è iniettiva}.
\end{Definizione}

\Esempi \nopagebreak

Preso lo spazio affine $\E^3$, determinare quando due punti $\bm{P_0,P_1 \in  \E^3}$ sono \textbf{linearmente indipendenti} \nopagebreak

Analizziamo due casi distinti:
\begin{itemize}
    \item $P_0 = P_1$ \\
    In questo caso, $\varphi_{P_0,P_1}$ è un'applicazione costante. Infatti, $\varphi_{P_0,P_1}(\lambda_0,\lambda_1) = \lambda_0 P_0 + \lambda_1 P_1 = P_0(\lambda_0 + \lambda_1)$, ma dato che per definizione $\lambda_0 + \lambda_1 = 1$, troviamo che $\varphi_{P_0,P_1}(\lambda_0,\lambda_1) = P_0$ indipendemente da $\lambda_0$ e $\lambda_1$.
    \item $P_0 \neq P_1$ \\
    In questo caso, $\varphi_{P_0,P_1} = \lambda_0 P_0 + \lambda_1 P_1$, e allo stesso modo \mbox{$\varphi_{P_0,P_1}(\mu_0,\mu_1) = \mu_0 P_0 + \mu_1 P_1$}. Se poniamo $\lambda_0 P_0 + \lambda_1 P_1 = \mu_0 P_0 + \mu_1 P_1$, troviamo che $P_0(\lambda_0 - \mu_0) + P_1(\lambda_1 - \mu_1) = \undzero$. Dato che \mbox{$\lambda_1 = 1 - \lambda_0$} e $\mu_1 = 1 - \mu_0$, troviamo che 
    \begin{align*}
    P_0(\lambda_0 - \mu_0) + P_1(1-\lambda_0 - 1 + \mu_0) &= \undzero  \\
    P_0(\lambda_0 - \mu_0) + P_1(\mu_0 - \lambda_0) &= \undzero \\
    \lambda_0(P_0 - P_1) - \mu_0(P_0 - P_1) &= \undzero \\
    (P_0 - P_1) - (\lambda_0 - \mu_0) &= \undzero
    \end{align*}
    Dato che $P_0 \neq P_1$, $(P_0 - P_1)(\lambda_0 - \mu_0) = \undzero$ è verificata se e solo se $\lambda_0 = \mu_0$, dunque anche $\lambda_1 = \mu_1$. \\
    Abbiamo quindi dimostrato che presi $\lambda_0,\lambda_1, \mu_0,\mu_1 \in \R$, $\varphi_{P_0,P_1} (\lambda_0, \lambda_1) = \varphi_{P_0,P_1}(\mu_0,\mu_1)$ è verificata se e solo se $(\lambda_0,\lambda_1) = (\mu_0,\mu_1)$, dunque $\varphi$ è un'applicazione iniettiva. 
\end{itemize}

Possiamo quindi dire che $P_0$ e $P_1$ sono linearmente indipendenti solo quando $\bm{P_0 \neq P_1}$.

\Separatore

Preso lo spazio affine $\E^3$, determinare quando tre punti $\bm{P_0,P_1,P_2 \in \E^3}$ sono \textbf{linearmente indipendenti}. \nopagebreak

Sappiamo che $\varphi_{P_0,P_1,P_2}(\lambda_0,\lambda_1,\lambda_2) = \lambda_0 P_0 + \lambda_1 P_1 + \lambda_2 P_2$. Dato che $\lambda_0 = 1 - \lambda_1 - \lambda_2$, troviamo che $\varphi_{P_0,P_1,P_2}(\lambda_0,\lambda_1,\lambda_2) = (1-\lambda_1-\lambda_2)P_0 + \lambda_1 P_1 + \lambda_2 P_2$. Allo stesso modo troviamo che $\varphi_{P_0,P_1,P_2}(\mu_0,\mu_1,\mu_2) = (1-\mu_1-\mu_2) P_0 + \mu_1 P_1 + \mu_2 P_2$. Se poniamo ora $\varphi_{P_0,P_1,P_2}(\lambda_0,\lambda_1,\lambda_2) = \varphi_{P_0,P_1,P_2}(\mu_0,\mu_1,\mu_2)$, troviamo che 
\begin{align*}
(1-\lambda_1 - \lambda_2)P_0 + \lambda_1 P_1 + \lambda_2 P_2 &= (1-\mu_1-\mu_2) P_0 + \mu_1 P_1 + \mu_2 P_2 \\
P_0(\mu_1 - \lambda_1 + \mu_2 - \lambda_2) + P_1(\lambda_1 - \mu_1) + P_2(\lambda_2 - \mu_2) &= \undzero \\
(P_1 - P_0)(\lambda_1 - \mu_1) + (P_2 - P_0)(\lambda_2 - \mu_2) &= \undzero 
\end{align*}

Dato che $P_1 - P_0 = \vv{P_0P_1}$ e allo stesso modo $P_2 - P_0 = \vv{P_0P_2}$, troviamo che \mbox{$\vv{P_0P_1}(\lambda_1 - \mu_1) + \vv{P_0P_2}(\lambda_2 - \mu_2) = \undzero$}. Affinché $\varphi_{P_0,P_1,P_2}$ sia iniettiva, è necessario che $\lambda_1 = \mu_1$ e che $\lambda_2 = \mu_2$. Per dimostrare ciò, ci basta dimostrare che $\vv{P_0P_1}$ e $\vv{P_0P_2}$ sinao lineramente indipendenti. Questo è vero se e solo se $P_0,P_1$ e $P_2$ non sono sulla stessa retta. Dunque concludiamo che $P_0,P_1$ e $P_2$ sono linearmente indipendenti solo quando \textbf{non si trovano sulla stessa retta}.

\begin{Osservazione}[]
\label{oss-punti-indipendenti}
$\bm{P_0,\ldots,P_d}$ sono punti \textbf{linearmente indipendenti} $\iff \vvbm{P_0P_1}\bm{, \ldots, }\vvbm{P_0P_d}$ sono \textbf{linearmente indipendenti}
\end{Osservazione}

Questa osservazione ci viene dall'esempio precedente con i 3 punti in $\E^3$, e ci semplifica molto il verificare se $d+1$ vettori sono linerarmente indipendenti o meno, dovendo solo verificare l'indipendenza di $d$ vettori.

\section{Sottospazi affini}

\begin{Definizione}[Sottospazio affine]
\label{def-sottospazio-affine}
Preso uno spazio affine $\S$, un sottoinsieme $\T \subset \S$ è un \textbf{sottospazio affine} se esiste un sottospazio vettoriale $\W \subset V(\S)$ e un punto $P_0 \in \S$ tali che $\bm{\T = P_0 + \W = \{P_0 + w \mid w \in \W\}}$.
\end{Definizione}

Ad esempio, $\T \subset A^n(\K)$ (dove $A^n(\K)$ indica lo spazio affine con gruppo di traslazioni $\K^n$) è un sottospazio affine se $\T = \overline{\X} + \W$, dove $\overline{\X} \in A^n(\K)$, e se $\W \subset \K^n$.

\begin{Proposizione}[]
\label{propos-sottospazio vettoriale e sottospazio affine}
Preso un sottospazio vettoriale $\V$ di $\K^n$, $\bm{\V}$ è anche un \textbf{sottospazio affine} di $A^n(\K)$.
\end{Proposizione}

\Dimostrazione \nopagebreak

La dimostrazione è evidente: infatti, se poniamo $\overline{\X} = \undzero$, troviamo che $A^n(\K) \supset \T = \undzero + \V = \V$. $\blacksquare$ \\
È importante però sottolineare che il contrario non è sempre vero: infatti, se prendiamo il sottospazio vettorile $\W = \{\undzero\}$ di $\K^n$, e $\overline{\X} \neq 0$, allora $\T = \overline{\X} + \V = \{\overline{\X}\} \not\ni \undzero$, dunque $\T$ non può essere un sottospazio vettoriale di $\K^n$. 

\begin{Proposizione}[]
\label{propos-sottospazio-affine}
Preso un sottospazio affine $\T \subset \S$, due punti $P_0,Q_0 \in \S$ e due sottospazi vettoriali $\W$ e $\U$ di $V(\S)$, $\bm{\T = P_0 + \W = Q_0 + \U} \implies \bm{\W = \U}$.
\end{Proposizione}

\Dimostrazione \nopagebreak

Per dimostrare che $\W = \U$, basta dimostrare che $\U \subset \W$ e che $\W \subset \U$. Dato che i due passaggi sono pressoché identici, dimostreremo solo che $\U \subset \W$. \\
Dato che $Q_0 \in \T$, sappiamo che esiste un vettore $w_0 \in \W$ tale che $P_0 + w_0 = Q_0$. Sappiamo anche che qualsiasi punto $P \in \T$ può essere scritto come $P = Q_0 + u$, dove $u$ è un vettore in $\U$. Dato che $Q_0 = P_0 + w_0$, troviamo che $P = (P_0 + w_0) + u = P_0 + (w_0 + u)$. Ma sappiamo anche che qualsiasi punto $P \in \T$ può essere riscritto come $P = P_0 + w$, dove $w$ è un vettore in $\W$. Dunque troviamo che $P_0 + (w_0 + u) = P_0 = + w$, ossia $u = w - w_0$. Abbiamo quindi dimostrato che un qualsiasi vettore di $\U$ è scrivibile come combinazione lineare di vettori di $\W$, ossia $\U \subset \W$. $\blacksquare$

\NB Solo da ora potremo indicare con $V(\T)$ il sottospazio vettoriale associato a $\T$, dato che abbiamo dimostrato che vi è un unico sottospazio che genera $\T$.

\begin{Proposizione}[]
\label{propos-sottospazio-affine-spazio}
Un sottospazio affine è anche uno \textbf{spazio affine}.
\end{Proposizione}

\Dimostrazione \nopagebreak

Se prendiamo un punto $P \in \T$ e un vettore $v \in V(\T)$, sappiamo che $P + v \in \T$. Possiamo quindi definire l'applicazione
\begin{align*}
\T \times V(\T) &\longrightarrow \T \\
(P,v) &\longmapsto P + v 
\end{align*} 
che rispetta tutte le Proprietà \ref{prop-somma-punto-vettore}. Dunque $\T$ è uno spazio affine. $\blacksquare$

\begin{Definizione}[Sottospazio affine generato]
\label{def-sottospazio-affine-generato}
Preso uno spazio affine $\S$ e i punti $P_0, \ldots, P_d \in \D$, il \textbf{sottospazio affine} di $\S$ \textbf{generato da tali punti} è 
\[\bm{\braket{P_0,\ldots,P_d}} := \left\{\sum_{i = 0}^d \lambda_i P_i \mid \sum_{i = 0}^d \lambda_i = 1\right\} = \left\{P_0 + \sum_{i = 1}^d \lambda_i \vv{P_0P_i}\right\} = \bm{P_0 + \langle} \vvbm{P_0P_1}\bm{,\ldots,} \vvbm{P_0P_d} \bm{\rangle}\] 
\NB Abbiamo potuto eliminare la condizione sui $\lambda$ perché, presi tutti i $\lambda_i$, ci basta prendere $\lambda_0 = 1 - \sum_{i = 1}^d \lambda_i$ per avere che $\sum_{i = 0}^d \lambda_i = 1$.
\end{Definizione}

\begin{Osservazione}[]
\label{oss-dimensione-sottospazio-affine-generato}
$\bm{\dim \langle P_0,\ldots,P_d \rangle \leq d}$ e in particolare \\ $\bm{\dim \langle P_0,\ldots,P_d \rangle = d} \iff P_0,\ldots,P_d$ sono \textbf{linearmente indipendenti}.
\end{Osservazione}

\Dimostrazione \nopagebreak

Sappiamo che $\dim \braket{P_0,\ldots,P_d} = \dim \braket{\vv{P_0P_1}, \ldots, \vv{P_0,P_d}}$ e dato che $\dim \braket{\vv{P_0P_1}, \ldots, \vv{P_0,P_d}} \leq d$, dove in particolare $\dim \braket{\vv{P_0P_1}, \ldots, \vv{P_0P_d}} = d \iff \vv{P_0P_1}, \ldots, \vv{P_0P_d}$ sono linearmente indipendenti, troviamo che $\dim \braket{P_0, \ldots, P_d} \leq d$. $\blacksquare$

È anche grazie a quest'osservazione che riusciamo a dimostrar ein modo formale l'Osservazione \ref{oss-punti-indipendenti}. Infatti, se $\vv{P_0P_1},\ldots,\vv{P_0P_d}$ sono linearmente indipendenti allora $\dim \braket{P_0,\ldots,P_d} = d$, e dunque $P_0,\ldots,P_d$ sono linearmente indipendenti, dato che $\ker \braket{\vv{P_0P_1},\ldots,\vv{P_0P_d}} = 0$ e dunque $\varphi$ è iniettiva. 

\section{Applciazioni affini e coordinate affini}
\subsection{Applicazioni affini}

\begin{Definizione}[Applicazione affine]
\label{def-applicazione-affine}
Un'applicazione $F: \S \longrightarrow \T$ è un'applicazione \textbf{affine} se esiste un'applicazione lineare $V(F): V(\S) \longrightarrow V(\T)$ tale che, presi due punti $P,Q \in \S$, $\vvbm{F(P) F(Q)} \bm{= V(F)(}\vvbm{PQ}\bm{)}$.
\end{Definizione}

In parole più semplici, $F$ è un'applicazione affine da $\S$ a $\T$ se esiste un'applicazione lineare tra i gruppi di traslazione $V(\S)$ e $V(\T)$ tale che $\vv{F(P)F(Q)} = V(F)(\vv{PQ})$.

\Esempio \nopagebreak

Preso $v \in V(\S)$, dimostrare che l'applicazione \nopagebreak

\begin{align*}
\tau_v: \S &\longrightarrow \S  \\
P &\longmapsto P + v
\end{align*}
è un'\textbf{applicazione affine} con $\bm{V(\tau_v) = Id_{V(\S)}}$. \nopagebreak

Dobbiamo quindi dimostrare che $\vv{\tau_v (P) \tau_v(Q)} = V(\tau_v)(\vv{PQ}) = Id_{V(\S)}(\vv{PQ}) = \vv{PQ}$. Sappiamo che $\tau_v(P) = P + v$ e che $\tau_v(Q) = Q + v$. Dunque troviamo che $\tau_v(P) + \vv{PQ} = \tau_v(Q)$ e che quindi $\vv{PQ} = \tau_v(Q) - \tau_v(P)$. Ma dato che $\tau_v (P) + \vv{\tau_v(P) \tau_v(Q)} = \tau_v(Q)$, troviamo che \mbox{$\tau_v(P) - \tau_v(Q) = \vv{\tau_v(P) \tau_v(Q)}$}, ossia troviamo che $\vv{PQ} = \vv{\tau_v (P) \tau_v(Q)}$. Dunque $\bm{\tau_v}$ è un'\textbf{applicazione affine}.

\subsection{Coordinate affini}
\begin{Definizione}[Riferimento affine]
\label{def-riferimento-affine}
In uno spazio affine $\S$ di dimensione $n$ un \textbf{riferimento affine} per $\S$ è dato da una coppia $\bm{R = RA(O;\B)}$, dove $O$ è un punto in $\S$ e $\B = \{v_1, \ldots, v_n\}$ è una base di $V(\S)$.
\end{Definizione}

Con un riferimento affine è possibile prendere un qualsiasi punto $P \in \S$ e calcolare il vettore $\vv{OP} = \lambda_1 v_1 + \ldots + \lambda_n v_n$, dove $(\lambda_1, \ldots, \lambda_n)$ rappresentano le \textbf{coordinate affini} del punto $P$ nel riferimento affine $R$.

\begin{Definizione}[Coordinate affini]
\label{def-coordinate-affini}
Preso uno spazio affine $\S$ di dimensione finita e un riferimento affine $R = RA(O; \B)$, definiamo l'applicazione
\begin{align*}
X_R: \S &\longrightarrow A^n(\K)  \\
 P &\longmapsto X_\B(\vv{OP}) 
\end{align*}
che restituisce le \textbf{coordinate affini} di un qualsiasi punto $P$ di $\S$.

\NB Dato che $\B$ è una base di $V(\S)$, per il Corollario \ref{propos-corol-basi} sappiamo che l'applicazione $X_R$ è biunivoca.
\end{Definizione}

\section{Equazioni parametriche e cartesiane di uno spazio affine}
\subsection{Equazioni parametriche}

\begin{Definizione}[Equazioni parametriche]
\label{def-equazioni-parametriche}
Preso uno spazio affine $\S$, una base $\B = \{w_1, \ldots, w_d\}$ di $\W \subset V(\S)$, un riferimento affine $R = RA(O,\B)$ e un sottospazio affine $\T = P_0 + \W = \{P_0 + \sum_{i = 1}^d t_i w_i \mid (t_1, \ldots, t_d) \in \K^d\}$ con gruppo delle traslazioni $\W$, le \textbf{equazioni parametriche} sono un sistema di equazioni lineari che ci permette di trovare, al variare dei parametri ($t_1, \ldots, t_d$), tutti i punti appartenenti al sottospazio affine $\T$.
\end{Definizione}

Se per esempio abbiamo il punto $P_0(1,2,3)$ (con $P_0(1,2,3)$ indichiamo un punto con coordinate affini $(1,2,3)$ nel nostro riferimento affine) e $\B = \{w\}$ con $w = (4,-1,0)$, troviamo che $\T = \{P_0 + tw \mid t \in \K\} = \{(1,2,3) + (4t,-t,0) \mid t \in \K\} = \{(4t+1,2-t,3) \mid t \in \K\}$. \\
Dunque le equazioni parametriche di $\T$ sono
\[
\begin{cases}
&x_1 = 4t + 1 \\
&x_2 = 2 - t  \\
&x_3 = 3
\end{cases}
\]

$\X = (x_1, x_2, x_3)$ sarà quindi un qualsiasi punto appartenente a $\T$.

Possiamo dire che le equazioni parametriche funzionano come un'applicazione che prende in input dei parametri e che restituisce solo punti appartenenti al sottospazio affine $\T$. Dunque, se chiamiamo quest'applicazione $f$, troviamo che $\bm{\T = Im(f)}$.

\subsection{Equazioni cartesiane}

\begin{Definizione}[Equazioni cartesiane]
\label{def-equazioni-cartesiane}
Preso un sottospazio affine $\S$ di dimensione $n$, un riferimento affine $R = RA(O,\B)$ e un sottospazio affine $\T = P_0 + \W$ di dimensione $d < n$, le \textbf{equazioni cartesiane} sono un sistema di $n-d$ equazioni lineari che ci permettono di verificare se un punto $\X = (x_1, \ldots, x_n) \in \S$ appartenga a $\T$.
\end{Definizione}

Riprendendo l'esempio visto in precedenza con le equazioni parametriche, abbiamo che $\T = \{(4t + 1, 2-t, 3) \mid t \in \K\}$. Dato che $\dim \S = 3$ e $\dim \T = 1$, il sistema avrà du eequazioni. Affinché $\X = (x_1, x_2, x_3) \in \T$, c'è bisogno di risolvere il sistema
\[
\begin{cases}
    x_1 = 1 + 4t \\
    x_2 = 2 - t \\
    x_3 = 3
\end{cases}
\]
che diventa
\[
\begin{cases}
 \bm{x_1 = 9 - 4x_2} \\
 \bm{x_3 = 3}
\end{cases}
\]

Le equazioni di quest'ultimo sistema sono proprio le nostre equazioni cartesiane, che ci dicono che un punto $\X \in \S$ appartiene anche a $\T$ se risolve questo sistema. 

Possiamo dire che le equazioni cartesiane funzionano come un'applicazione che prende in input un punto di $\S$ e restituisce $\undzero$ se tale punto appartiene anche a $\T$. Dunque, se chiamiamo quest'applicazione $g$, troviamo che $\bm{\T = \ker g}$.

\chapter{Forme quadratiche e spazi vettoriali euclidei}

\section{Nozioni preliminari}

Prima di parlare di forme quadratiche, forma bilineari simmetriche e spazi vettoriali euclidei, è fondamentale conoscere queste due definizioni, che ora vediamo riguardo lo spazio euclideo ma che poi estenderemo a tutti gli spazi vettoriali.

\begin{Definizione}[Prodotto scalare euclideo su $\bm{V(\E^3)}$]
\label{def-prodotto-scalare-euclideo-E^3}
Presi due vettori $v,w \in V(\E^3)$, definiamo il prodotto scalare euclideo di $v,w$ come $ \bm{\langle v,w \rangle =} \normbm{v} \bm{\cdot} \normbm{w} \bm{\cdot \cos(\theta)}$, dove $\norm{v}$ prende il nome di \textbf{norma} e indica la \textbf{lunghezza} di $v$ nel nostro sistema di riferimentom e $0 \leq \theta \leq \pi$ indica il più piccolo angolo compreso tra $v$ e $w$.
\end{Definizione}

Da questa definizione è facile convincersi che
\begin{itemize}
    \item $\braket{v,w} = 0 \iff \theta = \frac{\pi}{2}$, ossia $v \perp w$
    \item $\braket{v,w} > 0 \iff \theta < \frac{\pi}{2}$
    \item $\braket{v,w} < 0 \iff \theta > \frac{\pi}{2}$
    \item $\braket{v,v} = \norm{v}^2$
\end{itemize}

\begin{Definizione}[Base ortonormale in $\bm{V(\E^3)}$]
\label{def-base-ortonormale}
Una base $\B = \{v_1, \ldots, v_n\}$ di $V(\E^n)$ prende il nome di base \textbf{ortonormale} se 
\[
\begin{cases}
\normbm{v_i} \bm{= 1} &\enspace \forall i \in \{1, \ldots, n\}  \\
\bm{v_i \perp v_j} &\enspace \forall i \neq j 
\end{cases}
\]
Dove $v_i \perp v_j$ indica che $\theta(v_i, v_j) = \frac{\pi}{2}$ (indichiamo con $\theta(v_i,v_j)$ il più piccolo angolo $\theta$ tra i due vettori).
\end{Definizione}

\section{Forme quadratiche}

\begin{Definizione}[Forma quadratica]
\label{def-forma-quadratica}
Preso uno spazio vettoriale $\V$ su $\K$ finitamente generato, una \textbf{forma quadratica} su $\V$ è un'applicazione $Q: \V \longrightarrow \K$ tale che, in una base $\B = \{v_1, \ldots, v_n\}$,
\[
\bm{Q(x_1 v_1 + \ldots + x_n v_n) = \sum_{1 \leq i \leq j \leq n} a_{ij} x_i x_j}
\]
dove $a_{ij} \in \K$. Questo altro non è che un \textit{polinomio omogeneo} in $x_1, \ldots, x_n$ di \textit{grado} $2$, ossia un polinomio dove tutti i monomi sono di grado totale $2$.

\NB Cambiando base $\B$ si ottengono polinomi anche molto diversi tra loro, ma tutti sempre omogenei e di grado $2$.
\end{Definizione}

Un'esempio di forma quadratica è l'applicazione
\begin{align*}
Q: V(\E^3) &\longrightarrow R \\
v &\longmapsto \norm{v}^2
\end{align*}
In questo caso, $Q$ è una forma quadratica su $\V$, con base ortonormale $\B = \{v_1, v_2, v_3\}$. Infatti, sappiamo che ogni vettore $\X$ di $V(\E^3)$ può essere riscritto come $x_1 v_1 + x_2 v_2 + x_3 v_3$, con $x_1, x_2, x_3 \in \R$. Dunque $Q(\V) = \norm{v}^2 = \braket{v,v} = \braket{x_1v_1 + x_2v_2 + x_3 v_3, x_1v_1 + x_2v_2 + x_3v_3}$. A questo punto, anticipiamo una proprietà che vedremo tra poco e sfruttiamo la bilinearità del prodotto scalare. La rivedremo più avanti, ma per ora non ci serve sapere altro che \mbox{$\braket{v_1 + v_2, v_1 + v_2} = \braket{v_1, v_1} + \braket{v_1,v_2} + \braket{v_2, v_1} + \braket{v_2,v_2}$} e che $\braket{\lambda v, v} = \lambda \braket{v,v}$. Tornando ora al nostro esempio, applicando tante volte questa proprietà è possibile dire che
\begin{align*}
Q(\V) &= x_1^2\braket{v_1,v_1} + x_1x_2 \braket{v_1,v_2} + x_1 x_3 \braket{v_1,v_3}  \\
&+ x_2x_1 \braket{v_2,v_1} + x_2^2 \braket{v_2,v_2} + x_2 x_3 \braket{v_2,v_3} \\
&+ x_3x_3 \braket{v_3,v_1} + x_3x_2 \braket{v_3,v_2} + x_3^2 \braket{v_3,v_3} 
\end{align*}

Dato che la base $\B$ è ortonormale, tutti i termini $\braket{v_i, v_j}$ con $i \neq j$ sono pari a $0$, mentre tutti quelli con $i = j$ sono pari a $1$. Dunque troviamo che $Q(\V) = x_1^2 + x_2^2 + x_3^2$. Questa non è altro che una forma quadratica dove
\[
a_{ij} = 
\begin{cases}
0 & \enspace \text{se } i \neq j \\
1 & \enspace \text{se } i = j 
\end{cases}
\]

\begin{Definizione}[Insieme delle forme quadratiche]
\label{def-insieme-forme-quadratiche}
Chiamiamo  $\bm{Q(\V)}$ \textbf{l'insieme delle forme quadratiche} su $\V$.
\end{Definizione}

Definiamo le operazioni di \textbf{somma} e \textbf{prodotto per scalare} per $Q(\V)$.

\bigskip

\begin{minipage}[c]{0.45\textwidth}
    \centering
    {\Large\textbf{Somma}:}
    \begin{align*}
        \hspace{0.4cm}Q(\V) \times Q(\V) &\longrightarrow Q(\V) \\
        \hspace{0.4cm}(Q_1,Q_2) &\longmapsto Q_1 + Q_2 \\
        \hspace{0.4cm}\text{Dove } (Q_1 + Q_2)(v) &= Q_1(v) + Q_2(v)
    \end{align*}
\end{minipage}
\begin{minipage}[c]{0.45\textwidth}
    \centering
    {\Large\textbf{Prodotto per scalare}:}
    \begin{align*}
        \K \times Q(\V) &\longrightarrow Q(\V) \\
        (\lambda,Q) &\longmapsto \lambda Q \\
        \text{Dove } (\lambda Q)(v) &= \lambda \cdot Q(v)
    \end{align*}
\end{minipage}

\bigskip

\begin{Proposizione}[]
\label{propos-insieme-forme-quadratiche-spazio-vett}
Con le operazioni di somma e prodotto per scalare appena descritte, $\bm{Q(\V)}$ è uno \textbf{spazio vettoriale} su $\K$.
\end{Proposizione}

\Dimostrazione \nopagebreak

Sia $\B = \{v_1, \ldots, v_n\}$ una base di $\V$. Per prima cosa, dobbiamo assicurarci che l'applicazione nulla $\undzero$ si trovi in $Q(\V)$. Dunque dobbiamo verificare che \mbox{$\undzero(x_1 v_1 + \ldots + x_n v_n) = \sum_{1 \leq i \leq j \leq n} a_{ij} x_i x_j = 0$}. Questo è vero se prendiamo $a_{ij} = 0 \enspace \forall i,j$. Dunque $\undzero \in Q(\V)$ ($\undzero$ è anch'esso un polinomio di grado 2 omogeneo, semplicemente con i coefficienti tutti nulli). \\
Dobbiamo poi dimostrare che $Q(\V)$ è un insieme chiuso rispetto alla somma, ossia che prese due forme quadratiche $Q_1, Q_2 \in Q(\V)$, $(Q_1 + Q_2) \in Q(\V)$. Essendo $\B$ una base di $\V$, troviamo che \mbox{$(Q_1 + Q_2)(v) = Q_1(v) + Q_2(v) = Q_1(\sum_{i = 1}^n x_i v_i) + Q_2(\sum_{i = 1}^n x_i v_i)$}. Dato che sia $Q_1$ che $Q_2$ sono delle forme quadratiche, troviamo che \mbox{$(Q_1 + Q_2)(v) = \sum_{1 \leq i \leq n} b_{ij} x_i x_j + \sum_{1 \leq i \leq j  \leq n} c_{ij}x_i x_j = \sum_{1 \leq i \leq j \leq n} (b_{ij} + c_{ij} x_i x_j)$}. Ponendo ora $a_{ij} = b_{ij} + c_{ij}$, troviamo che $(Q_1 + Q_2)(v) = \sum_{1 \leq i \leq j \leq n} a_{ij} x_i x_j$, ossia $(Q_1 + Q_2)(v) \in Q(\V)$. \\
Dobbiamo infine dimostrare che $Q(\V)$ è un insieme chiuso rispetto al prodotto per scalare, ossia che presa una forma quadratica $Q \in Q(\V)$ e preso uno scalare $\lambda \in \K$, $\lambda Q \in Q(\V)$, $\lambda Q \in Q(V)$. Dato che \mbox{$\lambda Q(v) = \lambda Q(\sum_{i = 1}^n x_i v_i) = \lambda \sum_{1 \leq i \leq j \leq n} b_{ij} x_i x_j = \sum_{1 \leq i \leq j \leq n} (\lambda b_{ij}) x_i x_j$}. Ponendo ora $a_{ij} = \lambda b_{ij}$, troviamo che $\lambda Q(v) = \sum_{1 \leq i \leq j \leq n} a_{ij} x_i x_j$, ossia che $\lambda Q \in Q(\V)$.\\
Possiamo quindi concludere che $Q(\V)$ è uno spazio vettoriale. $\blacksquare$

\section{Forme bilineari simmetriche e forme quadratiche}
\subsection{Correlazione tra applicazioni bilineari simmetriche e forme quadratiche}
\NB Per il resto del capitolo consideremo di essere sempre in un campo $\K$ dove $2 \neq 0$.

\begin{Osservazione}[]
\label{oss-forma-bilineare-quadratica}
Preso uno spazio vettoriale $\V$ finitamente generato su $\K$ e presa una \textbf{forma bilineare} $F: \V \times \V \longrightarrow \K$, ossia un'applciazione bilineare, allora l'applicazione
\begin{align*}
\bm{q_F: \V} &\longrightarrow \bm{\K} \\
\bm{v} &\longmapsto \bm{F(v,v)}
\end{align*}
è una \textbf{forma quadratica}.
\end{Osservazione}

\Dimostrazione \nopagebreak

Prendiamo una base $\B = \{v_1, \ldots, v_n\}$ di $\V$. Dunque sappiamo che $q_F(v) = q_F(\sum_{i = 1}^n x_i v_i) = F(\sum_{i = 1}^n x_i v_i, \sum_{j = 1}^n x_j v_j)$. Essendo $F$ una forma bilineare, sappiamo che $F(\sum_{i = 1}^n x_i v_i, \sum_{j = 1}^n x_j v_j) = \sum_{1 \leq i,j \leq n} x_i x_j F(v_i, v_j)$. Questa non è altro che la definizione di un'applicazione bilineare. INfatti, nonostante ci siano dei termini in più, ossia quando $i > j$, basta semplicemente porre 
\[
a_{ij} = 
\begin{cases}
F(v_i,v_i) & \enspace \text{se } i = j \\
F(v_i,v_j) + F(v_j,v_i) & \enspace \text{se } i \neq j  
\end{cases}
\]
In questo modo, troviamo proprio che $q_F(v) = \sum_{1 \leq i \leq j \leq n} a_{ij} x_i x_j$. $\blacksquare$

\begin{Definizione}[Insieme delle forme bilineari]
\label{def-insieme-forme-bilineari}
Chiamiamo $\bm{\Bil(\V)}$ \textbf{l'insieme delle forme bilineari} che vanno da $\V \times \V$ a $\K$.
\end{Definizione}

Definiamo le operazioni di \textbf{somma} e \textbf{prodotto per scalare} per $\Bil(\V)$.

\bigskip

\begin{minipage}[c]{0.45\textwidth}
    \centering
    {\Large\textbf{Somma}:}
    \begin{align*}
        \hspace{0.4cm}\Bil(\V) \times \Bil(\V) &\longrightarrow \Bil(\V) \\
        \hspace{0.4cm}(F,G) &\longmapsto F + G \\
        \hspace{0.4cm}\text{Dove } (F + G)(v) &= F(v) + G(v)
    \end{align*}
\end{minipage}
\begin{minipage}[c]{0.45\textwidth}
    \centering
    {\Large\textbf{Prodotto per scalare}:}
    \begin{align*}
        \K \times \Bil(\V) &\longrightarrow \Bil(\V) \\
        (\lambda,F) &\longmapsto \lambda F \\
        \text{Dove } (\lambda F)(v) &= \lambda \cdot F(v)
    \end{align*}
\end{minipage}

\bigskip

\begin{Proposizione}[]
\label{propos-bil-spazio-vettoriale}
Con le operazioni di somma e prodotto per scalare appena descritte, $\bm{\Bil(\V)}$ è uno \textbf{spazio vettoriale} su $\K$.
\end{Proposizione}

\Dimostrazione \nopagebreak

Per prima cosa, dobbiamo assicurarci che l'applicazione nulla $\undzero$ si trovi in $\Bil(\V)$. Dunque dobbiamo verificare che $\undzero(\alpha u + v, w) = \alpha \undzero(u,w) + \undzero(v,w)$ (dovremmo verificare la stessa cosa anche per il secondo membro, ma il procedimento è identico e quindi lo omettiamo). Dato che $\undzero(v,w) = \undzero$ sempre, troviamo che $\undzero(\alpha u + v, w) = \undzero$ e che anche $\alpha \undzero(u,w) + \undzero(v,w) = \undzero$. Dunque $\undzero \in \Bil(\V)$. \\
Dobbiamo poi dimostrare che $\Bil(\V)$ è un insieme chiuso rispetto alla somma, ossia che prese due forme bilineari $F,G \in \Bil(\V)$, $(F+G) \in \Bil(\V)$. Sappiamo che
\begin{align*}
(F+G)(\alpha u + v, w) &= F(\alpha u + v, w) + G(\alpha u + v, w) \\
 &= \alpha F(u, w) + F(v,w) + \alpha G(u,w) + G(v,w) \\
 &= \alpha(F(u,w) + G(u,w)) + F(v,w) + G(v,w) \\
 &= \alpha (F+G)(u,w) + (F+G)(v,w)
\end{align*}
dunque troviamo che $(F+G) \in \Bil(\V)$. \\
Dobbiamo infine dimostrare che $\Bil(\V)$ è un insieme chiuso rispetto al prodotto per scalare, ossia che presa una forma bilineare $F \in \Bil(\V)$ e uno scalare $\lambda \in \K$, $\lambda F \in \Bil(\V)$. Dato che \mbox{$\lambda F(\alpha u + v, w) = \lambda \cdot F(\alpha u + v, w) = \lambda \cdot (\alpha F(u,w) + F(v,w)) = \alpha \lambda F(u,w) + \lambda F(v,w)$}, troviamo proprio che $\lambda F \in \Bil(\V)$. \\
Possiamo quindi concludere che $\Bil(\V)$ è uno spazio vettorile. $\blacksquare$

\begin{Definizione}[Forma bilineare simmetrica]
\label{def-forma-bilineare-simmetrica}
Presa una forma bilineare $F \in \Bil(\V)$, questa è una \textbf{forma bilineare simmetrica} se $\bm{F(v,w) = F(w,v)} \enspace \forall v,w \in \V$.
\end{Definizione}

Un esempio di forma bilineare simmetrica può essere proprio il prodotto scalare euclideo di $V(\E^3)$. Infatti, \mbox{$\braket{v,w} = \norm{v} \cdot \norm{w} \cdot \cos(\theta) = \norm{w} \cdot \norm{v} \cdot \cos(\theta) = \braket{w,v}$}. 

\begin{Definizione}[Forma bilineare simmetrica standard]
\label{def-forma-bilineare-simmetrica-standard}
La \textbf{forma bilineare simmetrica standard} su $\K^n$ è $\bm{F(X,Y) = X^t \cdot Y  = \sum_{i = 1}^n x_i y_i}$, con $X,Y \in M_{n1}$.
\end{Definizione}

\begin{Definizione}[Insieme delle forme bilineari simmetriche]
\label{def-insieme-forme-bilineari-simmetriche}
Chiamiamo $\bm{\Bil^+(\V)}$ \textbf{l'insieme delle forme bilineari simmetriche} che vanno da $\V \times \V$ a $\K$.
\end{Definizione}

\begin{Proposizione}[]
\label{propos-Bil+-spazio-vettoriale}
Con le stesse operazioni di somma e prodotto per scalare descritte per lo spazio vettoriale $\Bil(\V)$, anche $\bm{\Bil^+(\V)}$ è uno \textbf{spazio vettoriale} su $\K$.
\end{Proposizione}

\Dimostrazione \nopagebreak

Per prima cosa, dobbiamo assicurarci che l'applicazione nulla $\undzero$ si trovi in $\Bil^+(\V)$. Sappiamo già che $\undzero \in \Bil(\V)$, dunque dobbiamo solo verificare che $\undzero(v,w) = \undzero(w,v)$, ma questo è ovvio dato che in entrambi i casi il risultato è $\undzero$. Dunque concludiamo che $\undzero \in \Bil^+(\V)$. \\
Dobbiamo poi dimostrare che $\Bil^+(\V)$ è un insieme chiuso rispetto alla somma, ossia che prese due forme bilineari simmetriche $F,G \in \Bil^+(\V)$, $(F+G) \in \Bil^+(\V)$. Questo è evidente dato che $(F+G)(v,w) = F(v,w) + G(v,w) = F(w,v) + G(w,v) = (F+G)(w,v)$, dunque troviamo che $(F+G) \in \Bil^+(\V)$. \\
Dobbiamo infine dimostrare che $\Bil^+(\V)$ è un insieme chiuso rispetto al prodotto per scalare, ossia che presa una forma bilineare simmetrica $F \in \Bil^+(\V)$ e uno scalare $\lambda \in \K$, $\lambda F \in \Bil^+(\V)$. Dato che $\lambda F(v,w) = \lambda \cdot F(v,w) = \lambda \cdot F(w,v) = \lambda F(w,v)$, troviamo che $\lambda F \in \Bil^+(\V)$.\\
Possiamo quindi concludere che $\Bil(\V)$ è uno spazio vettoriale. $\blacksquare$

\begin{Proposizione}[]
\label{propos-isomorfismo-bilineari-quadratiche}
L'applicazione
\begin{align*}
\bm{\Phi: \Bil^+(\V)} &\longrightarrow \bm{Q(\V)}  \\
\bm{F} &\longmapsto \bm{q_F} 
\end{align*}
è un \textbf{isomorfismo} di spazi vettoriali.
\end{Proposizione}

\Dimostrazione \nopagebreak

Dobbiamo quindi dimostrare che $\Phi$ è un'applicazione lineare ed è biiettiva.

\begin{itemize}
    \item $\Phi$ è lienare \\
    Per dimostrare che $\Phi$ è lineare, dobbiamo dimostrare che, prese due forme bilineari simmetriche $F,G \in \Bil^+(\V)$ e due scalari $\lambda, \mu \in \K$, $\Phi(\lambda F + \mu G) = \lambda \Phi(F) + \mu \Phi(G) = \lambda q_F + \mu q_G$. \\
    Con un po' di manipolazione algebrica possiamo concludere proprio che
    \begin{align*}
    \Phi(\lambda F + \mu G)(v) &= q_{\lambda F + \mu G}(v) \\
     &= (\lambda F + \mu G)(v,v) \\
     &= \lambda F (v,v) + \mu G(v,v) \\
     &=  \lambda q_F(v) + \mu q_G(v) \quad \blacksquare
    \end{align*}
    \item $\Phi$ è biiettiva \\
    Per dimostrare che $\Phi$ è biiettiva, dobbiamo dimostrare che questa è sia iniettiva che suriettiva.
    \begin{itemize}
        \item $\Phi$ è iniettiva \\
        Essendo $\Phi$ lineare, ci basta dimostrare che $\ker \Phi = \{\undzero\}$. Sappiamo già che \mbox{$q_F(v+w) = F(v,v) + 2F(v,w) + F(w,w)$}, essendo $F$ una forma bilineare simmetrica. Dato che $F(v,v) = q_F(v)$ e $F(w,w) = q_F(w)$, troviamo che $q_F(v+w) = q_F(v) + 2F(v,w) + q_F(w)$. Se $\Phi(F) = \undzero$, allora $\Phi(F) = q_F(v) = \undzero$ sempre. Ma sappiamo che $q_F(v) = F(v,v)$, e dato che $q_F(v) = \undzero$, anche $F(v,v) = \undzero$. Troviamo quindi che $F(v+w,v+w) = 2F(v,w) = \undzero$ e quindi $F(v,w) = \undzero$, dato che siamo in un campo $\K$ dove $2 \neq 0$. Dunque abbiamo dimostrato che se $\Phi(F) = \undzero$, allora $F = \undzero$, ossia $\ker \Phi = \{\undzero\}$. Dunque $\Phi$ è iniettiva.
        \item $\Phi$ è suriettiva \\
        Dalla formula vista precedentemente \mbox{$q_F(v+w) = q_F(v) + 2F(v,w) + q_F(w)$}, è facile trovare che $F(v,w) = \frac{1}{2}(q_F(v+w) - q_F(v) - q_F(w))$. Sfruttando questa identità, prendiamo ora una forma quadratica $Q: \V \longrightarrow \K$ e prendiamo l'applicazione bilineare $F$ definita proprio come
        \begin{align*}
        F: \V \times \V &\longrightarrow \K  \\
        (v,w) &\longmapsto \frac{1}{2}(Q(v+w) - Q(v) - Q(w)) 
        \end{align*}
        Se riusciamo ora a verificare che $F \in \Bil^+(\V)$, allora abbiamo dimostrato che $q_F = Q$. Infatti, \mbox{$q_F = F(v,v) = \frac{1}{2} (Q(v+v) - Q(v) - Q(v)) = \frac{1}{2}Q(2v) - Q(v)$}, e dato che, presa una base $\B = \{v_1, \ldots, v_n\}$ di $\V$, $Q(v) = Q(\sum_{i = 1}^n x_i v_i) = \sum_{1 \leq i \leq j \leq n}^n a_{ij} x_i x_j$, essendo $Q$ una forma quadratica, troviamo che \mbox{$Q(2v) = \sum_{1 \leq i \leq j \leq n} a_{ij} 2x_i 2x_j = 4 \sum_{1 \leq i \leq j \leq n} a_{ij} x_i x_j = 4 Q(v)$}. Dunque troviamo che $q_F = \frac{1}{2} Q(2v) - Q(v) = 2Q(v) - Q(v) = Q(v)$. \\
        Per verificare che $F \in \Bil^+(\V)$, dobbiamo verificare che $F$ è bilineare e che è simmetrica. È evidente che questa sia simmetrica, dato che \mbox{$F(v,w) = \frac{1}{2}(Q(v+w) - Q(v) - Q(w)) = \frac{1}{2}(Q(w+v) - Q(w) - Q(v)) = F(w,v)$}. Per dimsotrare che è invece bilineare, prendiamo $F(v,w)$ e la base $\B$ usata in precedenza. Dato che $v = \sum_{i = 1}^n x_i v_i$ e $w = \sum_{i = 1}^n y_i v_i$, troviamo che 
        \begin{align*}
        Q(v+w) &= \sum_{1 \leq i \leq j \leq n} a_{ij}(x_i + y_i) (x_j + y_j)  \\
         &= \sum_{1 \leq i \leq j \leq n} a_{ij}(x_ix_j + x_i y_j + y_i x_j + y_i y_j) \\
         &= \sum_{1 \leq i \leq j \leq n} a_{ij} x_i x_j + \sum_{1 \leq i \leq j \leq n} a_{ij} (x_i y_j + y_i x_j) + \sum_{1 \leq i \leq j \leq n} a_{ij} y_i y_j \\
         &= Q(v) + \sum_{1 \leq i \leq j \leq n} a_{ij} (x_i y_j + y_i x_j) + Q(w)
        \end{align*}
        Troviamo quindi che 
        \begin{align*}
        F(v,w) &= \frac{1}{2}(Q(v + w) - Q(v) - Q(w))  \\
         &= \frac{1}{2}(Q(v) + \sum_{1 \leq i \leq j \leq n} a_{ij}(x_i y_j + y_i x_j) + Q(w) - Q(v) - Q(w)) \\
         &= \frac{1}{2} \sum_{1 \leq i \leq j \leq n} a_{ij} (x_i y_j + y_i x_j)
        \end{align*}
        Prendiamo ora $F(\alpha u + v, w) = \frac{1}{2} \sum_{1 \leq i \leq j \leq n} a_{ij} ((\alpha z_i + x_i) y_j + y_i (\alpha z_j + x_j))$, dove $z_i$ sono le coordinate del vettore $u$ nella base $\B$. Con un po' di calcoli concludiamo che
        \begin{align*}
        F(\alpha u + v, w) &=  \frac{1}{2} \sum_{1 \leq i \leq j \leq n} a_{ij}(x_i y_j + \alpha z_i y_j + y_i x_j + y_i \alpha z_j) \\
         &= \frac{1}{2} \sum_{1 \leq i \leq j \leq n} a_{ij}(\alpha (z_i y_j + y_i z_j) + (x_i y_j + y_i x_j)) \\
         &= \alpha \frac{1}{2} \sum_{1 \leq i \leq j \leq n} a_{ij} (z_i y_j + y_i z_j) + \frac{1}{2} \sum_{1 \leq i \leq j \leq n} a_{ij}(x_i y_j + y_i x_j) \\
         &= \alpha F(u,w) + F(v,w)
        \end{align*}
        Ripetendo lo stesso procedimento a secondo membro, dimsotriamo la bilinearità di $F$. \\
        Dunque $F \in \Bil^+(\V)$, e dunque $\Phi$ è suriettiva, dato che abbiamo dimostrato che per ogni $Q \in Q(\V)$ esiste una $F \in \Bil^+(\V)$ tale che $q_F = Q$.
    \end{itemize}
    Dunque $\Phi$ è biiettiva. $\blacksquare$
\end{itemize}

\subsection{Matrice di Gram}

\begin{Definizione}[Matrice di Gram]
\label{def-matrice-di-gram}
Preso uno spazio vettoriale $\V$ finitamente generato su $\K$, una base $\B = \{v_1, \ldots, v_n\}$ di $\V$ e una forma bilineare simmetrica $F \in \Bil^+(\V)$, la \textbf{matrice di Gram} associata a $F$ dalla scelta di base $\B$ è la matrice $\bm{M_\B(F)} \in M_{nn}(\K)$ la cui entrata $(i,j)$ è pari a $\bm{F(v_i,v_j)}$.
\end{Definizione}

Dunque troviamo che
\[
M_\B(F) = \begin{bmatrix}
    F(v_1,v_1) & \dots & F(v_1, v_n) \\
    \vdots & \ddots & \vdots \\
    F(v_n,v_1) & \dots & F(v_n,v_n)
\end{bmatrix}
\]

\begin{Proposizione}[]
\label{propos-matrice-gram-prodotto}
Preso uno spazio vettoriale $\V$ finitamente generato su $\K$, una base $\B$ di $\V$ e una forma bilineare simmetrica $F \in \Bil^+(\V)$, $\bm{F(v,w) = X_\B(v)^t \cdot M_\B(F) \cdot X_\B(w)} \enspace \forall v,w \in \V$.
\end{Proposizione}

\Dimostrazione \nopagebreak

Chiamiamo per comodità $X = X_\B(v)$, $Y = X_\B(w)$ e $A = M_\B(F)$. Sappiamo che \mbox{$F(v,w) = F(\sum_{i = 1}^n x_i v_i, \sum_{j = 1}^n y_j v_j)$}. Per la bilinearità di $F$, troviamo che $F(v,w) = \sum_{1 \leq i,j \leq n} x_i y_j F(v_i, v_j)$. Se chiamiamo $a_{ij} = F(v_i, v_j)$, troviamo che $F(v,w) = \sum_{1 \leq i,j \leq n} a_{ij} x_i y_j$. Se guardiamo ora il prodotto $X^t \cdot A \cdot Y$, troviamo proprio che 
\[
X^t \cdot A \cdot Y = X^t \cdot \begin{bmatrix}
    a_{11} & \dots & a_{1n} \\
    \vdots & \ddots & \vdots \\
    a_{n1} & \dots & a_{nn} 
\end{bmatrix}
\cdot \begin{bmatrix}
    y_1 \\
    \vdots \\
    y_n 
\end{bmatrix} = (x_1, \ldots, x_n) \cdot \begin{bmatrix}
    \sum_{j = 1}^n a_{1j} y_j \\
     \vdots \\
     \sum_{j = 1}^n a_{nj} y_j 
\end{bmatrix} = \sum_{1 \leq i,j \leq n} a_{ij} x_i y_j = F(v,w) \enspace \blacksquare
\]

\begin{Osservazione}[]
\label{oss-F bilineare simmetrica- MB simmetrica}
$\bm{F \in \Bil^+(\V)} \iff \bm{M_\B(F)}$ \textbf{è simmetrica}, ossia $M_\B(F) = M_\B(F)^t$.
\end{Osservazione}

\Dimostrazione \nopagebreak

\begin{enumerate}[label=\arabic*°]
    \item $F \in \Bil^+(\V) \implies M_\B(F)$ è simmetrica \\
    Dato che $F \in \Bil^+(\V)$ sappiamo che $F(v,w) = F(w,v)$. Dunque troviamo che 
    \[
    M_\B(F) = \begin{bmatrix}
        F(v_1, v_1) & F(v_1, v_2) & \dots & F(v_1,v_n) \\
        F(v_2, v_1) & F(v_2, v_2) & \dots & F(v_2, v_n) \\
        \vdots & \vdots &\ddots & \vdots \\
        F(v_n, v_1) & F(v_n, v_2) & \dots & F(v_n, v_n)
    \end{bmatrix}
    = \begin{bmatrix}
        m_{11} & m_{21} & \dots & m_{n1} \\
        m_{21} & m_{22} & \dots & m_{n2} \\
        \vdots & \vdots & \ddots & \vdots \\
        m_{n1} & n_{n2} & \dots & m_{nn}
    \end{bmatrix} = M_\B(F)^t \quad \blacksquare
    \]
    \item $F \in \Bil^+(\V) \impliedby M_\B(F)$ è simmetrica \\
    Chiamiamo per comodità $X = X_\B(v)$, $Y = X_\B(w)$ e $A = M_\B(F)$. Dato che $F \in \Bil^+(\V)$, sapiamo per la Proposizione \ref{propos-matrice-gram-prodotto} che $F(v,w) = X^t \cdot A \cdot Y$. Dato che $X^t \in M_{1n}(\K), A \in M_{nn}(\K)$ e $Y \in M_{n1}(\K)$, $X^t \cdot A \cdot Y \in M_{11}(\K)$. Dato che la matrice di Gram è simmetrica, sappiamo che $F(v_i, v_j) = F(v_i, v_j)^t$ sempre. Dato che $\B$ è una base di $\V$, troviamo che \mbox{$F(v,w) = F(\sum_{i = 1}^n x_i v_i, \sum_{j = 1}^n y_j v_j) = \sum_{1 \leq i, j \leq n} x_i y_j F(v_i, v_j) = \sum_{1 \leq i,j \leq n} x_i y_j F(v_i, v_j)^t = F(v,w)^t$}. Dunque
    \begin{align*}
    F(v,w) &= F(v,w)^t \\
    &= (X^t \cdot A \cdot Y)^t \\
    &= Y^t \cdot A \cdot X \\
    &= F(w,v)
    \end{align*}
    Concludiamo quindi che $F \in \Bil^+(\V)$. $\blacksquare$
\end{enumerate}

\begin{Definizione}[Matrice associata di una forma quadratica]
\label{def-matrice-associata-forma-quadratica}
Presa una forma quadratica $Q \in Q(\V) = \sum_{1 \leq i \leq j \leq n} a_{ij} x_i x_j$, una forma bilineare simmetrica $F \in \Bil^+(\F)$ tale che $q_F = Q$ e presa una base $\B$ di $\V$, la \textbf{matrice associata in} $\bm{\B}$ della forma quadratica $Q$, ossia la matrice di Gram della forma bilineare $F$, è la matrice
\[
\bm{M_\B(F) = } \begin{bmatrixb}
    a_{11} & \frac{1}{2} a_{12} & \dots & \frac{1}{2}a_{1n} \\
    \frac{1}{2} a_{12} & a_{22} & \dots & \frac{1}{2}a_{2n} \\
    \vdots & \vdots & \ddots & \vdots \\
    \frac{1}{2} a_{1n} & \frac{1}{2} a_{2n} & \dots & a_{nn}
\end{bmatrixb}
\]
\end{Definizione}

Cerchiamo di capire perché è proprio in questa forma. Sappiamo che
\[
M_\B(F) = \begin{bmatrix}
    F(v_1,v_1) & \dots & F(v_1, v_n) \\
    \vdots & \ddots & \vdots \\
    F(v_n,v_1) & \dots & F(v_n,v_n)
\end{bmatrix}
\]
Dato che 
\begin{align*}
Q(v) &= F(v,v) \\
 &= F(\sum_{i = 1}^n x_i v_i, \sum_{j = 1}^n x_jv_j) \\
 &= \sum_{\ \leq i,j \leq n} x_i x_j F(v_i, v_j)  \\
 &= \sum_{i = 1}^n x_i^2 F(v_i, v_j) + \sum_{1 \leq i < j \leq n} 2x_i x_j F(v_i, v_j)
\end{align*}
Scriviamo $2x_ix_j$ al posto di solo $x_ix_j$ perché appaiono sia nell'ordine $x_ix_jF(v_i,v_j)$ che nell'ordine $x_jx_i F(v_j,v_i)$, che però sono uguali dato che $F$ è simmetrica. È vero anche che \mbox{$Q(v) = \sum_{1 \leq i \leq j \leq n} a_{ij} x_i x_j = \sum_{i = 1}^n a_{ii}x_i^2 + \sum_{1 \leq i < j \leq n}  a_{ij} x_i x_j$}, dunque troviamo che \mbox{$F(v_i,v_i) = a_{ii}x_i^2$} e che \mbox{$F(v_i, v_j) = \frac{1}{2}a_{ij} x_i x_j$}. Dunque troviamo proprio che 
\[
M_\B(F) = \begin{bmatrix}
    a_{11} & \frac{1}{2} a_{12} & \dots & \frac{1}{2}a_{1n} \\
    \frac{1}{2} a_{12} & a_{22} & \dots & \frac{1}{2}a_{2n} \\
    \vdots & \vdots & \ddots & \vdots \\
    \frac{1}{2} a_{1n} & \frac{1}{2} a_{2n} & \dots & a_{nn}
\end{bmatrix}
\]

Se per esempio prendiamo la forma quadratica $Q(x_1, x_2) = x_1^2 + 3x_1x_2 + x_2^2$, troviamo che 
\[
M_\B(F) = \begin{bmatrix}
    1 & \frac{3}{2} \\
    \frac{3}{2} & -1
\end{bmatrix}
\]

\begin{Proposizione}[Cambiamento di base nella matrice di Gram]
\label{propos-cambiamento-base-matrice-Gram}
Presa una forma bilineare $F \in \Bil^+(\V)$ e due basi $\B$ e $\Bs$ di $\V$, \mbox{$\bm{M_\Bs(F) = M_\B^\Bs(Id)^t \cdot M_\B(F) \cdot M_\B^\Bs(Id)}$}.
\end{Proposizione}

\Dimostrazione \nopagebreak

Sappiamo per la Proposizione \ref{propos-matrice-gram-prodotto} che $F(v,w) = X_\B(v)^t \cdot M_\B(F) \cdot X_\B(w)$. Dato che per la Proposizione \ref{propos-matrice-associata-prodotto} $X_\B(v) = M_\B^\Bs(Id) \cdot X_\Bs(v)$, troviamo che
\begin{align*}
 F(v,w) &= (M_\B^\Bs(Id)  \cdot X_\Bs(v))^t \cdot M_\B(F) \cdot M_B^\Bs(Id) \cdot X_\Bs(w) \\
 &= X_\Bs(v)^t \cdot M_\B^\Bs(Id)^t \cdot M_\B(F) \cdot M_\B^\Bs(Id) \cdot X_\Bs(w) 
\end{align*}

e quindi $F(v,w) = X_\B(v)^t \cdot X_\Bs(v)^t \cdot M_\B^\Bs(Id)^t \cdot M_\B(F) \cdot M_\B^\Bs(Id) \cdot X_\Bs(w) \cdot X_\B(w)$. Sappiamo anche che $F(v,w) = X_\Bs(v)^t \cdot M_\Bs(F) \cdot X_\B(F)$, dunque
\begin{align*}
X_\Bs(v)^t \cdot M_\B^\Bs(Id)^t \cdot M_\B(F) \cdot M_\B^\Bs(Id) \cdot X_\Bs(w) &= X_\Bs(v)^t \cdot M_\Bs(F) \cdot X_\B(F) \\
M_\Bs(F) &= M_\B^\Bs(Id) \cdot M_\B(F) \cdot M_\B^\Bs(Id) \quad \blacksquare
\end{align*}

\begin{Definizione}[Matrici congruenti]
\label{def-matrici-congruenti}
Due matrici $A,B \in M^+_{nn}(\K)$ (dove $M^+_{nn}(\K)$ indica l'insieme delle matrici $n \times n$ simemtriche) sono \textbf{congruenti} se esiste una matrice $G \in GL_n(\K)$ tale che $\bm{A = G^t \cdot B \cdot G}$. 
\end{Definizione}

È quindi evidente come le matrici $M_\B(F)$ e $M_\Bs(F)$ siano matrici congruenti.

\begin{Osservazione}[]
\label{oss-matrici-congruenti-rango-uguale}
$\bm{A}$ e $\bm{B}$ \textbf{congruenti} $\implies \bm{rg(A) = rg(B)}$.
\end{Osservazione}

\Dimostrazione \nopagebreak

Sappiamo che esiste una matrice invertibile $G$ tale che $A = G^t \cdot B \cdot G$. Dato che \mbox{$rg(A) = \dim Im(A) = \dim \K^n - \dim \ker A = n - \dim \ker A$}, affinché $rg(A) = rg(B)$ ci basta dimostrare che $\dim \ker A = \dim \ker B$. \\
Sappiamo che $\ker A = \{\X \in \K^n \mid A \cdot \X = \undzero\}$, dunque $\X \in \ker A$ se e solo se $G^t \cdot B \cdot G \cdot \X = \undzero$. Dato che $G$ è invertibile, troviamo che $B \cdot G \cdot \X = (G^t)^{-1} \cdot \undzero = \undzero$. Dunque troviamo che $B \cdot G \cdot \X = \undzero$, ossia $G \cdot \X \in \ker B$. Abbiamo appena trovato un isomorfismo tra $\ker A$ e $\ker B$:
\begin{align*}
\Phi: \ker A &\longrightarrow \ker B \\
\X &\longmapsto G \cdot \X 
\end{align*}
Infatti, $\Phi$ è sicuramente lineare, è iniettiva dato che $\ker \Phi = \{\undzero\}$ ($\Phi(\X) = \undzero$ solo se $\X = \undzero$ perché $G$ è invertibile e quindi diverso da $\undzero$), ed è poi anche suriettiva dato che per ogni vettore $Y \in \ker B$ esiste un $\X \in \ker A$ tale che $Y = G \cdot \X$. \\
Dunque troviamo che $\dim \ker A = \dim \ker B$, ossia $rg(A) = rg(B)$. $\blacksquare$

\begin{Proposizione}[Corollario]
\label{corol-forme-equivalenti}
Preso uno spazio vettoriale finitamente generato $\V$, una forma bilineare $F \in \Bil(\V)$ e due basi $\B$ e $\Bs$ di $\V$, $\bm{rg(M_\B(F)) = rg(M_\Bs(F))}$.
\end{Proposizione}

Questo è evidente dato che $M_\B(F)$ e $M_\Bs(F)$ sono congruenti come abbiamo già visto, e per l'Osservazione \ref{oss-matrici-congruenti-rango-uguale} il loro rango è uguale.


\begin{Osservazione}[]
\label{oss-relazione-congruenza-equivalenz}
La relazione di congruenza è una \textbf{relazione di equivalenza}.
\end{Osservazione}

\Dimostrazione \nopagebreak

Per essere una relazione d'equivalenza la relazione di coniugio deve essere riflessiva, simmetrica e transitiva.
\begin{itemize}
    \item Riflessività \\
    $A$ è sempre congruente con $A$, dato che $A = 1_n^t \cdot A \cdot 1_n = 1_n \cdot A = A$.
    \item Simmetria \\
    Se $A$ è congruente con $B$, allora $B$ è congruente con $A$, essendo $G$ invertibile, e troviamo che $B = (G^{-1})^t \cdot A \cdot G^{-1}$. \\
    \item Transitività \\
    Se $A$ è congruente con $B$ e $B$ è congruente con $C$, allora abbiamo che \mbox{$A = G^t \cdot B \cdot G = G^t \cdot H^t \cdot C \cdot H \cdot G = (H \cdot G)^t \cdot C \cdot (H \cdot G)$}. Essendo $H$ e $G$ matrici invertibili, anche $H \cdot G$ lo è, dunque troviamo che $A$ e $C$ sono congruenti. 
\end{itemize}

Possiamo quindi concludere che la relazione di congruenza è una relazione d'equivalenza. $\blacksquare$

\section{Spazi vettoriali euclidei}

\begin{Definizione}[Forma bilineare definita positiva]
\label{def-forma-bilineare-definita-positiva}
Preso uno spazio vettoriale $\V$ su $\R$ e una forma bilineare simmetrica $F \in \Bil^+(\V)$, questa è \textbf{definita positiva} se, perso $v \neq \undzero$, $\bm{F(v,v) > 0}$.
\end{Definizione}

\begin{Definizione}[Prodotto scalare euclideo]
\label{def-prodotto-scalare-euclideo}
Preso uno spazio vettoriale $\V$ su $\R$, il \textbf{prodotto scalare euclideo su} $\bm{\V}$ è una forma bilineare simmetrica e definita positiva. Lo indichiamo con lo stesso simoblo usato per il prodotto scalare euclideo su $V(\E^3)$ $\bm{\langle v,w \rangle}$.
\end{Definizione}

\begin{Definizione}[Spazio vettoriale euclideo]
\label{def-spazio-vettoriale-euclideo}
Uno \textbf{spazio vettoriale eucldeo} è uno spazio vettoriale su $\R$ con prodotto scalare euclideo (è quindi uno spazio vettoriale dove è definito il prodotto scalare euclideo). Lo indichiamo con il simbolo $\bm{(\V, \langle , \rangle)}$.
\end{Definizione}

Abbiamo già visto che il prodotto scalare euclideo dello spazio euclideo è definito come \mbox{$\braket{v,w} = \norm{v} \cdot \norm{w} \cdot \cos(\theta)$}. Notiamo subito come, se $v \neq 0$, $\braket{v,v} > 0$ dato che $\norm{v} > 0$ (perché $\norm{v}$ rappresenta la lunghezza del vettore $v$ e quindi non può mai essere negativa) e $\cos(\theta) = \cos(0) = 1$.

Altri due esempi importanti da ricordare sono:
\begin{itemize}
    \item $\V = \R^n$ con $\braket{X,Y} := X^t \cdot Y = \sum_{i = 1}^n x_i y_i$. Questo prodotto in particolare prende il nome di \textbf{prodotto scalare standard}.
    \item $\V = C^0([a,b])$ (ossia l'insieme delle funzioni conitnue nell'intervallo $[a,b]$), dove $a,b \in \R$, con \mbox{$\braket{f,g} := \int_a^b f(t) g(t) \ dt$}.
\end{itemize}

\begin{Proposizione}[Corollario]
\label{corollario-norma-v}
Dato che $\braket{v,v} = \norm{v} \cdot \norm{v} \cdot \cos(0) = \norm{v}^2$, $\normbm{v} \bm{= \langle v,v \rangle^{\frac{1}{2}}}$.  \\
Allo stesso modo, $\normbm{\lambda v} = \braket{\lambda v, \lambda v}^{\frac{1}{2}} = \lambda^2 \braket{v,v}^{\frac{1}{2}} = \bm{\lvert \lambda \rvert \cdot} \normbm{v}$.
\end{Proposizione}

\begin{Proposizione}[Disuguaglianza di Cauchy-Schwarz]
\label{propos-disuguaglianza-cauchy-schwarz}
Preso uno spazio vettoriale euclideo $(\V, \braket{,})$, $\bm{v,w \in \V} \implies \bm{\lvert \langle v,w \rangle \rvert \leq} \normbm{v} \bm{\cdot} \normbm{w}$, e in particolare \mbox{$\abs{\braket{v,w}} = \norm{v} \cdot \norm{w} \iff v$ e $w$ sono \textbf{linearmente dipendenti}}.
\end{Proposizione}

\Dimostrazione \nopagebreak

Partiamo col dimostrare il caso in cui $\braket{v,w} = \norm{v} \cdot \norm{w} \cdot \cos(\theta)$, per poi dimostrare il caso generale. \\
In questo caso, $\abs{\braket{v,w}} = \abs{\norm{v} \cdot \norm{w} \cdot \cos(\theta)}$. Dato che $\abs{\cos(\theta)} \leq 1$, troviamo che \mbox{$\abs{\braket{v,w}} \leq \abs{\norm{v} \cdot \norm{w}} = \norm{v} \cdot \norm{w}$}, dato che $\norm{v} \geq 0$. Dunque troviamo che $\abs{\braket{v,w}} \leq \norm{v} \cdot \norm{w}$.

Passiamo ora al caso generale, dove distinguiamo due casi:
\begin{itemize}
    \item $v = 0$ \\
    In questo caso abbiamo $\braket{0,w} = \braket{0 + 0, w} = \braket{0, w} + \braket{0,w}$ e dunque $\braket{0,w} = 0$. Dunque troviamo che $\abs{\braket{0,w}} = 0 \leq \norm{0} \cdot \norm{w} = 0$, che è una disuguaglianza sempre vera. $\blacksquare$
    \item $v \neq 0$ \\
    In questo caso prendiamo una variabile $x \in \R$ e definiamo il polinomio \mbox{$p(x):= \braket{xv + w, xv + w} = \norm{v}^2 x^2 + 2 \braket{v,w} + \norm{w}^2$} per la bilinearità del prodotto scalare euclideo. Dato che $p(x) \geq 0$ sempre essendo il prodotto scalare euclideo definito positivo, troviamo che questo polinomio di secondo grado (secondo perché siamo sicuri che $v \neq 0$ e dunque $\norm{v} \neq 0$) può solo non avere radici reali o averne una sola doppia, essendo una parabola positiva verso l'alto con il vertice che al più si trova sull'asse $x$.
    Dunque troviamo che $\Delta = (2\braket{v,w})^2 - 4 \norm{v}^2 \cdot \norm{w}^2 \leq 0$. Da questa disequazione troviamo proprio la disequazione di Cauchy-Schwarz \mbox{$\abs{\braket{v,w}} \leq \norm{v}^2 \cdot \norm{w}^2$}. $\blacksquare$ \\
    Notiamo, infine, che $\Delta = 0$ se e solo se esiste una radice reale di $p(x)$, ossia esiste una $x$ tale che $xv + w = 0$, e dato che $v \neq 0$, troviamo che $v$ e $w$ sono linearmente dipendenti. $\blacksquare$
\end{itemize}

\begin{Proposizione}[Disuguaglianza triangolare]
\label{propos-disuguaglianza-triangolare}
Preso uno spazio vettoriale euclideo $(\V, \braket{v})$, $\bm{v,w \in \V} \implies \normbm{v + w} \bm{\leq} \normbm{v} \bm{+} \normbm{w}$ e in particolare $\norm{v + w} = \norm{v} + \norm{w} \iff v$ e $w$ sono \textbf{linearmente dipendenti}.
\end{Proposizione}

\Dimostrazione \nopagebreak

Questa proposizione è una diretta conseguenza della disuguaglianza di Cauchy-Schwarz. Prende il nome di disuguaglianza \textit{triangolare} perché, se prendiamo come spazio vettoriale lo spazio euclideo, se rappresentiamo i vettori $v,w$ e $v + w$ graficamente notiamo proprio come questi vadano a formare un triangolo.

\begin{minipage}[c]{0.35\textwidth}
\centering
\includegraphics[width=0.7\textwidth, valign=t]{Disuguaglianza triangolare.png}
\end{minipage}\hfill
\begin{minipage}[c]{0.60\textwidth}
È evidente infatti guardando l'immagine che \mbox{$\norm{v+w} \leq \norm{v} + \norm{w}$, e $\norm{v + w} = \norm{v} + \norm{w}$} solo se $v$ e $w$ sono sulla stessa retta, ossia nel caso in cui sono linearmente dipendenti.
\end{minipage}

Dimostriamolo ora algebricamente. \\
Sappiamo già $\norm{v} \geq 0$ sempre, dunque si può dimostrare che \mbox{$\norm{v+w} \leq \norm{v} + \norm{w}$} dimostrando che $\norm{v+w}^2 \leq (\norm{v} + \norm{w})^2$.
\begin{align*}
\norm{v+w}^2 &\leq (\norm{v} + \norm{w})^2 \\
\braket{v+w, v+w} &\leq \norm{v}^2 + 2 \norm{v} \cdot \norm{w} + \norm{w}^2 \\
\braket{v,v} + 2 \braket{v,w} + \braket{w,w} &\leq \norm{v}^2 + 2 \norm{v} \cdot \norm{w} + \norm{w}^2 \\
\norm{v}^2 + 2 \braket{v,w} + \norm{w}^2 &\leq \norm{v}^2 + 2 \norm{v}\cdot \norm{w} + \norm{w}^2 \\
\braket{v,w} &\leq \norm{v} \cdot \norm{w} 
\end{align*}
Questa è proprio la disuguaglianza di Cauchy-Schwarz, che abbiamo già dimostrato essere sempre vera. Dunque anche la disuguaglianza triangolare è dimostrata. $\blacksquare$ \\
Notiamo, infine, che l'uguaglianza è verificata solo quando i vettori sono dipendenti, dato che l'uguaglianza nella disuguaglianza di Cauchy-Schwarz è verificata solo quando sono dipendenti. $\blacksquare$

\begin{Definizione}[Distanza tra due punti]
\label{def-distanza-punti}
Preso uno spazio affine $\S$ su $\R$, un prodotto scalare euclideo $\braket{,}$ su $V(\S)$ e due punti $P,Q \in \S$, la \textbf{distanza} tra $P$ e $Q$ è $\bm{d(P,Q) = \lVert} \vvbm{PQ} \bm{\rVert}$.
\end{Definizione}

\begin{Proprieta}[Proprietà della distanza tra due punti]
\label{prop-distanza-tra-punti}
Preso uno spazio affine $\S$ su $\R$ e un prodotto scalare euclideo $\braket{,}$ su $V(\S)$
\begin{enumerate}
    \item $\forall P,Q \in \S \enspace \bm{d(P,Q) \geq 0}$ e $\bm{d(P,Q) = 0} \iff \bm{P = Q}$ 
    \item $\forall P,Q \in \S \enspace \bm{d(P,Q) = d(Q,P)}$ 
    \item $\forall P,Q,R \in \S \enspace \bm{d(P,Q) \leq d(P,R) + d(R,Q)}$
\end{enumerate}
\end{Proprieta}

\Dimostrazione \nopagebreak

\begin{enumerate}
    \item Banalmente, $d(P,Q) = \norm{\vv{PQ}}$ che per definizione è sempre $\geq 0$, e se $P = Q$ allora $d(P,Q) = \norm{\vv{PQ}} = \norm{0} = 0$. $\blacksquare$
    \item Geometricamente questo è evidente: la distanza da un punto $P$ a uno $Q$ è ovviamente uguale alla distanza da $Q$ a $P$. Algebricamente, si vede subito come \mbox{$d(P,Q) = \norm{\vv{PQ}} = \norm{-\vv{QP}} = \abs{-1} \norm{\vv{QP} = \norm{\vv{QP}}}$}. $\blacksquare$
    \item Sappiamo che in uno spazio affine $\vv{PQ} = \vv{PR} + \vv{RQ}$, dunque per la disuguaglianza triangolare troviamo che è vero che $\norm{\vv{PR} + \vv{RQ}} = \norm{\vv{PQ}} \leq \norm{\vv{PR}} + \norm{\vv{RQ}}$, e dato che \mbox{$d(P,Q) = \norm{\vv{PQ}}, d(P,R) = \norm{\vv{PR}}$} e $d(RQ) = \norm{\vv{RQ}}$, troviamo proprio che \mbox{$d(P,Q) \leq d(P,R) + d(R,Q)$}. $\blacksquare$
\end{enumerate}

\subsection{Ortogonalità}

\begin{Definizione}[Vettori ortogonali]
\label{def-vettori-ortogonali}
In uno spazio vettoriale euclideo $(\V, \braket{,})$ due vettori $v,w$ sono detti \textbf{ortogonali} ($\bm{v \perp w}$) se $\braketbm{v,w} \bm{= 0}$. 
\end{Definizione}

Ad esempio, nello spazio vettoriale $\V = C^0([-\pi, \pi])$, le funzioni $1, \sin t, \cos t, \sin 2t, \cos 2t, \ldots, \sin kt, \cos kt$ con $k \in \N$ sono funzioni a due a due ortogonali. Infatti, funzioni una pari e una dispari, il cui prodotto è sempre una funzione dispari. Dunque $\braket{f,g} = \int_{-\pi}^\pi f(t) g(t) \ dt = 0$.

\begin{Definizione}[Base ortonormale]
\label{def-base-ortonormale-totale}
Una base $\B = \{v_1, \ldots, v_n\}$ di $\V$ è una base \textbf{ortonormale} (\textbf{ON}) se 
\[
\braketbm{v_i,v_j} \bm{= \delta_{ij} =}
\begin{cases}
\bm{1} &\enspace \text{se } i = j \\
\bm{0} &\enspace \text{se }i \neq j 
\end{cases}
\]
Ossia $\norm{v_i} = 1 \enspace \forall i \in \{1, \ldots, n\}$ e $v_i \perp v_j \enspace \forall i \neq j$.
\end{Definizione}

\begin{Osservazione}[]
\label{oss-base-ortonormale}
Presa una base ortonormale $\B = \{v_1, \ldots, v_n\}$ di $\V$, $\braketbm{\sum_{i = 1}^n x_i v_i, \sum_{j = 1}^n y_j v_j} \bm{= \sum_{i = 1}^n x_i y_i}$.
\end{Osservazione}

\Dimostrazione \nopagebreak

Per bilinearità del prodotto scalare euclideo, troviamo facilmente che $\braket{\sum_{i = 1}^n x_i v_i, \sum_{j = 1}^n y_j v_j} = \sum_{1 \leq i,j\leq n} x_i y_j \braket{v_i, v_j}$. Dato che ci troviamo in una base ortonormale, se $i = j$ allora $\braket{v_i,v_i} = 1$, altrimenti $\braket{v_i,v_j} = 0$. Dunque troviamo che $\braket{\sum_{i = 1}^n x_i v_i, \sum_{j = 1}^n y_j v_j} = \sum_{i = 1}^n x_i y_i$. $\blacksquare$

\begin{Proposizione}[]
\label{propos-base-ON}
Ogni spazio vettoriale euclideo $(\V, \braket{,})$ finitamente generato \textbf{ha una base ON}.
\end{Proposizione}

\Dimostrazione \nopagebreak

Dimostriamo questa proposizione per induzione. 
\begin{itemize}
    \item Caso base: $\dim \V = 0$ \\
    In questo caso abbiamo $\V = \{\undzero\}$, e allora una base ON di $\V$ è banalmente la base vuota $\B = \{\}$. $\blacksquare$
    \item Ipotesi induttiva \\
    Assumiamo la proposizione vera per tutti gli spazi vettoriali di dimensione $n-1$.
    \item Passo induttivo \\
    Abbiamo uno spazio vettoriale $\V$ la cui dimensione è $n$. Prendiamo un vettore non nullo $v \in \V$. Se poniamo $v_1 = \lambda v$, sappiamo che \mbox{$\norm{v_1} = \norm{\lambda v} = \abs{\lambda} \norm{v}$}. Se poniamo ora $\lambda = \frac{1}{\norm{v}}$, che ha senso dato che $v \neq 0$, troviamo che \mbox{$\norm{v_1} = \left|{\frac{1}{\norm{v}}}\right| \norm{v} = 1$}. Dunque il vettore $v_1 = \frac{v}{\norm{v}}$ è un vettore con norma $1$. Prendiamo ora il sottospazio $v_1^\perp := \{v \in \V \mid \braket{v_1, v} = 0\}$ che rappresenta tutti i vettori ortogonali a $v_1$. Possiamo vederlo anche come il nucleo dell'applicazione
    \begin{align*}
    \Phi: \V &\longrightarrow \R \\
    v &\longmapsto \braket{v_1,v}
    \end{align*}
    Questa applicazione è lineare, ma soprattutto è suriettiva. Infatti, preso un qualsiasi scalare $\lambda \in R$, basta prendere $v = \lambda v_1 \in \V$ per trovare che $\braket{v_1, v} = \braket{v_1, \lambda v_1} = \lambda \braket{v_1, v_1} = \lambda$. Dunque qualsiasi valore in $\R$ è rappresentato da questa applicazione. \\
    Dato che $\dim \Phi = \dim \ker \Phi + \dim Im(\Phi) = \dim v_1^\perp + \dim Im(\Phi)$ essendo $\Phi$ lineare, e dato che $\dim \Phi = \dim \R = 1$, troviamo che $\dim v_1^\perp = n - 1$. Per ipotesi induttiva, esiste una base ON $\B = \{v_2, \ldots, v_n\}$ di $v_1^\perp$. Dato che $\V = \braket{v_1} \oplus v_1^\perp$, e dato che $\braket{v_1} \cap v_1^\perp = \{\undzero\}$, per la Proposizione \ref{propos-dim-spaz-vett-2} troviamo che $\Bs = \{v_1, \ldots, v_n\}$ è una base ON di $\V$. $\blacksquare$
\end{itemize}

\begin{Definizione}[Spazi vettoriali euclidei equivalenti]
\label{def-spazi-vettoriali-euclidei-equivalenti}
Due spazi vettoriali euclidei $(\V, \braket{,}_\V)$ e $(\W, \braket{,}_\W)$ con $\dim \V = \dim \W$ sono detti \textbf{equivalenti}, ossia esiste un isomorfismo che li collega.
\end{Definizione}

\section{Forme bilineari simmetriche equivalenti}

\begin{Definizione}[Forme quadratiche equivalenti]
\label{def-forme-quadratiche-equivalenti}
Due forme quadratiche $Q_1, Q_2 \in Q(\V)$ sono dette \textbf{equivalenti} se esistono due basi $\B = \{v_1, \ldots, v_n\}$ e $\Bs = \{w_1, \ldots, w_n\}$ di $\V$ tali che \mbox{$\bm{Q_1(\sum_{i = 1}^n x_i v_i) = Q_2(\sum_{i = 1}^n y_i w_i)}$}. 
\end{Definizione}

\begin{Definizione}[Forme bilineari simmetriche equivalenti]
\label{def-forme-bilineari-simmmetriche-equivalenti}
Due forme bilineari simmetriche $F_1, F_2 \in \Bil^+(\V)$ sono \textbf{equivalenti} se prese due basi $\B$ e $\Bs$ di $\V$, \mbox{$\bm{M_\B(F_1) = M_\Bs(F_2)}$}.
\end{Definizione}

\begin{Definizione}[Rango di una forma bilineare]
\label{def-rango-forma-bilienare}
Il \textbf{rango} di una forma bilineare è il rango di una sua matrice di Gram in una sua base qualsiasi $\B$, ossia $\bm{rg(F) = rg(M_\B(F))}$. \\
Ha senso dare questa definizione dato che per il Corollario \ref{corol-forme-equivalenti} il rango di tutte le matrici di Gram di una forma bilineare è lo stesso.
\end{Definizione}

\Esempio \nopagebreak

Trovare il \textbf{rango} della forma bilineare simmetrica \nopagebreak
\begin{align*}
\bm{F: \K^n \times \K^n} &\longrightarrow \bm{\K}  \\
\bm{(X,Y)} &\longmapsto \bm{\sum_{i = 1}^d x_i y_i} 
\end{align*}

Prendiamo la base standard $S = \{e_1, \ldots, e_n\}$ e studiamo la matrice di Gram di $F$ in tale base:
\[
M_S(F) = \begin{bmatrix}
    F(e_1,e_1) & F(e_1, e_2) & \dots & F(e_1, e_n) \\
    F(e_2, e_1) & F(e_2, e_2) & \dots & F(e_2, e_n) \\
    \vdots & \vdots & \ddots & \vdots \\
    F(e_n, e_1) & F(e_n, e_1) & \dots & F(e_n, e_n)
\end{bmatrix}
\]

In generale, 
\[
F(e_i,e_j) = 
\begin{cases}
1 & \enspace \text{se } i = j \text{ e } i \leq d\\
0 & \enspace \text{altrimenti}
\end{cases}
\]
Questo perché, se $i = j$ allora il prodotto $x_i y_i$ vale $1$, ma nella sommatoria questo prodotto è presente solo se $i \leq d$. Tutti gli altri prodotti, in qualsiasi altro caso, hanno almeno uno dei due fattori che è pari a $0$. \\
Dunque troviamo che
\[
M_S(F) = 
\begin{array}{@{}c@{}}
    \begin{bmatrix}
        1 & 0 & \dots & 0 & 0 & \dots & 0 \\
        0 & 1 & \dots & 0 & 0 & \dots & 0 \\
        \vdots & \vdots & \ddots & \vdots & \vdots & \ddots & \vdots \\
        0 & 0 & \dots & 1 & 0 & \dots & 0 \\
        0 & 0 & \dots & 0 & 0 & \dots & 0 \\
        \vdots & \vdots & \ddots & \vdots & \vdots & \ddots & \vdots \\
        0 & 0 & \dots & 0 & 0 & \dots & 0 
    \end{bmatrix} \\[-2ex]
    \begin{array}{@{}cc@{}}
        \underbrace{\rule{5.5em}{0pt}}_{d \text{ colonne}} & \rule{3.5em}{0pt}
    \end{array}
\end{array}
\]

E quindi $\bm{rg(F)} = rg(M_S(F)) = \bm{d}$. 

\begin{Proposizione}[Corollario]
\label{corol-rango-forme-equivalenti}
Prese due forme bilineari $F, G \in \Bil(F)$ equivalenti, $\bm{rg(F) = rg(G)}$.
\end{Proposizione}

Questo è ovvio dato che se $F$ e $G$ sono equivalenti allora esistono due basi $\B$ e $\Bs$ tali che $M_\B(F) = M_\Bs(F)$, e dato che $rg(F) = rg(M_\B(F))$ e $rg(G) = rg(M_\Bs(G))$, troviamo proprio che $rg(F) = rg(G)$.

\subsection{Vettori e basi F-ortogonali}

\begin{Definizione}[Vettori F-ortogonali]
\label{def-vettori-F-ortogonali}
Preso uno spazio vettoriale $\V$ e una forma bilineare simmetrica $F \in \Bil^+(\V)$, due vettori $v,w \in \V$ sono \textbf{F-ortogonali} ($\bm{v \perp w}$) se $\bm{F(v,w) = 0}$. \\
\NB Se $F$ nel contesto è omettibile, si può anche dire semplicemente che $v$ e $w$ sono ortogonali.
\end{Definizione}

\begin{Definizione}[Base F-ortogonale]
\label{def-base-F-ortogonale}
Preso uno spazio vettoriale $\V$ finitamente generato e una forma bilineare simmetrica $F \in \Bil^+(\V)$, una base $\B = \{v_1, \ldots, v_n\}$ di $\V$ è \textbf{F-ortogonale} se $\bm{F(v_i, v_j) = 0} \enspace \forall i \neq j$, ossia $\bm{v_i \perp v_j} \enspace \forall i \neq j$. \\
\NB Se $F$ nel contesto è omettibile, si può anche semplicemente dire che $\B$ è ortogonale.
\end{Definizione}

\begin{Osservazione}[]
\label{oss-MB-diagonale}
$\B = \{v_1, \ldots, v_n\}$ \textbf{base ortogonale} $\implies \bm{M_\B(F)}$ \textbf{è diagonale}.
\end{Osservazione}

\Dimostrazione \nopagebreak

Banalmente,
\[
M_\B(F) = \begin{bmatrix}
    F(v_1,v_1) & F(v_1,v_2) & \dots & F(v_1,v_n) \\
    F(v_2,v_1) & F(v_2,v_2) & \dots & F(v_2,v_n) \\
    \vdots & \vdots & \ddots & \vdots \\
    F(v_n,v_1) & F(v_n,v_2) & \dots & F(v_n,v_n)
\end{bmatrix} = \begin{bmatrix}
    F(v_1,v_1) & 0 & \dots & 0 \\
    0 & F(v_2, v_2) & \dots & 0 \\
    \vdots & \vdots & \ddots & \vdots \\
    0 & 0  & \dots & F(v_n, v_n)
\end{bmatrix}
\quad \blacksquare
\]

\begin{Proposizione}[]
\label{propos-base-ortogonale}
Ogni forma bilineare simmetrica $F \in \Bil^+(F)$ \textbf{ha una base F-ortogonale di} $\bm{\V}$. 
\end{Proposizione}

\Dimostrazione \nopagebreak

La dimostrazione è omessa perché pressoché identica alla dimostrazione della Proposizione \ref{propos-base-ON}.

\begin{Proposizione}[]
\label{propos-spazio-vettoriale-C-forme-bilineari-simmetriche-rango}
Preso uno spazio vettoriale $\V$ sul campo dei numeri complessi $\C$, \mbox{$\bm{F,G \in \Bil^+(\V)}$ \textbf{sono equivalenti} $\iff \bm{rg(F) = rg(G)}$}.
\end{Proposizione}

\Dimostrazione \nopagebreak

\begin{enumerate}[label=\arabic*°]
    \item $F,G \in \Bil^+(\V)$ sono equivalenti $\implies rg(F) = rg(G)$ \\
    Già dimostrato nel Corollario \ref{corol-rango-forme-equivalenti}.
    \item $F,G \in \Bil^+(\V)$ sono equivalenti $\impliedby rg(F) = rg(G)$ \\
    Sappiamo per la Proposizione \ref{propos-base-ortogonale} che per $F$ esiste una base F-ortogonale $\B = \{v_1, \ldots, v_n\}$ e che per $G$ anche esiste una base G-ortogonale $\Bs = \{w_1, \ldots, w_n\}$. Per l'Osservazione \ref{oss-MB-diagonale} che $M_\B(F)$ e $M_\Bs(G)$ sono entrambe matrici diagonali con coefficienti sulla diagonale pari rispettivamente a $\lambda_i = F(v_i,v_i) \neq 0$ e a $\mu_i = G(w_i,w_i) \neq 0 \enspace \forall i \in \{1, \ldots, r\}$ dove $r$ è il rango di $F$ e $G$. Prendiamo ora due nuove basi $\B' = \{v_1', \ldots, v_n'\}$ dove
    \[
    v_i' =
    \begin{cases}
    \frac{v_i}{\sqrt{\lambda_i}} & \text{se } i \leq r \\
    v_i & \text{se } i > r 
    \end{cases}
    \]
    e $\Bs = \{w_1', \ldots, w_n'\}$ dove
    \[
    w_i' =
    \begin{cases}
    \frac{w_i}{\sqrt{\mu_i}} & \text{se } i \leq r\\
    w_i & \text{se } i > r 
    \end{cases}
    \]
    Questo ha senso dato che $\lambda_i, \mu_i \neq 0$ e, trovandoci nel campo dei numeri complessi, $\sqrt{\lambda_i}$ e $\sqrt{\mu_i}$ sono sempre definiti. Queste basi rimangono ortogonali, dunque anche $M_{\B'}(F)$ e $M_{\Bs'}(G)$ sono diagonali, con entrate $F(v_i',v_i')$ e $G(w_i',w_i')$. In particolare,
    \[
    F(v_i', v_i') = 
    \begin{cases}
    F(\frac{v_i}{\sqrt{\lambda_i}}, \frac{v_i}{\sqrt{\lambda_i}}) = \frac{1}{\lambda_i} F(v_i,v_i) = \frac{\lambda_i}{\lambda_i} = 1 & \text{se } i \leq r\\
    F(v_i, v_i) = 0 & \text{se } i > r 
    \end{cases}
    \]
    e allo stesso modo
    \[
    G(w_i', w_i') = 
    \begin{cases}
    G(\frac{w_i}{\sqrt{\mu_i}}, \frac{w_i}{\sqrt{\mu_i}}) = \frac{1}{\mu_i} F(w_i,w_i) = \frac{\mu_i}{\mu_i} = 1 & \text{se } i \leq r\\
    G(w_i, w_i) = 0 & \text{se } i > r 
    \end{cases}
    \]
    Troviamo così che
    \[
    M_{\B'(F)} = \begin{array}{@{}c@{}}
    \begin{bmatrix}
        1 & 0 & \dots & 0 & 0 & \dots & 0 \\
        0 & 1 & \dots & 0 & 0 & \dots & 0 \\
        \vdots & \vdots & \ddots & \vdots & \vdots & \ddots & \vdots \\
        0 & 0 & \dots & 1 & 0 & \dots & 0 \\
        0 & 0 & \dots & 0 & 0 & \dots & 0 \\
        \vdots & \vdots & \ddots & \vdots & \vdots & \ddots & \vdots \\
        0 & 0 & \dots & 0 & 0 & \dots & 0 
    \end{bmatrix} \\[-2ex]
    \begin{array}{@{}cc@{}}
        \underbrace{\rule{5.5em}{0pt}}_{r \text{ colonne}} & \rule{3.5em}{0pt}
    \end{array}
\end{array}
= M_{\Bs'(G)}
    \]
    Dunque $F$ e $G$ sono equivalenti. $\blacksquare$
\end{enumerate}

\begin{Definizione}[Forma bilineare simmetrica definita positiva o negativa]
\label{def-forma-definita-positiva-negativa}
Presa una forma bilineare simmetrica $F \in \Bil^+(\V)$ questa è \textbf{definita positiva} se $\bm{F}$ è un \textbf{prodotto scalare euclideo}. È invece definita \textbf{negativa} se $\bm{-F}$ è un \textbf{prodotto scalare euclideo}.
\end{Definizione}

\begin{Definizione}[Segnatura positiva e negativa di $\bm{F}$]
\label{def-segnatura-positiva}
Presa una forma bilineare simmetrica $F \in \Bil^+(F)$, la \textbf{segnatura positiva} $\bm{\sigma_+(F)}$ è la \textbf{massima dimensione} di un sottospazio $\W \subset \V$ tale che la restrizione $F|_\W$ (la forma bilineare ottenuta da $F$ usando come dominio il sottoinsieme $\W \times \W$ al posto del sottoinsieme $\V \times \V$) è \textbf{defintia positiva}. \\
La \textbf{segnatura negativa} $\bm{\sigma_-(F)}$ è invece la \textbf{massima dimensione} di un sottospazio $\W \subset \V$ tale che la restrizione  $F|_\W$ sia \textbf{definita negativa}.
\end{Definizione}

\begin{Proposizione}[]
\label{propos-calcolare-segnatura}
Preso uno spazio vettoriale $\V$, una forma bilineare simmetrica $F \in \Bil^+(\V)$ e una base ortogonale $\B = \{v_1, \ldots, v_n\}$ di $\V$, sappiamo che $M_\B(F)$ è una matrice diagonale con entrate sulla diagonale $\lambda_1, \ldots, \lambda_n$. Allora
\[
\begin{cases}
\bm{\sigma_+(F) = \left|\{\lambda_i \mid \lambda_i > 0\}\right|} & \\
\bm{\sigma_-(F) = \left|\{\lambda_i \mid \lambda_i < 0\}\right|} &
\end{cases}
\]
È evidente quindi come $\bm{rg(F) = \sigma_+(F) + \sigma_-(F)}$.
\end{Proposizione}

\Dimostrazione \nopagebreak

Ci limiteremo a dimostrare che $\sigma_+(F) = \abs{\{\lambda_i \mid \lambda_i > 0\}}$, dato che la dimostrazione per $\sigma_-(F)$ è identica ma specchiata.

Chiamiamo $p = \{\lambda_i \mid \lambda_i > 0\}$ l'insieme delle entrate positive di $M_\B(F)$ e $n$ la dimensione di $\V$, e cerchiamo quindi di dimostrare che $\sigma_+(F) = p$.
\begin{enumerate}[label=\arabic*°]
    \item $\sigma_+(F) \geq p$ \\
    Prendiamo la base $\B$ e riordiniamola in modo tale da avere tutti i primi $p$ vettori come i vettori a cui sono associati i $\lambda_i$ positivi. Dunque fissiamo la base $\B' = \{v_1', \ldots, v_p', \ldots, v_n'\}$. Prendiamo ora il sottospazio $\W_+ = \braket{v_1',\ldots,v_p'}$. Sappiamo che $\dim \W_+ = p$. Preso un vettore $\X \in \W_+$, sappiamo che $F(\X,\X) = q_F(\X) = \lambda_1 x_1^2 + \ldots + \lambda_p x_p^2$. Dato che $\lambda_i > 0$ con $i \in \{1, \ldots, p\}$, troviamo che $F(\X, \X) > 0$ se prendiamo $\X \neq \undzero$, dunque $F$ è definita positiva, e dunque sicuramente $\sigma_+(F) \geq p$.
    \item $\sigma_+(F) \leq p$ \\
    Sia $\W \subset \V$ il sottospazio vettoriale di $\V$ dove $F(w,w) > 0 \enspace \forall w \in \W$ con $w \neq 0$. Sia invece $\U \subset \W$ il sottospazio vettoriale di $\V$ dove al contrario $F(u,u) \leq 0 \enspace \forall u \in \U$. Se riprendiamo la base vista prima $\B'$, possiamo dire che $\U = \braket{v_{p+1}', \ldots, v_n'}$, e che quindi $\dim \U = n - p$. Per la definizione stessa dei due sottospazi, troviamo che $\W \cap \U = \{\undzero\}$, e dunque $\dim (\U \cap \W) = 0$. Utilizzando la formula di Grassmann, troviamo che \mbox{$\dim (\W + \U) = \dim \W + \dim \U - \dim (\W \cap \U) = \dim \W + \dim \U = \dim \W + n - p$}. Poiché $\W + \U \subset \V$, troviamo che \mbox{$n \geq \dim(\W + \U) = \dim \W + n - p$}, e quindi $\dim \W \leq p$. Dato che $\dim \W = \sigma_+(F)$, troviamo proprio che $\sigma_+(F) \leq p$.    
\end{enumerate}
Dato che $\sigma_+(F) \geq p$ e allo stesso tempo $\sigma_+(F) \leq p$, troviamo proprio che $\sigma_+(F) = p$. $\blacksquare$

\begin{Proposizione}[]
\label{propos-finale}
Preso uno spazio vettoriale sul campo $\R$,
\[
\bm{F,G \in \Bil^+(F)} \textbf{ sono equivalenti } \iff
\begin{cases}
\bm{\sigma_+(F) = \sigma_+(G)} & \\ 
\bm{\sigma_-(G) = \sigma_-(G)} & \\ 
\end{cases}
\]
\end{Proposizione}
\end{document}